\bookpart{图像处理}{Image Processing}
\label{part:image-processing}

\BookPartGuide{%
一张照片可以用于判断对象类别、确定空间范围，也可以用于恢复受损内容或制作新图像。不同任务对输出和允许修改的内容有不同要求。本部分从任务目标出发，说明这些要求怎样决定计算过程与评价依据。

图像识别与检索章研究类别、区域、测量量和匹配关系；图像复原、增强与压缩章研究内容保真、呈现质量与比特成本；图像生成与编辑章研究条件创建、转换、局部修改和新视角合成。各章使用第一部分的基础知识和第二部分介绍的模型，保留输入、结果和来源之间的对应，区分教学算例与模型实验。}

\chapter{图像识别与检索}
\label{chap:vision-image-recognition}

一张路口照片中有两辆车，一辆红色货车靠近相机，一辆蓝色轿车停在公交站旁。说出“这是街道”，还没有说明车在哪里；给出两个框，还没有说明哪辆是红色、哪辆被遮挡、离相机多远；在图库中找到另一张相似的街景，也未必找到了同一辆车的其他照片。本章要回答的问题是：怎样把一幅图像变成这些逐渐具体、而且可以核验的结果，即语义判断、空间范围、测量量和跨图对应。

本章用到两类已有材料。第一部分讲过的经典方法给出了不经学习也能工作的做法：\ref{chap:vision-image-formation}章的相机模型说明像素与三维点怎样对应，\ref{chap:vision-image-transforms}章的金字塔与滤波说明怎样在多个尺度上观察图像，\ref{chap:vision-image-features}章的边缘、局部特征、经典分割和基于人工特征的识别说明怎样从像素中提出可比较的证据。\ref{chap:vision-image-video-models}章和\ref{chap:vision-multimodal-models}章则给出了可复用的视觉表示和图文表示。本章不重讲这些基础，而是研究在具体任务中还缺什么：判断针对整幅图、一个对象还是一个像素；监督约束哪一种输出；网络输出怎样还原成原图坐标、对象编号或米制距离；结果靠什么核验。每类任务先回顾经典路线怎样处理同一问题、在什么条件下失效，再讲学习方法改动了哪一步。

读完本章，读者应能说明各类识别与检索任务的输入、输出与评价指标，指出一个任务的主要困难来自外观、尺度、分辨率、几何还是开放类别，并能沿着本章的模型结构表复核一个检测器、分割器或几何模型的张量形状。第一节先把这些任务和困难摆清楚，后续各节再逐一回答。

街景中的两辆车贯穿对象与空间任务，正常零件与划痕用于说明没有完整缺陷类别表时怎样判断异常，建筑照片用于说明检索与几何恢复。以下算例均为教学构造，数值用于复核机制，不是模型实验结果。

\section{图像识别与检索的问题与挑战}
\label{sec:vision-recognition-challenges}

图像识别与检索不是一个任务，而是一组输出形式不同的任务。它们的输入都是一幅或多幅图像，有时附带类别表、文字、点击或图库；输出则从一个类别名，到一组框、一张逐像素的类别图、若干关节坐标、一张深度图，再到一份按相关程度排序的图库列表。输出形式不同，评价指标也不同，而指标决定了模型被要求做对什么。下面先把各任务的输入、输出和指标列清楚，再讨论使它们变难的五类挑战。

表\ref{tab:vision-recognition-task-metrics}汇总了本章各任务的输入、输出与主要指标。表中每个指标都在下文给出定义；同一名称在不同基准中的细节可能不同，比较结果时必须说明所用协议。

\begin{table*}[tbp]
 \centering
 \small
 \begin{tabularx}{\linewidth}{@{}lXXX@{}}
 \toprule
 任务 & 输入 & 输出 & 主要指标\\
 \midrule
 类别与场景判定 & 整图或对象裁剪，可附类别名 & 一个或多个类别 & 准确率、前$k$准确率、类别平均准确率\\
 异常判定 & 零件图像，只有正常样本可供参照 & 整图异常分数与异常热图 & 图像级与像素级AUROC\\
 边界框预测 & 图像，可附类别表或文字 & 框、类别与得分的集合 & IoU阈值下的AP与mAP\\
 像素类别预测 & 图像，或两期配准影像 & 逐像素类别图 & 类别IoU与mIoU\\
 对象掩码预测 & 图像，可附类别、点框或短语 & 实例掩码集合 & 掩码AP、PQ\\
 前景透明度预测 & 图像与三分图 & 逐像素$\alpha$ & SAD、MSE\\
 计数 & 图像与计入规则 & 数量，可附位置 & MAE、RMSE\\
 姿态估计 & 含人图像 & 二维或三维关节，或人体网格 & OKS下的AP、MPJPE\\
 深度与几何估计 & 一幅或多幅图像，可附相机 & 深度图、点图与相机 & AbsRel、$\delta<1.25$、旋转与平移误差\\
 图像检索 & 查询图与图库 & 图库图像的排序列表 & mAP、Recall@$K$\\
 \bottomrule
 \end{tabularx}
 \caption{本章各任务的输入、输出与主要评价指标。指标的协议细节（阈值、对齐方式、忽略项）随基准而变，报告时应一并给出。}
 \label{tab:vision-recognition-task-metrics}
\end{table*}

类别判定最直接的指标是\term{准确率}（accuracy），即预测类别$\hat y_i$等于标注$y_i$的样本比例$(1/n)\sum_i\mathbf 1\{\hat y_i=y_i\}$。类别很多且彼此相近时，常同时报告前$k$准确率：只要真实类别出现在得分最高的$k$个类别中就算正确。样本在类别间不均衡时，样本平均准确率会被大类主导，此时还要报告各类准确率的平均值，即类别平均准确率。异常判定的分数需要选阈值才能变成“正常／异常”，不同阈值给出不同的误报与漏报。为了在选阈值之前比较排序能力，使用ROC曲线下面积（area under the ROC curve, AUROC），它等于随机抽一个异常样本和一个正常样本、前者分数更高的概率。图像级AUROC按整图分数计算，像素级AUROC把每个像素当作一个样本，检查热图能否把缺陷像素排在正常像素之前。

检测和分割要先回答“预测与哪个真实对象对应”。两个区域的重叠用\term{交并比}（intersection over union, IoU）衡量，它将交集面积除以并集面积：
\begin{equation}
 \operatorname{IoU}(\bm{b},\bm{b}')
 =\frac{|\bm{b}\cap\bm{b}'|}
 {|\bm{b}|+|\bm{b}'|-|\bm{b}\cap\bm{b}'|}\text{。}
 \label{eq:vision-iou}
\end{equation}
$\bm{b},\bm{b}'$可以是矩形框，也可以是掩码覆盖的像素集合。取两个面积均为800的框$(10,10,50,30)$和$(20,10,60,30)$，交集面积为600，IoU为$600/1000=0.6$。检测评价先固定一个IoU阈值$t$，把预测按得分从高到低逐个与同类、尚未被匹配的真实对象配对，IoU不低于$t$的配对算真阳性，其余预测算假阳性，未被配上的真实对象算漏检。逐步降低得分门限，就得到一串精确率$P$（真阳性占全部预测的比例）和召回率$R$（真阳性占全部真实对象的比例），连成精确率—召回率曲线。\term{平均精度}（average precision, AP）是这条曲线下的面积；实际计算时先把精确率改成单调不增的包络$P_{\mathrm{interp}}(r)=\max_{r'\geq r}P(r')$，再在若干召回点上取平均。对所有类别的AP再求平均，就是\term{平均精度均值}（mean average precision, mAP）。COCO检测基准在召回率0至1之间均匀取101个点做插值，并把IoU阈值从0.50到0.95每隔0.05取一次、共10个阈值，对类别和阈值同时求平均；它不区分AP与mAP的写法，直接把这个平均值称为AP。高阈值上的AP奖励定位精细的框，0.5阈值上的AP则主要检查有没有找到对象。

逐像素输出的评价改成按类别累计像素。对类别$r$，令$\mathrm{TP}_r,\mathrm{FP}_r,\mathrm{FN}_r$分别为预测和标注都属于$r$、只有预测属于$r$、只有标注属于$r$的像素数，则类别IoU为$\mathrm{TP}_r/(\mathrm{TP}_r+\mathrm{FP}_r+\mathrm{FN}_r)$，对$K$个类别取平均得到\term{平均交并比}（mean intersection over union, mIoU）。它以类别为单位平均，因此路杆这类小类别与道路这类大类别权重相同。实例掩码沿用检测的AP，只是把框的IoU换成掩码的IoU。全景分割既有可数对象又有道路天空等区域，\term{全景质量}（panoptic quality, PQ）把匹配区域的IoU与对象错误合在一起：对同类、IoU超过0.5的区域建立匹配集合$T$，令$F_{\mathrm{p}},F_{\mathrm{n}}$为多报和漏报区域集合，
\begin{equation}
 \mathrm{PQ}=
 \frac{\sum_{(p,g)\in T}\operatorname{IoU}(p,g)}
 {|T|+\frac12|F_{\mathrm{p}}|+\frac12|F_{\mathrm{n}}|}\text{。}
 \label{eq:vision-pq}
\end{equation}
IoU超过0.5保证一个真实区域至多与一个预测区域匹配，所以匹配无须另作优化\scite{vision-1801-00868}。例如匹配区域IoU为0.8和0.6，另有一个漏报且无多报，则PQ为$1.4/2.5=0.56$；只报告匹配区域的平均IoU会得到0.7，掩盖了漏掉的区域。透明度预测在三分图的未知区内比较连续值，常用绝对误差和$\mathrm{SAD}=\sum_p|\hat\alpha_p-\alpha_p|$与均方误差。

测量任务的指标带单位。计数比较每幅图的预测数$\hat c_i$与真值$c_i$，平均绝对误差为$\mathrm{MAE}=(1/n)\sum_i|\hat c_i-c_i|$，均方根误差为$\mathrm{RMSE}=\sqrt{(1/n)\sum_i(\hat c_i-c_i)^2}$，后者对少数大误差更敏感。二维姿态沿用检测的AP，只是把IoU换成\term{对象关键点相似度}（object keypoint similarity, OKS）：
\begin{equation}
 \mathrm{OKS}=
 \frac{\sum_k\exp\!\left(-d_k^2/(2s^2\kappa_k^2)\right)\mathbf 1\{v_k>0\}}
      {\sum_k\mathbf 1\{v_k>0\}}\text{，}
 \label{eq:vision-oks}
\end{equation}
其中$d_k$是第$k$个关节的预测与标注之间的距离，$s$是该人的尺度（COCO取标注区域面积的平方根），$\kappa_k$是按关节类型设定的容差常数，$v_k>0$表示该关节有标注。OKS把每个关节的像素距离换算成0到1之间的相似度，再在有标注的关节上平均；COCO关键点基准同样在0.50到0.95的10个OKS阈值上平均AP。三维姿态用\term{平均关节位置误差}（mean per-joint position error, MPJPE），即各关节预测位置与真实位置三维距离的平均，单位通常是毫米。深度的常用指标来自早期单目深度工作\scite{vision-1406-02283}：对有效像素集合$\Omega$，绝对相对误差$\mathrm{AbsRel}=(1/|\Omega|)\sum_p|\hat z_p-z_p|/z_p$，阈值准确率$\delta<1.25$是满足$\max(\hat z_p/z_p,z_p/\hat z_p)<1.25$的像素比例，另有线性与对数尺度的RMSE。相机估计比较旋转与平移：两个旋转矩阵$\hat{\bm R},\bm R$之间的角度误差为$\arccos\{[\operatorname{tr}(\hat{\bm R}^\top\bm R)-1]/2\}$，平移在尺度未知时只比较方向夹角。

检索按相关性检查排名。若某个查询有$R$个相关项，排名$k$处的相关指示为$r_k$，截至该位置的精确率为$P(k)$，非插值平均精度为
\begin{equation}
 \mathrm{AP}=\frac1R\sum_k P(k)r_k\text{，}
 \label{eq:vision-retrieval-ap}
\end{equation}
对所有查询取平均即检索的mAP。三个正例出现在第1、3、5位时，AP为$(1+2/3+3/5)/3\approx0.756$。地点识别等任务更关心前几名里有没有正确结果，此时用Recall@$K$：前$K$个结果中至少含一个相关项的查询所占比例。AP既关心是否找到，也关心出现得早不早；Recall@$K$只关心前$K$名，适合候选还要交给后续步骤复核的场合。

有了输入、输出和指标，就可以看清这些任务为什么难。下面五个小节各讨论一类挑战。它们都不只属于某一个任务，但在不同任务上以不同形式出现。

\subsection{外观变化带来的类别歧义}
\label{sec:vision-recognition-challenge-appearance}

同一类别的物体，在图像上可以很不相同；不同类别的物体，又可能很相像。路口的红色货车在阴影里变暗，正对相机时和侧对相机时轮廓完全不同；同一型号的轿车喷成蓝色和红色，像素值几乎没有共同之处。反过来，起重机和鹤在英文里同名，两栋玻璃幕墙建筑在检索中很难分开。识别要求对视角、光照、形变和类内差异保持不变，同时对区分类别或个体的细节保持敏感，这两个要求方向相反。

经典方法用人工设计的特征缓和这一矛盾。\ref{sec:vision-features-template-matching}节的模板匹配直接比较像素，归一化互相关能抵消整体亮度和对比度的变化，却对视角和形变毫无办法。\ref{sec:vision-features-sift}节的SIFT和\ref{sec:vision-features-hog}节的HOG改为比较梯度方向的统计：梯度方向对光照的加性和乘性变化不敏感，局部直方图又容忍小的位移。\ref{sec:vision-features-bovw}节的视觉词袋再把一幅图的局部特征量化成“视觉词”的频率，用支持向量机等分类器判别，也用于在图库中检索同一地点。这些特征的不变性是设计者事先选定的，选定之外的变化（大视角、遮挡、细粒度差异）就无法处理；分类边界也只能在人工特征张成的空间里画。

学习方法把“对什么不变、对什么敏感”交给数据决定。卷积网络和视觉Transformer从大量图像中学习表示，自监督或图文对比训练又让同一表示可以被多个任务复用。代价是表示可能学到数据中的捷径，例如用背景判断类别，或在训练少见的领域失效。\ref{sec:vision-recognition-semantics}节在固定类别表与文字类别表两种条件下讨论类别判定，并把属性限制到对象区域内；\ref{sec:vision-recognition-retrieval}节讨论检索中的同一问题：区分同一实例与相似实例。

\subsection{尺度与遮挡带来的定位困难}
\label{sec:vision-recognition-challenge-localization}

定位要回答“对象在哪里”，而对象在图中的大小和可见程度变化极大。靠近相机的货车占画面三分之一，路口远处的轿车可能只有$20\times10$个像素；公交车站旁的车被行人挡住一半，相邻两车的框大面积重叠。一幅图里的对象数量也不固定，可以是零个，也可以是几十个。定位器必须在所有位置、所有尺度上搜索，又要给出一个不重复、不遗漏的集合。

经典检测器用穷举加剪枝的方式处理。\ref{sec:vision-features-viola-jones}节的Viola–Jones检测器在\ref{sec:vision-transforms-pyramid}节的图像金字塔各层上滑动一个固定大小的窗口，用积分图快速计算类Haar特征，再用级联分类器尽早拒绝明显不是人脸的窗口；尺度变化由金字塔覆盖，海量的背景窗口由级联的前几级过滤。\ref{sec:vision-features-dpm}节的可变形部件模型在HOG金字塔上用一个根滤波器描述整体、若干部件滤波器描述局部，并允许部件相对根位置有代价地移动，从而容忍形变和部分遮挡。两者都在得分图上用非极大值抑制去掉重复窗口。这条路线的问题是：窗口和部件的形状要人工设定，每个尺度要重算特征，被大面积遮挡或比例罕见的对象容易漏掉。

学习式检测器保留了“在所有位置和尺度上打分”的思想，但改变了实现方式。卷积网络在整幅图上一次前向就得到所有位置的预测，相当于把滑动窗口写成卷积；图像金字塔换成网络内部的特征金字塔，各层共享一次计算；级联对容易背景的剪枝，换成了训练时降低容易样本权重的损失；去重从后处理规则逐步变成训练目标。\ref{sec:vision-recognition-detection}节沿着这条线讲六个检测模型，\ref{sec:vision-recognition-mask-sets}节把同样的问题推到实例掩码。

\subsection{稠密输出与分辨率的矛盾}
\label{sec:vision-recognition-challenge-dense}

逐像素的输出要同时满足两个相互冲突的要求。判断道路中央的一块灰色是路面还是车顶，需要看到周围很大范围的上下文；画出路杆、车窗和发丝的细边界，又需要接近原图的分辨率。卷积网络靠逐层下采样扩大感受范围，代价正是丢掉位置细节；如果不下采样，计算和存储又会随分辨率迅速增长。对象掩码还要在像素层面分开相邻的同类对象，透明度预测则要给出连续的混合比例。

经典分割从低层线索把像素组织成区域。\ref{sec:vision-features-otsu}节的阈值法按灰度把像素分为两类，\ref{sec:vision-features-watershed}节的分水岭沿梯度的“山脊”切开区域，\ref{sec:vision-features-graph-segmentation}节的图割把分割写成能量最小化，用像素与标记的颜色相似度和相邻像素的平滑项共同决定边界，\ref{sec:vision-features-superpixel}节的超像素把颜色和位置相近的像素聚成小块。这些方法得到的是边界贴合的区域，却不知道区域属于什么类别；颜色相同的车顶和路面会被并到一起，纹理复杂的同一物体又会被切碎。

学习方法让每个像素的类别由网络在上下文中判断，再设法补回分辨率：空洞卷积在不降低分辨率的前提下扩大感受范围，编码器—解码器用浅层特征恢复边界，查询式分割让每个对象拥有一张整图掩码。经典线索并没有消失，图割、超像素和颜色传播的思想仍出现在交互分割与抠图中。\ref{sec:vision-recognition-segmentation}节到\ref{sec:vision-recognition-matting}节依次讨论像素类别、对象掩码和前景透明度三种稠密输出。

\subsection{单目观测带来的几何不确定}
\label{sec:vision-recognition-challenge-geometry}

一幅照片把三维场景投影到二维平面，深度信息在投影中丢失。\ref{sec:vision-formation-camera-intrinsics}节的针孔模型说明，一个像素只确定一条从相机中心出发的射线：同一像素可以对应4米外的点，也可以对应8米外的点。手腕在图上的位置相同，可能是向前伸也可能是向后伸；把整个场景连同相机平移放大一倍，所有照片保持不变。因此从图像得到深度、三维关节或相机位置，必须引入额外的观测或先验。

经典方法用多视图增加观测。\ref{sec:vision-formation-multi-view}节的双目立体在两台已知基线的相机之间沿极线搜索对应点，由视差反比得到深度；运动恢复结构则对一组无序照片先匹配特征点、估计两视图几何，再三角化三维点，最后用光束法平差联合优化所有相机和点的重投影误差。这些方法有明确的几何依据，在纹理丰富、重叠充分时很准，但依赖可靠的特征匹配：弱纹理、反光、重复图案和视图太少都会使流程中断，而且多步骤串联，前一步的错误会传到后面。单幅图像则无从三角化，经典方法只能借助消失点、已知物体尺寸等少量线索。

学习方法从两个方向改变了这一局面。一是用大规模数据学习单目先验，直接从一幅图预测相对深度或人体三维姿态；二是把多视图几何的计算步骤写成可学习的网络，例如在学习的特征上构造代价体，或让网络一次前馈输出所有相机和点图。前者给出完整但只确定到尺度的结果，后者需要多幅图但能利用跨视图约束。\ref{sec:vision-recognition-measurement}节按测量量依次讨论计数、姿态、深度和相机几何，并在每一处说明尺度从哪里来。

\subsection{开放类别与长尾带来的泛化困难}
\label{sec:vision-recognition-challenge-open}

训练时的类别表总是有限的，使用时却会遇到新东西。仓库画面里出现一个没有名字的搬运装置；零件上出现一种从未见过的划痕；用户要找“公交车旁的红色货车”，这个组合不在任何类别表中。即使在类别表之内，样本数量也极不均衡：轿车有成千上万张，少见车型只有几张。闭集模型对新类别只能二选一地误报成已知类或当作背景，对长尾类别则倾向于预测常见类别。

经典方法基本是闭集的。每个类别单独训练一个检测器或分类器，新类别要重新收集样本并训练；视觉词袋的词汇表在建库时固定，检索时只能比较这些词的频率，再用\ref{sec:vision-recognition-ransac}节的RANSAC几何验证排除外观相似但几何不一致的候选。几何验证不需要类别，这是经典路线中少有的“开放”环节。

学习方法提供了三种扩展方式。第一种用文字指定类别，例如CLIP、Grounding DINO与SAM 3把类别表写成文字，使用时再给出；第二种不命名类别，只判断“像不像一个对象”或“是否偏离正常”，例如OWOBJ的对象性与PatchCore的正常记忆；第三种在没有人工标注时自动构造监督，例如GeM用三维重建挖掘同一建筑的照片对。三者都扩大了可处理的范围，却也带来新的风险：文字条件可能被部分忽略，对象性会把背景提成候选，正常记忆分不清无害的外观变化。这些风险分散在\ref{sec:vision-recognition-semantics}节、\ref{sec:vision-recognition-regions}节和\ref{sec:vision-recognition-retrieval}节中讨论，\ref{sec:vision-recognition-evaluation}节再给出相应的检验协议。

五类挑战按输出逐步变具体的顺序在后续各节中得到回答。\ref{sec:vision-recognition-semantics}节先处理外观歧义下的“是什么”，并在异常判定中第一次面对开放类别；\ref{sec:vision-recognition-regions}节把判断落到框、像素和掩码上，同时回答尺度遮挡与稠密分辨率两类困难，开放词汇检测与未知对象发现也在这里出现；\ref{sec:vision-recognition-measurement}节在定位之上加入数量、关节和三维坐标，正面处理单目观测的几何不确定；\ref{sec:vision-recognition-retrieval}节把一幅图与图库中的其他图对应起来，需要的正是前面形成的全局表示、局部点和几何约束。每一处都先回顾经典路线的做法和失效条件，再讲学习模型改变了哪一步。最后一节把这些输出放在同一组评价条件下检验。

\section{语义与属性判定}
\label{sec:vision-recognition-semantics}

本节处理三种语义输出：类别说明“属于什么”，属性说明“具有什么特征”，异常判断说明“是否偏离规定的正常状态”。输入都是一幅图或一个裁剪区域，输出都是有限个标签或分数，难点却各不相同。类别判定的难点在于类别表从哪里来：训练时固定，还是使用时由文字临时给出。属性判定的难点在于归属：整图里同时有红色和蓝色，必须说清哪个颜色属于哪辆车。异常判定的难点在于缺少反例：缺陷类型无法穷举，训练数据主要是正常样本。

三个难点对应三种不同的设计选择。类别表固定时，可以在冻结的视觉表示上训练一个分类头，DINOv3是这条路线的代表；类别表在使用时才给出时，可以把类别名也编码成向量，用CLIP的图文相似度排序。属性判定沿用这两种接口，但先把表示限制在对象区域内。异常判定改变了监督形式：不学习“缺陷长什么样”，而是保存正常外观，再度量新观测离正常外观有多远。

\subsection{类别与场景的判定}
\label{sec:vision-recognition-classification}

类别判定的输入是整幅图像或对象裁剪，输出是类别表中的一项或多项。在\term{单标签分类}（single-label classification）中，一个样本只对应类别表中的一项；\term{多标签分类}（multi-label classification）允许多项同时成立。例如车辆裁剪图可以按主要车型给一个标签，整张街景却可以同时含公交车、轿车和行人。损失函数必须遵守这一区别：互斥类别用softmax在类别间归一化，可共存的对象则分别输出“是否存在”的得分。场景分类把整幅画面概括为街道或厨房，它与对象分类不是同一个问题，厨房这一场景名不能展开为一份完整的物品清单。

两条路线的差别在于类别表进入计算的方式。第一条路线把类别表写进分类头的输出维度：每个类别对应一行权重，用带标签的样本学习这些权重，视觉编码器可以冻结，也可以微调。第二条路线把类别表写成文字：每个类别对应一条描述，由文字编码器把描述变成向量，再与图像向量比较，不为当前类别表训练新参数。前者需要标注数据，但可以针对具体领域学到细致的判别边界；后者不需要标注，却受限于预训练图文数据覆盖了哪些概念、描述写得是否准确。下面分别用DINOv3和CLIP说明这两种用法。

\subsubsection{DINOv3}
\label{sec:vision-recognition-dinov3}

DINOv3的结构与自监督训练已在\ref{sec:vision-ivm-dinov3}节展开。这里只需要它的几个结论：DINOv3是在大规模无标注图像上自监督训练的视觉Transformer系列，输入图像被切成固定边长的图像块，网络输出一个汇总整图的类别词元和每个图像块各自的特征；训练中加入的格拉姆锚定（Gram anchoring）约束使长时间训练后的逐块特征仍保持局部区分能力。原论文强调，这些表示在多种任务上不经微调即可使用\scite{vision-2508-10104}。本节关心的是拿到这样的编码器后，怎样完成一个有标签的车辆或场景分类任务。

输入的组织先决定特征的形状。设所选检查点的图像块边长为16，输入缩放到$224\times224$，则网络得到$14\times14=196$个块特征和一个类别词元，每个向量为$d$维（$d$由检查点规模决定）。全局任务通常取类别词元输出作为$\bm{z}\in\mathbb{R}^{d}$，也可以把它与块特征的平均拼接；空间任务则保留$14\times14\times d$的块特征网格。缩放、裁剪和归一化必须按检查点的约定进行，并保证标签仍对应被保留的内容：若裁剪切掉了车头，“轿车”标签就可能失去依据。训练、调参和测试图像应按车辆、拍摄批次或采集条件分开，避免同一对象的近重复图同时出现在训练和测试中。

拿到特征后有三种使用方式，差别在于更新哪些参数。\term{线性探测}（linear probing）冻结编码器，只训练一个线性分类头；它把编码器当作固定的特征提取器，所以训练快、显存小，同一份特征还能供多个任务共用。$k$近邻分类连分类头也不训练，直接在带标签的训练特征中找最相近的样本投票。\term{微调}（fine-tuning）则同时更新编码器参数，通常使用比分类头小得多的学习率，或只解冻最后几层。设冻结编码器$f_{\bm{\theta}}$得到$\bm{z}=f_{\bm{\theta}}(\bm{X})$，分类头包含$\bm{W}\in\mathbb{R}^{k\times d}$与$\bm{b}\in\mathbb{R}^{k}$。对$k$个互斥类别，计算
\begin{equation}
 p(y=r\mid\bm{X})=
 \frac{\exp(\bm{w}_r^\top\bm{z}+b_r)}
 {\sum_{j=1}^{k}\exp(\bm{w}_j^\top\bm{z}+b_j)}\text{，}
 \label{eq:vision-linear-head}
\end{equation}
其中$\bm{w}_r^\top$是$\bm{W}$的第$r$行。训练目标为平均负对数似然$L=-n^{-1}\sum_i\log p(y_i\mid\bm X_i)$。线性探测只更新$\bm{W},\bm{b}$；微调还更新$\bm{\theta}$，所需梯度存储随之增加，训练数据较少时还可能破坏预训练表示中对当前任务有用的结构。

对一个样本，记类别得分为$a_r=\bm{w}_r^\top\bm{z}+b_r$、概率为$p_r$，则$\partial(-\log p_y)/\partial a_r=p_r-\mathbf{1}\{r=y\}$，头部梯度为$(p_r-\mathbf{1}\{r=y\})\bm{z}$。假设三类车型的得分为$(2,1,0)$，softmax概率依次约为0.665、0.245和0.090；若真实类别为第二类，得分梯度约为$(0.665,-0.755,0.090)$，下降更新提高第二类相对于其他类的得分。这些概率是在给定类别集合内归一化的结果，不能证明图像中一定存在可辨车辆。

选择哪种方式，取决于数据量、领域差距和资源。训练样本少、目标领域与预训练图像接近时，线性探测通常已足够，而且不会改动编码器；领域差距大（例如显微图像或遥感图像），或任务依赖预训练中少见的细节时，微调可能带来改进，但需要更多标注和计算，并要用验证集确认改进确实存在。调参时需要固定的超参数包括输入分辨率、所用特征层、头部学习率和训练轮数；微调还要加上编码器学习率和解冻层数。空间任务可以在每个块特征上接同一个线性头，得到$14\times14$的类别图，这正是\ref{sec:vision-recognition-dense-pixels}节像素类别预测的最简形式。

常见失败有三种。第一，若编码器读取整幅街景，判定左车车型的得分可能依赖道路背景；任务要针对某一辆车时，应先取得该车区域。第二，外观接近的细粒度类别依赖小部位：全局背景足以区分拍摄地点，却可能不足以区分鸟喙或叶缘。iNaturalist包含细粒度类别和不均衡样本，评价时应同时检查小部位分辨率和类别覆盖\scite{vision-1707-06642}。例如常见类别有90张且全对，少见类别有10张且全错，样本平均准确率为90\%，两个类别的平均准确率却为50\%。第三，使用时必须保留训练时的编码器版本、预处理、分类头和类别顺序；需要“未知”输出时，再加入由验证集确定的拒绝阈值。

\subsubsection{CLIP}
\label{sec:vision-recognition-clip}

如果使用时才给出“轿车、货车、公交车”，固定输出维度的分类头不能直接容纳新类别。CLIP的双编码器和对比训练已在\ref{sec:vision-models-contrastive}节展开，其结论是：图像编码器和文字编码器把匹配的图文对映射到单位球面上相近的位置，两者的内积可以比较图像与任意一条描述的对应程度\scite{radford2021}。\term{零样本分类}（zero-shot classification）利用这一点，把类别表写成文字，不为当前类别表训练分类头。

输入的组织分两侧。文字侧，把每个类别名填入模板，例如“一张{类别}的照片”，编码后归一化；同一类别可以使用多个模板，对其单位文字向量求平均后再归一化，CLIP原文在ImageNet上即对多条模板做了这种集成。这样得到的$k$个文字向量排成矩阵，起到线性分类头权重的作用，只是权重来自文字而非标注。图像侧，按检查点约定预处理后得到单位图像向量。令$\bar{\bm{z}}$与$\bar{\bm{t}}_r$分别为单位图像向量和第$r$类的文字向量，$s_r=\bar{\bm{z}}^\top\bar{\bm{t}}_r$，使用温度$\tau>0$计算候选内归一化得分
\begin{equation}
 q_r=\frac{\exp(s_r/\tau)}
 {\sum_{j\in C}\exp(s_j/\tau)}\text{，}
 \label{eq:vision-clip-classification}
\end{equation}
其中$C$是这次输入的非空候选索引集合。CLIP训练时把$1/\tau$作为可学习参数，并把它的上限设为100。分母随候选集合变化，所以$q_r$不是不依赖候选的存在概率。模板集合须在评价前固定，换模板就是换分类器。

CLIP默认不更新任何参数。若有少量标注，可以在冻结的CLIP图像向量上做线性探测，这时CLIP退化为第一条路线中的一个编码器；也可以用文字向量初始化分类头再微调。超参数主要是模板集合、类别名的写法和拒绝阈值。类别名有歧义时（例如“crane”既是鹤也是起重机），模板需要加入上下文。

设四项相似度为$0.31,0.28,0.19,0.12$，最后一项是不在图中的“帆船”。最大值只说明第一项在这四项中领先；把前三项全部换成不相关的词，模型仍会返回一个最大值。开放使用要另设含未知对象的验证协议、拒绝阈值或目标存在判断。整图相似度还无法保证“红色”绑定到货车而非旁边的轿车，也不擅长回答“图中有几辆车”这类计数问题。

两种用法可以在同一车型分类任务、同一测试集上比较，差别集中在类别来源、训练需求和失败方式上，见表\ref{tab:vision-recognition-classification-compare}。CLIP把类别表的指定从训练阶段推迟到使用阶段，换来了不需标注的灵活性，代价是判别边界由图文预训练决定，难以针对领域细节调整；DINOv3加分类头则反过来。二者并不互斥：CLIP图像编码器同样可以做线性探测，DINOv3的块特征也可以接入文字对齐的模块。

\begin{table}[htbp]
 \centering
 \small
 \begin{tabularx}{\linewidth}{@{}lXX@{}}
 \toprule
 比较项 & DINOv3加分类头 & CLIP零样本分类\\
 \midrule
 类别来源 & 训练时固定，写入头部输出维度 & 使用时以文字给出\\
 需要的数据 & 每类带标签的样本 & 无需当前任务标注\\
 更新的参数 & 分类头，或连同编码器微调 & 默认不更新\\
 空间输出 & 块特征可接逐位置分类头 & 主要为整图向量\\
 典型失败 & 背景捷径、细粒度部位不足 & 候选闭合、属性错绑、模板敏感\\
 \bottomrule
 \end{tabularx}
 \caption{同一车型分类任务下两种类别接入方式的比较。表中列出的是接口差别，不是两个模型在某一数据集上的精度结论。}
 \label{tab:vision-recognition-classification-compare}
\end{table}

\subsection{对象属性的判定}
\label{sec:vision-recognition-attributes}

要记录“红色货车”和“蓝色轿车”，只输出整图属性集合$\{\text{红色},\text{蓝色}\}$会丢掉归属。属性判定的输入因而不是整幅图，而是一个对象区域；输出是$(\text{对象编号},\text{属性},\text{取值})$记录。难点在于把属性证据限制在目标对象内：裁剪过小会切掉车灯，过大又会纳入旁车颜色。

路线有两条。一条是整图多标签预测，实现简单，但无法表达归属，只适合“图中是否出现红色物体”这类问题。另一条先定位对象，再对各区域提取表示并接属性头；区域表示可以来自DINOv3块特征在框内的汇总，也可以来自\ref{sec:vision-recognition-mask-rcnn}节的区域对齐特征。CLIP也可以用“一辆红色的车”作为文字查询，但应当对对象裁剪而不是整图计算相似度，否则同样会出现错绑。位置与属性共享对象编号，后续计数、搜索或人工核验才能指向同一个对象。

属性之间是否互斥由标签定义决定。“红色车身”和“车灯开启”可以同时成立，“完整车身的主色”则可以规定为单标签字段。对于可共存的$k$个二值属性，令区域向量$\bm{z}_i$经过每属性一个线性得分$h_{i,r}$，$p_{i,r}=1/(1+e^{-h_{i,r}})$为预测，$a_{i,r}\in\{0,1\}$为标注，$m_{i,r}$表示该属性是否有可靠标注，则
\begin{equation}
 L_{\mathrm{attr}}=
 -\frac{\sum_{i,r}m_{i,r}
 [a_{i,r}\log p_{i,r}+(1-a_{i,r})\log(1-p_{i,r})]}
 {\sum_{i,r}m_{i,r}}\text{。}
 \label{eq:vision-attribute-loss}
\end{equation}
分母要求至少有一个有效标注，全无标注的批次不计算此项。掩码$m_{i,r}$的作用是区分“未标注”和“不存在”：把未标注直接置为不存在，会训练出系统性的假负例。单个有效标注对得分的梯度是$p_{i,r}-a_{i,r}$。冻结区域编码器时只训练属性头；各属性阈值由独立验证记录选择，输出保留每车的属性分数和可见性。若主色被定义为互斥字段，则另用softmax头，不把多个二元输出强行当成唯一主色。

属性所需的部位不可见时，应输出“不可判定”而不是猜一个值。对象编号并不能修正错误的视觉证据，但它使属性错绑能够被单独定位和统计。

\subsection{正常与异常状态的判定}
\label{sec:vision-recognition-anomaly}

工业零件的输入是一幅零件图像，输出是整图的异常分数和标出可疑位置的异常热图。难点在于监督不对称：新的缺陷类型不断出现，训练集主要是正常样本，无法像车型分类那样为每种缺陷准备标签。\term{异常检测}（anomaly detection）因此改为学习或保存正常外观，再度量新观测与正常外观的偏离。

不同路线对“偏离”给出不同的度量。重建路线训练网络重建正常图像，用重建误差作为异常分数；它的困难是网络泛化过好时也能重建缺陷，泛化不足时又会把正常的细纹理重建模糊。分布路线为正常特征拟合一个密度，用低密度表示异常；单个高斯等简单分布可能无法覆盖正常外观中的多种模式。记忆路线保存正常样本的局部特征，用最近邻距离表示异常，不需要假设分布形式，代价是记忆的规模和检索时间。PatchCore属于第三条路线\scite{vision-2106-08265}。

PatchCore针对的具体问题是：整图特征比较会把一处细小划痕淹没在整幅零件的正常外观中，而直接用最深层特征又过于偏向ImageNet类别语义。它的主要想法是保存正常图像中\emph{每个局部位置}的中层特征，查询时逐位置找最近的正常片段；局部比较让细小划痕得到单独的分数，中层特征兼顾纹理细节和一定的上下文。原工作默认使用ImageNet预训练的WideResNet-50，取第2、3个残差阶段的特征。对每个位置，先在$3\times3$邻域内做自适应平均池化，把邻近上下文汇入该位置；再把第3阶段的特征双线性放大到第2阶段的空间尺寸并与之合并，得到每个位置一个固定维的片段向量。原文把图像缩放到$256\times256$后中心裁剪为$224\times224$，不做数据增强。WideResNet-50第2阶段的步幅为8，因此第2阶段的特征网格是$28\times28$，每幅正常图贡献784个片段向量；第3阶段步幅为16，网格$14\times14$，放大两倍后与第2阶段逐位置对齐。合并后的片段维数由自适应平均池化统一到一个预先设定的值，原文未给出这个数值。全部正常图的片段汇成特征集合$\bm{M}$。

PatchCore没有梯度训练，编码器始终冻结，“学习”就是建立记忆。全部片段可能非常多，因而从$\bm{M}$中保留代表子集$\bm{M}_{\mathrm{c}}$，即\term{核心集}（coreset），其覆盖目标为
\begin{equation}
 \min_{\substack{\bm{M}_{\mathrm{c}}\subseteq\bm{M}\\
 |\bm{M}_{\mathrm{c}}|=m}}
 \max_{\bm{z}\in\bm{M}}\min_{\bm{u}\in\bm{M}_{\mathrm{c}}}
 \lVert\bm{z}-\bm{u}\rVert_2\text{。}
 \label{eq:vision-patchcore-coreset}
\end{equation}
最内层最小值为每个正常片段寻找一个代表，外层最大值检查最难覆盖的正常片段。若只最小化平均距离，少量但正常的外观模式可能在记忆中完全缺席，查询时就会被误报为异常。精确求解需要穷举子集，原方法改用近似贪心选择：从一个起点开始，反复加入离当前子集最远的特征，并为每个正常特征维护到当前子集的最近距离；新加入$\bm{u}$后更新为原距离与$\lVert\bm{z}-\bm{u}\rVert_2$的较小者，达到预算$m$就停止。原工作用随机投影降低选择时的距离计算成本，并报告了保留25\%、10\%和1\%的设置。建立记忆的数据必须符合“正常”定义，异常混入记忆会把某些缺陷变成近邻参照。

查询图像位置$p$的局部分数为$d_p=\min_{\bm{u}\in\bm{M}_{\mathrm{c}}}\lVert\bm{z}_p-\bm{u}\rVert_2$。把这些分数按空间排列、双线性放大到输入分辨率并用$\sigma=4$的高斯核平滑，就得到异常热图。整图分数从最大局部距离$s_*=\max_p d_p$开始，再根据最异常片段$\bm{z}_*$的最近正常片段$\bm{u}_*$周围是否稠密重新加权。令$\bm{N}_b(\bm{u}_*)$包含$\bm{u}_*$在记忆中的$b$个近邻，则
\begin{equation}
 s=\left(1-
 \frac{\exp\lVert\bm{z}_*-\bm{u}_*\rVert_2}
 {\sum_{\bm{u}\in\bm{N}_b(\bm{u}_*)}
  \exp\lVert\bm{z}_*-\bm{u}\rVert_2}\right)s_*\text{。}
 \label{eq:vision-patchcore-score}
\end{equation}
邻域大小$b$须大于1，否则分子分母相同，分数退化为零。若最近正常片段较孤立，其他近邻的距离更大，括号中的系数更接近1，最大距离受到的压低更少；也就是说，靠近一个稀少正常模式的查询会被判得更可疑。

假设一维正常记忆为$\{0,1,4\}$，查询局部特征为$\{0.2,0.9,6\}$，对应距离是$0.2,0.1,2$，最大值由第三个片段给出。若$b=2$，取最近正常片段4及其近邻1，整图分数为$2[1-\exp(2)/(\exp(2)+\exp(5))]\approx1.905$。阈值应在独立验证条件下确定。若新设备使所有特征整体平移，正常零件也可能远离记忆；高分只能表示对当前正常参照的偏离，不能直接命名缺陷，也不自动等于业务不合格。

\section{对象定位与区域划分}
\label{sec:vision-recognition-regions}

上一节的判断都没有说明位置。定位要把判断接回图上，输入是图像以及“要找什么”的指定（固定类别表、一段文字或一个点框提示），输出按精细程度分为四种形式：边界框给出对象的大致矩形范围，像素类别图给出每个像素属于哪一类，对象掩码进一步把同类像素分给不同实例，前景透明度则为发丝等混合边界给出连续的贡献比例。四种输出回答同一个“在哪里”的问题，但保存的信息逐级增加，监督和计算代价也逐级增加。本节按输出形式划分小节，每个小节内再比较不同模型对同一输出的不同解法。

各小节共有两个难点。其一是目标数量不定：一幅图里可能有零辆、一辆或二十辆车，网络却只能输出固定形状的张量，因此需要规定哪些输出位置负责哪些目标，以及怎样去掉重复。其二是目标由谁指定：类别表固定时可以把类别写进输出维度，文字或点击指定目标时，需要让查询条件进入定位计算。下面各模型的差别，大多就是对这两个问题的不同回答。

\subsection{对象边界框的预测}
\label{sec:vision-recognition-detection}

\term{目标检测}（object detection）的输入是一幅图像，输出是若干$(\bm{b},c,s)$组成的集合，其中$\bm{b}$是框，$c$是类别，$s$是得分。\term{边界框}（bounding box）用一个矩形包围目标，以下采用连续边界$\bm{b}=(x_1,y_1,x_2,y_2)$，宽高分别为$x_2-x_1$和$y_2-y_1$。标签应说明框包围可见部分还是完整对象。两个框的重叠程度用式\eqref{eq:vision-iou}的IoU衡量。节首例子中IoU为0.6的两个框，若是对同一辆车的两个预测，应当去重；若属于两辆互相遮挡的车，直接去掉一个则会漏检。

在学习式检测器之前，经典路线已经确立了检测的基本框架：在所有位置和尺度上打分，再去掉重复。\ref{sec:vision-features-viola-jones}节的Viola–Jones检测器以$24\times24$左右的固定窗口在图像金字塔的每一层上滑动，每个窗口用积分图在常数时间内算出类Haar特征，再交给由AdaBoost训练的级联分类器；前几级只用少数特征就拒绝绝大多数背景窗口，只有少量窗口走完全部级联。它在正面人脸这类外观稳定的目标上实现了实时检测，但类Haar特征描述不了外观变化大的类别。\ref{sec:vision-features-hog}节的HOG加线性支持向量机把特征换成梯度方向直方图，对行人这类轮廓明确的目标更稳健；\ref{sec:vision-features-dpm}节的可变形部件模型再在HOG金字塔上组合一个粗分辨率的根滤波器和若干细分辨率的部件滤波器，部件位置可以偏离锚点、偏离量受二次代价约束，部件位置作为隐变量用隐支持向量机训练。三者的共同结构是“特征金字塔$\to$逐窗口打分$\to$非极大值抑制”，差别在特征和打分函数。它们的局限也在这三处：每个尺度都要重算人工特征，窗口和部件的形状要人为设定，每个类别单独训练一个检测器。

下面的学习式检测器保留了这套结构，但逐处替换。YOLO与RetinaNet用卷积网络在整幅图上一次前向得到所有位置的得分，滑动窗口变成卷积的逐位置输出，人工特征变成学习的特征，多个类别共用一个网络；RetinaNet还把图像金字塔换成网络内部的特征金字塔，各层共享主干计算，并用焦点损失降低大量容易背景的训练权重，在作用上对应级联对容易窗口的提前拒绝。可变形部件的“允许局部偏移”在卷积网络中没有显式对应，而是由多层非线性特征隐式吸收。去重则直到DETR才从非极大值抑制移进训练目标。

检测的难点集中在“哪个输出负责哪个目标”。本小节的六个模型可以看成对这个问题的一串回答，每一步只改动一个设计。YOLO把图像划成网格，让目标中心所在的格子负责它，用一次前向得到所有框，但每个格子容量有限，重复框要靠后处理去掉。RetinaNet在多尺度特征的每个位置放置一组预设的参考框（锚框），让与目标重叠足够大的锚框负责它；这样小目标和拥挤目标有了更多候选，代价是绝大多数锚框是背景，需要焦点损失平衡训练。FCOS去掉预设锚框，让落在目标内的特征位置直接预测到四条边的距离，减少了锚框尺寸和比例这类人工设定，但仍需规则分配和非极大值抑制。DETR进一步取消密集候选，用固定数量的查询向量和二分匹配直接输出一个无重复的集合，把去重写进了训练目标。Grounding DINO沿用DETR式的查询解码，但把类别表换成任意文字，使“要找什么”在使用时给出。OWOBJ则面对连名字都没有的对象，在检测器上增加对象性信号，使未知对象也能作为候选保留下来。

\subsubsection{YOLO}
\label{sec:vision-recognition-yolo}

YOLO之前的主流检测器先生成上千个候选区域，再对每个区域分别做分类，同一幅图要经过多次网络计算，速度难以满足实时视频。原始YOLO的想法是把检测写成一次回归：整幅图经过一个网络，直接输出所有格子的框和类别得分\scite{vision-1506-02640}。由于每个格子都在整图特征上作出预测，网络在判断时能看到全局上下文。

网络是24个卷积层加2个全连接层，逐层结构见表\ref{tab:vision-recognition-yolo-layers}。它借鉴GoogLeNet的思路，但不用Inception模块，而是在$3\times3$卷积前反复插入$1\times1$卷积：$1\times1$卷积先把通道压缩一半，$3\times3$卷积再在较少的通道上扩大感受范围，这样在加深网络的同时控制计算量。以第4组为例，单个子模块的计算顺序是“$1\times1$卷积降到256通道$\to$带泄漏的ReLU$\to3\times3$卷积回到512通道$\to$带泄漏的ReLU”，重复4次。网络用4个$2\times2$最大池化和2个步长为2的卷积把$448\times448$的输入缩小64倍，得到$7\times7\times1024$的特征；第一个全连接层把它展平后映射到4096维，其后接比例为0.5的随机失活，第二个全连接层输出$7\times7\times30$。除最后一层是线性输出外，各层都用斜率0.1的带泄漏的ReLU，即$x>0$时取$x$、否则取$0.1x$。原文还给出一个只有9个卷积层、滤波器更少的Fast YOLO，具体通道数论文未给出。

\begin{table*}[tbp]
 \centering
 \small
 \begin{tabular}{@{}lllrrl@{}}
 \toprule
 阶段或层组 & 层类型 & 卷积核与步长 & 输出通道 & 重复次数 & 输出尺寸\\
 \midrule
 输入 & — & — & 3 & — & $448\times448$\\
 第1组 & 卷积 & $7\times7$，步长2 & 64 & 1 & $224\times224$\\
       & 最大池化 & $2\times2$，步长2 & 64 & 1 & $112\times112$\\
 第2组 & 卷积 & $3\times3$ & 192 & 1 & $112\times112$\\
       & 最大池化 & $2\times2$，步长2 & 192 & 1 & $56\times56$\\
 第3组 & 卷积 & $1\times1$；$3\times3$ & 128；256 & 1 & $56\times56$\\
       & 卷积 & $1\times1$；$3\times3$ & 256；512 & 1 & $56\times56$\\
       & 最大池化 & $2\times2$，步长2 & 512 & 1 & $28\times28$\\
 第4组 & 卷积 & $1\times1$；$3\times3$ & 256；512 & 4 & $28\times28$\\
       & 卷积 & $1\times1$；$3\times3$ & 512；1024 & 1 & $28\times28$\\
       & 最大池化 & $2\times2$，步长2 & 1024 & 1 & $14\times14$\\
 第5组 & 卷积 & $1\times1$；$3\times3$ & 512；1024 & 2 & $14\times14$\\
       & 卷积 & $3\times3$ & 1024 & 1 & $14\times14$\\
       & 卷积 & $3\times3$，步长2 & 1024 & 1 & $7\times7$\\
 第6组 & 卷积 & $3\times3$ & 1024 & 2 & $7\times7$\\
 全连接 & 全连接 & — & 4096 & 1 & $1\times1$\\
 输出 & 全连接，重排 & — & 30 & 1 & $7\times7$\\
 \bottomrule
 \end{tabular}
 \caption{YOLO（第一版）的逐层结构。层配置取自原文图3，卷积共24层（第1至6组依次为1、1、4、10、6、2层）；输出尺寸按$448\times448$输入逐层推算，原文只在图中画出各组的张量形状。未注明步长的卷积步长为1并保持尺寸。“$1\times1$；$3\times3$”表示依次相连的两层，重复次数指这一对的重复次数。ImageNet预训练使用前20个卷积层，检测时再加上后4个卷积层和2个全连接层。}
 \label{tab:vision-recognition-yolo-layers}
\end{table*}

全连接层输出一个$S\times S\times(5B+C)$的张量；原文取$S=7$、每格$B=2$个框、$C=20$个类别，输出为$7\times7\times30$。每个格子的30个数中，前10个是两个框各自的$(u,v,w,h)$和置信度，后20个是该格子共享的条件类别概率$\mathbb{P}(c\mid\mathrm{obj})$。框中心$(u,v)$相对于格子定位，宽高相对于整图归一化。若格子列行为$(j,i)$，输入宽高为$(W_{\mathrm{in}},H_{\mathrm{in}})$，则中心还原为$((j+u)W_{\mathrm{in}}/S,(i+v)H_{\mathrm{in}}/S)$。

训练按“目标中心落在哪个格子，就由该格子负责”组织监督；在负责格子的$B$个框中，再选当前与真实框IoU最高的那个承担坐标监督。记格子为$a$、框槽为$b$，$I^{\mathrm{resp}}_{a,b}$标出负责框，$I^{\mathrm{other}}_{a,b}$标出其他框槽，$I^{\mathrm{obj}}_a$标出含目标中心的格子。预测中心、宽高和置信度分别为$(\hat x,\hat y)$、$(\hat w,\hat h)$和$\hat C$，真实值去掉帽子。原始目标全部用平方误差，先写位置项：
\begin{equation}
 \begin{aligned}
 e_{a,b}^{\mathrm{pos}}&=(\hat x_{a,b}-x_a)^2
                       +(\hat y_{a,b}-y_a)^2,\\
 e_{a,b}^{\mathrm{size}}&=(\sqrt{\hat w_{a,b}}-\sqrt{w_a})^2
                        +(\sqrt{\hat h_{a,b}}-\sqrt{h_a})^2,\\
 L_{\mathrm{coord}}&=\lambda_{\mathrm c}\sum_{a,b}
 I^{\mathrm{resp}}_{a,b}(e_{a,b}^{\mathrm{pos}}+e_{a,b}^{\mathrm{size}})\text{。}
 \end{aligned}
 \label{eq:vision-yolo-training-process}
\end{equation}
置信度和类别项为
\[
 \begin{aligned}
 L_{\mathrm{conf}}&=\sum_{a,b}I^{\mathrm{resp}}_{a,b}
 (\hat C_{a,b}-\operatorname{IoU}_{a,b})^2
 +\lambda_{\mathrm n}\sum_{a,b}
 I^{\mathrm{other}}_{a,b}\hat C_{a,b}^{\,2},\\
 L_{\mathrm{class}}&=\sum_a I^{\mathrm{obj}}_a\sum_r
 (\hat p_a(r)-\mathbf{1}\{r=c_a\})^2\text{。}
 \end{aligned}
\]
总目标是三项之和。原文取$\lambda_{\mathrm c}=5$、$\lambda_{\mathrm n}=0.5$：大多数框槽不含目标，若不降低它们的权重，“没有目标”的信号会压倒少数负责框的定位信号。宽高取平方根，使同样的像素误差在小框上产生更大的代价。训练分两个阶段：先在ImageNet上以$224\times224$输入预训练前20个卷积层，再加上其余卷积和全连接层，以$448\times448$输入用检测数据训练全部参数。

推理时，每框置信度表示$\mathbb{P}(\mathrm{obj})\operatorname{IoU}$，乘上共享的$\mathbb{P}(c\mid\mathrm{obj})$得到类别相关得分。例如格子$(2,3)$的中心偏移$(0.5,0.25)$还原为$(160,208)$；归一化宽高$(0.2,0.1)$对应像素宽高$(89.6,44.8)$；置信度0.8与货车条件概率0.75相乘，货车框分数为0.6。输入曾缩放或加边时，还须去掉加边并除缩放比例，才是原图上的框。然后按阈值过滤，再用\term{非极大值抑制}（non-maximum suppression, NMS）去重：按得分从高到低取一个框，删除同类中与它IoU超过阈值的候选，对余下框重复。对前述IoU为0.6、得分0.9与0.7的两个框，阈值0.5会删除第二个。

原始YOLO的限制来自它的分配方式。一个格子只预测一组类别概率和两个框，两辆车中心落在同一格时至少会漏掉一辆；7\ensuremath{\times}7的粗网格对远处小车尤其不利。NMS在拥挤场景中还可能删掉真实存在的相邻目标。后来的YOLO版本改变了预测、分配与去重设计，不能沿用此处全部细节。

\subsubsection{RetinaNet}
\label{sec:vision-recognition-fpn}
\label{sec:vision-recognition-retinanet}

YOLO的单一粗网格留下了两个问题：小目标在深层特征中可能已经消失，每个位置的候选又太少。RetinaNet从两处入手：用特征金字塔在多个尺度上预测，并在每个位置放置多个参考框。但候选一多，训练中背景候选就以数量压倒前景，RetinaNet为此提出焦点损失，使单阶段密集检测在训练上不再被大量容易背景主导\scite{vision-1708-02002}。

多尺度特征来自\term{特征金字塔网络}（feature pyramid network, FPN），它是RetinaNet的组成部分\scite{vision-1612-03144}。远处车辆在图上只占几个像素，深层特征步幅大，位置细节丢失；浅层特征保留位置，却缺少语义。经典检测器的做法是\ref{sec:vision-transforms-pyramid}节的图像金字塔：把输入缩放成多个尺寸，每个尺寸分别计算特征，各层语义水平相同，代价是特征计算重复多次。FPN改为只运行一次主干，利用卷积网络本身逐层下采样形成的特征层级，并补上一条自顶向下的路径，使高分辨率的层也带上深层语义。具体做法是把深层特征上采样，与相应浅层特征相加。记主干第$\ell$阶段输出为$\bm{C}_\ell$，自顶向下的中间特征为$\bm{M}_\ell$，则
$\bm{M}_\ell=\operatorname{Conv}_{1\times1}(\bm{C}_\ell)+\operatorname{Up}(\bm{M}_{\ell+1})$，再用$3\times3$卷积得到$\bm{P}_\ell$。$1\times1$横向卷积把通道数统一为256，最近邻上采样把空间尺寸放大两倍，否则两项无法相加。这些附加卷积层都不带非线性。RetinaNet使用$\bm{P}_3$到$\bm{P}_7$五层，$\bm{P}_\ell$的分辨率是输入的$1/2^\ell$，步幅从8到128。其中$\bm{P}_3$至$\bm{P}_5$由主干的$\bm{C}_3$至$\bm{C}_5$经上述自顶向下路径得到；$\bm{P}_6$由$\bm{C}_5$经一个步长为2的$3\times3$卷积得到，$\bm{P}_7$再对$\bm{P}_6$施加ReLU和一个步长为2的$3\times3$卷积；不使用步幅为4的$\bm{P}_2$，以减少计算。较密的层负责小目标，较疏的层负责大目标；若目标在输入缩小时已经丢失，特征融合不能补回观测。

整个检测器按主干、颈部和头部分段，各段的输入输出形状见表\ref{tab:vision-recognition-retinanet-structure}。主干的残差块结构在\ref{sec:vision-ivm-resnet}节已经讲过，这里只用到它各阶段输出的步幅和通道数。

\begin{table*}[tbp]
 \centering
 \small
 \begin{tabularx}{\linewidth}{@{}llXl@{}}
 \toprule
 部分 & 输入形状 & 结构与层数 & 输出形状\\
 \midrule
 主干ResNet-50 & $800\times800\times3$ & 去掉全连接层，取第3至5阶段输出$\bm C_3,\bm C_4,\bm C_5$ & $100^2\times512$，$50^2\times1024$，$25^2\times2048$\\
 颈部FPN & $\bm C_3$至$\bm C_5$ & 每层$1\times1$横向卷积到256通道，最近邻上采样两倍后相加，再经$3\times3$卷积；$\bm P_6$、$\bm P_7$各由一个步长2的$3\times3$卷积得到 & $\bm P_3$至$\bm P_7$：$100^2,50^2,25^2,13^2,7^2$，均为256通道\\
 锚框 & 每层每个位置 & 面积$32^2$（$\bm P_3$）至$512^2$（$\bm P_7$），每层尺度$\{2^0,2^{1/3},2^{2/3}\}$乘长宽比$\{1{:}2,1{:}1,2{:}1\}$，$A=9$ & 覆盖约32至813像素\\
 分类子网 & 每层$h\times w\times256$ & 4个$3\times3$卷积（256通道，各接ReLU），再接$3\times3$卷积与sigmoid；各层共享参数 & $h\times w\times KA$\\
 回归子网 & 每层$h\times w\times256$ & 结构同分类子网但参数独立，末层线性输出；与类别无关 & $h\times w\times4A$\\
 \bottomrule
 \end{tabularx}
 \caption{RetinaNet（ResNet-50-FPN）的分段结构。层数、通道与锚框配置取自原文；输出尺寸按$800\times800$输入推算（原文消融实验多用短边600像素，主结果用短边800像素）。$K$为类别数，$A$为每个位置的锚框数。}
 \label{tab:vision-recognition-retinanet-structure}
\end{table*}

在金字塔每个位置，RetinaNet放置$A=9$个\term{锚框}（anchor box），即3种尺度乘3种长宽比的参考框。两个小型全卷积子网在所有层间共享：分类子网经4层$3\times3$卷积后输出$KA$个通道，每个锚框对每个类别给一个独立sigmoid得分；回归子网结构相同，输出$4A$个通道，给出相对于锚框的修正。两个子网参数不共享，回归子网对所有类别给出同一组修正。设锚框为$(x_a,y_a,w_a,h_a)$，目标为$(x_*,y_*,w_*,h_*)$，回归标签为$((x_*-x_a)/w_a,(y_*-y_a)/h_a,\log(w_*/w_a),\log(h_*/h_a))$；给定预测修正$\bm{\delta}$，逆变换为
\begin{equation}
 \begin{aligned}
 \hat x&=x_a+w_a\delta_x,\qquad \hat y=y_a+h_a\delta_y,\\
 \hat w&=w_a e^{\delta_w},\qquad \hat h=h_a e^{\delta_h}\text{。}
 \end{aligned}
 \label{eq:vision-retinanet-decode}
\end{equation}
中心差除以参考尺寸、宽高比取对数，使标签无像素单位且能表达乘性缩放。教学锚框$(50,30,40,20)$配目标$(60,30,80,20)$，标签为$(0.25,0,\log2,0)$，逆算恢复原目标。

训练按锚框与真实框的IoU分配：IoU不低于0.5的锚框为前景，低于0.4为背景，介于两者之间的忽略。一幅图约有十万个锚框，绝大多数是很容易判断的背景。每个背景锚框的交叉熵很小，但数量巨大，总和仍可能主导梯度。令$p_t$为模型分给真实二值标签的概率，\term{焦点损失}（focal loss）为
\begin{equation}
 L_{\mathrm{focal}}=-\alpha_t(1-p_t)^\gamma\log p_t\text{，}
 \qquad\gamma\geq0\text{。}
 \label{eq:vision-focal-loss}
\end{equation}
其中$\alpha_t$调整正负样本权重，原文取$\gamma=2$、$\alpha=0.25$。$\gamma=0$时退回加权交叉熵；$\gamma=2$时，$p_t=0.99$的容易样本乘上$10^{-4}$，而$p_t=0.5$的困难样本乘上$0.25$。总目标是所有锚框的焦点损失加上正锚框的平滑$L_1$回归损失，再除以正锚框数。分类子网最后一层的偏置按先验前景概率0.01初始化，避免训练初期背景损失过大而不稳定。主干从ImageNet分类权重初始化，金字塔与子网由检测监督训练；原文用16张图的小批量做9万次随机梯度下降迭代。推理时先以0.05的得分阈值过滤，每层最多保留得分最高的1000个锚框并按式\eqref{eq:vision-retinanet-decode}解码，各层合并后以IoU阈值0.5做NMS，最后逆变换到原图。

焦点损失改变的是每个候选的训练贡献，不增加少见类别的观测，因此前景背景失衡与类别长尾仍须分别检查。锚框尺寸与比例是人工设定，与数据中目标形状不匹配时会降低召回，这正是下一个模型要去掉的设计。

\subsubsection{FCOS}
\label{sec:vision-recognition-fcos}

锚框带来了一组需要逐数据集调整的超参数：尺寸、比例、正负IoU阈值；它们还使每个位置的输出数量乘以$A$。FCOS的想法是让特征位置本身充当参考点：落在目标框内的位置直接预测它到框四条边的距离$(l,t,r,b)$，不再以锚框为坐标起点\scite{vision-1904-01355}。

FCOS沿用RetinaNet式的FPN（$\bm{P}_3$至$\bm{P}_7$，步幅8至128）和共享检测头，结构上只有三处改动。第一，$\bm{P}_6$、$\bm{P}_7$改由$\bm{P}_5$（而非$\bm{C}_5$）经步长为2的卷积逐级得到。第二，头部仍是分类与回归两个分支，各含4个256通道的$3\times3$卷积，新加的卷积层后使用组归一化（最后的预测层除外）；分类分支末端输出$h\times w\times K$，回归分支末端输出$h\times w\times4$。第三，与分类输出并列增加一个单层的中心度分支，输出$h\times w\times1$（原文脚注指出把它移到回归分支效果更好，但默认结构仍放在分类分支旁）。于是每个位置只输出一组预测：$K$个类别得分、4个距离和1个中心度，而RetinaNet每个位置要输出$A=9$组。步幅为$s$的特征位置$(i,j)$对应原图点$(x,y)=(\lfloor s/2\rfloor+js,\lfloor s/2\rfloor+is)$，预测框为$(x-l,y-t,x+r,y+b)$。各层共享同一个头，而不同层负责的距离范围相差很大，因此距离分支的输出$x$经$\exp(s_\ell x)$映射，$s_\ell$是每层一个可训练标量：指数保证距离为正，$s_\ell$让共享的头在不同层上自动调整输出量级。默认主干为ResNet-50，训练输入短边800像素、长边不超过1333像素。

监督仍需分配，只是规则换成了几何条件。落在某个真实框内的位置是正例候选；原方法再按层限定可回归的最大距离范围（$\bm{P}_3$至$\bm{P}_6$的上界依次为64、128、256、512像素，$\bm{P}_7$不设上界），把不同大小的目标分到不同层；同一位置落在多个框内时选面积最小的框。训练目标包括类别的焦点损失、正例的IoU损失和中心度的二元交叉熵。\term{中心度}（centerness）衡量位置是否接近框中心，其目标为
\begin{equation}
 c_*=\sqrt{\frac{\min(l,r)}{\max(l,r)}
                 \frac{\min(t,b)}{\max(t,b)}}\text{。}
 \label{eq:vision-fcos-centerness}
\end{equation}
中心点四侧距离均衡时$c_*=1$，靠近边界时其中一个比例趋于零。对于框$(20,10,80,50)$，中心$(50,30)$的距离为$(30,20,30,20)$，中心度为1；位置$(25,30)$的距离为$(5,20,55,20)$，中心度为$\sqrt{1/11}\approx0.302$。图\ref{fig:vision-box-points}把这两个位置与框画在同一坐标系中：两点都可以精确回归同一矩形，中心度却不同，因此它不是预测框IoU的直接测量。

\begin{figure}[htbp]
 \centering
 % 连续像素坐标下的FCOS教学例；所有长度与正文算例一致。
\begin{tikzpicture}[x=1mm,y=-1mm,font=\figurefont\small,
  dist/.style={-{Stealth[length=2mm]},draw=BookBlue,line width=0.9pt}]
 \draw[draw=BookMuted,line width=0.7pt,->] (10,0)--(90,0) node[right] {$x$};
 \draw[draw=BookMuted,line width=0.7pt,->] (10,0)--(10,58) node[below] {$y$};
 \draw[draw=BookTeal,fill=BookTeal!7,line width=1pt] (20,10) rectangle (80,50);
 \draw[draw=BookMuted,dashed,line width=0.7pt] (50,10)--(50,50);
 \draw[draw=BookMuted,dashed,line width=0.7pt] (20,30)--(80,30);
 \fill[BookInk] (50,30) circle (1mm);
 \fill[BookBlue] (25,30) circle (1mm);
 \node[anchor=north west] at (51,33) {中心：$c_*=1$};
 \node[anchor=south west] at (26,27) {边缘：$c_*\approx0.302$};
 \draw[dist] (25,18)--(20,18);
 \draw[dist] (25,18)--(80,18);
 \node[anchor=south] at (22.5,17) {$5$};
 \node[anchor=south] at (57,17) {$55$};
 \draw[draw=BookBlue,line width=0.7pt] (25,18)--(25,29);
 \node[anchor=north] at (20,51) {$(20,50)$};
 \node[anchor=south] at (80,9) {$(80,10)$};
\end{tikzpicture}

 \caption{同一真实框内的中心位置与边缘位置。实线箭头标出边缘位置到左右边界的距离，虚线标出中心轴；两点都可以回归同一矩形，但中心度目标分别为1与约0.302。坐标采用向右、向下为正的连续像素边界。}
 \label{fig:vision-box-points}
\end{figure}

推理时，类别得分乘以预测中心度作为排序分数，降低远离中心的低质量框，再做NMS。FCOS减少了锚框相关的超参数，每个位置的输出也少了$A$倍；但分配规则、尺度范围和NMS仍在，拥挤场景中的重复与漏检问题没有因去掉锚框而消失。

\subsubsection{DETR}
\label{sec:vision-recognition-detr}

前三个模型都先产生密集候选，再用NMS去重。NMS是训练之外的手工规则：阈值高了留下重复框，低了删掉相邻的真目标。DETR的想法是把检测直接视为集合预测：网络输出固定数量的预测，训练时用二分匹配让每个真实目标恰好对应一个预测，其余预测学习“无对象”；这样重复预测在训练中就受到惩罚，推理时不再需要NMS\scite{vision-2005-12872}。

DETR的结构由三段组成。主干（原文为ResNet-50）把$H\times W\times3$的图像变成$H/32\times W/32\times2048$的特征，$1\times1$卷积降到$d=256$通道，展平为$HW/1024$个词元并加正弦位置编码。6层Transformer编码器在这些位置之间交换上下文。解码器读入$N=100$个\term{对象查询}（object query），即可学习的$d$维向量，经6层自注意力和对图像特征的交叉注意力更新；查询之间的自注意力让各查询“知道”其他查询在预测什么，这是避免重复的结构基础。每个查询最后经共享的线性层输出$K+1$类的softmax（多出的一类是“无对象”），经共享的三层前馈网络输出归一化的中心宽高$(c_x,c_y,w,h)\in[0,1]^4$。查询是学习得到的向量，并非预先指定给某辆车的身份，预测序列也没有必须对应目标编号的顺序。Transformer结构本身可参考\BookChapterRef{chap:transformer}{卷五的Transformer章}。

DETR的Transformer部分沿用原始的后置层归一化写法。记编码器第$l$层的输入为$\bm Z^{(l)}\in\mathbb R^{n\times256}$，位置编码为$\bm E$，一个编码层依次计算
\[
 \begin{aligned}
 \bm Z'&=\operatorname{LN}\bigl(\bm Z^{(l)}+\operatorname{MSA}(\bm Z^{(l)}+\bm E,\ \bm Z^{(l)}+\bm E,\ \bm Z^{(l)})\bigr),\\
 \bm Z^{(l+1)}&=\operatorname{LN}\bigl(\bm Z'+\operatorname{FFN}(\bm Z')\bigr)\text{，}
 \end{aligned}
\]
其中$\operatorname{MSA}$的三个参数依次是查询、键和值：位置编码只加在查询和键上，值不含位置。这样位置只影响“读哪里”，不直接混入被读取的内容。解码器的输入在第一层初始化为零，对象查询作为“输出位置编码”在每个注意力层加到查询上；每个解码层在查询自注意力与前馈网络之间多一次交叉注意力，其键为编码器输出加空间位置编码，值为编码器输出。前馈网络是两层带ReLU的逐位置变换，原文未给出其中间宽度。表\ref{tab:vision-recognition-detr-config}按配置表的格式列出原文的四个规模。

\begin{table*}[tbp]
 \centering
 \small
 \begin{tabular}{@{}llccccl@{}}
 \toprule
 模型 & 主干与特征步幅 & 词元数 & 编码／解码层数 & 宽度 & 头数 & 参数量\\
 \midrule
 DETR & ResNet-50，32 & $HW/32^2$ & 6／6 & 256 & 8 & 41M\\
 DETR-DC5 & ResNet-50，16 & $HW/16^2$ & 6／6 & 256 & 8 & 41M\\
 DETR-R101 & ResNet-101，32 & $HW/32^2$ & 6／6 & 256 & 8 & 60M\\
 DETR-DC5-R101 & ResNet-101，16 & $HW/16^2$ & 6／6 & 256 & 8 & 60M\\
 \bottomrule
 \end{tabular}
 \caption{DETR的配置。DETR不切图像块，词元来自主干末层特征的每个位置，因此“图像块大小”一列换成特征步幅；DC5表示把主干最后一阶段的步幅去掉、改用空洞卷积，使特征分辨率加倍。对象查询数均为100，随机失活比例0.1；前馈网络的中间宽度论文未给出。原文说明DETR（R50）的4130万参数中，主干占2350万，Transformer占1780万。以$800\times1216$的输入为例，步幅32给出$25\times38=950$个词元，步幅16则为3800个。}
 \label{tab:vision-recognition-detr-config}
\end{table*}

训练的关键是匹配。标注有$m\leq N$个对象，第$i$个为$(c_i,\bm{b}_i)$，第$j$个预测为$(\hat{\bm{p}}_j,\hat{\bm{b}}_j)$。DETR寻找从真实对象到不同预测的单射$\sigma$，使总匹配代价最小：
\begin{equation}
 \hat{\sigma}=\argmin_\sigma
 \sum_{i=1}^{m}\left[
 -\hat p_{\sigma(i)}(c_i)
 +\lambda_1\lVert\bm{b}_i-\hat{\bm{b}}_{\sigma(i)}\rVert_1
 +\lambda_{\mathrm{g}}L_{\mathrm{GIoU}}
    (\bm{b}_i,\hat{\bm{b}}_{\sigma(i)})\right]\text{。}
 \label{eq:vision-detr-matching}
\end{equation}
这就是\term{二分匹配}（bipartite matching）：一侧是真实对象，一侧是预测，每条边的权重是匹配代价，要求每个真实对象恰好配一个不同的预测。它可以用匈牙利算法在多项式时间内精确求解。原文取$\lambda_1=5$、$\lambda_{\mathrm g}=2$。广义交并比$\operatorname{GIoU}=\operatorname{IoU}-|E\setminus(\bm b\cup\bm b')|/|E|$，其中$E$是两框的最小包围矩形，定位损失为$L_{\mathrm{GIoU}}=1-\operatorname{GIoU}$；不相交框的IoU都为零，额外的空白比例仍能区分隔得近和隔得远的框。

假设两个真实对象对三个查询的总代价矩阵为
\[
 \bm{C}=\begin{pmatrix}1&2&9\\1.2&5&8\end{pmatrix}\text{。}
\]
逐行独立取最小都会选第一个查询，违反唯一对应。二分匹配选择第一对象到第二查询、第二对象到第一查询，总代价为$3.2$，小于反过来分配的$6$。改变标注的排列只改变匹配的写法，不改变最佳集合代价。

匹配选定后视为本轮的固定监督，不对离散匹配本身求梯度。匹配的查询承担负对数分类损失和两项框损失，未匹配的查询学习“无对象”，该类的权重降为0.1，以免大量空查询压倒少数目标。每个解码层的输出都接同样的预测头并计算辅助损失，帮助深层解码器训练。主干从ImageNet权重初始化并冻结批归一化统计，其余部分由检测数据训练；优化器为AdamW，Transformer学习率$10^{-4}$、主干$10^{-5}$。基线模型训练300轮，在16块V100上约需3天；与Faster R-CNN比较时延长到500轮。推理时直接丢弃“无对象”和低分输出，不执行NMS。

DETR把去重从后处理移进了训练目标，代价也随之而来：全局注意力的计算随特征位置数平方增长，难以使用高分辨率特征，小目标效果较弱；每个查询初期没有空间先验，匹配在训练早期不稳定，导致收敛很慢；查询数$N$还限定了一次可输出的对象上限。后续的可变形注意力和以锚点初始化查询的改进主要针对这几点，下面的Grounding DINO即建立在这一系列改进之上。

\subsubsection{Grounding DINO}
\label{sec:vision-recognition-grounding-dino}

前四个模型的类别表都写在分类层的输出维度里。当目标变成“红色货车”或“公交车旁的货车”时，固定类别头无法表达；先用固定检测器找出所有车、再用CLIP对每个裁剪打分，又会把文字与视觉的交互推迟到定位之后，定位本身无法利用文字条件。Grounding DINO的想法是让文字从一开始就参与检测：文字与图像在特征阶段相互融合，用与文字最相关的图像位置初始化查询，解码时查询同时读取图像和文字，最后把类别分类换成查询与文字词元的对齐分数\scite{vision-2303-05499}。它的检测部分沿用DINO（一种DETR改进型检测器），因此集合匹配与无NMS输出都继承自DETR。

输入是一幅图像和一段文字。文字可以是用句点分隔的类别名列表，例如“轿车. 货车. 公交车.”，也可以是一个指代表达。文字经BERT-base编码为至多256个词元特征$\bm{T}\in\mathbb{R}^{n_t\times d}$；图像经Swin Transformer主干得到多尺度特征，展平为$\bm{V}\in\mathbb{R}^{n_v\times d}$。为了防止并列的类别名相互干扰，原文采用子句级文字表示：注意力掩码阻止不相关的类别名彼此注意，每个类别仍保留各自的逐词特征。

模型由三段计算组成。第一段是特征增强器，共6层，每层对图像特征做可变形自注意力、对文字特征做自注意力，再做图像到文字和文字到图像两个方向的交叉注意力，使车辆附近的图像特征带上“货车”的语义，文字特征也带上图中实际出现的内容。第二段是语言引导的查询选择：对每个图像位置$i$计算它与所有词元的最大相似度$\max_r\bm{v}_i^\top\bm{t}_r$，取得分最高的900个位置，用它们的位置初始化查询的参考框，查询的内容部分仍是可学习向量。此时选中的只是候选位置，不能当成已经确认的货车。第三段是6层跨模态解码器，每层依次做查询自注意力、对图像特征的可变形交叉注意力、对文字特征的交叉注意力和前馈网络，逐层修正框。输出是900个查询各自的归一化框，以及一个$900\times n_t$的对齐logit矩阵，其中第$(i,r)$项是查询$i$与词元$r$的点积。表\ref{tab:vision-recognition-gdino-structure}按分段格式列出这些组件。

\begin{table*}[tbp]
 \centering
 \small
 \begin{tabularx}{\linewidth}{@{}llXl@{}}
 \toprule
 部分 & 输入形状 & 结构与层数 & 输出形状\\
 \midrule
 图像主干 & $H\times W\times3$ & Swin-T（T版）或Swin-L（L版）；T版取步幅8至32的三个尺度，再把步幅32的特征下采样得到步幅64，共4个尺度；L版直接取步幅4至32的4个尺度 & 4个尺度，展平为$n_v\times256$\\
 文字主干 & 至多256个词元 & BERT-base，子句级注意力掩码 & $n_t\times256$\\
 特征增强器 & 图像与文字特征 & 6层；每层图像可变形自注意力、文字自注意力、图像到文字与文字到图像交叉注意力、前馈网络 & 形状不变\\
 查询选择 & 增强后的特征 & 按与词元的最大相似度取前900个图像位置初始化参考框 & 900个查询\\
 跨模态解码器 & 900个查询 & 6层；每层查询自注意力、图像可变形交叉注意力、文字交叉注意力、前馈网络 & 900个框与$900\times n_t$对齐logit\\
 \bottomrule
 \end{tabularx}
 \caption{Grounding DINO的分段结构，隐藏维数$d=256$。原文给出的总参数量为T版1.72亿、L版3.41亿。}
 \label{tab:vision-recognition-gdino-structure}
\end{table*}

训练数据混合了三类：带类别框的检测数据（类别名拼成文字输入）、带短语与框对应的定位数据，以及由已有模型为图文对产生伪框的图像描述数据。每条训练记录包含图像、文字、目标框以及框与文字中词元的对应，一个对象可以对应多个词元。训练沿用DETR式二分匹配，匹配代价中分类、$L_1$与广义IoU三项的权重为2.0、5.0和2.0。匹配后对每个查询—词元logit施加焦点损失，使匹配查询与其对应词元的分数高、与其他词元的分数低；框用$L_1$和广义IoU约束，三项损失权重依次为1.0、5.0和2.0，每个解码层和编码器输出都计算辅助损失。T版的预训练数据为Objects365，或再加上GoldG（Flickr30k Entities与Visual Genome的短语定位数据）和约四百万对伪标注图文数据；L版还加入了OpenImages、COCO和指代表达数据。图像和文字主干从各自的预训练权重初始化，融合与解码部分在上述数据上联合训练。使用时不再为每次查询更新权重。

推理时，对每个查询取它在各词元上的最大对齐分数作为框分数，保留超过框阈值的查询；再在该查询的词元分数中找出超过文字阈值的词元，拼回短语作为标签；最后把归一化中心宽高乘以原图尺寸并转换为四边。框阈值、文字阈值和短语恢复规则都会改变最终输出，评价时须固定。若模型只在“货车”词元上高分而“红色”得分弱，最终记录不能悄悄把整个属性短语写成已核实的事实。

能接收“红色货车”“公交车旁的货车”这类字符串，只说明接口容纳了条件，不能证明属性和关系已被正确解析。可以分别遮去参照物、更换颜色词或保留多个同类目标，检查预测是否按所需证据改变。完整的表达定位协议在\ref{chap:vision-multimodal-understanding}章讨论。

\subsubsection{OWOBJ}
\label{sec:vision-recognition-unknown-discovery}
\label{sec:vision-recognition-owobj}

仓库画面中出现一个尚未加入类别表的搬运装置。固定类别检测器可能把它认作货车，也可能当成背景；Grounding DINO允许用户写出装置名称，但用户此时未必知道应当查询什么。\emph{未知对象发现}（unknown object discovery）要求在没有名字时先提出“这里可能有一个对象”，保留区域供后续核验。它与开放词汇定位的差别在于名称是否已作为输入；它与拒识也不同，拒识只回答已有裁剪是否属于已知类，发现还要确定未知对象在哪里。

问题的根源在训练分配。DETR式训练中，没有匹配到已知标注的查询一律学习“无对象”，而这些查询中既有真实背景，也有恰好未被标注的对象；把它们全部用零监督，检测器就学会把一切未标注的东西判为背景。OWOBJ的想法是给检测器加一个\emph{对象性}（objectness）信号，即一个候选是否像一个独立对象，而不是它属于哪个已知类；未匹配查询不再一律接受零标签，而是按其特征与已知对象特征分布的接近程度获得软标签\scite{vision-owobj-2025}。

OWOBJ不另造检测器，而是作为插件接入已有的Faster R-CNN式或DETR式检测器，推理流程不变。在开放世界检测实验中，它接在两个以可变形DETR为基础的方法之上，主干为ResNet-101；少样本检测实验接在Faster R-CNN式的DeFRCN上，开放词汇检测实验接在以CLIP为主干的CORA上。以可变形DETR式的实现为例，主干给出多尺度特征，经可变形编码器和解码器，得到每个查询的嵌入$\bm{q}_i\in\mathbb{R}^d$。查询之后接三个分支：框分支是3层前馈网络；分类分支输出$Z+1$维，其中$Z$维对应已知类，多出的一维表示未知或对象性；对象性模块本身不含可学习的层，只维护查询嵌入的一个高斯分布，用来给多出的那一维生成软标签。查询数$N$和嵌入维数$d$沿用所接的基线检测器，论文正文未给出具体数值。用集合$C_{\mathrm{k}}$表示训练已标注类别，输出可以是已知记录$(\bm{b},c,s)$，也可以是未命名记录$(\bm{b},\mathrm{unknown},s_{\mathrm{o}})$。原方法由匹配查询特征的移动平均估计均值$\bm{\mu}$与正定对角协方差$\bm{\Sigma}$，用
\begin{equation}
 \begin{aligned}
 d_{\mathrm{M}}(\bm{q})=
 (\bm{q}-\bm{\mu})^\top\bm{\Sigma}^{-1}(\bm{q}-\bm{\mu}),\qquad
 s_{\mathrm{o}}(\bm{q})=\exp[-d_{\mathrm{M}}(\bm{q})]
 \end{aligned}
 \label{eq:vision-owobj-objectness}
\end{equation}
生成对象性软标签。\term{马氏距离}（Mahalanobis distance）按各特征方向的正常变动尺度比较偏离，避免数值尺度大的通道独占距离。$s_{\mathrm{o}}$在距离为零时为1，距离增大时下降；它是训练设计中的对象性信号，不是已经校准的真实对象概率。

训练在原检测器的类别与框损失之外加入三类约束。匹配查询继续接受已知类别和框监督，并通过减小$d_{\mathrm M}$使对象特征在分布中更集中；未匹配查询用$s_{\mathrm{o}}$监督对象性输出。原工作还用变分近似的先验匹配项约束分布，高斯先验的方差随估计后验调整并加入方差偏置$\beta^2$（原文消融中$\beta=0.5$最好），以处理少标注条件下的分布收缩；已知与未匹配查询之间的能量间隔项$\max(0,E_{\mathrm u}-E_{\mathrm k}+\delta)$（$\delta=0.2$）则改变分类得分的相对约束，这里的能量来自分类logit的对数指数和。移动统计按$\bm\mu^{(t+1)}=\rho\bm\mu^{(t)}+(1-\rho)\bar{\bm q}$更新，$0\leq\rho<1$控制旧批次的保留程度，对角方差同样更新并保持正数下界。训练从所接检测器的预训练权重出发；使用时沿用该检测器的框和对象性输出筛选未知候选，式\eqref{eq:vision-owobj-objectness}只在训练中提供软标签，不能拿新图的真实标注重新估计对象分布。

作为一维教学例子，若$\mu=0$、$\Sigma=1$，两个未匹配查询为0.2和2，则软对象性分别为$e^{-0.04}\approx0.961$和$e^{-4}\approx0.018$。给前者较强的对象监督，可以避免直接把未知候选抹掉；但若背景特征也靠近这个分布，它同样会被提出。因此对象性需要同时检查未知召回和背景误检，不能由“保留更多区域”推断发现已经正确。

发现之后可以暂存原图编号、框、得分和当前类别表，再由人工或后续证据命名新类，形成新增训练记录。把新类加入检测器属于\term{类别增量学习}（class-incremental learning），后续学习既要扩展输出集合，又要保留旧类能力，可参考\BookChapterRef{chap:ket-update-maintenance}{卷八“知识的更新和维护”章}；它不是这一次未命名检测已经完成的步骤。仓库装置被误报为已知货车、墙面被误报为未知对象、加入装置类后货车退化，是三种不同的错误。

六个模型的对比见表\ref{tab:vision-recognition-detection-compare}，比较固定在同一街景检测任务和同一输出形式$(\bm{b},c,s)$上。从YOLO到FCOS，改变的是候选的密度和参考方式，去重始终依靠NMS；DETR把去重移进训练，换来更简洁的推理，付出的是收敛速度和小目标；Grounding DINO与OWOBJ都保留DETR式集合输出，分别改变了类别从哪里来和未标注对象怎样监督。

\begin{table*}[tbp]
 \centering
 \small
 \begin{tabularx}{\linewidth}{@{}lXXXX@{}}
 \toprule
 模型 & 负责目标的输出 & 类别来源 & 去重方式 & 主要代价\\
 \midrule
 YOLO & 中心所在格子的框槽 & 固定类别表 & NMS & 粗网格，同格多目标漏检\\
 RetinaNet & 与目标IoU足够大的多尺度锚框 & 固定类别表 & NMS & 锚框设定，背景失衡需焦点损失\\
 FCOS & 落在框内且尺度相符的位置 & 固定类别表 & 中心度加NMS & 仍依赖分配规则与NMS\\
 DETR & 二分匹配选中的查询 & 固定类别表加无对象 & 训练中的唯一匹配 & 收敛慢，小目标弱，查询数上限\\
 Grounding DINO & 匹配查询，与词元对齐 & 使用时给出的文字 & 唯一匹配 & 阈值与短语恢复影响输出\\
 OWOBJ & 匹配查询加对象性分支 & 已知类加未命名对象 & 沿用所接检测器 & 未知召回与背景误检需权衡\\
 \bottomrule
 \end{tabularx}
 \caption{同一检测任务下六个模型在关键设计上的差别。每一行相对上一行只改动一到两项设计；表中不给出精度数值，因为原论文的数据集与设置各不相同。}
 \label{tab:vision-recognition-detection-compare}
\end{table*}

\subsection{像素类别的预测}
\label{sec:vision-recognition-segmentation}
\label{sec:vision-recognition-dense-pixels}

框内常含背景，也可能含另一辆车；需要准确范围时，就要给每个像素一个判断。\term{语义分割}（semantic segmentation）给每个像素分配类别，同类车辆可以连成一片；\term{实例分割}（instance segmentation）给不同车辆各自的掩码；\term{全景分割}（panoptic segmentation）同时覆盖道路、天空等不可数区域，使每个像素获得类别，可数对象的像素再获得实例编号。两辆车的并集可以是正确的车辆语义区域，却是错误的单个实例，因此三种标签不能未经说明相互转换。图\ref{fig:vision-segmentation-outputs}沿用同一个离散街景，显示增加实例编号与背景类别各带来什么信息。本小节讨论只输出类别、不区分实例的像素预测；区分实例的掩码留给下一小节。

\begin{figure}[htbp]
 \centering
 % 同一离散场景的三种分割输出；字符在灰度下保留类别和身份区别。
\begin{tikzpicture}[font=\figurefont\small,x=4.5mm,y=-4.5mm]
 \foreach \offset/\title in {0/语义分割,8/实例分割,16/全景分割}{
  \begin{scope}[xshift=\offset*4.5mm]
   \node[anchor=south] at (3,0) {\title};
   \fill[BookPaper] (0,0.4) rectangle (6,4.4);
   \draw[BookMuted,line width=0.7pt] (0,0.4) rectangle (6,4.4);
   \draw[BookRule,line width=0.5pt] (0,1.4)--(6,1.4);
   \ifnum\offset=8
    \node[text=BookMuted] at (3,3.6) {不标背景};
   \else
    \node[text=BookMuted] at (3,3.6) {道路};
    \node[text=BookMuted] at (3,0.9) {天空};
   \fi
   \filldraw[fill=BookTeal!18,draw=BookTeal,line width=0.8pt]
      (0.5,1.8) rectangle (2.8,2.8);
   \ifnum\offset=0
    \filldraw[fill=BookTeal!18,draw=BookTeal,line width=0.8pt]
      (3.2,1.8) rectangle (5.5,2.8);
    \node at (1.65,2.3) {车};
    \node at (4.35,2.3) {车};
   \else
    \filldraw[fill=BookAmber!16,draw=BookAmber,dashed,line width=0.8pt]
      (3.2,1.8) rectangle (5.5,2.8);
    \node at (1.65,2.3) {车1};
    \node at (4.35,2.3) {车2};
   \fi
  \end{scope}
 }
\end{tikzpicture}

 \caption{三种分割结果的区别。语义分割保留类别，实例分割分开可数对象，全景分割共同保留背景类别与对象编号；文字编号和边界线型使身份区别在灰度打印中仍可辨认。图中区域为离散教学构造，框状图元表示教学掩码，不是检测框。}
 \label{fig:vision-segmentation-outputs}
\end{figure}

像素类别预测的输入是$H\times W\times3$的图像，输出是$H\times W\times K$的类别得分图，逐像素取最大即得类别图。难点在于分辨率与上下文之间的矛盾：判断道路中央的一块灰色像素是路面还是车顶，需要很大范围的上下文；画出路杆和车窗的细边界，又需要接近原图的分辨率。卷积网络靠逐层下采样扩大感受范围，代价正是丢失位置细节，直接把最后的低分辨率特征放大只会增加采样位置，补不回细边界。

经典分割方法走的是另一条路：它们在原分辨率上工作，边界由像素之间的低层差异决定，但不判断类别。\ref{sec:vision-features-otsu}节的Otsu阈值法在灰度直方图上选一个使类间方差最大的门限，把像素分成两类，只在前景和背景灰度分布分得开时有效。\ref{sec:vision-features-watershed}节的分水岭把梯度幅值看成地形，从各个极小值“注水”，水域相遇处就是边界；它得到闭合轮廓，却对噪声敏感，容易过分割，通常要先用标记控制注水的起点。\ref{sec:vision-features-graph-segmentation}节的图割把每个像素看成图的顶点，一元项衡量像素颜色与前景、背景颜色模型的匹配程度，成对项惩罚相似的相邻像素被分开；对二值标记，这一能量的全局最小值可由最小割求得，GrabCut再用用户给的框迭代更新颜色模型。\ref{sec:vision-features-superpixel}节的SLIC超像素在颜色和位置的联合空间里做局部$k$均值，把图像切成大小相近、贴合边界的小块。这些方法共同的长处是边界精细、不需要训练；共同的短处是不知道“路面”和“车顶”是两类：颜色相近的不同类别会被合并，纹理复杂的同一物体会被切碎。学习方法补上的正是语义，而经典方法的边界线索仍以别的形式保留下来：DeepLab系列的早期版本曾在网络输出之后接一个带成对项的条件随机场来贴合边界，交互分割和抠图也继续使用图割与颜色传播的思想。

DeepLabv3+对这一矛盾的回答是在编码端用空洞卷积扩大上下文而不再继续下采样，在解码端引入浅层特征恢复边界。BIT处理的是同一矛盾的另一种形态：输入变成同一地点两个时期的影像，类别变成“变化／不变”，上下文不仅来自同一幅图的周围，还来自另一时期的对应区域；它用少量语义词元在两期之间交换上下文。两者的输出形式相同，都是逐像素类别图，差别在于类别的含义和上下文的来源。

\subsubsection{DeepLabv3+}
\label{sec:vision-recognition-deeplab}

在DeepLabv3+之前，DeepLabv3已经用空洞卷积和多尺度池化取得了较强的上下文，但它直接把输出步幅为16的得分图双线性放大16倍，物体边界因而粗糙；另一类编码器—解码器网络逐级恢复分辨率，边界较好，上下文却受编码器下采样限制。DeepLabv3+把两者结合：以DeepLabv3作为编码器提供上下文，再加一个简单的解码器，用浅层特征修正边界\scite{vision-1802-02611}。

\term{空洞卷积}（atrous convolution）在卷积核的相邻采样点之间留出间隔，空洞率$r$表示间隔为$r$。例如一维三点核读取$x_{i-r},x_i,x_{i+r}$，参数个数仍为三，覆盖范围却从3扩大到$2r+1$。它让网络在不降低特征分辨率、不增加参数的前提下读到更大范围。

网络按编码器与解码器分段，各段形状见表\ref{tab:vision-recognition-deeplab-structure}。主干（ResNet-101或改造的Xception）把最后几个阶段的下采样换成空洞卷积，输出步幅为16，即$H/16\times W/16\times2048$的特征。\term{空洞空间金字塔池化}（atrous spatial pyramid pooling, ASPP）模块在这一特征上并行使用五个分支：一个$1\times1$卷积、三个空洞率分别为6、12、18的$3\times3$卷积（原文图1），以及一个整图平均池化分支；各分支输出拼接后经$1\times1$卷积得到256通道的编码器输出。不同空洞率对应不同大小的邻域，整图池化提供全局场景信息。以输出步幅16、空洞率18为例，一个$3\times3$核的两个最远采样点相隔$2\times18=36$个特征位置，约合输入上的576像素。解码器先把编码器输出双线性放大4倍到$H/4\times W/4$；从主干步幅为4的浅层取出特征（ResNet-101取第2阶段残差块在步长2之前的最后一层），用$1\times1$卷积压到48通道，以免浅层通道数过多压过语义特征——原文在8至64通道之间比较，48通道最好；两者拼接后经两个256通道的$3\times3$卷积细化（一个或三个都稍差），再用$1\times1$卷积输出$K$个类别得分，最后双线性放大4倍回到输入尺寸。

\begin{table*}[tbp]
 \centering
 \small
 \begin{tabularx}{\linewidth}{@{}llXl@{}}
 \toprule
 部分 & 输入形状 & 结构与层数 & 输出形状\\
 \midrule
 主干 & $513\times513\times3$ & ResNet-101，最后一个阶段改用空洞卷积，输出步幅16；或改造的对齐Xception & 浅层$129\times129\times256$；深层$33\times33\times2048$\\
 ASPP & $33\times33\times2048$ & $1\times1$卷积；空洞率6、12、18的$3\times3$卷积；整图池化；拼接后$1\times1$卷积 & $33\times33\times256$\\
 解码器：上采样 & $33\times33\times256$ & 双线性放大4倍 & $129\times129\times256$\\
 解码器：浅层分支 & $129\times129\times256$ & $1\times1$卷积 & $129\times129\times48$\\
 解码器：融合 & $129\times129\times304$ & 拼接后2个$3\times3$卷积（256通道），$1\times1$卷积输出类别 & $129\times129\times K$\\
 输出 & $129\times129\times K$ & 双线性放大4倍 & $513\times513\times K$\\
 \bottomrule
 \end{tabularx}
 \caption{DeepLabv3+的分段结构。输出步幅、通道数和解码器层数取自原文；输出尺寸按原文训练裁剪$513\times513$推算，ResNet-101各阶段通道数沿用\ref{sec:vision-ivm-resnet}节。原文以ResNet-101为主干时，不带解码器的计算量约为810亿次乘加，加上解码器后约为1013亿次；参数量论文未给出。}
 \label{tab:vision-recognition-deeplab-structure}
\end{table*}

改造的Xception有三处变化：中间流的模块重复16次（原文图4），使网络更深；所有最大池化换成带步长的深度可分离卷积，使任意位置都可以改用空洞卷积来保持分辨率；每个$3\times3$深度卷积之后都加批归一化和ReLU。\term{深度可分离卷积}（depthwise separable convolution）把一个标准卷积拆成两步：先对每个通道单独做$3\times3$卷积，再用$1\times1$卷积混合通道，参数和计算量因此大幅下降。原文在ASPP和解码器中也改用深度可分离卷积，乘加次数下降约三至四成，精度基本不变。

训练使用逐像素交叉熵。对训练图$i$的有效位置集$\Omega_i$，类别标注为$y_{i,p}$，按类别softmax得到$p_{i,p,r}$，则
\begin{equation}
 L_{\mathrm{pixel}}=
 -\frac{\sum_i\sum_{p\in\Omega_i}\log p_{i,p,y_{i,p}}}
        {\sum_i|\Omega_i|}\text{。}
 \label{eq:vision-dense-pixel-training}
\end{equation}
无效标签的像素不进入$\Omega_i$。每个位置的得分梯度是预测类别分布减真实类别指示，再经上采样、解码器和ASPP传回主干。主干从分类权重初始化，ASPP和解码器由像素监督学习。缩放、裁剪图像时，标签须用最近邻等保持类别的采样，不能用插值制造不存在的类别编号。推理时先把类别得分图插值回原图，再逐像素取最大；直接插值整数类别编号没有概率含义。训练与测试可以使用不同的输出步幅（例如测试用8以换取更细的结果），这改变了计算量和精度，应共同记录。

道路内部的类别主要依赖上下文，细路杆依赖分辨率；若路杆在输入缩小时已经消失，解码器的浅层特征也无法补回。两车相邻时，即使边缘清楚，纯类别图也不分配实例身份。

\subsubsection{BIT}
\label{sec:vision-recognition-bit}

遥感应用中的目标常常是找出新增的建筑，而不是给单幅影像中的每栋楼分类。\term{变化检测}（change detection）的输入是同一地区两个时期、已经配准的影像，输出是逐像素的“变化／不变”图。常见的孪生卷积网络分别提取两期特征再比较，但卷积的感受范围有限，难以利用远处的上下文判断一块区域的变化是否有意义；对两期全部像素做自注意力，计算又随像素数平方增长。BIT（bitemporal image transformer）的想法是先把每期特征压缩成少量语义词元，在这个小空间里用Transformer建模两期之间的上下文，再把结果投回像素\scite{vision-2103-00208}。

两期影像$\bm{X}_1,\bm{X}_2$经共享权重的ResNet-18主干得到特征图，展平为$\bm{F}_1,\bm{F}_2\in\mathbb{R}^{n\times d}$，$n$是空间位置数。语义分词器为每期计算空间注意力$\bm{A}\in\mathbb{R}^{n\times k}$，每列在位置上做softmax归一化，得到$k$个加权汇总的词元$\bm{T}=\bm{A}^{\top}\bm{F}\in\mathbb{R}^{k\times d}$；原文取$k=4$。每个词元可以理解为“建筑区域”“道路区域”这类高层概念在该期影像中的汇总。两期词元拼成$2k$个，加位置编码后经Transformer编码器相互注意，使一期的概念能参考另一期的对应内容。编码后的词元再拆回两期，各自经Transformer解码器投回像素：每个像素特征作为查询，读取本期的$k$个词元，得到上下文增强的特征$\bm{F}_1',\bm{F}_2'$。最后把两期特征上采样到原图尺寸，计算$|\bm{F}_1'-\bm{F}_2'|$，由浅层卷积预测头输出两类得分。

表\ref{tab:vision-recognition-bit-structure}给出原文主模型的分段形状。ResNet-18原有5个阶段、各下采样2倍；BIT把后面阶段的步长改为1，主模型只保留前4个阶段，接一个输出32通道的逐点卷积压缩特征，再双线性放大，使特征的下采样倍数为4。Transformer很小：词元数$L=4$，编码器1层，解码器8层，每层8个头、每头8维，前馈网络为两层、中间宽度64并用GELU，各子层前置层归一化；只给词元加可学习位置编码，解码器只有交叉注意力和前馈网络，不在像素之间做自注意力。预测头是两个带批归一化的$3\times3$卷积，输出通道依次为32和2。整个模型只有约355万参数，$256\times256$输入下约43.5亿次浮点运算；同样主干、不加Transformer的基线为338万参数，可见语义词元带来的额外开销很小。

\begin{table*}[tbp]
 \centering
 \small
 \begin{tabularx}{\linewidth}{@{}llXl@{}}
 \toprule
 部分 & 输入形状（每期） & 结构与层数 & 输出形状（每期）\\
 \midrule
 主干 & $256\times256\times3$ & ResNet-18前4个阶段，第4阶段步长改为1；两期共享权重 & $32\times32\times256$\\
 压缩与上采样 & $32\times32\times256$ & 逐点卷积到32通道，双线性放大2倍 & $64\times64\times32$\\
 语义分词器 & $4096\times32$ & 逐点卷积得到4张空间注意力图，各自softmax后加权求和 & $4\times32$\\
 Transformer编码器 & 两期拼接$8\times32$ & 1层，8头 & $8\times32$，再拆回两期\\
 Transformer解码器 & 像素$4096\times32$与本期4个词元 & 8层，只有交叉注意力与前馈网络 & $64\times64\times32$\\
 预测头 & 两期特征放大到$256\times256\times32$后取绝对差 & 2个$3\times3$卷积（32、2通道），逐像素softmax & $256\times256\times2$\\
 \bottomrule
 \end{tabularx}
 \caption{BIT主模型（原文记为BIT\_S4）的分段结构。层数、宽度与参数量取自原文；张量形状按$256\times256$输入推算，ResNet-18第4阶段的256通道沿用其标准配置。}
 \label{tab:vision-recognition-bit-structure}
\end{table*}

训练使用配准影像与像素变化标签，在有效像素上以式\eqref{eq:vision-dense-pixel-training}的交叉熵学习“变化／不变”。每次训练对两期图像和变化标签施加同样的裁剪与翻转，保证同一位置仍指同一地表。主干从ImageNet权重初始化，分词器、Transformer和预测头由变化数据学习。推理时保持两期的顺序与配准，把分数恢复到原网格，再用阈值选出变化区域。

新建筑应计为变化，树木由绿变黄是否计为变化则取决于标注定义。若两期影像整体错开几个像素，差异分支也会在边缘产生响应，因此配准是任务前提。词元交换能够利用上下文关联，却不能把传感器、季节或光照差异自动解释成真实的建设活动。

两个模型的对比如表\ref{tab:vision-recognition-pixel-compare}。BIT相对DeepLabv3+改变了输入（一幅变两幅）和上下文的组织方式（空洞卷积的固定邻域变为少量全局词元），换来了两期之间的全局比较和较低的注意力成本；代价是只有$k$个词元时细小变化可能被汇总掉，且整个方法依赖配准。

\begin{table}[htbp]
 \centering
 \small
 \begin{tabularx}{\linewidth}{@{}lXX@{}}
 \toprule
 比较项 & DeepLabv3+ & BIT\\
 \midrule
 输入 & 单幅图像 & 同一地区两期配准影像\\
 类别含义 & 语义类别 & 变化／不变\\
 上下文来源 & 多空洞率卷积与整图池化 & 两期共享的少量语义词元\\
 边界恢复 & 解码器融合步幅4的浅层特征 & 词元投回像素后上采样\\
 主要前提与风险 & 细结构须在输入中保留 & 配准误差与非目标变化\\
 \bottomrule
 \end{tabularx}
 \caption{两种像素类别预测模型的比较。两者输出都是逐像素类别图，差别在类别含义与上下文的组织。}
 \label{tab:vision-recognition-pixel-compare}
\end{table}

\subsection{对象掩码的预测}
\label{sec:vision-recognition-mask-sets}

对象掩码的输出是一组二值图，每张图对应一个实例，并附类别或得分。它比像素类别图多出的信息是身份：两辆相邻的车必须分在两张掩码里。难点因此有两层：一是确定有哪些对象，即检测问题中“哪个输出负责哪个目标”在像素层面的重现；二是确定要哪些对象，即目标由类别表、用户点击还是概念文字指定。

四个模型对这两层问题给出了不同的组合。Mask R-CNN先检测框，再在每个框内预测掩码，把实例分开的任务交给检测器，掩码分支只需在框内区分前景和背景。Mask2Former取消了框，沿用DETR的集合匹配思路，让每个查询直接预测一张整图掩码，同一结构可以输出语义、实例和全景结果。SAM改变了目标指定方式：不再依赖类别表，而是由用户的点或框指定一个对象，输出不带类别名的掩码。SAM 3再把指定方式从一个局部点击扩展为一个概念短语，要求找出图中所有符合概念的实例。

\subsubsection{Mask R-CNN}
\label{sec:vision-recognition-mask-rcnn}

两阶段检测器Faster R-CNN已经能为每个对象给出框和类别，自然的扩展是在框内再预测一张掩码。但它从特征图上截取区域特征时，会把框的连续坐标两次取整到特征格点；这对框分类影响不大，对像素级掩码却会造成可见的错位。Mask R-CNN的想法是在检测分支旁并行增加一个小型掩码分支，并用区域对齐代替取整，使框内位置与特征准确对应\scite{vision-1703-06870}。

网络的第一阶段是\term{区域建议网络}（region proposal network, RPN）：在FPN各层的每个位置放置锚框，预测对象性得分和框修正，解码并去重后得到约一千个候选框。第二阶段对每个候选用\term{区域对齐}（RoIAlign）取特征：把候选框等分为固定数目的小格，在每格中按连续坐标取若干采样点，用双线性插值读取特征值，再平均或取最大，得到固定尺寸的区域特征。设四个相邻格点的特征值为$f_{0,0},f_{0,1},f_{1,0},f_{1,1}$，采样点在水平方向偏移$a$、竖直方向偏移$b$，二者均在$[0,1]$内，则
\begin{equation}
 f(a,b)=(1-a)(1-b)f_{0,0}
 +a(1-b)f_{0,1}+(1-a)bf_{1,0}+abf_{1,1}\text{。}
 \label{eq:vision-roialign}
\end{equation}
用数值$(0,4,2,6)$与$(a,b)=(0.25,0.5)$可得$f=2$，既保留横向的四分之一偏移，也保留纵向中点；若直接取最近格点，这些偏移信息就会消失。原文在每个小格中取4个规则分布的采样点，再取平均，并指出结果对采样点的具体位置和个数不敏感。以ResNet-FPN为主干时，两个头的结构见表\ref{tab:vision-recognition-maskrcnn-heads}：框分支读取$7\times7\times256$的区域特征，经两个1024维全连接层输出类别和框修正；掩码分支读取$14\times14\times256$的区域特征，经4个$3\times3$卷积和一个$2\times2$、步长2的反卷积，输出$28\times28\times K$，每个类别一个二值掩码通道。原文另有以ResNet第4阶段特征为主干的C4版本：区域特征为$7\times7\times1024$，先经第5阶段残差块得到$7\times7\times2048$，框分支在其平均池化后输出，掩码分支经一次反卷积得到$14\times14\times K$，分辨率比FPN版低一半。

\begin{table*}[tbp]
 \centering
 \small
 \begin{tabular}{@{}lllrrl@{}}
 \toprule
 阶段或层组 & 层类型 & 卷积核与步长 & 输出通道 & 重复次数 & 输出尺寸\\
 \midrule
 框分支 & 区域对齐 & — & 256 & 1 & $7\times7$\\
        & 全连接 & — & 1024 & 2 & $1\times1$\\
        & 全连接 & — & $K$；$4K$ & 1 & $1\times1$\\
 掩码分支 & 区域对齐 & — & 256 & 1 & $14\times14$\\
          & 卷积 & $3\times3$ & 256 & 4 & $14\times14$\\
          & 反卷积 & $2\times2$，步长2 & 256 & 1 & $28\times28$\\
          & 卷积 & $1\times1$ & $K$ & 1 & $28\times28$\\
 \bottomrule
 \end{tabular}
 \caption{Mask R-CNN（ResNet-FPN）的两个区域头，取自原文图4。隐藏层均接ReLU；COCO上$K=80$。框分支最后一层并列输出类别得分与每类的框修正。}
 \label{tab:vision-recognition-maskrcnn-heads}
\end{table*}

训练时，与某个真实框IoU不低于0.5的候选为正例，并记住匹配的实例编号。总目标为$L=L_{\mathrm{cls}}+L_{\mathrm{box}}+L_{\mathrm{mask}}$，再加上RPN的对象性与框损失。掩码损失只在正例上、只对其匹配真实类别的那个通道计算逐像素sigmoid二元交叉熵；目标是匹配实例的掩码裁剪到候选框并缩放为$28\times28$。这样各类别的掩码通道不在同一像素上竞争，类别判断交给框分支。一个正候选若匹配错了实例，即使“车辆”类别正确，也会学到另一辆车的掩码。主干从分类权重初始化，各部分联合训练。原文输入短边为800像素，FPN版每幅图采样512个区域，正负比例1比3；8块GPU、每块2幅图，共16万次迭代。

推理时FPN版保留1000个候选，经框分支筛选、去重后，只对得分最高的100个框运行掩码分支，取预测类别对应的掩码通道，以0.5为阈值二值化，并缩放放回原图。若一个$28\times28$掩码对应原框$(20,10,80,50)$，归一化区域位置$(\xi,\eta)$回到$(20+60\xi,10+40\eta)$。恢复后的掩码、类别和框共享实例编号。区域对齐解决采样量化误差，不能补回候选阶段漏掉的车辆；$28\times28$的掩码分辨率对大对象的细边界也有限。

\subsubsection{Mask2Former}
\label{sec:vision-recognition-mask2former}

Mask R-CNN依赖框：框漏了，掩码就没有；道路、天空这类不可数区域又不适合用框表示，语义分割和全景分割只好另用一套网络。Mask2Former的想法是沿用DETR的集合预测，但每个查询输出的不是框而是一张整图掩码加一个类别，这样语义、实例和全景三种任务都可以写成“预测一组带类别的掩码”\scite{vision-2112-01527}。它相对早期掩码分类方法的关键改动是掩码注意力：每个查询只读取自己当前预测的前景区域。

网络由三部分组成。主干产生多尺度特征；像素解码器（默认为6层多尺度可变形注意力）在$1/32$、$1/16$、$1/8$三个尺度上交换信息，再经一次上采样和横向连接得到$1/4$分辨率的逐像素嵌入$\bm{f}_p\in\mathbb{R}^{d}$。Transformer解码器维护$N$个查询（默认100），共9层，按$1/32$、$1/16$、$1/8$的顺序循环读取三个尺度的特征。每层先做掩码交叉注意力，再做查询间自注意力和前馈网络。设第$i$个查询解码成掩码向量$\bm{e}_i\in\mathbb{R}^d$，它在位置$p$的掩码概率为$\sigma(\bm{e}_i^\top\bm{f}_p)$。\term{掩码注意力}（masked attention）把上一层预测的掩码缩放到当前尺度并以0.5二值化，在注意力得分中加入屏蔽项$M_{i,j}$：前景位置为0，其余为$-\infty$，再作softmax。这样每个查询只聚合自己可能对应区域内的特征，而不是被整图背景稀释。若某个查询的掩码为空，所有位置都被屏蔽，分母失效，实现需恢复可读位置。与早期掩码分类方法相比，Mask2Former还把查询特征改为可学习的初始值，并在每层给特征加正弦位置编码和可学习的尺度编码，去掉了解码器中的随机失活。

\begin{table}[htbp]
 \centering
 \small
 \begin{tabular}{@{}lcccc@{}}
 \toprule
 主干 & 查询数 & 解码层数 & 参数量 & 计算量\\
 \midrule
 ResNet-50 & 100 & 9 & 44M & 226G\\
 ResNet-101 & 100 & 9 & 63M & 293G\\
 Swin-L & 200 & 9 & 216M & 868G\\
 \bottomrule
 \end{tabular}
 \caption{Mask2Former在COCO全景分割上的配置，取自原文。像素解码器均为6层多尺度可变形注意力；计算量为每幅图的浮点运算次数。Transformer的宽度、头数与前馈网络宽度在正文中只以符号出现，论文未给出具体数值。}
 \label{tab:vision-recognition-mask2former-config}
\end{table}

表\ref{tab:vision-recognition-mask2former-config}列出原文报告的三个规模。参数增长主要来自主干；解码部分在三种配置中结构相同，只有Swin-L的全景与实例模型把查询数增加到200。

训练沿用DETR式二分匹配，匹配代价和损失都由类别项和掩码项组成。Dice项度量软预测与真实前景的重叠，对于预测$p_j$和真实标签$m_j$，一种带平滑量$\epsilon>0$的形式为
\begin{equation}
 L_{\mathrm{Dice}}=
 1-\frac{2\sum_j p_jm_j+\epsilon}
 {\sum_j p_j+\sum_j m_j+\epsilon}\text{。}
 \label{eq:vision-dice}
\end{equation}
二元交叉熵逐像素约束预测，Dice从区域整体约束重叠，二者对小对象和大对象的敏感程度不同。原文的类别、二元交叉熵和Dice权重分别为2、5、5，“无对象”类权重为0.1。掩码损失不在全图计算，而是在$K=12544$个（即$112\times112$个）采样点上计算，匹配时均匀采样，最终损失时按重要性采样偏向预测不确定的点，以降低高分辨率掩码的显存开销；原文报告这使每幅图的训练显存从18GB降到6GB。每层解码输出都计算辅助损失。COCO上的训练使用AdamW，学习率$10^{-4}$（主干乘0.1），批大小16，训练50轮，输入经大尺度抖动后裁剪为$1024\times1024$。

推理时按任务组织输出。语义输出将各查询的类别概率与掩码概率相乘，对查询求和，再逐像素选类；实例输出保留查询编号，按类别得分与掩码质量筛选；全景输出还要处理区域重叠、无对象和不可数类别。两个车辆查询在语义图中可以合并，在实例图中不能合并。原工作按任务与数据集分别训练，统一架构不等于一份参数对所有任务通用；早期掩码错误也可能限制后续读取，小对象和被遮挡对象需要单独检查。

\subsubsection{SAM}
\label{sec:vision-recognition-sam}

前两个模型都只能分割训练类别表中的对象。用户有时只想选中图上的一个区域，甚至说不出它的类别名，例如一块污渍或车门上的贴纸。SAM（segment anything model）定义了\term{可提示分割}（promptable segmentation）任务：给定图像和一个提示（点、框或粗略掩码），输出提示所指对象的有效掩码；提示有歧义时，至少给出一个合理的结果\scite{vision-2304-02643}。它的想法是把昂贵的图像编码与廉价的提示解码分开，同一幅图编码一次后，用户每改一次提示只重跑轻量解码器，从而支持实时交互。

三个组件的配置见表\ref{tab:vision-recognition-sam-config}，其中ViT-H的层数与宽度沿用ViT原文的Huge配置\scite{dosovitskiy2020vit}。图像编码器是以掩码自编码预训练的ViT-H，图像块边长16，输入缩放并填充到$1024\times1024$，因此得到$64\times64=4096$个词元。对4096个词元做全局注意力代价很高，SAM的编码器大多数块只在$14\times14$的窗口内做注意力，另有4个均匀分布的块做全局注意力，让远处信息仍能交换。编码器之后接一个颈部：$1\times1$卷积把通道降到256，再接一个256通道的$3\times3$卷积，两者之后各有层归一化，输出$64\times64\times256$的图像嵌入。

提示编码器把点的坐标经位置编码映射为256维向量，并加上“正点”“负点”或“框左上角”“框右下角”类型的学习嵌入；粗略掩码以图像嵌入4倍的分辨率输入，经两个$2\times2$、步长2的卷积（输出4和16通道）和一个$1\times1$卷积变成256通道，与图像嵌入逐位置相加，没有掩码提示时加一个学习的“无掩码”嵌入。掩码解码器是两层改造的Transformer，宽度256、8个头、前馈网络中间宽度2048，交叉注意力中查询、键、值降到128维以节省计算。输出词元（3个掩码词元加1个IoU词元）与提示词元拼在一起，每层依次做：词元自注意力、词元到图像的交叉注意力、逐词元前馈网络、图像到词元的交叉注意力。最后一次的“图像读词元”让图像嵌入也吸收了提示信息。解码后，图像嵌入经两个$2\times2$、步长2的反卷积（输出64和32通道）上采样4倍到$256\times256$，每个掩码词元经3层MLP得到32维的动态权重，与上采样特征逐位置做内积，得到一张$256\times256$的掩码logit。

\begin{table*}[tbp]
 \centering
 \small
 \begin{tabular}{@{}lccccccl@{}}
 \toprule
 组件 & 图像块大小 & 词元数 & 层数 & 宽度 & 头数 & MLP倍率 & 参数量\\
 \midrule
 图像编码器（ViT-H） & 16 & 4096 & 32 & 1280 & 16 & 4 & 约6.3亿\\
 提示编码器 & — & 每点或角点1个 & — & 256 & — & — & 论文未单列\\
 掩码解码器 & — & 输出与提示词元 & 2 & 256 & 8 & 8 & 论文未单列\\
 \bottomrule
 \end{tabular}
 \caption{SAM的配置。ViT-H的层数、宽度、头数与MLP倍率沿用ViT原文的Huge配置，SAM原文图13标注其编码器约636M参数；SAM编码器中4个块为全局注意力、其余为$14\times14$窗口注意力。掩码解码器的MLP倍率按前馈宽度2048与宽度256折算。}
 \label{tab:vision-recognition-sam-config}
\end{table*}

一个提示可能有多种合理解释：落在车门上的点可以指车门，也可以指整辆车。SAM为单个提示同时输出三张掩码（大致对应局部、部件、整体），另有一个头为每张掩码预测它与目标的IoU，用于排序。训练时只对三张中与真实掩码损失最小的那一张反传掩码损失，这让不同输出分工到不同粒度，而不是被平均成模糊的中间结果；掩码损失是焦点损失与Dice损失按20比1线性组合，IoU头用均方误差监督。训练模拟交互：先从真实掩码中采一个点或框，之后在当前预测的错误区域采新点，并把上一轮的掩码logit作为提示输入，每个掩码共模拟11轮。训练数据来自作者的数据引擎，模型辅助人工标注，再逐步转为全自动标注，最终形成约1100万张图像、11亿个掩码的SA-1B数据集。最终模型用AdamW训练9万次迭代，批大小256，在256块A100上约需68小时。

使用时先缓存图像嵌入，点或框改变只重跑提示编码和解码；掩码logit上采样到输入尺寸，去掉填充并恢复原图尺寸，再以0为阈值二值化。用户在漏掉的车尾补正点、在误选的路面补负点，可以逐步消除歧义；IoU预测可以推荐候选，但选择车门还是整车仍由用户意图决定。报告结果时应说明提示来源和点击次数：用真实标注框作提示与由检测器自动给框，具有不同的信息和错误条件。SAM的掩码不附带类别名称。

\subsubsection{SAM 3}
\label{sec:vision-recognition-sam3}

SAM的一个点只选择一个对象。“找出所有校车”则需要先判断这个概念是否出现，再枚举所有符合概念的实例；用SAM完成这件事，必须由用户逐个点击，或者先用开放词汇检测器给框。SAM 3把指定方式从局部提示扩展为概念：输入一个简短名词短语和若干带正负标签的图像示例框，输出所有匹配实例的框和掩码。原报告称这一任务为可提示概念分割（promptable concept segmentation, PCS）\scite{vision-2511-16719}。

模型的核心想法是把“概念是否出现”与“哪个实例匹配”分开。普通DETR式检测器让每个查询同时回答这两个问题：一辆普通客车的查询既要判断自己是车，又要判断全图是否存在校车，局部信息不足以回答后者。SAM 3增加一个共享的存在词元，与对象查询一同进入解码器，由存在预测头得到$p_{\mathrm{pres}}$；查询$i$的局部分数$s_{i\mid\mathrm{pres}}$只表达“若概念存在，我是否为其一例”，最终分数为$s_i=p_{\mathrm{pres}}s_{i\mid\mathrm{pres}}$。

检测器的结构沿用DETR式查询解码。文字短语经文字编码器得到词元特征，示例框经编码后同时携带位置、框内视觉内容和正负标签，二者共同组成概念条件。图像特征来自与文字编码器联合预训练的视觉编码器，经融合编码器读取概念条件；对象查询在解码器中细化框，掩码头再为每个查询产生实例掩码。负示例告诉模型哪些相似对象不属于当前概念，而不是把该区域的像素删除。

原报告附录给出了各组件的规模，见表\ref{tab:vision-recognition-sam3-structure}。视觉编码器是在54亿图文对上对比训练的感知编码器（Perception Encoder），共32层，其中4层做全局注意力，其余层在窗口内做注意力：$1008\times1008$的输入被分成$3\times3$个336像素见方的窗口，每个窗口$24\times24$个词元。融合编码器的每个块依次做图像自注意力、对提示词元的交叉注意力和前馈网络；检测解码器沿用DETR式结构，并加入以框为参照的位置偏置、逐层修正框等改进。存在词元与对象查询一同进入解码器，输出经MLP头给出所有查询共享的$p_{\mathrm{pres}}$。

\begin{table*}[tbp]
 \centering
 \small
 \begin{tabularx}{\linewidth}{@{}llXl@{}}
 \toprule
 部分 & 输入 & 结构与层数 & 规模或输出\\
 \midrule
 视觉编码器 & $1008\times1008$图像 & 感知编码器，32层（4层全局注意力，其余窗口注意力） & 约4.5亿参数\\
 文字编码器 & 名词短语 & 因果Transformer，上下文至多32个词元 & 约3亿参数\\
 融合编码器 & 图像特征与提示词元 & 6个块：自注意力、对提示的交叉注意力、前馈网络 & 条件化图像特征\\
 检测解码器 & 200个对象查询加存在词元 & 6个块，DETR式 & 每查询的框、局部分数；共享的存在分数\\
 分割头 & 解码后的查询与多尺度特征 & 改编自MaskFormer，语义与实例分割共用 & 每查询一张实例掩码\\
 跟踪器 & 视频帧 & 沿用SAM~2的记忆注意力与记忆库，与检测器共用视觉编码器 & 跨帧掩码\\
 \bottomrule
 \end{tabularx}
 \caption{SAM~3的组件，取自原报告附录。检测器与跟踪器合计约1亿参数，全模型约8.5亿参数；图像块边长原文未直接写出。}
 \label{tab:vision-recognition-sam3-structure}
\end{table*}

训练分阶段进行。先用大规模图文对预训练视觉和文字编码器；再以图像、短语和掩码组成的数据训练检测器及两侧编码器，数据既有伪标注也有人工核验的标注；高质量阶段加入存在头，并根据漏检采正示例、根据高分误检采负示例，累加示例以模拟交互纠错。存在头用二元交叉熵覆盖含概念的正图和难负概念；概念存在时，实例按类别和范围代价做二分匹配，再以局部得分、框的$L_1$与广义IoU、掩码的焦点与Dice损失训练查询。概念缺席的图像不向局部分类分支施加逐查询的负梯度，这与普通DETR的无对象训练不同。视频跟踪阶段冻结主干训练跟踪器，不应倒推为前面的图像训练也一直冻结。使用时更换“校车”文字只改变条件，不重新学习权重。

若存在分数为0.9，两个实例的局部分数为0.8和0.6，最终分数为0.72和0.54；若存在分数为0.05，同样的局部分数只剩0.04和0.03。保留超过阈值的查询，将各掩码放回原图并保持独立实例编号。共享存在分支的一个错误会同时改变所有查询：把“校车”误判为不存在，多个真实实例会一起被压低，形成相关的漏检。评价因此分别检查概念有无、每实例召回和掩码边界。简短概念和示例框扩大了目标的指定方式，但不足以推断任意关系嵌套或长指令都已解析；视频中的身份传播在\ref{chap:vision-video-recognition}章继续讨论。

四个模型的对比见表\ref{tab:vision-recognition-mask-compare}。从Mask R-CNN到Mask2Former，改变的是实例的来源（先有框，到查询直接给掩码），换来了三类分割的统一与不依赖框，代价是训练和显存更重；从Mask2Former到SAM，改变的是目标的指定方式（类别表到用户提示），换来了类别无关的交互分割，代价是失去类别名；SAM 3再把提示从一个点扩展到一个概念，恢复了语义，代价是存在判断错误会牵连所有实例。

\begin{table*}[tbp]
 \centering
 \small
 \begin{tabularx}{\linewidth}{@{}lXXXX@{}}
 \toprule
 模型 & 目标的指定 & 实例的来源 & 输出是否带类别 & 主要代价或风险\\
 \midrule
 Mask R-CNN & 固定类别表 & 先检测框，框内预测掩码 & 是 & 框漏检即掩码缺失，掩码分辨率有限\\
 Mask2Former & 固定类别表 & 查询经集合匹配直接给整图掩码 & 是 & 训练成本高，早期掩码错误限制读取\\
 SAM & 点、框或粗掩码 & 每个提示一组歧义候选 & 否 & 依赖提示质量，需交互消歧\\
 SAM 3 & 概念短语与正负示例 & 查询枚举全部匹配实例 & 是（概念） & 存在判断错误牵连全部实例\\
 \bottomrule
 \end{tabularx}
 \caption{同一街景掩码任务下四个模型的比较。表中只比较接口与机制，不给出不同论文设置下的精度。}
 \label{tab:vision-recognition-mask-compare}
\end{table*}

\subsection{前景透明度的预测}
\label{sec:vision-recognition-matting}

发丝和半透明边缘的一个像素可能同时收到前景与背景的光，二值掩码只能选择保留或删除，无法表达这种混合。\term{图像抠图}（image matting）估计每个像素的前景贡献比例$\alpha\in[0,1]$。在线性颜色空间的简化合成模型中，每个像素满足
\begin{equation}
 \bm{x}=\alpha\bm{f}+(1-\alpha)\bm{b}\text{，}
 \label{eq:vision-matting}
\end{equation}
其中$\bm{x},\bm{f},\bm{b}\in\mathbb{R}^3$是观测、前景与背景颜色。每个像素只有三个颜色观测，却有前景、背景和$\alpha$共七个未知数，问题本身欠定，需要额外条件或先验。路线的差别就在于额外条件从哪里来：传统方法假设颜色在局部平滑变化，由用户画出的\term{三分图}（trimap）给出确定前景、确定背景和未知区域三类标记，再在未知区域内传播颜色；学习方法仍可使用三分图作为条件，但把“怎样从图像和条件推出$\alpha$”交给网络，从大量合成样本中学习。不使用三分图的方法则让网络自行判断前景，代价是对目标指定的歧义更敏感。下面以Deep Image Matting说明有三分图条件的学习路线。

此前的抠图方法主要依赖局部颜色统计：当前景和背景颜色相近，或背景纹理复杂时，颜色本身不足以区分，结果在发丝处出现明显误差；另一个限制是缺少足够大的真实$\alpha$数据集来训练学习方法。Deep Image Matting从两方面入手：构造大规模合成训练集，并用深度网络从图像的结构和语义信息中推断$\alpha$，而不只依赖颜色\scite{vision-1703-03872}。

网络分两个阶段。第一阶段是以VGG-16为基础的编码器—解码器，输入是RGB图像与三分图拼接成的$H\times W\times4$张量，编码器逐级下采样提取特征，解码器逐级上采样输出$H\times W\times1$的初步$\hat\alpha$。第二阶段是一个小型全卷积细化网络，输入是原图与初步$\hat\alpha$，预测残差以锐化边缘细节，输出最终$\alpha$。确定前景和确定背景区域直接取1和0，学习的重点是未知带中的混合比例。

逐层结构见表\ref{tab:vision-recognition-matting-layers}。编码器有14个卷积层和5个最大池化层，取VGG-16的前14个卷积层（第14层由全连接层fc6改写为卷积）初始化，第四个输入通道的权重初始化为零。解码器比编码器小得多：6个卷积层和5个反池化层，最后接一个输出单通道的预测层；每个反池化层把特征放大回上一级的分辨率，与编码器的5次池化对应。细化网络是4个$3\times3$卷积，前3层64通道并接ReLU，最后一层输出1通道；它不做下采样，输入的第四通道（初步$\hat\alpha$，缩放到$[0,1]$）直接加到输出上，网络只学习残差。

\begin{table*}[tbp]
 \centering
 \small
 \begin{tabular}{@{}lllrrl@{}}
 \toprule
 阶段或层组 & 层类型 & 卷积核与步长 & 输出通道 & 重复次数 & 输出尺寸\\
 \midrule
 编码器 & 卷积与池化 & VGG-16前14层 & 至512 & 14＋5 & $H/32\times W/32$\\
 解码器 & 卷积 & $1\times1$ & 512 & 1 & $H/32\times W/32$\\
        & 反池化与卷积 & $5\times5$ & 512，256，128，64，64 & 5 & 逐级至$H\times W$\\
        & 预测层 & $5\times5$ & 1 & 1 & $H\times W$\\
 细化网络 & 卷积 & $3\times3$ & 64 & 3 & $H\times W$\\
          & 卷积 & $3\times3$ & 1 & 1 & $H\times W$\\
 \bottomrule
 \end{tabular}
 \caption{Deep Image Matting的逐层结构。编码器与解码器的层数取自原文正文，解码器各层的核大小与通道数取自原文图3；输入为$H\times W\times4$（RGB加三分图），训练时$H=W=320$。细化网络的输入为原图与初步$\hat\alpha$，共4通道。}
 \label{tab:vision-recognition-matting-layers}
\end{table*}

训练数据由合成得到：把已有精确$\alpha$和前景颜色的对象，按式\eqref{eq:vision-matting}合成到大量不同的背景图上；原文用493个前景对象各配100张来自COCO和VOC的背景，得到49300张训练图，测试集则由50个对象各配20张背景组成；训练时以未知像素为中心随机裁剪$320\times320$的块；三分图由真实$\alpha$的未知区域随机膨胀生成，使网络适应不同宽度的未知带。第一阶段的目标在未知区$U$上组合两项：
\[
 \begin{aligned}
 L={}&w_\ell\frac1{|U|}\sum_{p\in U}\sqrt{(\hat\alpha_p-\alpha_p)^2+\epsilon^2}\\
 &+(1-w_\ell)\frac1{|U|}\sum_{p\in U}
 \sqrt{\lVert\hat\alpha_p\bm f_p+(1-\hat\alpha_p)\bm b_p-\bm x_p\rVert_2^2+\epsilon^2}\text{，}
 \end{aligned}
\]
其中$\epsilon$是避免零点不可导的小常数，原文取$w_\ell=0.5$。第一项直接约束贡献比例，第二项是合成损失，用已知的前景和背景重新合成，检查该比例的颜色后果；两项通过同一混合公式连接。编码器从分类网络权重初始化，先训练第一阶段，再固定它训练细化网络，最后联合微调。三分图可以由目标$\alpha$构造以供训练，但使用时必须实际取得提示，不能把真实$\alpha$隐藏在输入中。

使用时输出$\alpha$；要把对象放到新背景$\bm{b}'$上，还需估计或取得前景颜色$\bm{f}$，再合成$\alpha\bm{f}+(1-\alpha)\bm{b}'$。例如灰度前景为0.8、原背景为0.2、$\alpha=0.25$，观测为0.35。换成0.6的背景，正确合成为$0.25\times0.8+0.75\times0.6=0.65$；若二值化为背景，结果为0.6，发丝的贡献消失；若直接把观测0.35当成前景颜色，又会混入旧背景。三分图若误把发丝标成确定背景，网络会被错误条件限制。合成训练数据与真实照片在光照和噪声上存在差距，边界、贡献比例和最终合成质量应分别核验。

\section{数量、姿态与空间测量}
\label{sec:vision-recognition-measurement}

“有几辆车”“人的手腕在哪里”“路牌离相机多远”“这几张照片是从哪里拍的”分别要求数量、带部位语义的位置、深度和相机位姿。它们都可以借用上一节的对象范围，但输出已经是测量量：有单位、有坐标系，误差也要按测量量计算，不能由检测或分割指标代替。本节按测量量划分为四个小节。数量只需要规定个体和计入规则；姿态要给每个关节一个位置，并说明它属于哪个人；深度要给每个像素一个距离；相机与场景几何则要同时确定拍摄位置和三维结构。从数量到几何，所需的几何约定越来越多，下面先回顾后三个小节共用的相机模型。

要把手腕或路牌放到三维空间中，需要\ref{sec:vision-formation-camera}节建立的相机模型。这里只回顾后文要用的结论和记号，推导与标定方法见\ref{sec:vision-formation-camera-intrinsics}节和\ref{sec:vision-formation-calibration}节。本章忽略镜头畸变（实际使用前先用标定参数校正），像素坐标向右、向下增加，相机光轴向前。

世界点$\bm{P}_{\mathrm w}\in\mathbb{R}^3$经外参$\bm{R},\bm{t}$变到相机坐标$\bm{P}_{\mathrm c}=(x_{\mathrm c},y_{\mathrm c},z)^\top=\bm{R}\bm{P}_{\mathrm w}+\bm{t}$，其中$\bm{R}$是旋转矩阵；$\bm{t}$是世界到相机变换的平移，相机中心的世界坐标则是$-\bm{R}^{\top}\bm{t}$。内参矩阵$\bm{K}$含像素焦距$f_x,f_y$与主点$(c_x,c_y)$，于是像素$(u,v)$满足
\begin{equation}
 z\begin{pmatrix}u\\v\\1\end{pmatrix}
 =\bm{K}(\bm{R}\bm{P}_{\mathrm w}+\bm{t}),\qquad
 \bm{K}=\begin{pmatrix}f_x&0&c_x\\0&f_y&c_y\\0&0&1\end{pmatrix}\text{。}
 \label{eq:vision-camera-projection}
\end{equation}
$z>0$表示点在相机前方。本章反复用到这个式子的逆向读法：令$\tilde{\bm{p}}=(u,v,1)^\top$，沿这条像素射线的相机点都是$\bm{P}_{\mathrm c}(z)=z\bm{K}^{-1}\tilde{\bm{p}}$，$z>0$待定，所以像素不能单独确定深度。光轴深度$z$与到相机中心的直线距离$\rho$也不是同一个量：
\begin{equation}
 \rho=\lVert\bm{P}_{\mathrm c}\rVert_2
 =z\sqrt{1+\left(\frac{u-c_x}{f_x}\right)^2
            +\left(\frac{v-c_y}{f_y}\right)^2}\text{。}
 \label{eq:vision-depth-range-measurement}
\end{equation}
教学中取$f_x=f_y=800$像素、主点$(320,240)$和像素$(520,240)$。若$z=4$米，反投影得到$(1,0,4)$米，直线距离为$\sqrt{17}\approx4.123$米；若$z=8$米，点$(2,0,8)$仍投到同一像素。两次投影相同，测量量却不同。

多视图增加了射线约束，但还需区分\emph{相对尺度}（relative scale）与\emph{度量尺度}（metric scale）：前者只确定长度比例，后者还确定一单位对应多少米。把全部世界点和相机平移同时乘$a>0$，相机点与$z$也乘$a$，式\eqref{eq:vision-camera-projection}的像素不变。已知相机基线、已知物体长度或带单位的深度观测可以确定这个比例；只把首相机设为原点，固定的是参照位置和朝向，没有测出尺度。网络从训练分布中学到的尺度是附加先验，须在当前设备与场景中验证。

同一个世界点进入第二幅图时，消去世界点，可得\ref{sec:vision-formation-two-view}节的两相机相对位姿：
\begin{equation}
 \begin{aligned}
 \bm{P}_{\mathrm c,2}
 &=\bm{R}_{2,1}\bm{P}_{\mathrm c,1}+\bm{t}_{2,1},\\
 \bm{R}_{2,1}&=\bm{R}_2\bm{R}_1^\top,\qquad
 \bm{t}_{2,1}=\bm{t}_2-\bm{R}_2\bm{R}_1^\top\bm{t}_1\text{。}
 \end{aligned}
 \label{eq:vision-relative-camera-measurement}
\end{equation}
将第一幅图的反投影点代入，再用第二幅图的$\bm{K}_2$投影，就得到一个可检查的跨图位置。裁剪和缩放也必须进入内参：若横纵缩放为$a_x,a_y$，裁剪原点为$(u_0,v_0)$，则新焦距为$(a_xf_x,a_yf_y)$，新主点为$(a_x(c_x-u_0),a_y(c_y-v_0))$。把新像素配旧内参会直接改变射线方向。落在图像内且位于相机前方只是可见的必要条件；\emph{遮挡}（occlusion）指同一射线上更近的表面挡住更远的表面，跨图比较一个投影点时还要检查深度是否相容。以下姿态、深度和几何恢复都沿用这些约定。

\subsection{对象数量的估计}
\label{sec:vision-recognition-counting}

计数的输入是图像和一个计数规则（数什么类别、在哪个区域、被遮挡一半的算不算），输出是一个数，通常还希望附带每个对象的位置。难点随密度变化：路口两三辆车时，逐个找出来再数最可靠；广场上几百个人头相互遮挡时，单个人的范围难以标注，检测器也难以分开。

两条路线的差别在于是否先形成个体。基于检测实例汇总的计数先用上一节的检测器得到对象清单，再按规则数清单中的项；它保留每个对象的位置和身份，可以逐个核验，但在密集遮挡中会漏检和重复。基于密度图积分的计数不形成个体，而是预测每个位置分摊多少计数，再对整幅图求和；它适合密集人群，代价是失去个体清单。

\subsubsection{基于检测实例汇总的计数}
\label{sec:vision-recognition-instance-counting}

这条路线不需要新模型：检测器可以是\ref{sec:vision-recognition-detection}节的任何一个，计数完全由汇总规则决定。给定去重后的检测集合，按类别、得分和区域筛选，再把满足条件的指示量相加。若区域为$R$，以框中心$\bm{o}_i$是否在$R$内为计入规则，则
\begin{equation}
 \hat n_R=\sum_i
 \mathbf{1}\{c_i\in C_{\mathrm{target}}\}
 \mathbf{1}\{s_i\geq\eta\}
 \mathbf{1}\{\bm{o}_i\in R\}\text{，}
 \label{eq:vision-detection-count}
\end{equation}
其中$c_i$是类别，$s_i$是得分，$\eta$是得分阈值。“中心落入”与“任意可见部分落入”是两种不同规则，边界上的对象会得到不同计数，规则须事先规定。

阈值$\eta$直接决定计数的偏差：阈值高则漏数，阈值低则多数，它应在验证集上按计数误差而不是检测精度选择。拥挤场景中NMS可能删掉相邻的真实对象，DETR式检测器则受查询数上限约束。例如真实有5辆车，漏掉1辆，另1辆被重复预测，总数仍为5；只看计数误差会掩盖两种错误，因此还应检查对象匹配。跨帧统计时，同一辆车会在多帧中出现，需由\ref{chap:vision-video-recognition}章的跟踪维持身份。

\subsubsection{基于密度图积分的计数}
\label{sec:vision-recognition-mcnn}

人群过密时，标注每个人的框既昂贵又不准确，检测器在严重遮挡下也会大量漏检；同一画面中，近处的人头可能有几十像素，远处的只有几个像素。MCNN用头部点标注代替框，学习一张\term{密度图}（density map），即每个位置分摊多少计数贡献；对密度图求和就是人数\scite{vision-mcnn2016}。它针对尺度变化的想法是用几个感受范围不同的卷积分支并行处理同一幅图，各自适应不同大小的人头。

首先需要把点标注变成密度监督。若标注点为$\bm{p}_i$，用归一化核构造
\begin{equation}
 D(\bm{p})=\sum_{i=1}^{n}G_{\sigma_i}(\bm{p}-\bm{p}_i),
 \qquad
 \sum_{\bm{p}}G_{\sigma_i}(\bm{p}-\bm{p}_i)=1\text{。}
 \label{eq:vision-density-target}
\end{equation}
每个核的和为1，所以整幅密度图之和等于标注人数。原工作采用几何自适应核：若第$i$点最近$k$个邻点的平均距离为$\bar d_i$，取$\sigma_i=\beta\bar d_i$，原文取$\beta=0.3$。拥挤处邻点近，核就窄；稀疏处核较宽，以此用点间距近似透视造成的人头大小变化。孤立点或点数不足时需回退为固定核宽。

网络有三列结构相同的卷积网络，每列都是“卷积—池化—卷积—池化—卷积—卷积”共4个卷积层和2个$2\times2$最大池化，激活函数为ReLU；三列只在卷积核大小和滤波器个数上不同，逐层配置见表\ref{tab:vision-recognition-mcnn-layers}。大核列从$9\times9$起始，适合近处的大人头；小核列从$5\times5$起始，适合远处的小人头。原文有意让大核的列使用较少的滤波器，以控制参数量。两次池化使三列的输出都是输入的$1/4$分辨率，于是可以在通道维拼接成$8+10+12=30$个通道，再经一个$1\times1$卷积映射为单通道密度图$\hat D$。这个$1\times1$卷积相当于为三列学习融合权重，它取代了全连接层，因此网络可以接收任意尺寸的输入，原文也不把图像缩放到统一尺寸，以免缩放扭曲密度。训练目标是与目标密度的平方误差$L=(2n)^{-1}\sum_{i=1}^n\sum_p(\hat D_i(p)-D_i(p))^2$，标签密度也须下采样到$1/4$分辨率，并保持总和不变。原文先分别预训练每一列（直接把该列第4个卷积层的输出映射为密度图），再用这些权重初始化三列并联合微调；数据增强是从每幅图裁出9个、各为原图四分之一大小的块。推理时对$\hat D$求和得到人数，结果可以是非整数。

\begin{table*}[tbp]
 \centering
 \small
 \begin{tabular}{@{}lllrrl@{}}
 \toprule
 阶段或层组 & 层类型 & 卷积核与步长 & 输出通道 & 重复次数 & 输出尺寸\\
 \midrule
 大核列 & 卷积，池化 & $9\times9$；$2\times2$，步长2 & 16 & 1 & $H/2\times W/2$\\
        & 卷积，池化 & $7\times7$；$2\times2$，步长2 & 32 & 1 & $H/4\times W/4$\\
        & 卷积 & $7\times7$ & 16，8 & 2 & $H/4\times W/4$\\
 中核列 & 卷积，池化 & $7\times7$；$2\times2$，步长2 & 20 & 1 & $H/2\times W/2$\\
        & 卷积，池化 & $5\times5$；$2\times2$，步长2 & 40 & 1 & $H/4\times W/4$\\
        & 卷积 & $5\times5$ & 20，10 & 2 & $H/4\times W/4$\\
 小核列 & 卷积，池化 & $5\times5$；$2\times2$，步长2 & 24 & 1 & $H/2\times W/2$\\
        & 卷积，池化 & $3\times3$；$2\times2$，步长2 & 48 & 1 & $H/4\times W/4$\\
        & 卷积 & $3\times3$ & 24，12 & 2 & $H/4\times W/4$\\
 融合 & 拼接，卷积 & $1\times1$ & 1 & 1 & $H/4\times W/4$\\
 \bottomrule
 \end{tabular}
 \caption{MCNN的逐层结构，取自原文图3。输入为$H\times W$的人群图像；“16，8”表示两个连续卷积层的输出通道依次为16和8。卷积保持空间尺寸，下采样只来自两次最大池化。原文未给出参数总量。}
 \label{tab:vision-recognition-mcnn-layers}
\end{table*}

保持总和是密度标签的关键约束。将每个密度单元分成四个，四个新单元应分享原来的质量，不能每个都复制原值；连续情形下，缩放$(x',y')=a(x,y)$须取$D'(x',y')=D(x'/a,y'/a)/a^2$，变量代换才给出$\int D'=\int D$。裁剪训练时先按中心规则确定哪些点计入，再为这些点在裁剪域构造归一化核；直接把原密度图裁一块，会把边外对象的核尾也计入。三个点各贡献1，整图密度之和为3；若一个边缘点的核有一半落到图外，直接截断后该点只贡献约0.5。

总和正确仍可能把本属左半边的人数移到右半边，应同时检查局部计数和空间误差。新的拍摄高度和密度范围可能超出训练条件；密度图也没有恢复每个人的身份或完整位置。两条路线在同一路口与广场图像上的取舍可以概括为：检测汇总在稀疏场景中给出可核验的清单，在密集场景中系统性漏数；密度积分在密集场景中更稳定，却只能核验区域总量，不能列出个体。

\subsection{人体姿态的估计}
\label{sec:vision-recognition-pose}

\term{姿态估计}（pose estimation）的输入是含人的图像，输出是每个人的关节位置。人体关节不同于图像中任意可匹配的角点：左腕、右肘具有固定的部位含义，多人场景还必须知道这些部位属于谁。输出可以是图像平面上的二维关节、相机坐标中的三维关节，也可以是完整的人体表面网格。难点随输出维数增加：二维要解决多人归属，三维要面对节首的深度歧义，网格还要确定关节之间的身体形状。

三个模型对应三级输出。OpenPose在整图上预测二维关节，并用部位亲和场把关节组装成不同的人；Martinez基线把二维关节作为输入，用一个简单网络回归三维关节，将二维定位误差与三维提升误差分开；SAM 3D Body则从图像直接恢复参数化人体网格，并接受关节点或掩码提示来纠正结果。

\subsubsection{OpenPose}
\label{sec:vision-recognition-openpose}

多人姿态的常见做法是自顶向下：先检测每个人，再在每个人的框内定位关节。它的计算量随人数线性增长，人与人重叠时检测框和裁剪质量又会直接影响关节结果。OpenPose采用自下而上路线：在整图上一次性找出所有关节候选，再判断哪些候选属于同一个人\scite{vision-1611-08050}。组装所依据的线索是\term{部位亲和场}（part affinity field, PAF）：对每一种肢段（例如前臂），在真实肢段经过的带状区域内写入从一端关节指向另一端关节的单位向量，其他位置为零。两个候选关节之间若真是同一人的前臂，连线经过的像素上的PAF方向就应与连线方向一致。

网络先用VGG-19的前10层把图像变成特征$\bm{F}$（空间步幅为8）。之后是多阶段的双分支卷积网络：第一阶段从$\bm{F}$同时预测$J$张关节置信图$\bm{S}^{(1)}$和$C$个肢段的PAF $\bm{L}^{(1)}$（每个肢段两通道，共$2C$通道）；第$t$阶段把$\bm{F}$、$\bm{S}^{(t-1)}$、$\bm{L}^{(t-1)}$拼接作为输入，重新预测两组输出。后续阶段因此能利用前一阶段的关节和肢段结果修正彼此，例如由清晰的肘部和前臂方向推断模糊的手腕。
各阶段的层配置见表\ref{tab:vision-recognition-openpose-stages}。第一阶段每个分支是3个$3\times3$卷积加2个$1\times1$卷积；从第二阶段起，每个分支换成5个$7\times7$卷积加2个$1\times1$卷积，更大的卷积核扩大了每层的感受范围，使后续阶段能在步幅8的特征上看到整条肢段两端的证据。原文共用6个阶段，每个阶段的两组输出都接受监督。测试时把$1080\times1920$的视频帧缩放到$368\times654$，在GTX 1080上九个人的画面中卷积网络约需99.6毫秒，后续的组装只需约0.58毫秒。原文的期刊扩展版把每个$7\times7$卷积换成3个连续的$3\times3$卷积并拼接其输出，先跑若干只预测PAF的阶段，再跑预测置信图的阶段；这里按会议版讲解。

\begin{table*}[tbp]
 \centering
 \small
 \begin{tabular}{@{}lllrrl@{}}
 \toprule
 阶段或层组 & 层类型 & 卷积核与步长 & 输出通道 & 重复次数 & 输出尺寸\\
 \midrule
 特征提取 & VGG-19前10层 & $3\times3$，含池化 & 论文未给出 & 1 & $H/8\times W/8$\\
 阶段1（每分支） & 卷积 & $3\times3$ & 论文未给出 & 3 & $H/8\times W/8$\\
                 & 卷积 & $1\times1$ & 末层为$J$或$2C$ & 2 & $H/8\times W/8$\\
 阶段$t\geq2$（每分支） & 卷积 & $7\times7$ & 论文未给出 & 5 & $H/8\times W/8$\\
                 & 卷积 & $1\times1$ & 末层为$J$或$2C$ & 2 & $H/8\times W/8$\\
 \bottomrule
 \end{tabular}
 \caption{OpenPose（会议版）的结构，取自原文图3与正文。阶段$t\geq2$的输入是$\bm F$、$\bm S^{(t-1)}$与$\bm L^{(t-1)}$的拼接，共6个阶段。两个分支结构相同、参数独立，分别输出$J$张置信图和$C$个肢段的PAF；中间层通道数和$J,C$的具体取值在原文正文中未给出。}
 \label{tab:vision-recognition-openpose-stages}
\end{table*}

训练从关键点标注构造两类目标。关节置信图在每个标注位置放一个高斯峰，同类多人的峰取逐点最大值；PAF在肢段带内写入单位方向，多人肢段重叠处取平均。每个阶段都对两组输出计算加权平方误差，权重$W(\bm{p})$在未标注的人身上为零，以免把漏标的人当作负例；每个阶段都接受监督，缓解深层网络的梯度消失。

推理时先在置信图上做非极大值抑制得到关节候选。对候选肘部$\bm{a}$和腕部$\bm{b}$，令方向$\bm{u}=(\bm{b}-\bm{a})/\lVert\bm{b}-\bm{a}\rVert_2$，沿连线对PAF积分：
\begin{equation}
 s(\bm{a},\bm{b})=
 \int_0^1
 \bm{L}((1-\lambda)\bm{a}+\lambda\bm{b})^\top
 \bm{u}\,\mathrm{d}\lambda\text{。}
 \label{eq:vision-paf-score}
\end{equation}
实现中在有限个点上采样求平均，$\bm{a}=\bm{b}$的退化候选应排除。对每一种肢段，所有肘部候选与腕部候选构成一个二分图，以$s$为边权求最大权匹配，保证一个关节在该肢段中只连一次；再按人体骨架树把共享关节的肢段拼成完整的人。假设五个采样点的方向内积为$(1,1,0.8,1,0.7)$，平均为0.9；连到另一人的手腕时变为$(0.1,0,-0.2,0,0.1)$，平均为0。两个手腕都定位正确，也仍需这一关联步骤才能把手臂接对。投影重合、严重遮挡和交叉肢段会使证据混杂，因此关节定位与人体组装应分开评价。

\subsubsection{Martinez基线}
\label{sec:vision-recognition-lifting}

从单幅图像估计三维姿态的端到端网络，误差可能来自两处：二维上没看准关节，或者看准了却无法推出深度。端到端训练把这两种误差混在一起，难以判断瓶颈。Martinez等人的基线把问题拆开：假定二维关节已经给出，只用一个简单的全连接网络把二维坐标“提升”为三维坐标\scite{vision-1705-03098}。实验结果表明这种简单提升已能取得有竞争力的误差，说明相当一部分困难在于二维检测，而不在三维推断本身。

输入是按固定部位顺序排列的$j$个二维关节，展平为$2j$维向量$\bm{x}_i$，并用训练集统计量做零均值单位方差规范化；输出是相机坐标中相对根关节（骨盆）的三维关节，展平为$3j$维向量$\bm{y}_i$。网络先用线性层把输入映射到1024维，接两个残差块，每块包含两组“线性层、批归一化、ReLU、随机失活”，残差连接为$\bm{h}'=\bm{h}+F_{\bm{\theta}}(\bm{h})$，最后一层线性层输出$3j$维。全网共6个线性层，约400万至500万个可训练参数；每层权重的范数被约束不超过1，随机失活比例为0.5，原文报告残差连接使误差下降约10\%。逐层结构见表\ref{tab:vision-recognition-martinez-layers}。单个残差块内的计算顺序是“线性层$\to$批归一化$\to$ReLU$\to$随机失活$\to$线性层$\to$批归一化$\to$ReLU$\to$随机失活$\to$与块输入相加”。网络很小，原文在Titan Xp上一次前向约2毫秒。

\begin{table*}[tbp]
 \centering
 \small
 \begin{tabular}{@{}lllrrl@{}}
 \toprule
 阶段或层组 & 层类型 & 卷积核与步长 & 输出通道 & 重复次数 & 输出尺寸\\
 \midrule
 输入 & 二维关节 & — & $2j$ & — & 向量\\
 输入层 & 线性 & — & 1024 & 1 & 向量\\
 残差块 & 线性（见正文） & — & 1024 & 2块\ensuremath{\times}2组 & 向量\\
 输出层 & 线性 & — & $3j$ & 1 & 向量\\
 \bottomrule
 \end{tabular}
 \caption{Martinez基线的结构，取自原文图1与正文。网络只有全连接层，“卷积核与步长”一列不适用；每个残差块含两组“线性、批归一化、ReLU、随机失活”。关节数$j$在原文中以符号出现。}
 \label{tab:vision-recognition-martinez-layers}
\end{table*}

训练以均方误差$(1/n)\sum_i\lVert g_{\bm{\theta}}(\bm{x}_i)-\bm{y}_i\rVert_2^2$学习（Adam，初始学习率0.001并指数衰减，批大小64，200轮），数据是带动作捕捉三维真值的室内数据集，二维输入可以是真值投影，也可以是二维检测器的输出。使用时必须复现训练时的规范化，再把输出逆规范化；关节顺序错位会把一个部位的输入当成另一个部位，维度却仍然匹配，因而不会报错。输出是相对根关节的坐标，不含人到相机的距离，要得到后者还需另估根平移。

节首的像素$(520,240)$可以对应$(1,0,4)$或$(2,0,8)$。网络用训练数据中的人体比例和姿态分布在这些解释中选一种，向前或向后弯曲的肢段也可能有相近的二维投影，选择同样来自先验。评价要说明是否允许旋转、平移和尺度对齐：相似变换对齐后的误差较小，不意味着原始输出已有正确朝向或距离。已知相机的多视图关节可以通过射线求交估计三维点，增加的是观测，而单目提升网络增加的只是先验。

\subsubsection{SAM 3D Body}
\label{sec:vision-recognition-sam3d-body}

两个三维骨架可以有相同的关节位置，却对应不同的肩宽、躯干厚度和手指形状；若要把人体轮廓投回图像比对，或驱动一个虚拟人，关节点还不够。\emph{人体网格恢复}（human mesh recovery, HMR）估计由顶点和面片组成的人体表面，同时解释姿态、体形与相机。已有方法在遮挡、罕见姿态和手部细节上容易失败，而且结果一旦出错，用户无法指出错误。SAM 3D Body的想法是把SAM式的提示接口引入人体网格恢复：模型可以直接预测，也可以接受二维关键点或掩码提示，按提示修正结果\scite{vision-sam3d-body-2026}。

输出采用参数化人体表示Momentum Human Rig（MHR），它将骨架结构、关节姿态和表面形状分开控制。记这些参数为$\bm{\beta}_{\mathrm{sk}},\bm{\theta}_{\mathrm{pose}},\bm{\beta}_{\mathrm{surf}}$，人体模型资产分别产生关节和顶点：
\[
 \bm{J}=\mathcal{J}(\bm{\beta}_{\mathrm{sk}},
                   \bm{\theta}_{\mathrm{pose}}),\qquad
 \bm{V}=\mathcal{V}(\bm{\beta}_{\mathrm{sk}},
                   \bm{\theta}_{\mathrm{pose}},
                   \bm{\beta}_{\mathrm{surf}})\text{，}
\]
其中$\bm{J}\in\mathbb{R}^{j\times3}$、$\bm{V}\in\mathbb{R}^{v\times3}$，$j,v$由模型资产规定。这两个式子只说明输出分工，不是把MHR的全部变形计算写成一个线性式。网络只需预测少量参数，输出就自动遵守人体连接结构；把关节或顶点经相机投影，才得到可与原图比对的位置和轮廓。

模型先取得人体裁剪及裁剪到原图的变换，把裁剪缩放到$512\times512$，由共享图像编码器得到空间特征；相机内参由现成的视场角估计器提供。原文给出两个规模：以ViT-H为编码器的版本约6.32亿参数，以DINOv3为编码器的版本约8.4亿参数。身体解码器的查询包含初始MHR与相机估计、辅助二维和三维关节查询，以及可选的提示。二维关键点提示写成$(x,y,\ell)$，$\ell$指出关节身份，位置编码与学习的关节嵌入使“左腕在这里”不同于一个无部位含义的点；掩码提示经卷积编码后加到图像特征上。具体地，解码器的输入词元依次是1个MHR与相机词元、$N$个关键点提示词元、$J_{2\mathrm D}$个二维和$J_{3\mathrm D}$个三维辅助关节词元，以及2个手部位置词元，宽度均为$D$；它们经交叉注意力读取图像特征后，第一个词元的输出经MLP回归姿态、形状、骨架和相机参数，再经MHR资产恢复网格。解码器的层数、宽度$D$以及MHR各参数向量的维数，论文均未给出。手部在整人裁剪中占的像素很少，模型另设可选的手部裁剪和独立手部解码器（共用编码器权重）；为避免直接替换手部参数连带改变肘部，原方法用手部解码得到的腕位置和身体解码得到的肘位置再次提示身体分支，细化后按人体连接树合并。

训练数据约700万幅图像，来自SA-1B、合成数据和多个手部与多视图数据集；三维依据来自多阶段标注：人工关键点、可微拟合、多视图几何和密集关键点检测以不同组合产生监督，另有合成和多视图数据。二维关节在裁剪坐标中监督；身体三维关节减去骨盆、手关节减去手腕后比较，只有已标注的手部参与对应目标。关节使用带学习不确定性权重的$L_1$项，姿态与形状参数使用$L_2$项，另加关节限位和手框误差。训练中随机进行多轮提示以模拟纠错，并提高用户给定关节的监督权重。

使用时人体框可由外部检测器给出，框误差属于完整流程的一部分。用户点先变到裁剪坐标，预测的关节和网格投影后再用裁剪逆变换回原图。教学中设左腕被遮住，无提示预测投回原图$(120,150)$，用户根据额外依据给点$(128,146)$，像素残差为$\sqrt{8^2+4^2}\approx8.94$。加入提示后应重跑解码、恢复网格并重新投影，检查这个残差以及肘腕连接、手形和轮廓是否共同改变；这里只定义需要检查的量，不假定提示后的误差必然下降。轮廓相符的网格仍可能把手腕放在不同深度，未知真实身高时也不能把模板的身体尺寸当成该人的测量结果。连续姿态与身份关联由\ref{chap:vision-video-recognition}章承接。

三个模型的对比见表\ref{tab:vision-recognition-pose-compare}。Martinez基线相对OpenPose改变的是输出维数，它把OpenPose这类二维结果当作输入，换来三维关节，付出的是对人体先验的依赖和单人输入的限制；SAM 3D Body再把输出从关节换成网格，并引入提示纠错，换来表面与手部，付出的是对人体模型资产、多阶段标注和裁剪流程的依赖。

\begin{table}[htbp]
 \centering
 \small
 \begin{tabularx}{\linewidth}{@{}lXXX@{}}
 \toprule
 比较项 & OpenPose & Martinez基线 & SAM 3D Body\\
 \midrule
 输入 & 整幅多人图像 & 单人二维关节 & 人体裁剪与可选提示\\
 输出 & 二维关节与人体归属 & 相对根关节的三维关节 & 参数化人体网格与相机\\
 关键机制 & 置信图加PAF组装 & 全连接残差回归 & 提示查询解码与MHR\\
 深度来源 & 不估计 & 人体姿态先验 & 图像与人体先验\\
 主要局限 & 交叉肢段与遮挡 & 依赖二维输入与对齐方式 & 尺度歧义、依赖模型资产\\
 \bottomrule
 \end{tabularx}
 \caption{三种姿态输出的比较。三者回答的不是同一个问题，比较时应固定关节集合、坐标与对齐方式。}
 \label{tab:vision-recognition-pose-compare}
\end{table}

\subsection{场景深度的估计}
\label{sec:vision-recognition-depth}

\term{深度估计}（depth estimation）的输入是一幅或多幅图像，输出是参考图每个像素的光轴深度$z$。节首的式\eqref{eq:vision-camera-projection}已经说明，单个像素只确定一条射线，射线上的$z$需要额外信息。信息来源决定了路线：单目路线只有一幅图，靠物体形状、透视、遮挡关系和训练中学到的先验推断深度，通常只能可靠地给出相对远近；多视图路线有多幅已知相机的图像，同一表面在不同图中的位置提供跨射线约束，在相机基线已知时可以得到度量深度。Depth Anything V2和MVSNet分别代表这两条路线；两者的输出都要说明是深度、逆深度还是相对量，以及是否已经确定度量尺度。

多视图路线有一段很长的经典前史，MVSNet正是在它上面做的改动，这里先回顾要点，完整推导见\ref{sec:vision-formation-multi-view}节。最简单的情形是平行双目：两台内参相同的相机沿水平方向相隔基线$b$，同一点在左右图中只在横坐标上相差视差$\Delta u$，由相似三角形得到$z=fb/\Delta u$，其中$f$是像素焦距。取$f=800$像素、$b=0.2$米、$\Delta u=40$像素，得$z=4$米；视差误差1像素约引起0.1米的深度误差，而深度越远视差越小，同样的视差误差带来的深度误差越大。经典立体匹配因此把问题化为沿极线搜索对应：对左图每个像素，在右图同一行的若干候选视差上计算匹配代价（例如窗口内灰度差的平方和或归一化互相关），得到一个“像素$\times$视差”的代价体，再加上相邻像素视差应当平滑的约束，选出代价最小的视差。多于两台相机、而且不平行时，平面扫描法把候选视差换成参考相机中的一组深度平面，每个深度平面通过单应把其他图像变换到参考视角，再比较外观是否一致。这条路线在纹理丰富处精确，但手工的匹配代价在弱纹理、反光和遮挡处不可靠，平滑约束又会抹掉细结构。MVSNet保留了平面扫描与代价体这两个经典结构，把手工代价换成学习的特征和三维卷积。

\subsubsection{Depth Anything V2}
\label{sec:vision-recognition-depth-anything}

单目深度模型需要大量带深度标签的训练图，真实数据的深度却来自深度传感器或多视图重建，往往有空洞、噪声，物体边缘尤其粗糙；用这些标签训练，模型会学到模糊的边界。合成图像有精确深度，但外观与真实照片存在差距。Depth Anything V2的想法是把两种数据的长处分开使用：只用合成图像训练一个大教师，得到精细的几何；再让教师给大量无标注真实图像产生伪标签，学生模型从这些伪标签学习真实外观\scite{vision-2406-09414}。伪标签学习的一般依据可参考\BookChapterRef{chap:data-derived-knowledge-acquisition}{卷八“基于数据本身的知识获取”章}。

网络由DINOv2编码器和密集预测变换器（dense prediction transformer, DPT）解码器组成。DINOv2的图像块边长为14\scite{vision-2304-07193}，以输入$518\times518$为例，编码器输出$37\times37$个词元；DPT从编码器的四个不同层取出词元，各自重组为二维特征图并缩放到不同分辨率，再自深而浅逐级融合、上采样，最后输出$H\times W\times1$的逐像素相对逆深度。DPT把多层空间特征融合成逐位置数值，而不是把全图汇总向量当作深度图。原文按DINOv2的四个规模训练了四个学生模型，配置见表\ref{tab:vision-recognition-dav2-config}。

\begin{table*}[tbp]
 \centering
 \small
 \begin{tabular}{@{}lccccccl@{}}
 \toprule
 模型 & 图像块大小 & 词元数 & 层数 & 宽度 & 头数 & MLP倍率 & 参数量\\
 \midrule
 Depth Anything V2-S（ViT-S） & 14 & 1369 & 12 & 384 & 6 & 未列出 & 约25M\\
 Depth Anything V2-B（ViT-B） & 14 & 1369 & 见注 & 768 & 12 & 未列出 & 论文未给出\\
 Depth Anything V2-L（ViT-L） & 14 & 1369 & 24 & 1024 & 16 & 未列出 & 约335M\\
 Depth Anything V2-G（ViT-g） & 14 & 1369 & 40 & 1536 & 24 & 未列出 & 约1.3B\\
 \bottomrule
 \end{tabular}
 \caption{Depth Anything V2的配置。参数量取自Depth Anything V2原文（含DPT解码器）；编码器的块数、宽度与头数取自DINOv2原文表17，其中ViT-B的块数在该表中印为18，与ViT-B通常的12层不同，此处不作推断。DINOv2原文只说明ViT-S/B/L用普通MLP、ViT-g用SwiGLU作前馈网络，未列出前馈宽度。词元数按$518\times518$输入、块边长14计为$37\times37$。}
 \label{tab:vision-recognition-dav2-config}
\end{table*}

训练分三步。教师（最大规模的编码器）只用约59.5万张合成RGB和精确深度训练，批大小64、16万次迭代；教师为大量真实无标注图像预测伪深度；不同规模的学生只用这些约6200万张真实图与伪标签对训练，批大小192、48万次迭代，编码器与解码器的学习率分别为$5\times10^{-6}$和$5\times10^{-5}$。基础检查点预测尺度、平移不变的相对逆深度：比较预测与标签之前，先消去两者之间的整体尺度和偏移。设预测为$\hat q_p$、目标为$q_p$，一个教学用的平方对齐目标为
\begin{equation}
 \min_{a,b}\frac1{|\Omega|}
       \sum_{p\in\Omega}(a\hat q_p+b-q_p)^2\text{，}
 \label{eq:vision-relative-depth-alignment}
\end{equation}
其中$\Omega$是有效监督位置，$a,b$是比较时的对齐量。令导数为零，若预测有非零方差，得$a=\operatorname{Cov}(\hat q,q)/\operatorname{Var}(\hat q)$、$b=\bar q-a\bar{\hat q}$；全图预测为常数时尺度无从估计。这个平方例子解释对齐消去了什么信息，原论文的尺度平移不变损失在具体形式上与此不同。训练还加入梯度匹配项，比较相邻位置残差的差，防止整体对齐合适而边界过于平滑。尺度平移不变项与梯度匹配项的权重比为1比2。对伪标签样本，原工作忽略损失最大的10\%区域以降低噪声影响。例如预测$(1,2,3)$与目标$(3,5,7)$在$a=2,b=1$下完全重合；对齐误差为零并不说明原预测的数值已经是逆米单位。

使用相对检查点时，按模型尺寸和归一化处理图像，解码出$\hat q$，插值回原图有效区域，并保留其相对逆深度身份。若三个位置的相对逆深度为$(0.8,0.4,0.2)$，可以判断第一个最近，却不能未经校准说深度为$(1.25,2.5,5)$米；颜色图只是显示映射，也不能从色值读米。度量适配用带单位的深度标签微调已有网络，使输出在该数据、最大深度和相机条件下对应实际距离；即便如此，换相机或场景仍可能带来尺度误差。

\subsubsection{MVSNet}
\label{sec:vision-recognition-mvsnet}

单目深度依赖先验，遇到训练中少见的场景结构就没有依据。如果有多张重叠照片且相机已知，同一表面在不同图中的位置可以直接约束深度。传统多视图立体方法用手工设计的相似度比较图像块，在弱纹理和反光区域表现差，也难以融合多张图的证据。MVSNet的想法是在学习的特征上做这件事：对参考图的每个像素枚举一组候选深度，把各源图的特征按候选深度变换到参考视角，用方差度量多视图是否一致，形成代价体，再用三维卷积从中推断深度\scite{vision-1804-02505}。

每个候选深度$d_k$对应参考相机中$z=d_k$的一个平面。参考像素$\tilde{\bm p}_r$在该深度下为$d_k\bm K_r^{-1}\tilde{\bm p}_r$，进入源相机得到$\bm R_{s,r}d_k\bm K_r^{-1}\tilde{\bm p}_r+\bm t_{s,r}$。由于$\bm e_3^\top\bm K_r^{-1}\tilde{\bm p}_r=1$，其中$\bm e_3=(0,0,1)^\top$，除去齐次比例$d_k$得到
\begin{equation}
 \begin{aligned}
 \tilde{\bm p}_s&\sim\bm H_{s,r}(d_k)\tilde{\bm p}_r,\\
 \bm H_{s,r}(d_k)=
 \bm K_s\left(\bm R_{s,r}
 +\frac{\bm t_{s,r}\bm e_3^\top}{d_k}\right)\bm K_r^{-1}\text{。}
 \end{aligned}
 \label{eq:vision-mvsnet-homography-process}
\end{equation}
这个映射对整张平面成立，因而可以用一个$3\times3$矩阵变换整幅源特征图，并对$d_k$可微。特征图尺度不同于图像时，先按步幅调整内参，再用双线性插值读取源特征；越界或落到相机后方的位置不得当成有效观测。

网络结构如下。$n_{\mathrm v}$幅图（1幅参考、其余为源图）经共享的8层二维卷积网络，各自得到$H/4\times W/4\times32$的特征图。对$D$个候选深度，把每幅源特征图按式\eqref{eq:vision-mvsnet-homography-process}变换到参考视角，得到$n_{\mathrm v}$个$D\times H/4\times W/4\times32$的特征体。将参考特征和变换后的各源特征记为$\bm{f}_r$，逐元素计算方差：
\begin{equation}
 \begin{aligned}
 \bar{\bm{f}}&=\frac1{n_{\mathrm{v}}}\sum_{r=1}^{n_{\mathrm{v}}}\bm{f}_r,\\
 \bm{c}(d_k,u,v)&=\frac1{n_{\mathrm{v}}}\sum_{r=1}^{n_{\mathrm{v}}}
   (\bm{f}_r-\bar{\bm{f}})\odot(\bm{f}_r-\bar{\bm{f}})\text{，}
 \end{aligned}
 \label{eq:vision-mvs-cost}
\end{equation}
其中$\odot$表示逐元素相乘。方差对视图数没有限制，任意多幅源图都可以合成一个\term{代价体}（cost volume），即形状为$D\times H/4\times W/4\times32$、记录每个位置在每个候选深度下多视图有多不一致的张量。三维U形卷积网络在空间和深度三个方向上整理代价体，输出单通道的$D\times H/4\times W/4$得分，沿深度轴做softmax得到概率$\pi_k$，以期望$\hat d=\sum_k\pi_kd_k$回归深度。这种期望回归对概率可微，结果也不必等于某个离散候选。最后一个浅层细化网络读取参考图像和初步深度，预测残差以修正边界。各段的结构与形状见表\ref{tab:vision-recognition-mvsnet-structure}。

\begin{table*}[tbp]
 \centering
 \small
 \begin{tabularx}{\linewidth}{@{}llXl@{}}
 \toprule
 部分 & 输入形状 & 结构与层数 & 输出形状\\
 \midrule
 特征提取（每图） & $512\times640\times3$ & 8层二维卷积，第3、6层步长为2，每个尺度2层；除末层外均接批归一化与ReLU；各视图共享权重 & $128\times160\times32$\\
 单应变换与代价 & $n_{\mathrm v}$张特征图 & 对$D$个深度平面按式\eqref{eq:vision-mvsnet-homography-process}变换，按式\eqref{eq:vision-mvs-cost}求方差 & $D\times128\times160\times32$\\
 代价整理 & 代价体 & 四个尺度的三维U形网络，第一层三维卷积把通道从32降到8，每个尺度2个卷积层，末层输出1通道 & $D\times128\times160$\\
 深度回归 & 得分体 & 沿深度轴softmax，求期望深度 & $128\times160$\\
 细化 & 初步深度与缩放的参考图（4通道） & 3个32通道二维卷积加1个单通道卷积，输出残差 & $128\times160$\\
 \bottomrule
 \end{tabularx}
 \caption{MVSNet的分段结构。层数与通道取自原文正文；形状按原文DTU训练设置推算：输入裁剪为$640\times512$，$n_{\mathrm v}=3$，深度假设从425毫米到935毫米、间隔2毫米，共$D=256$个。特征网络各中间层的通道数在正文中未给出。}
 \label{tab:vision-recognition-mvsnet-structure}
\end{table*}

同一像素的教学标量特征在正确假设下为$(1,1.1,0.9)$，方差约0.0067；在错误假设下读到$(1,2,3)$，方差为$2/3$。若候选深度为2、4、6米，概率为0.1、0.7、0.2，回归得到4.2米。双峰概率会把结果放在两个表面之间，因而概率的集中程度也应作为置信度检查。

训练数据包含参考图、若干重叠源图、已知相机和有效参考深度（原文使用DTU室内物体数据集，约2.7万个训练样本，测试时用5幅图、分辨率$1600\times1184$）。在有效像素上以$L_1$损失比较初步深度和细化深度与真值，两项相加，更新二维特征网络、三维代价整理网络和细化网络。使用时相机由当前图集提供；得到深度后按概率和多视图几何一致性过滤，再反投影合并点云。代价体的存储为$O(Dh'w'd)$，随候选数和分辨率增长，这是高分辨率重建的主要瓶颈。与本小节开头的平行双目一样，已知的相机基线正是MVSNet输出度量尺度的来源。

\begin{figure}[htbp]
 \centering
 % 针孔投影的几何关系；示意深度候选，不按物理长度绘制。
\begin{tikzpicture}[font=\figurefont\small,x=1cm,y=1cm,
 ray/.style={draw=BookBlue,line width=0.9pt},
 alt/.style={draw=BookAmber,dashed,line width=0.9pt}]
 \coordinate (a) at (0,0);
 \coordinate (b) at (7,0);
 \coordinate (p) at (2,3);
 \coordinate (q) at (3,4.5);
 \draw[ray] (a)--(q);
 \draw[ray] (b)--(p);
 \draw[alt] (b)--(q);
 \draw[draw=BookInk,line width=1.2pt] (-0.3,1)--(2.2,1);
 \draw[draw=BookInk,line width=1.2pt] (4.5,1)--(7.5,1);
 \fill[BookInk] (a) circle(1.1mm);
 \fill[BookInk] (b) circle(1.1mm);
 \fill[BookTeal] (p) circle(1.1mm);
 \draw[BookAmber,fill=white,line width=0.9pt] (q) circle(1.1mm);
 \fill[BookBlue] (0.667,1) circle(1mm);
 \fill[BookBlue] (5.333,1) circle(1mm);
 \draw[BookAmber,fill=white,line width=0.9pt] (6.111,1) circle(1mm);
 \node[below] at (a) {参考相机};
 \node[below] at (b) {源相机};
 \node[above right] at (0.667,1) {同一参考像素};
 \node[below] at (5.333,0.9) {$u'_1$};
 \node[below] at (6.111,0.9) {$u'_2$};
 \node[left] at (p) {深度$d_1$};
 \node[right] at (q) {深度$d_2$};
 \node[right] at (7.5,1) {源图平面};
 \draw[draw=BookMuted,<->,line width=0.7pt] (0,-0.8)--(7,-0.8)
   node[midway,below] {已知相机相对位置};
\end{tikzpicture}

 \caption{已知相机下，参考像素沿射线形成不同深度假设，分别投到源图的不同位置。图为针孔几何教学示意，只表示对应关系，不按物理长度绘制。跨图比较可排除不一致的假设，但遮挡或重复纹理会使判断不唯一。}
 \label{fig:vision-depth-hypotheses}
\end{figure}

图\ref{fig:vision-depth-hypotheses}展示两个深度假设如何查询不同的源图位置。弱纹理使多个候选代价相近，遮挡使源图看到另一表面，反光破坏外观一致性，相机误差使正确深度也无法对齐。两条路线在同一街景上的比较可以概括为：Depth Anything V2只要一幅图，覆盖所有像素，边界细致，但默认输出相对量，尺度来自先验或额外适配；MVSNet需要多幅图和已知相机，在相机基线已知时给出度量深度，但在弱纹理、遮挡区域可能无有效结果，代价体的存储也限制分辨率。

\subsection{相机与场景几何的估计}
\label{sec:vision-recognition-geometry}

MVSNet把相机当作已知条件。若手头只有一组普通照片，相机未知，就要同时估计每张图的内外参和场景的三维结构。输入是一组RGB图像，输出是每图的相机、深度或点云，以及它们所在的公共坐标系。难点在于相机和结构相互依赖：知道相机才能三角化点，知道点才能求相机。经典的回答是\ref{sec:vision-formation-multi-view}节的运动恢复结构，这里回顾它的步骤和失效条件。第一步在每幅图上检测SIFT等局部特征，在图对之间匹配，并用RANSAC估计基础矩阵或本质矩阵来剔除错误匹配；第二步从一对重叠充分的图开始，分解本质矩阵得到相对位姿，三角化出初始三维点；第三步逐幅加入新图，用已有三维点与新图特征的对应求新相机的位姿，再三角化新点。每加入若干幅图，就做一次光束法平差：对相机参数$\{\bm K_i,\bm R_i,\bm t_i\}$和三维点$\{\bm X_j\}$联合最小化重投影误差$\sum_{(i,j)}\lVert\bm p_{i,j}-\pi(\bm K_i,\bm R_i,\bm t_i,\bm X_j)\rVert_2^2$，其中$\bm p_{i,j}$是第$j$个点在第$i$幅图中的观测，$\pi$按式\eqref{eq:vision-camera-projection}投影。这个流程的每一步都有明确的几何意义，在纹理丰富、重叠充分时可以达到很高的精度；但它是一条串联的流水线，任何一步失败都会向后传递。弱纹理使特征点太少，重复图案使匹配出错，视图太少或基线太短使三角化病态，动态物体又违反静态场景的假设。

本小节的五个模型都用前馈网络直接预测几何，差别在于一次处理多少视图、信息怎样到达。DUSt3R一次只看两张图，直接输出两张图在同一坐标系中的点图，多张图则靠逐对预测加全局优化拼起来。VGGT让任意多张图在一次前向中相互注意，直接输出所有相机和深度，省去了逐对再优化，但注意力成本随视图数平方增长。VGG-T$^3$针对这一成本，把全局注意力中随视图增长的键值压缩到一个固定大小的网络里。InfiniteVGGT面对另一种条件：视图逐帧到达，不能等全集，它为流式处理维护一个有限的缓存。FF3R则在几何之外同时恢复可查询的语义，使三维结构能回答“椅子在哪里”。

\subsubsection{DUSt3R}
\label{sec:vision-recognition-dust3r}

上面的经典流程中，特征点匹配不足，相机就估不出来，后续一切都无从谈起。DUSt3R的想法是跳过显式的相机估计，让网络对一对图像直接回归\term{点图}（point map）：每个有效像素对应一个三维点，两张点图都表达在第一幅图的相机坐标中\scite{vision-2312-14132}。点图同时隐含了深度、两图的相对位姿和像素对应，相机可以事后从点图中恢复。

两幅图各自切块后经共享的ViT编码器，得到两组词元。两个Transformer解码分支每层先做自注意力，再对另一分支的词元做交叉注意力，从而反复交换跨图信息。每个分支最后由DPT式头输出$H\times W\times3$的点图和$H\times W\times1$的置信度；第二张图的点图同样表达在第一相机坐标中，这一约定迫使网络在内部推断两图的相对位姿。编码器和解码器从CroCo跨视图补全预训练初始化。原文的编码器为ViT-Large规模，解码器为ViT-Base规模，头部最终采用DPT；块边长、解码块数和参数总量论文未给出。训练分三段：先在$224\times224$上用线性头训练，再改用最长边512像素、长宽比随批次随机变化的图像，最后换成DPT头。训练数据约850万对图像，来自室内、室外、物体和驾驶等八个数据集。一对图像的前向在H100上约需40毫秒。

训练记录含同一场景的图像对，以及由深度、内参和外参得到的真实点图和有效位置；使用时不输入相机，并不意味着训练不曾使用几何标签。原工作分别用有效点到原点的平均距离规范化预测和目标，得到尺度无关的点误差$e_i$。置信加权目标为$\sum_i(c_i e_i-\alpha\log c_i)$，$\alpha>0$，置信度以$c_i=1+\exp h_i>1$参数化。对固定$e_i>0$，该项对$c_i$的导数为$e_i-\alpha/c_i$，误差大的点倾向于较小的置信，而对数项使网络不能把所有置信都压到最小。以$\alpha=1$为例，$e=0.1$时驻点为$c=10$；$e=2$时驻点0.5不在允许范围，便取下边界。这是可靠性权重，不是校准过的正确率。天空等无效区域按标签排除。

多于两张图时，每一对都以自己的首图为参考，不能直接拼接。DUSt3R把有足够重叠的图对连成图，边$e$保存两张点图$\bm X^{v,e}$及置信；对每幅视图求共同点图$\bm\chi^v$，对每条边求尺度$s_e>0$与刚体变换$(\bm R_e,\bm t_e)$，使$\sum_{e,v,p}c_p^{v,e}\lVert\bm\chi_p^v-s_e(\bm R_e\bm X_p^{v,e}+\bm t_e)\rVert_2$最小，并固定尺度规范以防所有点缩到零。这是对场景变量的迭代优化，网络权重不变。相机内参可以从点图拟合：例如假设主点在中心，由$(u-c_x,v-c_y)\approx f(x/z,y/z)$对高置信点拟合$f$。若甲乙图预测建筑宽度为2个单位，乙丙图为5个单位，对齐要先估计约2.5的相对尺度；没有额外度量依据时，共同场景只有约定尺度。逐对预测、全局对齐两个阶段使处理大量图像时计算重复，弱重叠和动态物体也会干扰对齐。

\subsubsection{VGGT}
\label{sec:vision-recognition-vggt}

DUSt3R一次只看两张图，$n_{\mathrm v}$张图需要处理许多对，再做一次全局优化；每一对看不到其他视图的证据，三张图共同才能消除的歧义在逐对阶段无法利用。VGGT的想法是让所有视图在一次前向中共同参与计算，直接输出每张图的相机、深度、点图和跟踪特征，不再依赖后续优化\scite{vision-2503-11651}。

每张图经DINOv2初始化的编码器切成带位置的图像块词元，并附一个相机词元和若干辅助寄存词元；首图使用与其他图不同的学习词元，标记它为坐标参照。主干由交替注意力组成：帧内层只在同一张图的词元之间注意，保留各图的结构；全局层在所有视图的全部词元之间注意，交换跨视角证据。两种层交替堆叠，相机词元因此始终与所属图像关联。最后，相机头读取相机词元，输出九维编码（四维旋转四元数、三维平移和横纵视场角）；DPT式稠密头把多层图像特征恢复为空间网格，输出深度、点图、置信和跟踪特征。视场角$\varphi_x,\varphi_y$在中心主点条件下给出$f_x=w/[2\tan(\varphi_x/2)]$、$f_y=h/[2\tan(\varphi_y/2)]$，因此得到的是当前网络输入网格上的内参。深度$D_i$是第$i$相机自身的光轴深度，点图和外参则用首相机作为公共参照，二者可以互相换算：
\begin{equation}
 \bm{P}_{1,i}(u,v)=
 \bm R_i^\top\left[
 D_i(u,v)\bm K_i^{-1}(u,v,1)^\top-\bm t_i\right]\text{。}
 \label{eq:vision-vggt-unprojection-process}
\end{equation}
这里$\bm R_i,\bm t_i$沿用世界到第$i$相机的约定，公共世界就是首相机。

表\ref{tab:vision-recognition-ff-geometry-config}给出VGGT的配置。主干共24个注意力块，每块含一个帧内自注意力层和一个全局自注意力层，即帧内与全局各24层；宽度1024、16个头，沿用DINOv2中ViT-L的设置，并使用QK归一化与LayerScale，全网没有交叉注意力。单个块的计算可以按ViT的块写法表示为两次Transformer更新（原文未说明层归一化的具体位置，这里按ViT的前置写法）：先把$n_{\mathrm v}$幅图的词元按图分组，对每组做$\bm z'=\bm z+\operatorname{MSA}(\operatorname{LN}(\bm z))$、$\bm z''=\bm z'+\operatorname{MLP}(\operatorname{LN}(\bm z'))$；再把全部图的词元拼成一个序列，做同样的两步。帧内层与全局层的差别只在注意力的范围。相机头在相机词元上再加4个自注意力层和一个线性层；DPT头取第4、11、17、23块的词元。每帧除图像块词元外还有1个相机词元和4个寄存词元，寄存词元的输出被丢弃。

\begin{table*}[tbp]
 \centering
 \small
 \begin{tabular}{@{}lccccccl@{}}
 \toprule
 模型 & 图像块大小 & 每图词元数 & 层数 & 宽度 & 头数 & MLP倍率 & 参数量\\
 \midrule
 DUSt3R编码器 & 论文未给出 & — & ViT-Large & — & — & — & 论文未给出\\
 DUSt3R解码器 & — & — & ViT-Base（两分支） & — & — & — & 论文未给出\\
 VGGT主干 & 14 & $37^2+5$ & 帧内24＋全局24 & 1024 & 16 & 未列出 & 约12亿（全模型）\\
 VGG-T$^3$ & 14 & 同VGGT & 同VGGT & 1024 & — & 快权重4 & 论文未给出\\
 \bottomrule
 \end{tabular}
 \caption{前馈几何模型的配置。DUSt3R原文只给出编码器与解码器的规模名称；VGGT的“每图词元数”按最长边518、正方形输入计为$37\times37$个图像块词元加1个相机词元与4个寄存词元。VGG-T$^3$沿用VGGT的层数，只把每个全局注意力层换成测试时训练模块，故不再有注意力头；其快权重是输入输出1024维、中间层4倍宽的SwiGLU网络。}
 \label{tab:vision-recognition-ff-geometry-config}
\end{table*}

训练记录是同一场景的多张RGB、每图相机、有效深度和公共坐标点图，跟踪监督还包含查询点的跨图位置及可见性；这些几何来自传感器、合成引擎或运动恢复结构。训练先把全部量变到首相机坐标，再用有效点到原点的平均距离除相机平移、点坐标和深度。与DUSt3R在比较时才归一化预测不同，VGGT要求网络直接输出这套规范化的结果。相机编码用Huber误差监督（小残差为平方、大残差为线性），深度和点图用置信加权的位置误差与空间梯度误差监督，跟踪用位置误差和二元可见性误差监督。主干与任务头联合训练，每个场景随机采样2至24幅图，跟踪损失的权重为0.05；优化器为AdamW，共16万次迭代，在64块A100上训练约9天。

使用时，一次前向得到全部相机和稠密量。教学中首相机为单位旋转、零平移，主点$(320,240)$、焦距800、像素$(520,240)$与深度4个规范单位恢复点$(1,0,4)$，仍不能将4个单位写成4米。可用式\eqref{eq:vision-vggt-unprojection-process}从深度构造公共点，也可直接读取点图；两种输出不应未经核验混拼。跟踪接口把原图查询点变到处理网格，在查询图特征中插值取其描述，再与各图跟踪特征相关，估计跨图位置和可见性。全局注意力比较$(n_{\mathrm v}p)^2$个词元对（$p$为每图词元数），图像数增多时计算和显存都迅速增长，这是下一个模型要解决的问题。

\subsubsection{\texorpdfstring{VGG-T$^3$}{VGG-T3}}
\label{sec:vision-recognition-vgg-t3}

建筑的照片从三张增到数百张时，VGGT的全局注意力成为瓶颈：$n_{\mathrm v}$幅图、每图$p$个词元，比较次数为$(n_{\mathrm v}p)^2$，键和值本身也占$O(n_{\mathrm v}pd)$的空间。高效注意力实现可以不保存完整得分矩阵，却不能减少比较次数。VGG-T$^3$的想法是在使用时把每个全局层中随视图增长的键值集合压缩到一个固定大小的多层感知机（multi-layer perceptron, MLP）中：这个MLP把键映射为对应的值，其参数就保存了当前场景；之后查询只需送入MLP读取，计算量随词元数线性增长\scite{vision-vgg-t3-2026}。

\term{测试时训练}（test-time training, TTT）指在处理当前输入时临时更新一小组参数。对一个全局层，设当前图集产生$N=n_{\mathrm v}p$个键$\bm{k}_i$、查询$\bm{q}_i$和目标值$\bm{v}'_i$，维数均为$d$；场景权重为$\bm{\phi}$，MLP记为$T_{\bm{\phi}}$。内层目标和读取为
\begin{equation}
 \begin{aligned}
 \bm{\phi}_*&\approx\argmin_{\bm{\phi}}
   \sum_{i=1}^{N}
   \ell\bigl(T_{\bm{\phi}}(\bm{k}_i),\bm{v}'_i\bigr),\\
 \bm{o}_i&=T_{\bm{\phi}_*}(\bm{q}_i)\text{。}
 \end{aligned}
 \label{eq:vision-vgg-t3-fast-weights}
\end{equation}
其中$\ell$是逐词元的拟合损失，衡量MLP从键读出的结果与目标值是否一致；近似号表示实际只做规定次数的更新。原文实现沿用一种已有的大块测试时训练设计，把$\ell$取为$T_{\bm\phi}(\bm k_i)$与$\bm v'_i$的内积形式；下面的教学例为了便于手算改用平方误差，两者都是对各词元求和的损失，这一点是后文分批计算的依据。目标值$\bm{v}'_i$中图像块部分先在各图网格上经过短二维卷积混合邻域，相机和寄存词元不经过这项卷积；邻域混合使拟合目标带有局部上下文，避免只是复现同一词元的两个线性投影。查询和键还做$L_2$归一化。

模型以VGGT权重初始化，保留图像编码、帧内层和预测头，把全局注意力层换成上述测试时训练模块，并适配其投影和新增参数。每层先从训练得到的初始化取得场景权重，按键值残差计算梯度，更新后把查询送入MLP，所得词元进入下一层；全部层完成后由相机、深度和点图头解码。快权重MLP是输入输出均为1024维、中间层宽度为4倍的SwiGLU网络；图像块的值在各图网格上经一个$3\times3$的二维卷积混合邻域。原文内层采用Muon优化器（5次Newton–Schulz迭代、内层学习率0.1），训练时内层更新一步，测试默认更新两步，更多步数没有进一步改善。外层只训练全局层的查询、键、值和输出投影以及新增参数，VGGT原有的其他参数全部冻结；在8块A100上训练10万步，原文估计约为VGGT训练成本的12\%。外层训练沿用VGGT的几何监督，使初始化和投影学会“适合快速拟合”的表示。

求和形式的内层目标允许把全集分批拟合。固定一次更新开始时的$\bm{\phi}$，把词元划成批次$B_s$，则$\nabla_{\bm{\phi}}L=\sum_s\sum_{i\in B_s}\nabla_{\bm{\phi}}\ell(T_{\bm{\phi}}(\bm{k}_i),\bm{v}'_i)$。这一等式来自有限和的微分线性性，要求各批在同一组权重上计算梯度并累加后再更新；逐批改权重就不是同一次全集梯度。未处理的批可以放在主机内存或分到多个设备，降低单卡显存。用标量线性模型$T_\phi(k)=\phi k$和平均平方误差作教学简化，两个键值样本为$(1,2)$、$(2,4)$；在$\phi=0$处梯度为$-10$，步长0.1更新到$\phi=1$，再更新一次到1.5，查询1.5读出2.25，尚未达到精确映射给出的3。有限次拟合与精确保存全部键值之间的差距，就是固定容量带来的近似误差。原文在其基准中指出相机预测较不精确，并采用直接点图头评估点图；原图、编码特征和输出点图的存储仍随视图数增长。在A100上，原文摘要报告重建1000幅图约需54秒，比使用全局注意力的基线快11.6倍。

\subsubsection{InfiniteVGGT}
\label{sec:vision-recognition-infinitevggt}

前两个模型都假定完整图集已经到齐。沿走廊行走时，第$t$幅图已经到达，未来的画面尚未发生，当前输出不能借未来视图纠正自身。\emph{因果几何预测}（causal geometry prediction）限制当前帧只读取自身与已到达的历史。直接保留每层所有历史帧的键值可以复用旧特征，但缓存随$t$无限增长；只保留最近几帧，又会删去初始参照与回访线索。InfiniteVGGT的想法是维护一个固定预算的滚动缓存，按键的多样性决定保留哪些历史，并始终保留首帧作为锚点\scite{vision-infinitevggt-2026}。

InfiniteVGGT建立在已有的因果几何模型StreamVGGT上，不重新训练权重，因此结构与VGGT式模型相同（表\ref{tab:vision-recognition-ff-geometry-config}），只是全局注意力改成只读过去的因果注意力：图像仍经帧内和因果跨帧注意力，输出相机、深度、点图和跟踪特征；变化只在于哪些历史键值可读。各层各头把第一帧的键值固定为锚点，因为几何输出以首相机为参照；其余历史与新到词元构成可裁减的候选集。相邻画面含大量重复内容，原方法用归一化键的方向差异选择保留内容，从而不必计算完整注意力矩阵。对一个层和头，令候选键的单位向量为$\hat{\bm{k}}_i$，平均向量为$\bm{\mu}=m^{-1}\sum_i\hat{\bm{k}}_i$，多样性评分为
\begin{equation}
 s_i=-\frac{\hat{\bm{k}}_i^\top\bm{\mu}}
                 {\lVert\bm{\mu}\rVert_2}\text{。}
 \label{eq:vision-infinitevggt-diversity}
\end{equation}
与平均方向差别大的键得分高，与其值一起保留。这个评分近似识别冗余，不直接度量三维信息量；平均向量为零时需要数值退化处理。不同层的多样性不同，原文把层平均评分$\bar s_\ell$经softmax变成预算比例$b_\ell=\exp(\bar s_\ell/\tau)/\sum_j\exp(\bar s_j/\tau)$，按总预算$B$分配各层保留数，再按分数选取键值，最后加回首帧锚点。每到一帧，执行一次追加、评分、裁减、保留；整个过程不更新网络权重。每帧约新增1000个词元，不加裁减时100帧内就累积约$10^5$个。原文消融中每层每头的初始预算取1万、2.5万和5万个词元，在2.5万附近效果趋于饱和；温度$\tau$的取值论文未给出。在单块A100上，用余弦相似度评分时每帧约0.17秒、峰值显存约14.5GB，改用注意力权重评分则每帧多约120毫秒。

教学中令候选键方向为$(1,0)$、$(1,0)$、$(0,1)$，平均为$(2/3,1/3)$；前两者与平均方向的余弦为$2/\sqrt5\approx0.894$，第三者为$1/\sqrt5\approx0.447$，第三者评分更高。预算只容两项时，它和一个重复方向留下。这说明裁减能保留不同方向，却没有证明被删的方向在之后回访时不重要。固定缓存控制的是工作状态，每帧输出的点图若全部写盘，总存储仍随时间增长。走廊返回起点时，应比较短程相机误差、累计漂移和回访位置重合；保留首帧有助于维持参照，但不是显式的回环检测与全图校正。原文为此构建的Long3D基准有5个序列，长度从约2100帧到约9500帧；这些有限序列可以提供长程测试，不能证明无限时长下的稳定。

\subsubsection{FF3R}
\label{sec:vision-recognition-ff3r}

房间两幅照片都标出“椅子”，还不能知道它们是否指向同一张椅子的同一表面。前四个模型只恢复几何，要回答“椅子在哪里”，常见做法是先重建几何，再把二维语义分割结果投到三维；但几何错位会把椅面语义放到墙上，逐场景优化又很慢。FF3R的想法是在一次前馈中同时预测几何、外观和开放词汇语义，并用跨视角的语义一致性反过来约束几何\scite{vision-ff3r-2026}。

输入是至少两幅多视图图像。几何分支沿用VGGT式的交替注意力词元；语义分支用冻结的CLIP-LSeg编码器取得逐位置、与文字对齐的语义特征。只在几何主干的最后一层，几何词元作为查询、语义词元作为键和值做一次交叉注意力，使几何解码获得语义上下文。解码器输出深度和相机，并为每个像素预测一个三维高斯：中心由深度反投影确定，协方差描述空间扩展，颜色参数描述外观，不透明度描述遮挡贡献，另附语义向量$\bm{f}_i$。高斯是可投影、可聚合的场景表示，不应把其中心直接当作精确的物体表面。

为用图像监督这些三维属性，要把它们投回观察相机。每个高斯的中心用节首的投影得到二维位置，协方差经投影的局部线性近似成为像平面上的椭圆；按深度先后计算对像素的贡献。若第$i$项在该像素的合成系数为$a_i\in[0,1]$，前方透过比例为$T_i=\prod_{j<i}(1-a_j)$，则颜色和特征分别按$\sum_iT_i a_i\bm c_i$与$\sum_iT_i a_i\bm f_i$汇总。两者共用贡献权重，因此语义特征的监督也会反向影响所预测的几何。完整的新视角渲染在\ref{sec:vision-generation-novel-view}节展开。

若两视图看到同一表面，它们的语义特征应在几何对应下相容。使用式\eqref{eq:vision-relative-camera-measurement}，将目标图像像素$\tilde{\bm{p}}_t$在预测深度$D_t$下反投影，再投到上下文图$c$：
\begin{equation}
 \bm{p}_c=
 \Pi\!\left[
 \bm{K}_c\left(
 \bm{R}_{c,t}D_t(\bm{p}_t)\bm{K}_t^{-1}\tilde{\bm{p}}_t
 +\bm{t}_{c,t}\right)\right]\text{，}
 \label{eq:vision-ff3r-warp}
\end{equation}
其中$\Pi$用第三坐标除前两坐标。按$\bm{p}_c$插值采样源特征，计算$1-\cos(\bm{f}_t,\bm{f}_{c\to t})$，只在投影落入图像且深度相容的位置求平均，再交换两图求另一方向；边界和遮挡筛选避免把不曾看见同一表面的位置强行拉到一起。

训练从同一场景采样上下文视图，几何主干和深度头由VGGT权重初始化，新增的融合与输出部分随机初始化，CLIP-LSeg保持冻结。完整目标包含渲染RGB的像素与感知误差、渲染特征与CLIP-LSeg特征的余弦误差、上述双向跨视角特征误差，以及以VGGT输出为伪标签的相机和深度蒸馏项（深度只在高置信位置比较）。原文的总目标为$L_{\mathrm{rgb}}+0.1L_{\mathrm{feat}}+0.1L_{\mathrm{warp}}+L_{\mathrm d}+10L_{\mathrm p}$，其中$L_{\mathrm{rgb}}$是$L_1$误差加0.05倍的感知误差，$L_{\mathrm p}$是相机位姿的Huber蒸馏误差。训练用DL3DV-10K中的约6500个场景、$448\times448$的图像，每次至多16个视图，在8块A100上约两天；高斯头的输出通道、语义特征维数和参数总量论文未给出。因此“无需人工相机、深度和语义标注”不等于只有RGB残差，也不等于没有预训练监督。视图密集时，可按三维位置把高斯分到体素，用几何置信$C_i$和与体素平均语义的余弦距离$d_i^{\mathrm{sem}}$共同加权：
\begin{equation}
 \begin{aligned}
 w_{i\to s}=
 \frac{\exp(C_i-\lambda d_i^{\mathrm{sem}})}
 {\sum_{j:\,v(j)=s}\exp(C_j-\lambda d_j^{\mathrm{sem}})},\qquad
 \bar{\bm{a}}_s=\sum_{i:\,v(i)=s}w_{i\to s}\bm{a}_i\text{，}
 \end{aligned}
 \label{eq:vision-ff3r-voxel}
\end{equation}
其中$v(i)$是体素编号，$\bm{a}_i$可以是位置、颜色或语义。若两项置信均为1、语义距离为0与1、$\lambda=1$，权重为$e/(e+1)\approx0.731$与0.269，偏离项贡献下降但未消失。使用时，把“椅子”编码到同一语义空间，与三维或投影后的特征比较，再按阈值恢复区域。仅颜色渲染自然，不能证明椅子的语义或尺度正确；未见的背面仍只有先验预测。作者项目页标注其为CVPR 2026 Findings论文。

五个模型的对比见表\ref{tab:vision-recognition-geometry-compare}，比较固定在同一组建筑或房间照片上。VGGT相对DUSt3R改变了一次处理的视图数，换来单次前向和多视图共同推断，代价是注意力随视图数平方增长；VGG-T$^3$把全局注意力换成测试时拟合的固定容量映射，换来线性增长，代价是近似误差和测试时的额外优化步；InfiniteVGGT改变的是信息到达方式，换来流式处理和固定工作内存，代价是被裁减的历史无法再读取；FF3R在几何上增加语义，换来可查询的三维表示，代价是更复杂的训练目标并依赖VGGT伪标签。

\begin{table*}[tbp]
 \centering
 \small
 \begin{tabularx}{\linewidth}{@{}lXXXX@{}}
 \toprule
 模型 & 视图的处理方式 & 主要输出 & 跨视图信息的保存 & 主要代价或局限\\
 \midrule
 DUSt3R & 逐对前馈，再全局对齐 & 成对点图与置信 & 每对的交叉注意力 & 对数随视图数增长，需额外优化\\
 VGGT & 全部视图一次前馈 & 相机、深度、点图、跟踪 & 全局注意力的全部键值 & 计算随视图数平方增长\\
 VGG-T$^3$ & 全部视图，测试时拟合 & 同VGGT，点图为主 & 每层一个固定容量MLP & 拟合近似误差，相机较弱\\
 InfiniteVGGT & 逐帧到达，因果处理 & 每帧相机与点图 & 首帧锚点加多样性缓存 & 删去的历史不可恢复，长程漂移\\
 FF3R & 多视图一次前馈 & 几何、三维高斯与语义 & 几何读取语义词元 & 训练目标复杂，依赖伪标签\\
 \bottomrule
 \end{tabularx}
 \caption{五种前馈几何模型的比较。尺度均为规范化的约定尺度，若无额外度量依据，都不直接给出米制结果。}
 \label{tab:vision-recognition-geometry-compare}
\end{table*}

\section{图像匹配与检索}
\label{sec:vision-recognition-retrieval}

前几节把一幅图像变成类别、范围和测量。检索的输入是一幅查询图像（或其中一块区域）和一个大图库，输出是按相关程度排序的图库图像列表。它面临两个相互牵制的难点。第一，什么算“相关”要事先规定：找相似的街景和找同一辆车的其他照片，正例完全不同。第二，图库可能有百万张图，逐张做精细的局部匹配与几何验证不可承受；只用一个紧凑向量比较又很快，却容易被外观相似的其他实例混淆。

常用的解决方式是按步骤分工。先规定检索相关性，确定正例、忽略项和评价协议；再用代价低的表示从全库中召回少量候选，宁可多召回，不可漏掉；最后只对这些候选做代价高的局部匹配和几何验证，重新排序。三个步骤的顺序反映了信息由粗到细、成本由低到高的安排，后一步只能在前一步留下的候选中纠错，不能找回前一步漏掉的图。

\subsection{检索相关性的规定}
\label{sec:vision-recognition-relevance}

“找类似的红色汽车”和“找同一辆车的其他照片”使用不同的正例。前者容许不同实例，只要外观或语义相近；后者需要车身划痕、车牌等区分个体的证据，同时允许颜色因光照而变化。\term{实例检索}（instance retrieval）指后一种任务：正例是同一个物体或地点的照片，而不是同类物体。若没有预先规定相关性，高相似度排序可能朝错误的目标优化：一个按类别语义训练的表示会把所有红色货车排在前面，这对实例检索恰恰是干扰。

查询的范围也要规定。整幅照片包含建筑、树木和天空，局部查询只指定一块招牌，两者容许的背景、遮挡和视角变化不同。图库中的近重复图会使任务过于容易，干扰图的规模则决定排序有多难。Revisited Oxford和Paris数据集修订了地标实例检索的标注，按难度把正例分为容易和困难两档，并把难以判断的图像标为忽略项，还提供了约一百万张干扰图\scite{vision-1803-11285}。评价前应固定查询范围、正例定义、忽略项、图库版本和干扰图规模；文字查询图像的检索由\ref{chap:vision-multimodal-understanding}章继续讨论。

\subsection{候选图像的召回}
\label{sec:vision-recognition-candidates}

召回的输入是查询图像和预先处理好的图库，输出是前$k$个候选及其图号。难点在于既要快，又不能漏掉视角差异大的正例。图库图像的表示在建库时离线计算一次，查询时只计算查询图的表示，然后在库中搜索。

表示有两种形式。全局表示把整幅图汇总为一个固定维向量，比较只需一次内积，配合近似最近邻索引可以在百万级图库中毫秒级返回候选；GeM是一种为实例检索学习全局向量的方法。局部表示为每幅图保留几百到几千个关键点及其描述向量，信息更细，可以描述招牌一角这样的局部，但直接逐对比较代价高；它在建库时一并提取存储，既可经视觉词汇聚合后用于投票式召回，也是下一步重排的输入。SuperPoint是一种学习得到的局部点检测与描述方法。近似索引带来的召回损失与表示本身的区分能力应分开检查。

这两种表示都有经典前身，而且学习方法正是沿着经典流程的结构逐处替换的。\ref{sec:vision-features-bovw}节的视觉词袋先在每幅图上检测SIFT关键点（\ref{sec:vision-features-sift}节），得到几百到几千个128维描述子；再在大量描述子上用$k$均值聚出一个视觉词汇表，把每个描述子量化到最近的“视觉词”。一幅图于是变成一个视觉词频率向量，按TF-IDF加权：在少数图中出现的词更有区分力，权重更大。检索时用倒排索引只访问与查询共享视觉词的图库图，因而可以扩展到百万级图库。VLAD进一步保留了量化丢掉的信息：它不只数每个视觉词出现几次，而是把落到第$c$个词的描述子与该词中心的残差$\bm x-\bm\mu_c$累加，拼成一个$K\times128$维向量，再归一化，得到更紧凑也更有区分力的全局描述。召回之后，经典流程用\ref{sec:vision-recognition-ransac}节的RANSAC做几何验证：只有能被同一单应或基础矩阵解释的匹配才算数。这条流水线的弱点在前两步：SIFT在大光照变化和弱纹理处检测不稳定，量化把相近但不同的描述子合并或把同一结构拆开，词频统计又看不到空间排列。GeM把“局部特征加聚合”换成卷积特征加可学习的池化，并用三维重建自动产生训练对；SuperPoint把手工的关键点检测和描述换成网络；LightGlue把最近邻匹配换成带上下文的注意力匹配。几何验证保留了下来，因为它不依赖任何外观假设。

\subsubsection{GeM}
\label{sec:vision-recognition-gem}

卷积网络的最后一层特征是一个空间网格，需要汇总成一个向量。平均池化平等计入每个位置，背景中的天空和路面会稀释目标；最大池化只保留每个通道最强的响应，又对单个噪声响应过于敏感。另一个问题是监督：分类网络学到的是类别语义，而实例检索需要区分同类中的不同个体，人工标注“哪些照片是同一建筑”成本很高。GeM针对这两点，提出可学习指数的广义均值池化，并利用三维重建自动挖掘训练图对\scite{vision-1711-02512}。

\term{广义均值池化}（generalized-mean pooling, GeM）对第$k$通道的非负特征$x_{p,k}$计算
\begin{equation}
 g_k=\left(\frac1n\sum_{p=1}^{n}x_{p,k}^{\,q_k}\right)^{1/q_k},
 \qquad q_k>0\text{。}
 \label{eq:vision-gem}
\end{equation}
$q_k=1$为平均，$q_k$趋于无穷时趋近最大值，介于其间时强响应权重更大，但其他位置仍有贡献。对响应$(1,1,1,4)$，平均为1.75，$q=3$时为$(67/4)^{1/3}\approx2.559$，最大值为4。实现时在幂运算前给特征加下界，避免零点的数值问题；指数可以固定，也可以与编码器一同学习。

完整的流程是：全卷积主干去掉最后的池化和全连接层，把图像变成$h\times w\times C$的非负特征（ReLU之后），GeM得到$C$维向量，可选地接一个白化变换，再做$L_2$归一化得到单位向量$\bar{\bm g}$。原文比较了三种主干：AlexNet给出256维、VGG16给出512维、ResNet-101给出2048维描述。以ResNet-101和$1024\times1024$的测试图为例，末层特征约为$32\times32\times2048$，GeM把这1024个位置汇成一个2048维向量。指数$p$可以每通道一个，也可以全部通道共享；原文最终采用共享且可学习的单个$p$，初始化为3，VGG16微调后收敛到约2.92。白化是在训练集描述上学到的线性判别投影，作为后处理接在GeM之后，默认不降维。单位向量满足$\lVert\bm{u}-\bm{v}\rVert_2^2=2-2\bm{u}^{\top}\bm{v}$，按欧氏距离或内积排序结果相同。

训练图对来自运动恢复结构：对大量无标注的地标照片做三维重建，共享足够多三维点的图对作为正对，不同重建模型中的图像作为负对；负对中优先选取当前描述最相近的困难负例，并随训练用当前参数重新挖掘。原文从约740万张网络照片重建出713个三维模型、16.3万余张图，其中551个模型（约13.3万张图）用于训练、162个用于验证，并去掉与Oxford、Paris测试地标重叠的模型；每个训练元组含1个查询、1个正例和5个困难负例，训练图最长边不超过362像素。因此“不需要人工标注”并不等于没有监督依据。一条训练图对$(\bm X,\bm X')$的同实例标签为$y\in\{0,1\}$，经共享编码、GeM和归一化后距离为$d=\lVert\bar{\bm g}-\bar{\bm g}'\rVert_2$，对比目标为
\begin{equation}
 \ell=\tfrac12 y d^2+
 \tfrac12(1-y)[\max(0,m-d)]^2\text{。}
 \label{eq:vision-gem-contrastive-process}
\end{equation}
正对的梯度推动距离缩短；负对只在$d<m$时受到推开的梯度，超过间隔$m$就不再计入。原文的间隔$m$对AlexNet、VGG16和ResNet-101分别取0.7、0.75和0.85。教学中$m=1$，正对距离0.4的目标为0.08，负对距离0.4的目标为0.18，负对距离1.2为零。编码器从分类权重初始化，训练共同更新编码器与池化指数。

使用时为每幅图库图像保存归一化描述与图号，对查询做完全相同的处理，按内积取前$k$个候选。指数、白化和归一化须一致，只对图库做白化而忘了查询，会改变距离定义。同一建筑的两张照片可能因视角差异而距离较远，另一栋同样玻璃幕墙的建筑却可能很近；全局描述只提供候选，不能证明局部细节一致。已有的DINOv3向量也可以作为召回基线，比较时应固定输入、归一化与图库。

\subsubsection{SuperPoint}
\label{sec:vision-recognition-superpoint}

局部特征需要两样东西：在不同照片中能重复检测到的位置，以及能把对应位置配起来的描述。SIFT等手工方法用设计好的规则完成这两件事\scite{vision-lowe2004}，在模糊、光照变化和弱纹理处常常失效；用学习方法改进时，困难在于“兴趣点”没有自然的人工标注，人很难一致地标出所有可重复的角点。SuperPoint的想法是用自监督产生兴趣点标签：先在具有已知角点的合成几何图形上训练检测器，再对真实图像施加多个已知的单应变换，把各次检测结果变换回原图并汇总，作为更稳定的伪标签；检测与描述共用一个编码器联合训练\scite{vision-1712-07629}。这里的兴趣点没有“左腕”那样的对象部位语义，只要求适合跨图重复辨认。

输入是$H\times W$的灰度图，VGG式编码器经三次下采样得到$H/8\times W/8\times128$的特征。编码器有8个$3\times3$卷积层，通道依次为64、64、64、64、128、128、128、128，每两层之后接一个$2\times2$最大池化（共三个）。两个头各是一个256通道的$3\times3$卷积加一个$1\times1$卷积，检测头的$1\times1$卷积输出65通道，描述头输出256通道；全网各层都接ReLU和批归一化（输出层除外）。逐层结构见表\ref{tab:vision-recognition-superpoint-layers}。

\begin{table*}[tbp]
 \centering
 \small
 \begin{tabular}{@{}lllrrl@{}}
 \toprule
 阶段或层组 & 层类型 & 卷积核与步长 & 输出通道 & 重复次数 & 输出尺寸\\
 \midrule
 编码器1 & 卷积；池化 & $3\times3$；$2\times2$，步长2 & 64 & 2 & $H/2\times W/2$\\
 编码器2 & 卷积；池化 & $3\times3$；$2\times2$，步长2 & 64 & 2 & $H/4\times W/4$\\
 编码器3 & 卷积；池化 & $3\times3$；$2\times2$，步长2 & 128 & 2 & $H/8\times W/8$\\
 编码器4 & 卷积 & $3\times3$ & 128 & 2 & $H/8\times W/8$\\
 检测头 & 卷积 & $3\times3$；$1\times1$ & 256；65 & 1 & $H/8\times W/8$\\
 描述头 & 卷积 & $3\times3$；$1\times1$ & 256；256 & 1 & $H/8\times W/8$\\
 \bottomrule
 \end{tabular}
 \caption{SuperPoint的逐层结构，取自原文第6节。重复次数指每组中$3\times3$卷积的个数；两个头共享编码器。检测头输出经softmax、去掉“无点”通道并重排后得到$H\times W$的点分数图，描述头输出经双三次插值和$L_2$归一化得到每点的256维描述，这两步都没有可学习参数。}
 \label{tab:vision-recognition-superpoint-layers}
\end{table*}

检测头输出$H/8\times W/8\times65$：每个粗格对应原图一个$8\times8$区域，64个通道表示点落在其中哪个像素，第65个通道表示“该区域没有点”。对65个通道做softmax，去掉无点通道，再把64个通道重排回$8\times8$的空间位置，就得到原图分辨率的点分数图，不需要反卷积。描述头输出$H/8\times W/8\times256$的粗描述，在关键点位置做双三次插值采样，并$L_2$归一化。最终保存每个点的$(x,y,\text{得分},\bm{d})$。

训练分三步：在合成形状上训练检测头（这一阶段的模型称为MagicPoint）；用单应适应为真实图像生成伪标签；用伪标签和已知单应同时训练检测和描述。检测损失是每个粗格上65类的交叉熵。描述损失利用已知单应：一幅图与其变换图之间，对应的粗格描述内积应大于正间隔$m_p=1$，不对应的粗格应小于负间隔$m_n=0.2$，形成铰链损失；由于不对应的粗格对远多于对应的，正对项乘以$\lambda_d=250$来平衡，描述损失整体再乘0.0001后与检测损失相加。单应适应对每幅图采样$N_h=100$个随机单应，原文做了两轮。真实图像训练使用约8万张COCO图像，缩放为$240\times320$的灰度图；在Titan X上处理$480\times640$的图像约需13毫秒。单应适应提供的只是平面变换下的监督，不能覆盖任意三维视差。

使用时按阈值和局部峰值抑制保留点，并可按分数限制点数。若输入缩小为原图的一半，网络点$(100,60)$还原为原图$(200,120)$，描述向量本身不变。保存点坐标所在网格和图号，后续匹配器才能在同一坐标系中做几何验证；没有点时应返回空集合，不虚造对应。窗格重复时，检测器可以在许多窗角上稳定响应，但描述可能也很相似；“点能重复检测”与“点能匹配到正确窗格”是两个不同要求，后者留给重排步骤。

两种召回表示的对比见表\ref{tab:vision-recognition-recall-compare}。它们不是互相替代的方案，而是在召回步骤中为后续计算准备的两份不同信息。

\begin{table}[htbp]
 \centering
 \small
 \begin{tabularx}{\linewidth}{@{}lXX@{}}
 \toprule
 比较项 & GeM & SuperPoint\\
 \midrule
 每图存储 & 一个$C$维单位向量 & 数百至数千个点及256维描述\\
 比较成本 & 一次内积，可用近似索引 & 点集之间的匹配，代价高\\
 监督来源 & 三维重建挖掘的正负图对 & 合成形状与单应适应的伪标签\\
 擅长 & 大图库中快速召回 & 局部对应，供几何验证\\
 易错情形 & 外观相似的不同实例 & 重复纹理、大视差\\
 \bottomrule
 \end{tabularx}
 \caption{召回步骤中两种表示的比较。前者决定哪些图像进入候选，后者为候选的重排准备局部证据。}
 \label{tab:vision-recognition-recall-compare}
\end{table}

\subsection{候选图像的重排}
\label{sec:vision-recognition-verification}

重排的输入是查询图和召回的前$k$个候选，每张图都附有局部点与描述；输出是重新排序的列表，并附每张候选与查询之间的对应点和几何变换。难点在于外观相似的其他实例：同样的玻璃幕墙、同样的窗格会产生许多描述相近的点对。重排分两步：先在两图的点之间建立对应，再检查这些对应能否被一个几何变换共同解释。LightGlue负责第一步，用注意力利用上下文给出点对；RANSAC负责第二步，在含大量错误的点对中寻找被多数支持的几何模型。几何通过之后，最终的相关性仍要回到查询定义判断。

\subsubsection{LightGlue}
\label{sec:vision-recognition-lightglue}

独立的最近邻匹配只比较两个描述，无法利用邻近点怎样共同移动的信息，重复窗格因此很容易配错。用注意力让两图的点相互交流可以解决这一问题，但此前的注意力匹配器对每一对图像都执行固定层数的计算，简单图对与困难图对花费相同。LightGlue的想法是在注意力匹配的基础上让计算量自适应：容易的图对提前结束，已经确定无法匹配的点提前丢弃\scite{vision-2306-13643}。它接收已有的关键点和描述（例如SuperPoint的输出），不负责从像素检测点。

两图分别有$n,m$个点，每个点的输入是归一化坐标和描述。网络有9个相同的层，宽度$d=256$，每个注意力单元4个头；每层先做图内自注意力，再做跨图交叉注意力。旋转位置编码只用在自注意力中：它按两点坐标之差旋转查询和键，使注意力得分只依赖相对位置，点因此能感知邻近点的空间排列；交叉注意力不加位置编码，因为两图之间没有共同的坐标系。配置见表\ref{tab:vision-recognition-lightglue-config}。

\begin{table}[htbp]
 \centering
 \small
 \begin{tabular}{@{}lcccccc@{}}
 \toprule
 模型 & 词元 & 层数 & 宽度 & 头数 & MLP倍率 & 参数量\\
 \midrule
 LightGlue & 每个关键点1个 & 9 & 256 & 4 & 论文未给出 & 论文未给出\\
 \bottomrule
 \end{tabular}
 \caption{LightGlue的配置。LightGlue不切图像块，词元是两图的关键点，数量随输入而变（原文微调时每图取2048个点）；每层含一次自注意力和一次交叉注意力。}
 \label{tab:vision-recognition-lightglue-config}
\end{table}每层之后都可以接匹配头：把更新后的描述投影为$d$维，以$1/\sqrt d$缩放两两内积得到$S_{i,j}$；另由线性层和sigmoid给出每个点的可匹配性$\mu_i,\nu_j\in[0,1]$。匹配概率为
\begin{equation}
 P_{i,j}=\mu_i\nu_j
 \frac{\exp S_{i,j}}{\sum_{j'=1}^{m}\exp S_{i,j'}}
 \frac{\exp S_{i,j}}{\sum_{i'=1}^{n}\exp S_{i',j}}\text{。}
 \label{eq:vision-lightglue-score}
\end{equation}
第一个分式要求$i$在对面所有候选中偏好$j$，第二个要求$j$在本侧所有候选中偏好$i$，可匹配性则允许被遮挡或越界的点放弃匹配。若$\mu_i=0.9,\nu_j=0.8$，两个方向的归一化分数为0.7和0.6，则$P_{i,j}=0.3024$。

训练数据是带真实点对的图像对，点对由深度和相机重投影得到，或由已知单应得到；遮挡和越界的点必须标为无对应。若真实匹配集合为$\mathcal M$，两侧无对应点集为$U,V$，目标为
\begin{equation}
 \begin{aligned}
 L_{\mathrm{match}}={}&-\frac1{|\mathcal M|}
 \sum_{(i,j)\in\mathcal M}\log P_{i,j}\\
 &-\frac1{|U|}\sum_{i\in U}\log(1-\mu_i)
 -\frac1{|V|}\sum_{j\in V}\log(1-\nu_j)\text{。}
 \end{aligned}
 \label{eq:vision-lightglue-training-process}
\end{equation}
空集合对应的项应跳过。这一目标在每一层都计算，使任何一层都能输出可用的匹配，这是提前退出的前提。之后另训练一个小的置信分类器，预测某个点在当前层的决定是否已与最终层一致，其梯度不回传到主干，避免为早退判断改动描述。匹配器的权重与所用的描述类型绑定，为一种特征训练的权重不能无条件用于另一种描述。

推理时每层结束后检查置信分类器。第$\ell$层的置信阈值为$\lambda_\ell=0.8+0.1e^{-4\ell/L}$，浅层要求更高的置信；当至少95\%的点超过阈值时提前退出。置信已超过阈值、可匹配性又低于0.01的点被判为不可匹配，从后续层中删去。原文报告容易的图对2至3层后就退出；训练先在约17万张图生成的合成单应图对上预训练，再在MegaDepth上微调。最终取$P$中行最大与列最大互相指向且超过阈值的点对，输出两幅原图上的坐标和匹配得分。任一图没有点时应返回空对应，不能为了让后续几何估计运行而强行填点。$P_{i,j}$是模型的配对评分，不是几何正确性的证明；比较速度时还需固定特征类型、图像尺寸、硬件和图对难度分布。

\subsubsection{RANSAC}
\label{sec:vision-recognition-ransac}

LightGlue的点对中仍会有错误，重复纹理可以产生看似可信但错误的对应。最小二乘法用全部点对拟合几何模型，少数严重错误就会把结果拉偏。\term{随机采样一致性}（random sample consensus, RANSAC）的想法是反过来：每次只随机抽取恰好能确定模型的最少点对，估计一个候选模型，再数有多少点对与它相符，保留获得最多支持的模型\scite{vision-ransac-1981}。只要某一次抽到的全是正确点对，就能得到好模型，而错误点对再多也只是不支持它。

几何模型取决于场景。平面上的点在两图间由\term{单应变换}（homography）联系，齐次坐标满足$\tilde{\bm{p}}'\sim\bm{H}\tilde{\bm{p}}$；一般三维场景的两视图则遵守\term{极线约束}（epipolar constraint）$\tilde{\bm{p}}'^\top\bm{F}\tilde{\bm{p}}=0$，$\bm{F}$为基础矩阵，它只约束对应点落在另一图的一条线上。以单应为例，消去齐次比例可写$\tilde{\bm p}'\times\bm H\tilde{\bm p}=\bm0$，每个点对提供两条独立线性约束，四对约束$\bm H$除尺度外的八个自由度；共线或重复的点会使估计退化。先把坐标平移缩放到合适的数值范围再求解，再变回原图。候选模型的对称转移误差为
\[
 e_i^2=\lVert\bm p_i'-\Pi(\bm H\tilde{\bm p}_i)\rVert_2^2
       +\lVert\bm p_i-\Pi(\bm H^{-1}\tilde{\bm p}_i')\rVert_2^2\text{，}
\]
要求$\bm H$可逆；误差低于像素阈值的点对为内点，其余为外点。

RANSAC不需要训练，它的“学习”是在当前点对上的搜索。算法从无模型开始，反复执行抽样、估计、计内点；若新模型的内点更多，或内点相同而残差更小，就替换当前最佳；最后用全部内点重新估计模型。抽样次数决定成功的可能：每次抽取$s$个点对，内点比例为$w$，全内点样本的概率约为$w^s$，$k$次都没抽到的概率为$(1-w^s)^k$。要求它不超过$1-p$，取对数得
\begin{equation}
 k\geq\frac{\log(1-p)}{\log(1-w^s)}\text{，}
 \label{eq:vision-ransac-iterations}
\end{equation}
分母为负，除法时不等号方向随之改变。取$w=0.5,s=4,p=0.99$，至少需要72次；内点率降低时次数迅速增加，$w=0$时不存在有限次数保证。实现常根据当前最佳模型的内点率动态更新$k$，这依赖估计可信，错误的高支持模型也可能使搜索过早结束。

原始RANSAC只用残差判断内点，忽略了点在空间上的相邻关系。Graph-Cut RANSAC在局部优化阶段为每个点对设二值标签$\ell_i$（内点或外点），把局部目标写成$\sum_iD_i(\ell_i;e_i)+\lambda\sum_{(i,j)\in E}w_{i,j}\mathbf1\{\ell_i\ne\ell_j\}$：一元项按残差惩罚不合适的标记，邻接项倾向于让空间上相邻的点对标签一致\scite{vision-1706-00984}。在满足二值可割条件时，图割求得这一次标记的最优解，再用选出的内点重估几何。它解决的是给定模型附近的局部标记问题，不是对所有模型和标签组合的全局最优；邻近的错误匹配也可能外观一致，仍须检查模型假设。

重排时以内点数量、内点覆盖查询范围的比例和残差为依据重新排序，排序规则应事先固定。建筑正面近似平面时，单应适用；有明显前后视差时，同一个单应无法对齐所有表面，应改用基础矩阵或分区建模。纯旋转、近共线点和重叠太少都会改变估计条件。一小块重复窗格获得的高内点数，也不能替代整幅查询的相关性。图\ref{fig:vision-retrieval-evidence}把全局召回、局部配对与几何验证分开列出，避免把一个阶段的高分误作整条证据链已经完成。

\begin{figure}[htbp]
 \centering
 % 流程节点间箭头均连接左右中心；下方为各阶段尚未排除的错误。
\begin{tikzpicture}[font=\figurefont\small,
 block/.style={draw=BookTeal,fill=BookTeal!8,minimum width=28mm,
   minimum height=12mm,align=center,inner sep=1.5mm},
 flow/.style={-{Stealth[length=2mm]},draw=BookInk,line width=0.9pt},
 note/.style={align=center,text width=28mm,font=\figurefont\footnotesize,
   text=BookMuted}]
 \node[block] (global) at (0,0) {全局描述\\筛选候选};
 \node[block] (local) at (3.7,0) {局部对应\\匹配细节};
 \node[block] (geo) at (7.4,0) {几何支持\\核验重排};
 \draw[flow] (global.east)--(local.west);
 \draw[flow] (local.east)--(geo.west);
 \node[note] at (0,-1.6) {可能只同类\\并非同一实例};
 \node[note] at (3.7,-1.6) {重复窗格\\可能相互错配};
 \node[note] at (7.4,-1.6) {几何匹配仍需\\满足查询定义};
 \foreach \x in {0,3.7,7.4}
   \draw[draw=BookRule,line width=0.7pt] (\x,-0.65)--(\x,-1);
\end{tikzpicture}

 \caption{图像检索中逐步保留更具体的证据。上方沿主流程减少候选，下方说明每步仍可能留下的错误。几何一致可以帮助重排，却仍要检查所得实例是否满足查询的相关性定义。}
 \label{fig:vision-retrieval-evidence}
\end{figure}

两种重排方法的对比见表\ref{tab:vision-recognition-rerank-compare}。LightGlue用学习得到的上下文消解描述的歧义，给出更多正确点对；RANSAC不学习任何参数，而是用几何模型筛掉剩余的错误。二者串联使用：没有好的点对，RANSAC需要的抽样次数会急剧增加；没有几何验证，LightGlue的高分点对也可能整体指向一个相似的错误实例。

\begin{table}[htbp]
 \centering
 \small
 \begin{tabularx}{\linewidth}{@{}lXX@{}}
 \toprule
 比较项 & LightGlue & RANSAC\\
 \midrule
 输入 & 两图的点与描述 & 两图的候选点对\\
 输出 & 点对及匹配概率 & 几何模型与内点集合\\
 依据 & 学习的图内与跨图上下文 & 单应或极线几何的一致性\\
 是否训练 & 需要带真实对应的图对 & 无需训练，按次数搜索\\
 失败情形 & 重复纹理、描述类型不符 & 内点率过低、退化配置、模型不符\\
 \bottomrule
 \end{tabularx}
 \caption{重排步骤中两种方法的分工。前者提高对应质量，后者检验对应能否被同一几何关系解释。}
 \label{tab:vision-recognition-rerank-compare}
\end{table}

\section{识别与检索结果的检验}
\label{sec:vision-recognition-evaluation}

本章的各种输出最终要组合使用，检验因此不能只看单项分数。组合系统需要让类别、框、掩码、属性和测量指向同一对象：若车辆检测正确，属性却来自旁车，逐模块的总体准确率仍可能很高。街景应逐实例保留标识并核验空间对应；工业零件应同时检查异常热图、整图得分与规定的缺陷，避免正常材质变化被业务阈值放大。

\ref{sec:vision-recognition-challenges}节已经给出各项指标的定义，这里关心的是报告这些指标时必须固定什么。检测的AP取决于插值方式、IoU阈值、尺寸分组和每图检测数量上限：PASCAL VOC式的单阈值AP与COCO式在10个阈值上平均的AP不是同一个量，COCO还规定每图最多计入100个得分最高的检测，并按面积小于$32^2$、介于$32^2$与$96^2$之间、大于$96^2$像素分出小、中、大三组。因此不能把任意一个“AP”视为同一指标，跨论文比较前要先核对协议。

区域任务有不同的分母。mIoU按类别平均，小类别与大类别同权，适合检查路杆这类细结构；按像素计的总体准确率则被道路和天空主导。实例分割的掩码AP先要匹配对象；式\eqref{eq:vision-pq}的PQ同时受范围与实例错误影响，前文的例子已说明只报告匹配区域的平均IoU会掩盖漏报。

检索的式\eqref{eq:vision-retrieval-ap}在相关项全部可见时才是完整的AP；只返回前若干项时，计算截断指标须说明分母和遗漏规则，无相关项的查询应按协议排除或单独评价。Revisited Oxford与Paris的“忽略项”在计算时既不算正例也不算负例，若误当负例，排在前面的难以判断的图就会无端压低AP。

开放条件需要单独的检验。未知发现在隐藏新类名称的条件下计算未知召回，同时记录未知被误报为已知的数量、背景假阳性，以及后续加类前后旧类的AP。教学中有5个未知物体，正确发现3个、误当已知1个、完全漏掉1个，未知召回为$3/5=0.6$；这不能由已知货车的AP代替。文字定位和概念分割要固定候选短语、正负示例及提示成本，使用真实框作提示取得的优势不能直接归于无提示系统。

测量量按单位和坐标检验。姿态比较固定关节集合、单位和坐标，以MPJPE计算$(1/j)\sum_k\lVert\hat{\bm J}_k-\bm J_k\rVert_2$。根对齐只去平移，相似变换对齐还去旋转和尺度，二者回答的问题不同。人体表面另比对应顶点误差、轮廓重投影和手足局部，少量关节误差不能证明全部网格准确。单目深度需区分相对排序、边界和米制误差，允许对齐的结果不能与未对齐的度量结果混为一表。离线场景比较固定视图集合、分辨率、首相机参照和拟合步数，分别报告相机、点图、重投影、完整度及总成本；流式场景只允许已到达的图，记录缓存预算、局部误差与长程漂移，一次全集对齐后的轨迹误差会隐藏部分累计误差。三维语义先在有效几何下找到同一表面，再比较各视角的查询归属；RGB渲染、语言相似度和几何一致各有误差来源，三者之一正确不能替其他两项通过。表\ref{tab:vision-recognition-protocols}汇总了各类输出需要固定的条件。

\begin{table}[htbp]
 \centering
 \small
 \begin{tabularx}{\linewidth}{@{}lXX@{}}
 \toprule
 输出 & 需要固定的条件 & 不能由它单独证明的内容\\
 \midrule
 类别、属性 & 类别表、未知项、对象范围 & 对象存在及属性归属都正确\\
 框、掩码 & 匹配规则、提示、可见范围 & 未观测部分及三维形状正确\\
 数量、姿态 & 计入规则、关节与对齐方式 & 实例清单、身份及原始朝向正确\\
 深度、点图 & 尺度、坐标、有效区域 & 距离带实际单位且适合测量\\
 检索排序 & 正例、图库、忽略项及难度 & 对象属性、文字和距离正确\\
 \bottomrule
 \end{tabularx}
 \caption{不同视觉成果所需的评价条件。右列列出仍需另行核验的推论，不表示这些任务必然无法提供相应能力。}
 \label{tab:vision-recognition-protocols}
\end{table}

\section{本章小结}
\label{sec:vision-recognition-summary}

回到开篇的路口照片。“这是街道”由场景分类给出，类别表固定时用冻结表示加分类头，类别在使用时给出时用图文相似度；“哪辆是红色”要先把属性限制在对象区域内；“车在哪里”按需要的精细程度给出框、像素类别、实例掩码或透明度；“有几辆、离多远”在定位之上引入计数规则、相机模型和尺度；“另一张照片里是不是同一辆车”则由全局召回、局部配对和几何验证逐步回答。每一步都把预测接回明确的对象、坐标和单位，这正是本章要建立的从像素到可核验结果的计算。

贯穿各节的是同一类设计选择：哪个输出负责哪个目标，目标由谁指定，结果在什么坐标和尺度下成立。检测从网格、锚框、逐位置距离走到集合匹配，再把类别表换成文字或对象性；掩码从框内预测走到查询掩码、点提示和概念提示；几何从逐对点图走到多视图联合、固定容量拟合和流式缓存。每一次改动都换来一种能力，也带来新的代价，表格中的比较说明了这些交换。

第一节提出的五类挑战，在本章得到的回答程度并不相同。外观变化带来的类别歧义，经典方法用梯度直方图和视觉词袋这类人工不变量缓和，学习方法用大规模预训练表示加分类头或文字类别处理，但背景捷径和细粒度部位仍需单独检查。尺度与遮挡带来的定位困难，从滑动窗口、图像金字塔和级联，演变为卷积的逐位置预测、特征金字塔和焦点损失，再到集合匹配；拥挤遮挡下的漏检并没有消失。稠密输出与分辨率的矛盾，由空洞卷积、编码器—解码器和查询式掩码缓解，而阈值、分水岭、图割和超像素给出的边界线索仍在交互分割和抠图中延续。单目观测带来的几何不确定无法由网络消除，只能由多视图观测或训练先验补足，本章在每一处都标明了尺度的来源；双目视差、平面扫描与光束法平差的结构，分别在MVSNet的代价体和前馈几何模型的全局对齐中延续下来。开放类别与长尾带来的泛化困难，靠文字类别、对象性、正常记忆和自动挖掘的训练对得到部分缓解，几何验证则是经典与学习两条路线共用的、与类别无关的核验手段。各模型的逐层结构表把这些设计落到了具体的层数、通道和张量形状上，读者可以据此复核每一步的输入输出。

仍未解决的部分同样清楚。候选高分不等于目标存在，区域正确不等于身份正确，相对深度不等于米制距离，几何支持也不等于任意业务相关性。开放概念增加了目标歧义，长尾类别需要更细的覆盖分析，遮挡和单目投影限制了可恢复的空间量；有限缓存不保证长期回访没有漂移，三维语义只有与位置和可见表面共同核验才可用。LVIS把大词表和少见实例纳入分割评价，MVTec AD把正常训练与测试缺陷定位分开，为这些问题提供了不同的观察条件\scite{vision-1908-03195,vision-mvtec2019}。连续帧中的身份、事件和检索留给\ref{chap:vision-video-recognition}章，文字与图像的联合定位和问答留给\ref{chap:vision-multimodal-understanding}章。

\paragraph{扩展阅读}
检测方面，Deformable DETR针对DETR收敛慢、小目标弱的问题，让每个查询只在参考点附近采样少量位置做注意力，从而能读取多尺度特征图；Grounding DINO中的可变形交叉注意力和OWOBJ实验所接的检测器都以它为基础\scite{vision-2010-04159}。DINO检测器在此之上加入含正负样本的去噪训练、混合查询选择和逐层修正参考框，Grounding DINO正是把这一检测器扩展到文字类别\scite{vision-2203-03605}。

分割、姿态与深度方面，PSPNet在主干末端对特征图做多种尺度的池化再拼接，用金字塔池化汇集全局上下文，可与DeepLabv3+以空洞卷积扩大感受野的做法对照\scite{vision-1612-01105}。HRNet不采用先降分辨率再恢复的结构，而是全程并行保持高分辨率分支，并在各分辨率之间反复交换信息，适合姿态和语义分割这类需要精确位置的输出\scite{vision-1908-07919}。ViTPose表明普通ViT主干加轻量解码器回归关节热图即可在COCO关键点上取得很强的结果；它采用先检测人、再逐人估计的自顶向下流程，与OpenPose自底向上组装肢体的方式不同\scite{vision-2204-12484}。MiDaS用对尺度和平移不变的损失混合多个标注口径不同的深度数据集训练，输出相对深度，与Depth Anything V2同样面向相对深度；ZoeDepth在相对深度预训练之上加入按场景类型选择的米制深度头，尝试把相对深度换算为以米为单位的距离\scite{vision-1907-01341,vision-2302-12288}。

匹配检索与基准方面，DELG在一个网络中同时输出GeM池化的全局描述子和经注意力筛选的局部描述子，把召回和重排所需的两类特征合到一次前向计算中\scite{vision-2001-05027}。SuperGlue是LightGlue的前身，它用图神经网络在两组关键点之间交替做自注意力和交叉注意力，再用Sinkhorn迭代求解带“未匹配”选项的最优传输\scite{vision-1911-11763}。LoFTR不检测关键点，先在粗网格上用Transformer建立稠密对应、再逐块细化，在弱纹理区域更有优势\scite{vision-2104-00680}。COCO的原始论文交代了本章多次使用的检测和实例分割标注怎样采集、类别怎样设计\scite{vision-1405-0312}；ADE20K提供150类的场景解析基准，常用于比较语义分割与全景分割模型\scite{vision-1608-05442}。

\chapter{图像复原、增强与压缩}
\label{chap:vision-image-restoration}

夜间拍到的招牌偏暗，字画周围有噪点，放大后又出现模糊和方块。我们希望读出招牌上的字，也希望把照片保存成较小的文件发给别人。把照片处理得明亮、锐利，并不能单凭外观判断字已经恢复正确；文件变小了，也不能只看重建图像就知道实际节省了多少存储。本章要回答三个相互关联的问题：受损的观测能恢复到什么程度，依据是什么；不改变内容时，怎样调整亮度和颜色使图像适合观看；怎样用尽量少的比特保存图像，同时说清丢掉了什么。

设希望恢复的图像为$\bm{X}\in\mathbb{R}^{h\times w\times c}$，实际观测为$\bm{Y}$，其中$h,w$是高和宽，$c$是通道数。一个便于分析的成像模型是
\begin{equation}
 \bm{Y}=\mathcal{Q}\bigl(\mathcal{S}(\bm{K}*\bm{X})+\bm{N}\bigr)\text{，}
 \label{eq:vision-image-degradation}
\end{equation}
其中$\bm{K}$描述模糊，$*$表示卷积，$\mathcal{S}$表示采样，$\bm{N}$表示噪声，$\mathcal{Q}$概括量化等后处理。它是一个简化模型：传感器噪声可能发生在颜色形成之前，图像也可能经历多次缩放和压缩，不能把这些操作一律交换顺序。RGB图像通常已经经过颜色插值、白平衡和色调映射；RAW保留更接近传感器读出的数据，两者可供反演的信息不同，成像管线的各个步骤见\ref{sec:vision-formation-isp}节。

式\eqref{eq:vision-image-degradation}把三个目的区分开来。\term{图像复原}（image restoration）依据受损观测估计$\bm{X}$，即尽量撤销$\bm{K}$、$\mathcal{S}$、$\bm{N}$造成的损失；\term{图像增强}（image enhancement）不追求某个“原图”，而是调整亮度、对比或颜色，使呈现适合指定用途；\term{图像压缩}（image compression）主动引入新的$\mathcal{Q}$，用更少比特表达图像，再控制它造成的失真。本章用到的滤波、频域运算、离散余弦变换与小波都已在\ref{chap:vision-image-transforms}章建立，这里只回顾所需的结论；\ref{chap:vision-image-video-models}章的卷积网络和Transformer在复原中充当映射，\ref{chap:vision-multimodal-models}章的潜空间扩散模型为严重退化提供生成先验，这些模型基础也不在本章重讲。下文的像素、信号和码流数字，凡未注明来源者均为教学构造。

\section{图像复原、增强与压缩的问题与挑战}
\label{sec:vision-restoration-challenges}

三类处理的输入和输出各不相同。复原的输入是受损图像$\bm{Y}$，有时附带掩码或已知的模糊核，输出是同尺寸或更高分辨率的估计$\widehat{\bm{X}}$。增强的输入是一幅RGB图像或一份RAW数据，输出是适合观看或测量的图像，通常没有唯一正确的参考。压缩的输入是图像$\bm{X}$，编码端输出一段比特流，解码端再由比特流得到重建$\widehat{\bm{X}}$。评价随之分成三类：重建与参考有多接近，看起来是否自然，用了多少比特。

最直接的接近程度是逐元素误差。对含$n$个有效标量元素的对齐图像，均方误差和\term{峰值信噪比}（peak signal-to-noise ratio, PSNR）为
\begin{equation}
 \operatorname{MSE}=\frac1n\sum_{i=1}^{n}(x_i-\hat x_i)^2,
 \qquad
 \operatorname{PSNR}
 =10\log_{10}\frac{v_{\max}^2}{\operatorname{MSE}}\text{，}
 \label{eq:vision-psnr}
\end{equation}
$v_{\max}$是约定的峰值，归一化图像通常取1，8位整数通常取255；单位为分贝（dB），越大越接近。逐元素误差对所有位置一视同仁，不区分误差落在平坦区还是边缘上。\term{结构相似性}（structural similarity, SSIM）改为在对应的局部窗口中比较亮度、对比和结构，单个窗口的值为
\begin{equation}
 \operatorname{SSIM}(\bm{x},\hat{\bm{x}})
 =
 \frac{(2\mu_x\mu_{\hat x}+c_1)(2\sigma_{x,\hat x}+c_2)}
 {(\mu_x^2+\mu_{\hat x}^2+c_1)
  (\sigma_x^2+\sigma_{\hat x}^2+c_2)}\text{，}
 \label{eq:vision-ssim}
\end{equation}
其中$\mu$是窗口内的加权均值，$\sigma^2$是相应方差，$\sigma_{x,\hat x}$是协方差，$c_1,c_2>0$稳定近常量窗口的计算；整图值由各窗口平均得到，越接近1越相似\scite{vision-ssim2004}这两项只衡量与参考的数值或结构差别。\term{学习感知图像块相似度}（learned perceptual image patch similarity, LPIPS）把两幅图送入预训练网络，比较多层归一化特征的加权距离，权重用人的相似性判断校准，越小越接近；它的定义与训练在\ref{sec:vision-restoration-lpips}节展开。

压缩还要计量比特。设一幅$h\times w$图像的完整码流有$B$比特，\term{每像素比特数}（bits per pixel, bpp）为
\begin{equation}
 \operatorname{bpp}=\frac{B}{hw}\text{，}
 \label{eq:vision-bpp}
\end{equation}
分母按像素位置计，不乘通道数。一个编解码器在不同压缩强度下得到一串（码率，PSNR）点，连成率失真曲线；比较两个编解码器，要比较整条曲线，而不是某一个点。常用的汇总量是\term{BD码率}（Bjøntegaard delta rate, BD-rate），由Bjøntegaard提出\scite{vision-bjontegaard2001}它的计算分三步：对编解码器A和B，分别用三次多项式把对数码率$r=\ln R$拟合成PSNR值$q$的函数$r_{\mathrm{A}}(q)$、$r_{\mathrm{B}}(q)$，在两条曲线共同覆盖的质量区间$[q_{\mathrm{L}},q_{\mathrm{H}}]$上求平均差，再换回比例：
\begin{equation}
 \Delta r=\frac{1}{q_{\mathrm{H}}-q_{\mathrm{L}}}
 \int_{q_{\mathrm{L}}}^{q_{\mathrm{H}}}
 \bigl(r_{\mathrm{B}}(q)-r_{\mathrm{A}}(q)\bigr)\,\mathrm{d}q,
 \qquad
 \text{BD-rate}=e^{\Delta r}-1\text{。}
 \label{eq:vision-bd-rate}
\end{equation}
它回答的问题是：在相同质量下，B平均比A多用或少用多少比特。若B在每个质量点都只需A的0.8倍码率，则$r_{\mathrm{B}}-r_{\mathrm{A}}=\ln0.8$恒定，$\Delta r=\ln0.8$，BD-rate为$-20\%$，负值表示节省。在对数码率上取平均，是为了让低码率和高码率区段按比例而不是按绝对比特数计入；两条曲线的质量区间若重叠很少，这个平均就缺少依据。

有了这些量，我们来看三类处理面对的主要困难。下面五个挑战按“复原—增强—复原与压缩共有的取舍—压缩”的顺序排列，每个都用夜间招牌说明，并交代经典方法和学习方法分别怎样处理。

\subsection{观测信息丢失带来的解不唯一}
\label{sec:vision-restoration-challenge-illposed}

把两个相邻像素平均成一个时，$(0,1)$和$(0.5,0.5)$都得到0.5；模糊核在某些频率上的响应为零时，原图在这些频率上的成分无论取什么值，模糊后都一样。复原因此常是\term{病态问题}（ill-posed problem）：解可能不存在、不唯一，或者对观测的微小扰动极为敏感。招牌上一个笔画被模糊成一团时，“8”和“3”都可能与观测相符。

经典方法用显式的先验在多个候选中选一个。正则化反卷积惩罚相邻像素的剧烈变化，维纳滤波按信号与噪声的功率谱分配各频率的增益，二者都在\ref{chap:vision-image-transforms}章的频域运算上进行。学习方法把先验交给训练数据：网络在大量“受损—清晰”配对上学到哪些图像更常见，但以平方误差训练时，它给出的是多个候选的平均（式\eqref{eq:vision-conditional-mean}），细节会被抹平。\ref{sec:vision-restoration-restoration}节按噪声、模糊、分辨率不足和缺失四种退化回答这个挑战：每种退化丢失信息的方式不同，补回信息的依据也不同。

\subsection{退化未知且混合带来的模型失配}
\label{sec:vision-restoration-challenge-mismatch}

网上转发的招牌照片可能先被手抖模糊，再被缩小、保存成JPEG，又被手机自动锐化。式\eqref{eq:vision-image-degradation}中的每个算子都不知道，而且按未知的顺序叠加。按单一高斯噪声或单一已知模糊核设计的方法，遇到这种输入时，假设的退化与实际退化不符，这称为模型失配：方法会把压缩方块当作纹理保留，或把锐化白边当作字画加强。

经典方法先估计退化参数，再套用对应的复原公式：例如从平坦区域估计噪声方差，或在盲反卷积中交替估计模糊核和图像。估计一旦出错，后续步骤全部建立在错误的模型上。学习方法改为扩大训练时的退化分布，让一个网络见过多种强度和组合：DnCNN-B在一个范围内随机取噪声强度，Real-ESRGAN把“模糊—缩放—加噪—压缩”重复两轮。合成退化覆盖得越广，模型失配越少，但它仍然不能穷尽真实相机和传播链路。\ref{sec:vision-restoration-denoising}节、\ref{sec:vision-restoration-restormer}节和\ref{sec:vision-restoration-real-esrgan}节分别在噪声、模糊和复合退化下讨论这个问题。

\subsection{低照度与目标不唯一带来的增强取舍}
\label{sec:vision-restoration-challenge-enhancement}

夜景中墙面很暗，招牌很亮。把墙面提亮到能看清时，暗处的噪声也被一起放大，因为线性提亮同时放大信号和噪声；而一条全局曲线拉开了墙面的层次，又会压缩招牌的层次。增强没有唯一正确的目标，每次调整都要说明希望看清什么、哪些数值关系必须保留。

经典方法用点运算和直方图方法重新分配灰度：伽马曲线、直方图均衡和受限局部直方图均衡化（\ref{sec:vision-transforms-intensity}节）都只改变数值映射，不能补回被噪声淹没的信号。学习方法在低照度下回到更早的成像阶段，直接读取RAW数据，联合完成去噪和颜色形成。\ref{sec:vision-restoration-enhancement}节先讨论暗光下的提亮，再讨论内容充足时全局映射和局部映射的取舍。

\subsection{细节补全带来的保真与感知冲突}
\label{sec:vision-restoration-challenge-perception}

把低清招牌放大四倍时，平方误差最优的输出是所有可能高清图像的平均，笔画边缘模糊；若让模型输出锐利自然的笔画，它就必须在多个候选中选定一个，选中的可能是“3”而不是原来的“8”。前者PSNR较高却看着模糊，后者看着清楚却可能与原图不符。这种冲突不是某个模型的缺陷：可以证明，在一定失真预算下，输出分布与真实图像分布的接近程度存在下限（\ref{sec:vision-restoration-lpips}节）。

经典方法大多以平方误差为准则，维纳滤波正是最小均方误差意义下的最优线性滤波器，因此结果偏平滑。学习方法加入感知损失、对抗损失或预训练扩散先验，让输出落在真实图像的分布上，代价是补出的细节缺少观测依据。\ref{sec:vision-restoration-super-resolution}节的三个超分辨率模型给出逐步加强的回答，\ref{sec:vision-restoration-evaluation}节说明怎样同时报告保真、感知与下游用途的结果。

\subsection{比特预算带来的码率与失真权衡}
\label{sec:vision-restoration-challenge-rate}

假设招牌照片必须压到0.5 bpp。比特花在平滑的天空上，字画边缘就会出现方块和振铃；花在字画上，天空可能出现色带。码率越低失真越大，问题是在给定码率下把失真放在用途最不敏感的地方，并让解码端只凭比特流复现结果。

经典编码器用固定的离散余弦变换或小波集中能量，用人工设计的量化表把更多误差放到人眼不敏感的高频，再用Huffman编码或算术编码去掉符号的统计冗余。各个模块分别设计，码率由质量因子间接控制。学习方法把变换、量化和概率模型放进同一个目标“码率$+\lambda\times$失真”联合训练，概率模型还可以随图像内容逐位置调整。\ref{sec:vision-restoration-compression}节先讲JPEG与JPEG 2000，再讲学习式压缩改变了哪些环节。

这五个挑战决定了本章的推进顺序。复原一节先处理第一个挑战，并在噪声和模糊两种退化下讨论第二个挑战：每种退化先给出经典方法依据的成像模型，再看学习方法放宽了哪些假设；分辨率不足时第四个挑战变得突出，三个模型依次改变补出细节的来源。增强一节回答第三个挑战，暗光下的提亮需要复原的手段，内容充足时只需选择映射。压缩一节回答第五个挑战，经典编码的各个环节正好为学习式压缩提供了逐项对照的框架。最后一节回到第四个挑战，说明保真、感知、用途和成本为什么要分开核验。

\section{图像复原}
\label{sec:vision-restoration-restoration}

复原的输入是一幅受损图像，输出是同尺寸或更高分辨率的估计图像。难点在于观测丢失了什么：噪声扰动每个像素，模糊把邻域混在一起，下采样把多个位置合并成一个，缺口则完全没有相应观测。这四种退化丢失信息的方式不同，所以本节按退化条件组织，而不是按网络结构组织。

面对同一个困难，即“观测不足以唯一确定原图”，各条路线回答的方式不同。若退化过程已知且可写成公式，可以显式比较候选图像经过退化后是否接近输入，再用一个正则项或信号统计在多个候选中挑选，这是噪声模型下的空域滤波和已知模糊核下的逆滤波、维纳滤波与约束最小二乘的做法。若退化参数未知，经典方法先估计参数再套用公式，学习方法则从大量“退化—清晰”配对中学习一个映射，让网络从图像本身判断该怎样修正；DnCNN、Restormer和部分卷积网络都属于后一类，它们的差别在于网络能看到多大范围、怎样处理无效像素。\term{盲复原}（blind restoration）指使用时不知道全部退化参数，需要从图像及训练先验中共同估计；这里的“盲”描述可用信息，不表示已经覆盖所有真实相机和损坏形式。当缺失信息太多时，平均意义下最优的估计会变得平滑，Real-ESRGAN和DiffBIR进一步借助对抗训练或预训练生成模型补出合理细节，这时就需要区分“观测支持的内容”和“先验补出的内容”。

\subsection{噪声干扰下的复原}
\label{sec:vision-restoration-denoising}
\label{sec:vision-restoration-dncnn}

细小字画与噪声都表现为局部亮暗变化。要保留前者、去掉后者，先要说明噪声怎样产生。这一小节的输入是带噪图像，输出是同尺寸的估计图像；难点是噪声的类型和强度可能未知，而且噪声与纹理在局部统计上相似。路线由“对噪声知道多少”决定：知道噪声的分布形式时，可以为它选一个对应的滤波器，或先把它变换成更易处理的形式；只知道噪声来自某个范围时，就从合成的带噪配对中学习一个去噪网络。下面先写出常见的噪声模型，再回顾与之对应的经典滤波，最后讲DnCNN。

最简单的模型是加性噪声$\bm{Y}=\bm{X}+\bm{N}$，噪声与图像内容无关，各位置独立同分布。不同物理来源给出不同的分布，表\ref{tab:vision-noise-models}列出常用的五种。\term{高斯噪声}（Gaussian noise）来自读出电路和热扰动，$n_{i,j}\sim\mathcal{N}(0,\sigma^2)$。椒盐噪声是一种脉冲噪声（\ref{sec:vision-transforms-nonlinear-denoising}节）：像素以概率$P_a$被置为最小值$a$（“椒”），以概率$P_b$被置为最大值$b$（“盐”），否则保持原值，常见于坏点和传输错误；它不是加性的，被污染的像素完全失去原值。均匀噪声$n\sim\mathcal{U}(a,b)$的方差为$\int_a^b(t-\tfrac{a+b}{2})^2\,\mathrm{d}t/(b-a)=(b-a)^2/12$，后文压缩一节用它近似量化误差。\term{周期噪声}（periodic noise）来自电源或机械干扰，形如$A\sin\bigl(2\pi(u_0i/h+v_0j/w)+\phi\bigr)$，它在空间上铺满全图，在二维傅里叶频谱上却只是一对位于$\pm(u_0,v_0)$的尖峰。

\begin{table*}[tbp]
 \centering
 \small
 \begin{tabularx}{\linewidth}{@{}p{16mm}p{40mm}p{28mm}XX@{}}
  \toprule
  噪声 & 模型 & 均值与方差 & 典型来源 & 经典处理\\
  \midrule
  高斯 & $y=x+n$，$n\sim\mathcal{N}(0,\sigma^2)$ & $0$；$\sigma^2$ & 读出电路、热噪声 & 均值或高斯平滑、维纳滤波\\
  泊松 & $y\sim\operatorname{Poisson}(x)$，$x$为期望光子数 & $x$；$x$ & 光子到达的随机性 & 方差稳定化后按高斯处理\\
  椒盐 & 以概率$P_a,P_b$取$a,b$，否则$y=x$ & 依赖$x$ & 坏点、传输错误 & 中值、自适应中值滤波\\
  均匀 & $y=x+n$，$n\sim\mathcal{U}(a,b)$ & $\frac{a+b}{2}$；$\frac{(b-a)^2}{12}$ & 量化误差的近似 & 均值平滑\\
  周期 & $y=x+A\sin(2\pi(u_0i/h+v_0j/w)+\phi)$ & 全图均值为零 & 电源、机械干扰 & 频域陷波滤波\\
  \bottomrule
 \end{tabularx}
 \caption{常用噪声模型。前四种按像素独立定义，周期噪声在空间上相关、在频域上集中；均值与方差一栏对泊松噪声指$y$本身，对加性噪声指$n$。}
 \label{tab:vision-noise-models}
\end{table*}

真实传感器的噪声介于高斯与泊松之间：\ref{sec:vision-formation-sensor}节已说明，光子到达的随机性产生散粒噪声，读出电路再叠加与信号无关的读出噪声。这里把它写成读数层面的模型。传感器把$x/\alpha$个期望光子转换成读数，光子数服从泊松分布，读出电路再叠加方差为$\sigma^2$的高斯噪声，于是读数可写成
\begin{equation}
 y=\alpha p+n,\qquad
 p\sim\operatorname{Poisson}(x/\alpha),\quad
 n\sim\mathcal{N}(0,\sigma^2)\text{，}
 \label{eq:vision-poisson-gaussian}
\end{equation}
其中$\alpha>0$是每个光子对应的读数增益，$p$与$n$独立。泊松分布的均值和方差都等于参数（可参考\BookChapterRef{chap:random-variables}{卷二“随机变量”章}），所以$\mathbb{E}[y]=\alpha\cdot x/\alpha=x$，$\operatorname{Var}[y]=\alpha^2\cdot x/\alpha+\sigma^2=\alpha x+\sigma^2$。方差随亮度线性增长，这就是\term{泊松—高斯噪声}（Poisson-Gaussian noise）模型\scite{vision-foi2008}暗处的信噪比因此特别低：取$\alpha=1$、$\sigma=2$，亮度4处的标准差为$\sqrt{8}\approx2.83$，信噪比约1.4；亮度100处的标准差约10.2，信噪比约9.8。

方差随均值变化，会使只为固定$\sigma$设计的去噪器在亮处去噪过度、在暗处去噪不足。\term{方差稳定化变换}（variance-stabilizing transform, VST）寻找一个单调函数$f$，使$f(y)$的方差近似为常数。依据是Delta方法：若$y$集中在均值$\mu$附近，一阶展开$f(y)\approx f(\mu)+f'(\mu)(y-\mu)$给出$\operatorname{Var}[f(y)]\approx f'(\mu)^2\operatorname{Var}[y]$（可参考前文已引的卷二“随机变量”章）。要让它等于1，就要$f'(\mu)=1/\sqrt{\alpha\mu+\sigma^2}$，积分得
\begin{equation}
 f(y)=\frac{2}{\alpha}\sqrt{\alpha y+\sigma^2}\text{。}
 \label{eq:vision-vst}
\end{equation}
纯泊松情形$\alpha=1,\sigma=0$时，它就是$2\sqrt{y}$；Anscombe把它修正为$2\sqrt{y+3/8}$，使小计数时的方差更接近1\scite{vision-anscombe1948}以纯泊松为例：均值4时$\operatorname{Var}[y]=4$，$f'(4)=1/2$，变换后方差约$4\times\frac14=1$；均值100时$\operatorname{Var}[y]=100$，$f'(100)=1/10$，变换后方差仍约1。使用流程是先做变换，用任意针对单位方差高斯噪声的去噪器处理，再做逆变换。一阶近似在计数很小时偏差较大，而且$f$是非线性的，直接代数求逆会带来系统偏差，需要专门设计的无偏逆变换。

已知噪声类型后，可以选择对应的滤波器。\ref{sec:vision-transforms-nonlinear-denoising}节已经讲过：均值和高斯滤波对独立的高斯噪声有效，平均$k$个独立样本使方差降为$1/k$，代价是边缘一起被平滑；中值滤波输出窗口内的中位数，一个窗口中只要被污染的像素少于一半，中位数就仍来自未被污染的像素，因此适合椒盐噪声；双边滤波和非局部均值按像素值或图块的相似程度加权，能在平滑噪声时保留边缘。周期噪声在频域上集中成尖峰，可以用陷波滤波器只去掉那几个频率（\ref{sec:vision-transforms-selective-filtering}节）。

固定窗口的中值滤波有两个限制。椒盐噪声的密度较高时，窗口中被污染的像素可能超过一半，中位数本身就是噪声值；而在噪声密度较低时，它仍会替换每个像素，包括没有被污染的像素，细小笔画因此变细或断开。\term{自适应中值滤波}（adaptive median filter）让窗口大小随局部情况变化，并且只替换被判为脉冲的像素\scite{vision-hwang-haddad1995,vision-gonzalez-woods2018}记窗口$W$内的最小值、中位数和最大值为$z_{\min},z_{\mathrm{med}},z_{\max}$，当前像素值为$z$，算法\ref{alg:vision-adaptive-median}分两个阶段。

\begin{algorithm}[自适应中值滤波（单个像素）]
 \label{alg:vision-adaptive-median}
 \begin{algorithmic}[1]
  \Require 带噪图像，当前位置的像素值$z$，初始窗口边长$s_0$，最大边长$s_{\max}$
  \Ensure 当前位置的输出值
  \State $s\gets s_0$
  \Loop
   \State 计算边长为$s$的窗口内的$z_{\min},z_{\mathrm{med}},z_{\max}$
   \If{$z_{\min}<z_{\mathrm{med}}<z_{\max}$}
    \Comment{阶段一：中位数不是脉冲}
    \If{$z_{\min}<z<z_{\max}$}
     \State \Return $z$ \Comment{阶段二：当前像素不是脉冲，保留}
    \Else
     \State \Return $z_{\mathrm{med}}$
    \EndIf
   \EndIf
   \State $s\gets s+2$
   \If{$s>s_{\max}$}
    \State \Return $z_{\mathrm{med}}$ \Comment{窗口已达上限}
   \EndIf
  \EndLoop
 \end{algorithmic}
\end{algorithm}

阶段一检查中位数是否落在最小值与最大值之间。若中位数等于最小值或最大值，说明窗口里的脉冲太多，中位数不可靠，就扩大窗口重算。阶段二检查当前像素本身：它若严格介于最小值与最大值之间，就被判为未受污染而原样保留，否则用中位数替换。以$3\times3$窗口为例，窗口值为$\{0,255,120,118,z,122,119,0,121\}$。若中心$z=255$，排序后$z_{\min}=0$、$z_{\mathrm{med}}=120$、$z_{\max}=255$，阶段一通过；阶段二中$z=z_{\max}$，判为“盐”，输出120。若中心$z=121$，排序后$z_{\min}=0$、$z_{\mathrm{med}}=119$、$z_{\max}=255$，阶段一通过；阶段二中$0<121<255$，输出121，而普通中值滤波会把它改成119。若窗口里有五个0，中位数等于$z_{\min}$，阶段一不通过，窗口扩到$5\times5$再判断。这个规则假设脉冲取极值，高斯噪声下没有这种可分离的结构，阶段二几乎总会保留原值。

这些经典滤波器的共同特点是：滤波规则由噪声模型推出，参数只有窗口大小、$\sigma$等几个；它们对指定的噪声有效，却不知道图像中哪些结构是字画、哪些是纹理，噪声与细节频带重叠时只能二者一起处理。学习方法要改变的正是这一点。

在深度网络之前，判别式去噪器通常为每个噪声强度单独训练一个模型；遇到强度未知的图像，要先估计$\sigma$，再选用对应模型，估计误差会直接传到结果中。另一方面，让很深的卷积网络直接输出清晰图像，网络既要复制大部分图像内容，又要改掉噪声，训练时容易收敛缓慢。DnCNN针对这两点提出：让网络只预测观测与参考之差，即噪声残差，并用批归一化稳定深层网络的训练\scite{vision-dncnn2017}

残差预测的直觉是这样的：带噪图像与清晰图像的大部分内容相同，二者之差是一张接近零均值、结构简单的噪声图。网络若只需输出这张差图，就不必在每一层保存并复制整幅图像。这一想法与ResNet的残差连接相关，但位置不同：ResNet在网络内部逐块学习残差（见\ref{sec:vision-ivm-resnet}节），DnCNN则把整个网络的输出定义为残差，再在网络外做一次减法。

DnCNN的结构是一串同尺寸卷积，表\ref{tab:vision-dncnn-layers}逐层列出。输入为$h\times w\times c$的带噪图像（灰度$c=1$，彩色$c=3$）。第一层用64个$3\times3\times c$卷积加ReLU，得到$h\times w\times64$的特征；第2层到第$D-1$层都按“64个$3\times3\times64$卷积 \ensuremath{\to} 批归一化 \ensuremath{\to} ReLU”的顺序计算，特征形状保持$h\times w\times64$；第$D$层用$c$个$3\times3\times64$卷积回到$h\times w\times c$，不加激活函数，因此输出可正可负。每层卷积前都做零填充，使空间尺寸不变，相减能逐位置执行。每经过一层$3\times3$卷积，一个输出像素所依赖的输入范围向四周各扩大一个像素，深度为$D$的网络的感受野因此为$(2D+1)\times(2D+1)$。作者参照当时主流去噪方法实际使用的图块大小来定深度：已知噪声强度的DnCNN-S取$D=17$，感受野为$35\times35$；盲去噪的DnCNN-B和多任务的DnCNN-3取$D=20$，感受野为$41\times41$。

\begin{table}[htbp]
 \centering
 \small
 \begin{tabularx}{\linewidth}{@{}p{14mm}Xp{14mm}p{10mm}p{16mm}p{18mm}@{}}
  \toprule
  层组 & 层类型 & 卷积核与步长 & 输出通道 & 重复次数 & 输出尺寸\\
  \midrule
  第1层 & 卷积，ReLU & $3\times3$，1 & 64 & 1 & $h\times w\times64$\\
  中间层 & 卷积，批归一化，ReLU & $3\times3$，1 & 64 & $D-2$（15或18） & $h\times w\times64$\\
  第$D$层 & 卷积 & $3\times3$，1 & $c$ & 1 & $h\times w\times c$\\
  网络外 & 观测减去残差 & — & $c$ & 1 & $h\times w\times c$\\
  \bottomrule
 \end{tabularx}
 \caption{DnCNN的逐层结构。DnCNN-S取$D=17$，DnCNN-B与DnCNN-3取$D=20$；所有卷积均零填充以保持空间尺寸，输入可为任意$h\times w$。论文未给出参数总量。}
 \label{tab:vision-dncnn-layers}
\end{table}

记这一残差网络为$r_{\bm{\theta}}$，由$n$对带噪图像与清晰参考训练：
\begin{equation}
 L(\bm{\theta})=\frac{1}{2n}\sum_{i=1}^{n}
 \bigl\lVert r_{\bm{\theta}}(\bm{Y}_i)
             -(\bm{Y}_i-\bm{X}_i)\bigr\rVert_{\mathrm{F}}^2,
 \qquad
 \widehat{\bm{X}}=\bm{Y}-r_{\bm{\theta}}(\bm{Y})\text{。}
 \label{eq:vision-dncnn}
\end{equation}
这里的Frobenius范数表示所有像素与通道的平方和。训练配对由清晰图像加合成高斯噪声得到，灰度模型用伯克利分割数据集中400幅$180\times180$的图像。DnCNN-S分别在$\sigma=15,25,50$下各训一个模型，裁取$128\times1600$个$40\times40$的图块；DnCNN-B在$\sigma\in[0,55]$内随机取噪声强度，裁取$128\times3000$个$50\times50$的图块，只训一个模型。优化采用带动量0.9、权重衰减$10^{-4}$的随机梯度下降，批量128，训练50轮，学习率从$10^{-1}$按指数衰减到$10^{-4}$。训练更新卷积参数和批归一化参数；使用时固定这些参数，输入一幅带噪图像，一次前向得到残差，再按式\eqref{eq:vision-dncnn}相减，不再需要清晰参考，DnCNN-B也不需要事先估计$\sigma$。

以招牌一行的四个灰度值为例，参考为$(0.2,0.2,0.8,0.8)$，加噪后得到$(0.25,0.17,0.84,0.74)$，残差标签就是$(0.05,-0.03,0.04,-0.06)$。若网络预测$(0.04,-0.02,0.03,-0.05)$，相减得到$(0.21,0.19,0.81,0.79)$。若网络把明暗跳变也解释成噪声，输出会抹掉字画；残差形式本身没有提供“噪声”和“真实细节”的语义区分，这一区分来自训练数据及损失。

平方误差还决定了输出的倾向。对固定观测$\bm{Y}=\bm{y}$和任意候选$\bm{a}$，有
\begin{equation}
 \mathbb{E}\!\left[\lVert\bm{X}-\bm{a}\rVert_{\mathrm{F}}^2
                  \mid\bm{Y}=\bm{y}\right]
 =
 \mathbb{E}\!\left[\lVert\bm{X}-\bm{\mu}\rVert_{\mathrm{F}}^2
                  \mid\bm{Y}=\bm{y}\right]
 +\lVert\bm{a}-\bm{\mu}\rVert_{\mathrm{F}}^2\text{，}
 \label{eq:vision-conditional-mean}
\end{equation}
其中$\bm{\mu}=\mathbb{E}[\bm{X}\mid\bm{Y}=\bm{y}]$，并假设二阶矩有限。把$\bm{X}-\bm{a}$写成$(\bm{X}-\bm{\mu})+(\bm{\mu}-\bm{a})$后展开平方，交叉项$2\langle\bm{X}-\bm{\mu},\bm{\mu}-\bm{a}\rangle$在给定$\bm{Y}=\bm{y}$时的期望为零，因为$\bm{\mu}-\bm{a}$是常量，而$\bm{X}-\bm{\mu}$的条件均值为零。所以理想的平方误差预测就是条件均值。若一个细节存在多个可能位置，平均结果会较平滑；这一点在噪声较轻时影响不大，在后文超分辨率中会变成主要问题。

DnCNN-S与DnCNN-B的对照说明了“已知条件”的作用。前者为每个噪声强度单独训练，在该强度下得到最好的结果，但使用前必须知道$\sigma$；后者用一个模型覆盖$[0,55]$，省去了噪声估计，代价是要在一个参数集合里兼顾多种强度。与经典滤波相比，DnCNN同样依赖噪声模型，只是这个模型体现在训练数据中而不是滤波公式中：两个版本都只见过加性高斯噪声。真实相机的噪声服从式\eqref{eq:vision-poisson-gaussian}一类与亮度相关的分布，经过去马赛克和压缩后又有空间相关性，超出训练分布时，网络可能把噪声当成纹理保留，也可能把细纹理当成噪声抹去。一个折中的做法是先用式\eqref{eq:vision-vst}把泊松—高斯噪声稳定成近似单位方差，再交给按高斯噪声训练的网络，但相关性问题仍然存在。把多种退化统一写成$\bm{Y}-\bm{X}$只统一了输出接口，不能推得它们具有同一分布。

\subsection{模糊观测下的复原}
\label{sec:vision-restoration-deblurring}

模糊不是逐像素独立的扰动。相机运动或失焦使一个场景点的光落到多个像素上，观测中每个像素都是邻域的加权和。描述这种扩散的核称为点扩散函数，即一个理想亮点在图像上留下的光斑。这一小节的输入是模糊图像，输出是清晰估计；难点是要把混在一起的贡献拆开，而某些频率成分在模糊后几乎为零，直接反演会把噪声成倍放大。

能否知道模糊核，决定了路线。若模糊核已知且空间不变，例如由陀螺仪记录的相机运动或标定过的镜头，恢复问题可以写成线性方程，在傅里叶域逐频率求解，不需要训练数据；逆滤波、维纳滤波和约束最小二乘滤波的差别只在于怎样处理弱频率。若模糊核未知，经典的盲反卷积先从图像中估计核，再回到已知核的情形；学习方法则从大量“模糊—清晰”配对中直接学习映射，让网络从图像本身的边缘和纹理判断该怎样修正。前一条路线的知识全在成像公式里，后一条路线的知识越来越多地来自训练数据，下面两部分分别展开。

\subsubsection{已知模糊核下的复原}
\label{sec:vision-restoration-deconvolution}

先考虑一维信号，假设模糊空间不变，并采用周期边界。卷积可写为$\bm{y}=\bm{H}\bm{x}+\bm{n}$，其中$\bm{x},\bm{y}\in\mathbb{R}^{m}$，$\bm{H}\in\mathbb{R}^{m\times m}$是由模糊核生成的循环矩阵。离散傅里叶变换将循环矩阵对角化，即卷积定理：周期卷积在频域中变成逐频率相乘（可参考\BookChapterRef{chap:signal-analysis}{卷一“信号”章}，二维情形见\ref{sec:vision-transforms-dft2d}节）。于是每个频率$\omega$满足
\begin{equation}
 Y_\omega=H_\omega X_\omega+N_\omega\text{，}
 \label{eq:vision-blur-frequency}
\end{equation}
大写字母表示相应向量的离散傅里叶系数。二维图像的推导完全相同，只是频率换成二维下标$(u,v)$。

最直接的做法是\term{逆滤波}（inverse filtering），逐频率除以$H_\omega$：
\[
 \widehat X_\omega=\frac{Y_\omega}{H_\omega}=X_\omega+\frac{N_\omega}{H_\omega}\text{。}
\]
第二个等号说明，逆滤波在无噪声时精确恢复原信号，在有噪声时把噪声分量除以$H_\omega$。模糊核是低通的，高频处$|H_\omega|$很小，噪声在这些频率被大幅放大；$H_\omega=0$处则根本没有唯一的反演。逆滤波失败的原因不是计算不准，而是它把“观测一致”当作唯一要求，没有告诉算法在弱频率上该相信什么。

\term{维纳滤波}（Wiener filtering）补上的是信号和噪声的统计。设$\bm{x}$和$\bm{n}$是零均值、互不相关的随机信号，并且在周期边界下是平稳的，即协方差只依赖位置差。这时它们的协方差矩阵也是循环矩阵，同样被离散傅里叶变换对角化，于是不同频率的系数互不相关，各频率的方差
\[
 S_x(\omega)=\mathbb{E}\bigl[|X_\omega|^2\bigr],\qquad
 S_n(\omega)=\mathbb{E}\bigl[|N_\omega|^2\bigr]
\]
称为信号与噪声的\term{功率谱}（power spectrum），描述能量在各频率上的分布。白噪声的功率谱是常数$S_n(\omega)=m\sigma^2$；自然图像的功率谱通常随频率升高而下降。

维纳滤波在线性估计$\widehat X_\omega=W_\omega Y_\omega$中选择$W_\omega$，使均方误差$\mathbb{E}|X_\omega-W_\omega Y_\omega|^2$最小；不同频率互不相关，所以总误差可以逐频率分别最小化。由式\eqref{eq:vision-blur-frequency}和不相关假设，
\[
 \mathbb{E}\bigl[X_\omega\overline{Y_\omega}\bigr]=\overline{H_\omega}S_x,
 \qquad
 \mathbb{E}\bigl[|Y_\omega|^2\bigr]=|H_\omega|^2S_x+S_n=:P_\omega\text{，}
\]
这里省略了$S_x,S_n$的自变量$\omega$，第一个式子用到$\mathbb{E}[X_\omega\overline{N_\omega}]=0$。展开平方并配方：
\[
 \begin{aligned}
 \mathbb{E}|X_\omega-W_\omega Y_\omega|^2
 &=S_x-W_\omega H_\omega S_x-\overline{W_\omega}\,\overline{H_\omega}S_x+|W_\omega|^2P_\omega\\
 &=P_\omega\Bigl|W_\omega-\frac{\overline{H_\omega}S_x}{P_\omega}\Bigr|^2
  +\frac{S_xS_n}{P_\omega}\text{。}
 \end{aligned}
\]
最后一步把含$W_\omega$的项写成一个非负的平方，可以展开右边逐项核对。第二项与$W_\omega$无关，所以误差在平方项为零时最小：
\begin{equation}
 W_\omega=\frac{\overline{H_\omega}S_x}{|H_\omega|^2S_x+S_n}
 =\frac{\overline{H_\omega}}{|H_\omega|^2+S_n/S_x}\text{，}
 \qquad
 \text{最小误差}=\frac{S_xS_n}{|H_\omega|^2S_x+S_n}\text{。}
 \label{eq:vision-wiener}
\end{equation}

式\eqref{eq:vision-wiener}的分母比逆滤波多了噪声与信号的功率比$S_n/S_x$。信噪比高的频率上，$S_n/S_x\ll|H_\omega|^2$，维纳滤波接近逆滤波；信噪比低的频率上，分母由$S_n/S_x$主导，增益趋于零，滤波器宁可输出零（即该频率的先验均值），也不放大噪声。没有噪声时$S_n=0$，它退化为逆滤波；$H_\omega=0$时增益为零，不再出现除以零。若$\bm{x}$和$\bm{n}$进一步服从高斯分布，线性估计就是条件均值，维纳滤波也就是式\eqref{eq:vision-conditional-mean}意义下所有估计中的最优者；没有高斯假设时，它只在线性估计中最优。

实际使用时，$S_x$未知，因为要恢复的正是$\bm{x}$。常见的近似是把比值换成一个常数$K$，得到$W_\omega=\overline{H_\omega}/(|H_\omega|^2+K)$，再在带参考的样本上挑选$K$。这个近似忽略了自然图像功率谱随频率下降的事实：真实的$S_n/S_x$在高频更大，常数$K$在高频抑制不足、在低频又抑制过度。

\term{约束最小二乘滤波}（constrained least squares filtering）换一种方式提供先验：不需要信号的功率谱，只要求结果平滑，并要求残差与噪声水平相符\scite{vision-gonzalez-woods2018}设$\bm{C}$是一个衡量粗糙程度的差分算子，二维时常取离散拉普拉斯核$\bigl(\begin{smallmatrix}0&-1&0\\-1&4&-1\\0&-1&0\end{smallmatrix}\bigr)$，一维时取周期一阶差分$\bm{D}$。问题写成
\begin{equation}
 \min_{\bm{x}}\ \lVert\bm{C}\bm{x}\rVert_2^2
 \quad\text{满足}\quad
 \lVert\bm{y}-\bm{H}\bm{x}\rVert_2^2=\lVert\bm{n}\rVert_2^2\text{，}
 \label{eq:vision-cls}
\end{equation}
即在所有“残差大小恰好等于噪声大小”的候选中，选最平滑的一个。约束的意思是：残差比噪声还小，说明结果把噪声也拟合进去了；残差比噪声大，说明信号没有被充分利用。引入拉格朗日乘子，问题的驻点满足
\[
 \widehat{\bm{x}}=\argmin_{\bm{x}}\
 \lVert\bm{H}\bm{x}-\bm{y}\rVert_2^2+\gamma\lVert\bm{C}\bm{x}\rVert_2^2\text{，}
\]
$\gamma>0$由约束确定。这个目标是二次函数，令梯度为零得到法方程$(\bm{H}^{\top}\bm{H}+\gamma\bm{C}^{\top}\bm{C})\widehat{\bm{x}}=\bm{H}^{\top}\bm{y}$。$\bm{H}$和$\bm{C}$都是循环矩阵，可由同一傅里叶基对角化，$\bm{H}^{\top}$对应复共轭$\overline{H_\omega}$；在分母非零时，每个频率独立求解为
\begin{equation}
 \widehat X_\omega=
 \frac{\overline{H_\omega}}
 {|H_\omega|^2+\gamma|C_\omega|^2}\,Y_\omega\text{。}
 \label{eq:vision-regularized-inverse}
\end{equation}
与式\eqref{eq:vision-wiener}对照，约束最小二乘就是把$S_n/S_x$换成$\gamma|C_\omega|^2$的维纳滤波。差分算子在高频处$|C_\omega|$大，相当于假设信号功率谱随频率升高而下降，这正好弥补了常数$K$的缺陷。在$H_\omega$较大的频率，分母由$|H_\omega|^2$主导，式子接近逆滤波；在$H_\omega$很小的频率，分母由$\gamma|C_\omega|^2$主导，增益被压低。若$H_\omega=C_\omega=0$，该频率仍没有唯一解，需要另外规定约束。这一形式也就是反卷积中常用的二次正则化；\term{反卷积}（deconvolution）泛指撤销卷积模糊的计算，正则项的强度在这里由噪声水平确定，而不是人为指定。

剩下的问题是怎样找到满足约束的$\gamma$。记残差$\bm{r}(\gamma)=\bm{y}-\bm{H}\widehat{\bm{x}}(\gamma)$，由式\eqref{eq:vision-regularized-inverse}，
\[
 R_\omega(\gamma)=Y_\omega-H_\omega\widehat X_\omega
 =\frac{\gamma|C_\omega|^2}{|H_\omega|^2+\gamma|C_\omega|^2}\,Y_\omega\text{。}
\]
系数$\gamma|C_\omega|^2/(|H_\omega|^2+\gamma|C_\omega|^2)$在$H_\omega,C_\omega\neq0$时随$\gamma$单调增加，所以残差平方和$\phi(\gamma)=\lVert\bm{r}(\gamma)\rVert_2^2$随$\gamma$单调不减（由Parseval关系，空间与频域的平方和只差常数因子）。单调性保证可以用调整$\gamma$的方式逼近目标。噪声能量$\lVert\bm{n}\rVert_2^2$本身未知，用它的期望$m(\sigma^2+\mu_n^2)$代替，其中$\mu_n,\sigma^2$是噪声的均值和方差，可以从平坦区域估计。算法\ref{alg:vision-cls-gamma}给出一种简单的调整方式。

\begin{algorithm}[约束最小二乘滤波的参数选择]
 \label{alg:vision-cls-gamma}
 \begin{algorithmic}[1]
  \Require 观测$\bm{y}$，模糊核，差分算子$\bm{C}$，噪声能量估计$\eta=m(\sigma^2+\mu_n^2)$，容差$a>0$，初值$\gamma^{(0)}$
  \Ensure 满足$|\phi(\gamma)-\eta|\leq a$的$\gamma$及复原结果
  \State $t\gets0$
  \Loop
   \State 按式\eqref{eq:vision-regularized-inverse}计算$\widehat{\bm{x}}(\gamma^{(t)})$与$\phi(\gamma^{(t)})$
   \If{$|\phi(\gamma^{(t)})-\eta|\leq a$}
    \State \Return $\gamma^{(t)}$与$\widehat{\bm{x}}(\gamma^{(t)})$
   \ElsIf{$\phi(\gamma^{(t)})>\eta+a$}
    \State 减小$\gamma$，例如取已知可行区间的中点 \Comment{残差太大，平滑过度}
   \Else
    \State 增大$\gamma$ \Comment{残差太小，拟合了噪声}
   \EndIf
   \State $t\gets t+1$
  \EndLoop
 \end{algorithmic}
\end{algorithm}

用一个八点例子把三种滤波器放在一起。取边缘$\bm{x}=(0,0,0,0,1,1,1,1)$，模糊核在中心和左右邻点的权重为$(0.6,0.2,0.2)$，再加入交替的$\pm0.04$噪声，得到观测
\[
 \bm{y}=(0.24,-0.04,0.04,0.16,0.84,0.96,1.04,0.76)\text{。}
\]
这个核的频率响应为$H_\omega=0.6+0.4\cos\omega$，最高频$\omega=\pi$处$H_\omega=0.2$，逆滤波的增益为5，把幅值0.04的交替噪声放大到0.2，输出为$(0.2,-0.2,0.2,-0.2,1.2,0.8,1.2,0.8)$。取常数比$K=0.05$的维纳滤波，$\omega=\pi$处增益为$0.2/(0.04+0.05)\approx2.22$，噪声仍被放大一倍多。约束最小二乘取$\bm{C}=\bm{D}$，周期差分在该处$|D_\omega|^2=2-2\cos\pi=4$；取$\gamma=0.05$时，增益为$0.2/(0.04+0.05\times4)\approx0.833$，噪声分量幅值降为约0.033。图\ref{fig:vision-filter-gains}画出三种增益在各频率上的变化：常数比维纳滤波只在高频把逆滤波的增益削去一半左右，差分约束则把高频增益压到1以下。

\begin{figure}[htbp]
 \centering
 % 八点周期模糊核(0.2,0.6,0.2)，H(w)=0.6+0.4cos(w)。三条曲线为逐频率增益的解析值：
% 逆滤波1/H；常数比维纳 H/(H^2+0.05)；一阶差分约束 H/(H^2+0.05(2-2cos w))。宽约95mm。
\begin{tikzpicture}[x=72mm,y=8.5mm,font=\figurefont\small]
 \draw[BookMuted,line width=.7pt,->] (0,0)--(1.08,0) node[right] {$\omega$};
 \draw[BookMuted,line width=.7pt,->] (0,0)--(0,5.5) node[above] {增益};
 \foreach \y in {1,2,3,4,5} {
  \draw[BookRule,line width=.5pt] (0,\y)--(1,\y);
  \node[anchor=east] at (-.01,\y) {$\y$};
 }
 \node[anchor=east] at (-.01,0) {$0$};
 \foreach \x/\lab in {0/0,.5/{\pi/2},1/\pi} {
  \draw[BookMuted,line width=.5pt] (\x,0)--(\x,-.1);
  \node[below] at (\x,-.1) {$\lab$};
 }
 \draw[BookAmber,dashed,line width=1pt] plot coordinates {
  (0,1)(.0833,1.014)(.1667,1.057)(.25,1.133)(.3333,1.25)(.4167,1.421)
  (.5,1.667)(.5833,2.014)(.6667,2.5)(.75,3.153)(.8333,3.943)(.9167,4.681)(1,5)};
 \draw[BookBlue,line width=1pt] plot coordinates {
  (0,.952)(.0833,.964)(.1667,1.001)(.25,1.064)(.3333,1.159)(.4167,1.291)
  (.5,1.463)(.5833,1.675)(.6667,1.905)(.75,2.106)(.8333,2.218)(.9167,2.234)(1,2.222)};
 \draw[BookTeal,line width=1pt] plot[mark=*,mark size=1.2pt] coordinates {
  (0,1)(.0833,1.01)(.1667,1.041)(.25,1.092)(.3333,1.159)(.4167,1.236)
  (.5,1.304)(.5833,1.333)(.6667,1.29)(.75,1.169)(.8333,1.011)(.9167,.882)(1,.833)};
 \draw[BookAmber,dashed,line width=1pt] (.04,6.2)--(.12,6.2);
 \node[anchor=west] at (.13,6.2) {逆滤波};
 \draw[BookBlue,line width=1pt] (.34,6.2)--(.42,6.2);
 \node[anchor=west] at (.43,6.2) {常数比维纳};
 \draw[BookTeal,line width=1pt] (.71,6.2)--(.79,6.2);
 \node[anchor=west] at (.8,6.2) {差分约束};
\end{tikzpicture}

 \caption{八点例中三种滤波器在各频率上的增益。横轴为频率$\omega\in[0,\pi]$，模糊核的频率响应$H_\omega=0.6+0.4\cos\omega$在高频降到0.2。虚线为逆滤波$1/H_\omega$；蓝实线为常数比维纳滤波，$K=0.05$；带点的实线为一阶差分约束，$\gamma=0.05$。增益大于1的频率被放大，三者的差别集中在高频。}
 \label{fig:vision-filter-gains}
\end{figure}

再用算法\ref{alg:vision-cls-gamma}选$\gamma$。噪声能量为$8\times0.04^2=0.0128$，取容差$a=0.001$。从$\gamma^{(0)}=0.05$开始，残差平方和为0.0223，超过$0.0138$，说明平滑过度；减半到$\gamma^{(1)}=0.025$，残差平方和为0.0132，落入$[0.0118,0.0138]$，停止。此时输出为$(0.177,-0.094,0.021,0.062,0.938,0.979,1.094,0.823)$，与原信号的均方误差约为0.0111；$\gamma=0.05$时为0.0159，逆滤波为0.04。在这个例子中，按噪声水平选出的$\gamma$接近使误差最小的取值（逐点搜索得到约0.018，误差约0.0104），但这依赖噪声能量估计准确，以及差分算子与信号结构相符。

图\ref{fig:vision-deconvolution}显示$\gamma=0.05$时的全部八点结果。正则化后的过冲变小，但跃迁也变缓，且仍有轻微振铃。这里保留了负值与大于1的值以展示反演行为，没有在绘图前裁剪。计算上，对$h\times w$的图像，二维版本对每个通道做一次快速傅里叶变换，逐频率相乘，再做逆变换，复杂度为$O(hw\log(hw))$。真实图像若采用反射边界，不能继续照搬循环卷积公式；若模糊核错误或随位置变化，正则化只能缓解数值不稳定，不能修正错误的成像模型。

\begin{figure}[htbp]
 \centering
 % Eight-point circular blur, kernel (.2,.6,.2), alternating noise .04.
% DFT ridge inverse uses first differences and lambda=.05; values are un-clipped.
\begin{tikzpicture}[x=11mm,y=30mm,font=\figurefont\small]
 \draw[BookMuted,line width=.7pt,->] (-.35,-.3)--(7.5,-.3)
   node[right] {$i$};
 \draw[BookMuted,line width=.7pt,->] (-.35,-.3)--(-.35,1.38)
   node[above] {灰度};
 \foreach \y in {0,0.5,1} {
  \draw[BookRule,line width=.5pt] (-.35,\y)--(7.2,\y);
  \node[anchor=east] at (-.45,\y) {$\y$};
 }
 \foreach \i in {0,...,7} {
  \draw[BookMuted,line width=.5pt] (\i,-.3)--(\i,-.34);
  \node[below] at (\i,-.34) {$\i$};
 }
 \draw[BookInk,line width=1.2pt]
  plot coordinates {(0,0)(1,0)(2,0)(3,0)(4,1)(5,1)(6,1)(7,1)};
 \draw[BookAmber,dashed,line width=1pt]
  plot[mark=square*,mark size=1.5pt] coordinates
  {(0,.2)(1,-.2)(2,.2)(3,-.2)(4,1.2)(5,.8)(6,1.2)(7,.8)};
 \draw[BookTeal,line width=1pt]
  plot[mark=*,mark size=1.5pt] coordinates
  {(0,.1997)(1,-.0766)(2,-.01)(3,.133)(4,.867)(5,1.01)(6,1.0766)(7,.8003)};
 \draw[BookInk,line width=1.2pt] (0,1.65)--(.6,1.65);
 \node[anchor=west] at (.7,1.65) {参考};
 \draw[BookAmber,dashed,line width=1pt] (2.1,1.65)--(2.7,1.65);
 \node[anchor=west] at (2.8,1.65) {直接求逆};
 \draw[BookTeal,line width=1pt] (4.8,1.65)--(5.4,1.65);
 \node[anchor=west] at (5.5,1.65) {正则化};
\end{tikzpicture}

 \caption{八点周期边缘的教学反卷积。实线为清晰参考，带方块的虚线为逆滤波，带圆点的实线为一阶差分约束（$\gamma=0.05$）的结果。两端因周期边界相接；纵轴未裁剪，便于观察过冲与平滑的取舍。}
 \label{fig:vision-deconvolution}
\end{figure}
\FloatBarrier

三种滤波器回答同一个问题：弱频率上该相信什么。逆滤波只相信观测，维纳滤波相信信号与噪声的功率谱，约束最小二乘相信“结果应当平滑且残差与噪声相当”。后两者在数学上是同一形式，只是先验的来源不同。这一思想在学习方法中继续存在：学习式复原网络的训练数据隐含地给出了比任何人工功率谱都细致的图像先验，而后文DiffBIR的复原引导，仍是在“与观测一致”和“符合先验”之间作加权。

\subsubsection{未知模糊核下的复原}
\label{sec:vision-restoration-restormer}

手持拍摄的招牌，模糊方向和长度事先不知道，画面中前景与背景的模糊程度还可能不同，式\eqref{eq:vision-regularized-inverse}中的$H_\omega$无从写出。经典的回答是\term{盲反卷积}（blind deconvolution）：同时估计模糊核$\bm{k}$和清晰图像$\bm{x}$。最自然的写法是联合最大后验估计
\begin{equation}
 (\widehat{\bm{x}},\widehat{\bm{k}})=\argmin_{\bm{x},\bm{k}}\
 \lVert\bm{k}*\bm{x}-\bm{y}\rVert_2^2
 +\lambda\sum_{i}\rho\bigl((\nabla\bm{x})_i\bigr)
 \quad\text{满足}\quad k_j\geq0,\ \sum_jk_j=1\text{，}
 \label{eq:vision-blind-map}
\end{equation}
其中$\rho$惩罚图像梯度，非负与归一化约束排除了核的负值和整体缩放。这个问题的困难在于解的多义性：取$\bm{k}$为只在中心取1的单位冲激、$\bm{x}=\bm{y}$，数据项同样为零，这是“没有模糊”的解释。

先验能否排除这个解释，取决于$\rho$。仍用八点边缘，无噪声时模糊结果为$(0.2,0,0,0.2,0.8,1,1,0.8)$，其周期差分为$(-0.2,0,0.2,0.6,0.2,0,-0.2,-0.6)$；清晰边缘的周期差分只有一个$+1$和一个$-1$。若$\rho(g)=g^2$，清晰解的先验代价为2，模糊解只有$4\times0.04+2\times0.36=0.88$，联合估计会偏向“没有模糊”。若改用重尾的$\rho(g)=|g|^{1/2}$，清晰解为2，模糊解为$4\sqrt{0.2}+2\sqrt{0.6}\approx3.34$，清晰解胜出。这说明盲反卷积需要能区分清晰边缘与模糊过渡的稀疏梯度先验\scite{vision-fergus2006}但Levin等人分析指出，在含大量细小纹理的真实图像上，联合估计仍大多偏向无模糊的解释；由于模糊核的未知数远少于像素数，先把图像积分掉、只对核做最大后验估计，问题受到的约束要强得多，能较准确地恢复核\scite{vision-levin2009}实际系统因此常先在由粗到细的图像金字塔上估计核，再固定核，用上一小节的非盲方法恢复图像。它的局限也随之而来：核估计依赖足够的强边缘，且大多假设整幅图共用一个空间不变的核，而Levin等人收集的真实模糊数据显示，这个假设经常不成立。

学习式去模糊绕开显式核，直接学习从模糊图像到清晰图像的映射。卷积网络是自然的起点，但它有两个限制：感受野有限，一个像素只能读取固定大小的邻域，而运动模糊的拖影可能横跨几十个像素；卷积权重在推理时固定，不随当前图像内容调整。自注意力能让每个位置读取全图并按内容加权，但逐像素注意力的计算量随像素数平方增长，一幅$256\times256$图像就有约$6.6\times10^4$个位置；把注意力限制在$8\times8$等小窗口内又回到了局部范围。

Restormer的主要想法是把注意力从空间维换到通道维：不再计算“哪些像素与哪些像素相关”，而是计算“哪些通道在整幅特征图上的响应彼此相关”，用这种通道间的关系矩阵混合特征\scite{vision-restormer2022}关系矩阵的大小只取决于通道数，计算量对位置数近似线性，因而能直接处理高分辨率图像；全图统计又使通道组合随输入变化。空间上的局部信息由注意力前的深度卷积和多尺度结构补充。

整体结构是一个四级对称编码器—解码器，表\ref{tab:vision-restormer-layers}列出各级配置。输入$\bm{Y}\in\mathbb{R}^{h\times w\times3}$先经一层卷积得到$h\times w\times C$的浅层特征，$C=48$。编码器从第一级到第四级依次用像素反重排把空间尺寸减半、通道加倍，特征形状依次为$h\times w\times48$、$\frac{h}{2}\times\frac{w}{2}\times96$、$\frac{h}{4}\times\frac{w}{4}\times192$和$\frac{h}{8}\times\frac{w}{8}\times384$；四级分别含4、6、6、8个Transformer块，注意力头数分别为1、2、4、8。解码器用像素重排逐级上采样，与同级编码特征沿通道拼接，再用$1\times1$卷积把通道减半；第一级例外，保留拼接后的$h\times w\times96$。随后在全分辨率上用4个块细化，最后一层卷积输出残差图$\bm{R}\in\mathbb{R}^{h\times w\times3}$，复原结果为$\widehat{\bm{X}}=\bm{Y}+\bm{R}$。注意这个符号约定与式\eqref{eq:vision-dncnn}相反：DnCNN预测噪声后相减，Restormer预测修正量后相加。

\begin{table*}[tbp]
 \centering
 \small
 \begin{tabularx}{\linewidth}{@{}p{22mm}Xp{24mm}p{12mm}p{12mm}p{28mm}@{}}
  \toprule
  阶段 & 层类型 & 卷积核与步长 & 输出通道 & 重复次数 & 输出尺寸\\
  \midrule
  浅层嵌入 & 卷积 & 论文未给出 & 48 & 1 & $h\times w\times48$\\
  编码第1级 & Transformer块（1头） & — & 48 & 4 & $h\times w\times48$\\
  下采样 & 像素反重排 & 空间减半 & 96 & 1 & $\frac h2\times\frac w2\times96$\\
  编码第2级 & Transformer块（2头） & — & 96 & 6 & $\frac h2\times\frac w2\times96$\\
  编码第3级 & 下采样后Transformer块（4头） & 空间减半 & 192 & 6 & $\frac h4\times\frac w4\times192$\\
  第4级 & 下采样后Transformer块（8头） & 空间减半 & 384 & 8 & $\frac h8\times\frac w8\times384$\\
  解码第3级 & 像素重排，拼接，$1\times1$卷积，Transformer块（4头） & $1\times1$，1 & 192 & 6 & $\frac h4\times\frac w4\times192$\\
  解码第2级 & 同上（2头） & $1\times1$，1 & 96 & 6 & $\frac h2\times\frac w2\times96$\\
  解码第1级 & 像素重排，拼接，不降通道，Transformer块（1头） & — & 96 & 4 & $h\times w\times96$\\
  细化 & Transformer块 & — & 96 & 4 & $h\times w\times96$\\
  输出 & 卷积，与输入相加 & 论文未给出 & 3 & 1 & $h\times w\times3$\\
  \bottomrule
 \end{tabularx}
 \caption{Restormer的逐级结构。块数、头数与通道数取自原文；论文只说明首尾各为“一层卷积”，未给出其卷积核大小，也未在正文中给出完整模型的参数总量。每个Transformer块由通道注意力和门控前馈两部分组成，门控前馈的扩展倍率$\gamma=2.66$。}
 \label{tab:vision-restormer-layers}
\end{table*}

每个Transformer块内的计算顺序为：层归一化 \ensuremath{\to} 通道注意力 \ensuremath{\to} 残差相加 \ensuremath{\to} 层归一化 \ensuremath{\to} 门控前馈 \ensuremath{\to} 残差相加。在一个注意力头内，设位置数$n=hw$，该头通道数$d$。归一化特征先经$1\times1$卷积混合通道，再经$3\times3$深度卷积混合邻域，得到按“通道$\times$位置”排列的$\bm{Q},\bm{K},\bm{V}\in\mathbb{R}^{d\times n}$。沿位置维对每行做$L_2$归一化得到$\bar{\bm{Q}},\bar{\bm{K}}$后，
\begin{equation}
 \bm{A}=\operatorname{softmax}_{\mathrm{row}}
       (\beta\,\bar{\bm{Q}}\bar{\bm{K}}^\top),
 \qquad
 \bm{O}=\bm{A}\bm{V}\text{，}
 \label{eq:vision-restormer-attention}
\end{equation}
其中$\beta$是每头学习的缩放量。$\bm{A}\in\mathbb{R}^{d\times d}$的第$(i,j)$项比较第$i$个通道的查询与第$j$个通道的键在整幅图上的响应模式；$\bm{O}$的每一行是若干通道的加权组合，在每个位置上用同一组权重混合。它不是每个像素直接对所有其他像素加权求和。

计算量可以具体比较。取$n=256\times256$、第一级$d=48$：逐位置注意力需为一个头建立$n^2\approx4.3\times10^9$个关系，通道注意力矩阵只有$48^2=2304$项。仅考虑注意力中的矩阵乘法，通道注意力的时间量级为$O(nd^2)$，关系矩阵存储为$O(d^2)$；特征本身仍需$O(nd)$存储，卷积和其余部分另有计算。说它对位置数近似线性，前提是固定通道数和其他结构。

门控前馈模块把归一化特征送入两个并行支路，每支都是$1\times1$卷积把通道扩大到$\gamma$倍（$\gamma=2.66$）再接$3\times3$深度卷积；一支经GELU后与另一支逐元素相乘，再用$1\times1$卷积投影回原通道数。乘法让一个支路调制另一个支路的响应，决定哪些特征继续向后传递；门值不限于$[0,1]$，不应理解成概率。

训练使用模糊—清晰配对，损失为逐像素$L_1$距离$\lVert\widehat{\bm{X}}-\bm{X}\rVert_1$，更新全部参数。优化器为AdamW（$\beta_1=0.9,\beta_2=0.999$，权重衰减$10^{-4}$），学习率从$3\times10^{-4}$按余弦退火降到$10^{-6}$，共30万次迭代。作者还采用渐进式训练：开始时裁取$128\times128$的图块、批量64，在第9.2万、15.6万、20.4万、24万和27.6万次迭代时依次改为$160^2$、$192^2$、$256^2$、$320^2$、$384^2$的图块，批量相应降为40、32、16、8、8。小块训练快，大块让网络在后期看到更大范围的上下文，接近推理时的整图输入。运动去模糊模型只在GoPro数据集上训练，再直接用于HIDE和RealBlur测试；论文报告它在GoPro上的PSNR为32.92 dB，高于此前的MPRNet的32.66 dB。推理时输入整幅模糊图像，一次前向得到结果，不需要模糊核。去噪、去雨、运动模糊和失焦各自使用不同的训练数据和检查点，同一个架构名不能代替检查点的适用范围。

三条路线可以放在同一张模糊招牌上对照，表\ref{tab:vision-deblur-compare}列出主要差别。非盲反卷积的全部假设写在公式里，结果可以逐频率解释，但核一旦错误或随位置变化，公式本身就错了；盲反卷积把核的估计交给图像先验，仍受空间不变假设的限制；Restormer不需要核，能处理非均匀模糊，代价是训练数据决定了它认识哪些模糊，超出GoPro一类运动模糊时没有保证，并且网络可能补出看似合理但观测并不支持的笔画。三种结果都应做同一项检查：把估计图像按已知或估计的退化重新模糊，看是否与观测一致。

\begin{table*}[tbp]
 \centering
 \small
 \begin{tabularx}{\linewidth}{@{}p{16mm}XXX@{}}
  \toprule
  比较维度 & 已知核下的约束最小二乘 & 盲反卷积 & Restormer\\
  \midrule
  所需信息 & 准确的模糊核、边界条件与噪声水平 & 足够的强边缘与梯度先验 & 与目标退化相近的模糊—清晰配对\\
  可处理的模糊 & 空间不变、核已知 & 空间不变、核未知 & 训练分布内的非均匀模糊\\
  参数与训练 & 只需选择$\gamma$ & 先验权重与金字塔层数 & 全部网络参数，30万次迭代\\
  推理计算 & 每通道一次FFT与逆FFT & 迭代估计核，再非盲复原 & 一次深层网络前向\\
  主要失效 & 核错误、零响应频率、振铃 & 偏向无模糊解、核随位置变化 & 分布外模糊、补出无观测支持的细节\\
  \bottomrule
 \end{tabularx}
 \caption{模糊观测下三条复原路线的比较。三列面对同一个难点，即弱频率成分无法稳定反演，区别在于约束来自显式成像公式、手工图像先验，还是训练配对。}
 \label{tab:vision-deblur-compare}
\end{table*}

\subsection{分辨率不足下的复原}
\label{sec:vision-restoration-super-resolution}

\term{超分辨率}（super-resolution）根据较低分辨率的输入估计更高分辨率的图像。输入是$h\times w\times3$的图像，输出是$sh\times sw\times3$的图像，$s$为放大倍数。放大网格增加了输出位置，却没有自动增加独立观测。例如把两个像素平均成一个时，$(0,1)$与$(0.5,0.5)$都会得到0.5。输出中的细线和文字怎样出现，取决于参考监督、退化假设与图像先验，而非输出像素数本身。最简单的经典做法是双三次插值（\ref{sec:vision-formation-interpolation}节），它只用相邻采样的加权和填补新网格，不补任何新信息。

这里的难点比去噪更尖锐：一个低清观测对应大量可能的高清图像，按式\eqref{eq:vision-conditional-mean}，像素平方误差最优的估计是它们的平均，细节会被抹平。三个模型对这一难点给出了逐步加强的回答。DnCNN的超分辨率变体沿用像素损失，在简单的双三次退化下回归条件均值，回答是“接受平滑，只求平均误差小”。Real-ESRGAN扩大合成退化的范围，并加入感知损失和对抗损失，让网络输出真实图像中常见的纹理，回答是“用判别器约束输出像真实图像”。DiffBIR先用复原模块去掉退化，再借助在大规模图像上预训练的扩散模型生成细节，回答是“从一个强大的图像先验中采样”。三者补出的细节越来越多，与观测的关系也越来越需要单独核验。

\subsubsection{DnCNN超分辨率变体}
\label{sec:vision-restoration-dncnn-sr}

复用DnCNN的残差网络前，需要先让输入与目标处在同一网格。由高分辨率参考$\bm{X}$合成低分辨率图像$\bm{Y}=\mathcal{S}(\bm{X})$，再用双三次插值$\mathcal{U}$放大为$\bm{B}=\mathcal{U}(\bm{Y})$，$\bm{B}$与$\bm{X}$尺寸相同。网络结构与表\ref{tab:vision-dncnn-layers}相同，取$D=20$：第一层卷积加ReLU，18层“卷积—批归一化—ReLU”，最后一层卷积输出$c$通道残差，感受野$41\times41$。网络学习$\bm{B}-\bm{X}$，使用时得到$\widehat{\bm{X}}=\bm{B}-r_{\bm{\theta}}(\bm{B})$。残差现在含插值和平滑带来的结构差异，主要集中在边缘和纹理处，不再是独立高斯噪声。所有计算都在高分辨率网格上进行，放大$s$倍时计算量也随之增大约$s^2$倍，这是先插值再处理的代价。

作者把这一变体与盲去噪、JPEG去块效应合在一起训练成DnCNN-3。训练图像为91幅常用超分辨率图像加上伯克利分割数据集的200幅训练图像；超分辨率输入由双三次下采样再双三次上采样得到，倍数取2、3、4；JPEG输入的质量因子取5到99；噪声输入的$\sigma\in[0,55]$。三类输入共裁取$128\times8000$个$50\times50$的图块对，参数从DnCNN-B初始化，损失仍为式\eqref{eq:vision-dncnn}。使用时不需要告诉网络当前属于哪类退化或哪个倍数，但必须先做与训练相同的双三次放大。

为了看清监督的含义，暂用两点平均和常数插值作教学简化。原细线为$(0,1)$，观测为0.5，插值得到$(0.5,0.5)$，残差标签为$(0.5,-0.5)$。若训练数据中细线多在右侧，网络学到的预测残差能恢复这一排列；若真实原图为$(1,0)$，相同观测需要相反残差。唯一低分辨率观测无法区分两者，平方误差下网络会输出两种残差的加权平均，也就是保留部分平滑。

像素监督衡量输出与训练参考的差。若应用还要求输出服从特定退化过程，应计算$\mathcal{S}(\widehat{\bm{X}})$并与$\bm{Y}$比较，这称为对观测的回投影核验。上述两个排列都能通过平均采样的回投影，所以观测一致性是必要依据，却不足以确认缺失细节。这一变体还有一个更直接的局限：训练只用了干净的双三次下采样，而网上的照片往往先模糊、再缩小、又多次JPEG压缩，输入偏离训练分布时，网络学到的残差不再适用。

\subsubsection{Real-ESRGAN}
\label{sec:vision-restoration-real-esrgan}

网上流传的招牌照片可能先被缩小、保存成JPEG，再截取、放大和重新压缩，还可能被手机自动锐化，在字边出现白边和振铃。只用一次双三次下采样训练的网络遇到这种输入，要么保留这些伪影，要么把它们放大。Real-ESRGAN要解决两件事：让训练退化尽量覆盖真实照片的退化，并在细节无法唯一确定时输出清晰而非平滑的纹理\scite{vision-real-esrgan2021}

第一件事靠\term{高阶退化模型}（high-order degradation model），即把“模糊—缩放—加噪—JPEG压缩”这一经典退化过程重复执行，每轮参数随机抽取。原方法用两轮，一对训练样本写成$\bm{Y}=\mathcal{D}_2(\mathcal{D}_1(\bm{X}))$与$\bm{X}$。每轮的模糊核以0.7、0.15、0.15的概率取自各向同性或各向异性高斯核、广义高斯核和平台形核，核尺寸在7到21之间，标准差在$[0.2,3]$内（第二轮为$[0.2,1.5]$），另以0.1的概率改用sinc核，第二轮还以0.2的概率跳过模糊。缩放可以向上或向下，插值方式随机。噪声以各半的概率取高斯噪声或泊松噪声，以0.4的概率加成灰度噪声。JPEG质量因子取自$[30,95]$。sinc滤波器是截断高频的理想低通滤波器，在锐利边缘附近产生振铃和过冲，用来模拟锐化及压缩带来的白边；最后一步以0.8的概率施加sinc滤波，它与最后一次JPEG压缩的先后顺序随机交换。这里的“高阶”指重复退化过程，不表示对图像求高阶导数。

第二件事靠感知损失与对抗损失。\term{感知损失}（perceptual loss）不逐像素比较，而是把两幅图都送入一个固定的预训练分类网络，比较中间层特征的差别，从而对纹理和结构敏感、对轻微位移不那么敏感。对抗损失来自生成对抗网络：判别器学习区分真实高清图与生成结果，生成器则学习让判别器分辨不出；生成对抗网络的训练原理可参考\BookChapterRef{chap:neural-generative-models}{卷五“基于神经网络的生成模型”章}。这两项让网络不再满足于条件均值，而是输出真实图像中常见的锐利纹理。

生成器沿用ESRGAN的结构\scite{vision-esrgan2018}表\ref{tab:vision-esrgan-generator}列出各段。它保留SRResNet的总体布局：浅层卷积之后堆叠基本块，整个主干在低分辨率上计算，末端上采样后用两层卷积重建。ESRGAN把基本块换成残差中的残差密集块（residual-in-residual dense block, RRDB），并去掉全部批归一化层，因为批归一化在训练与测试统计不同时会引入伪影。一个RRDB由3个密集块串联而成，每个密集块含5层卷积，前4层后接LeakyReLU，每层的输入是块内所有前面层输出的拼接；密集块的输出乘以残差缩放系数0.2后加回块的输入，3个密集块之外再套一层同样缩放的残差连接，“残差中的残差”即指这两级残差。较深的模型用23个RRDB。论文没有给出各卷积的通道数、增长通道数、卷积核大小和上采样的具体实现。对$\times4$放大，输入$h\times w\times3$直接进入主干，输出$4h\times4w\times3$。为了用同一主干做$\times2$和$\times1$，Real-ESRGAN先用像素反重排把输入变成$\frac{h}{2}\times\frac{w}{2}\times12$（$\times2$）或$\frac{h}{4}\times\frac{w}{4}\times48$（$\times1$），使大部分计算发生在较小的空间尺寸上。

\begin{table}[htbp]
 \centering
 \small
 \begin{tabularx}{\linewidth}{@{}p{14mm}Xp{12mm}p{10mm}p{18mm}@{}}
  \toprule
  阶段 & 层类型 & 卷积核与步长 & 重复次数 & 输出尺寸（$\times4$）\\
  \midrule
  输入整理 & 像素反重排（仅$\times2$、$\times1$） & — & 1 & 见正文\\
  浅层 & 卷积 & 未给出 & 1 & $h\times w$\\
  主干 & RRDB：3个密集块，每块5层卷积与4个LeakyReLU，残差缩放0.2 & 未给出 & 23 & $h\times w$\\
  主干末 & 卷积 & 未给出 & 1 & $h\times w$\\
  上采样 & 上采样$\times4$ & 未给出 & 1 & $4h\times4w$\\
  重建 & 卷积 & 未给出 & 2 & $4h\times4w\times3$\\
  \bottomrule
 \end{tabularx}
 \caption{ESRGAN与Real-ESRGAN生成器的结构。块数和块内组成取自ESRGAN原文的较深配置；“未给出”指原文未写明卷积核与通道数。}
 \label{tab:vision-esrgan-generator}
\end{table}

判别器由ESRGAN的VGG式结构改为U形网络：先逐级下采样提取特征，再逐级上采样并通过跳跃连接引入同级特征，输出与输入同尺寸的单通道真实度图，对每个像素给出反馈，而不是只给整幅图一个分数。复杂退化和U形结构都会使训练不稳，因此判别器的卷积层加了谱归一化，限制每层权重的谱范数。论文只用示意图给出这一结构，没有列出下采样级数和各层通道数。

训练分两个阶段。第一阶段只用$L_1$像素损失训练生成器，得到的模型称为Real-ESRNet，它从ESRGAN的参数微调而来，训练100万次迭代，学习率$2\times10^{-4}$。第二阶段以Real-ESRNet初始化生成器，联合训练生成器与判别器40万次迭代，学习率$10^{-4}$，生成器的目标为
\begin{equation}
 L_G=\bigl\lVert G(\bm{Y})-\bm{X}\bigr\rVert_1
 +\sum_{\ell=1}^{5}\eta_\ell
  \bigl\lVert\phi_\ell(G(\bm{Y}))-\phi_\ell(\bm{X})\bigr\rVert_1
 +0.1\,L_{\mathrm{adv}}\text{，}
 \label{eq:vision-real-esrgan-loss}
\end{equation}
其中$G$是生成器，$\phi_\ell$是预训练VGG19第$\ell$段卷积（conv1至conv5）在激活前的特征，权重$(\eta_1,\dots,\eta_5)=(0.1,0.1,1,1,1)$，$L_{\mathrm{adv}}$是生成器的对抗损失，由U形判别器给出的逐像素真实度计算，促使生成图在各位置都被判为真实；三项的总权重为1、1、0.1。像素项把输出拉向参考，感知项约束纹理结构，对抗项推动输出落在真实图像分布上。训练数据为DIV2K、Flickr2K和OutdoorSceneTraining，高清图块256，总批量48，用Adam优化，退化在训练时即时合成，参数更新采用指数滑动平均。作者还发现，训练时先锐化参考图能让输出更清晰而不明显增加伪影，这个版本称为Real-ESRGAN+。推理时只保留生成器，输入一幅低清图，一次前向得到放大结果，不需要知道退化参数。

用多次JPEG保存的招牌检查结果时，应分别看方块是否减弱、笔画是否连贯、文字是否改变。重复退化扩大了合成覆盖，感知和对抗目标又鼓励自然细节；这些条件也允许模型补出看似合理的纹理，例如把模糊的“8”画成笔画清楚的“3”。它们既没有穷尽所有相机退化，也没有为补出的字画提供新的观测。

\subsubsection{DiffBIR}
\label{sec:vision-restoration-diffbir}

严重退化下，Real-ESRGAN这类一次前向的生成器能补出的细节有限，人脸、文字等结构复杂的内容容易出现不自然的纹理。预训练的文生图扩散模型见过大量清晰图像，掌握了丰富的图像先验，自然的想法是用它来生成细节。但作者观察到，若直接把退化图像作为条件交给扩散模型，不同类型的退化会在结果中引出不同的伪影：生成过程无法区分哪些是图像内容、哪些是退化。DiffBIR因此把处理拆成两个阶段：第一阶段的复原模块只负责去掉退化，不生成新内容；第二阶段的生成模块以复原结果为条件，只负责补充细节。两阶段分开优化，于是第一阶段可以换用各任务现成的复原模型，第二阶段对所有任务共用\scite{vision-diffbir2024}本节依据2024年3月的第二版预印本。

DiffBIR使用的生成先验是Stable Diffusion 2.1-base，一种在自编码器潜空间中去噪的扩散模型。它的编码器$\mathcal{E}$把图像压缩为潜变量，解码器$\mathcal{D}$还原图像，U-Net去噪网络$\bm{\epsilon}_{\bm{\theta}}$在潜空间中预测所加噪声。这类潜空间扩散模型的结构与训练见\ref{sec:vision-models-ldm}节；扩散模型前向加噪和反向去噪的一般原理，可参考\BookChapterRef{chap:pgm-generative}{卷六“概率生成模型”章}。本节只用到两点：去噪网络在每一步给出噪声预测，由此可以算出当前对清晰潜变量的估计；去噪网络可以通过额外的条件支路接受控制。表\ref{tab:vision-diffbir-modules}按数据流顺序列出各模块的组成和训练状态。

\begin{table*}[tbp]
 \centering
 \small
 \begin{tabularx}{\linewidth}{@{}p{24mm}Xp{34mm}p{18mm}@{}}
  \toprule
  模块 & 组成 & 输入与输出（$512\times512$） & 训练状态\\
  \midrule
  训练用复原模块 & 改造的SwinIR：像素反重排下采样8倍后进入SwinIR主体，再做3次“最近邻插值—卷积—LeakyReLU”上采样 & $512^2\times3\to512^2\times3$ & 单独训练后丢弃\\
  推理用复原模块 & 盲超分辨率BSRNet，人脸复原SwinIR，盲去噪SCUNet-PSNR & 同上 & 现成模型，固定\\
  条件编码器 & Stable Diffusion的VAE编码器$\mathcal{E}$ & $512^2\times3\to64^2\times4$ & 固定\\
  条件网络 & 复制去噪U-Net的编码器与中间块；首层为拼接多出的4个通道新增零初始化参数 & $64^2\times8\to$多尺度特征 & 训练\\
  特征调制 & 零初始化卷积，加到中间块特征与跳跃连接特征 & 与U-Net各级特征同形 & 训练\\
  去噪网络 & Stable Diffusion 2.1-base的U-Net & $64^2\times4\to64^2\times4$ & 固定\\
  解码器 & VAE解码器$\mathcal{D}$ & $64^2\times4\to512^2\times3$ & 固定\\
  \bottomrule
 \end{tabularx}
 \caption{DiffBIR的模块组成。组成与训练状态取自原文正文及附录；Stable Diffusion内部U-Net的层数与通道沿用原模型，DiffBIR论文未重新列出，也未给出条件网络的参数量。}
 \label{tab:vision-diffbir-modules}
\end{table*}

第一阶段在训练和使用时扮演不同角色。训练生成模块时，作者另训练一个以像素平方误差为目标的复原模块$\mathrm{RM}$。它在一阶经典退化模型下合成训练输入：高斯模糊、下采样、加高斯噪声、JPEG压缩，再缩放回原尺寸；模糊和下采样造成的信息损失最大，所以这两项的范围取得很宽，模糊标准差取自$[0.1,12]$，下采样倍数取自$[1,12]$。这样得到的条件图像$\bm{X}_{\mathrm{c}}=\mathrm{RM}(\bm{Y})$多样且退化较少，复原模块训练15万次迭代、批量96。这个模块只为生成模块准备训练条件，处理不了复杂的真实退化，推理时被丢弃；推理时改接各任务现成的复原模型。把两者统称为“第一阶段”会掩盖这个差别。

第二阶段的条件网络称为IRControlNet。条件图像$\bm{X}_{\mathrm{c}}$先经固定的编码器$\mathcal{E}$得到条件潜变量$\bm{c}=\mathcal{E}(\bm{X}_{\mathrm{c}})$，它与当前带噪潜变量$\bm{z}^{(t)}$的形状相同，对$512\times512$的图像都是$64\times64\times4$。二者沿通道拼接成$64\times64\times8$，送入一个复制自去噪U-Net编码器和中间块的可训练支路；因拼接多出的输入通道，其权重初始化为零。支路输出的多尺度特征经零初始化卷积，加到固定U-Net的中间块特征和跳跃连接特征上。零初始化保证训练开始时支路的输出为零，整个模型与原扩散模型完全一致，随机初始化的条件支路不会一开始就破坏预训练的去噪能力。作者的消融实验说明了各设计的作用：若改用一个从头训练的小编码器处理条件图像，输出难以保持输入的颜色，PSNR平均下降约3 dB；若条件支路不读取$\bm{z}^{(t)}$，训练损失始终更高，保真指标虽略好，输出却偏平滑、缺少纹理；若支路不从去噪网络复制初始化，训练难以收敛。

训练只更新条件支路和连接参数，去噪U-Net、编码器和解码器保持预训练参数不变。对清晰参考$\bm{X}$取潜变量$\bm{z}=\mathcal{E}(\bm{X})$，随机选取步$t$和标准高斯噪声$\bm{\epsilon}$，得到$\bm{z}^{(t)}=\sqrt{\bar\alpha_t}\bm{z}+\sqrt{1-\bar\alpha_t}\bm{\epsilon}$，目标为
\begin{equation}
 L_{\mathrm{GM}}=
 \mathbb{E}_{\bm{z},t,\bm{\epsilon},\bm{c}}
 \Bigl[\bigl\lVert\bm{\epsilon}-
 \bm{\epsilon}_{\bm{\theta}}(\bm{z}^{(t)},t,\bm{c})\bigr\rVert_2^2\Bigr]\text{，}
 \label{eq:vision-diffbir-loss}
\end{equation}
其中$\bar\alpha_t$由噪声日程给定。它与原扩散模型的训练目标形式相同，只是去噪多了条件$\bm{c}$；学到的是“在给定复原结果的前提下，清晰图像的潜变量应该怎样”。训练数据是从LAION-2B英文子集中筛选出的约1500万幅高质量图像，在$512\times512$分辨率上用Adam训练8万次迭代、批量256，学习率前3万次为$10^{-4}$，之后为$10^{-5}$。

推理从随机潜变量开始，采用50步的间隔采样日程，正向文字提示置空，负向提示取“low quality”“blurry”一类描述。每步由条件支路与固定去噪网络预测噪声$\hat{\bm{\epsilon}}$，再按$\tilde{\bm{z}}^{(t)}=(\bm{z}^{(t)}-\sqrt{1-\bar\alpha_t}\,\hat{\bm{\epsilon}})/\sqrt{\bar\alpha_t}$得到对清晰潜变量的当前估计，最后解码成图像。生成先验倾向于在纹理、边缘处补出细节，在天空、墙面等平坦区域也可能补出原图没有的内容。为此DiffBIR提供一个无需训练的复原引导：每步得到$\tilde{\bm{z}}^{(t)}$后，减少解码图像与第一阶段结果的加权距离，
\begin{equation}
 \begin{aligned}
 L_{\mathrm{g}}(\tilde{\bm{z}})
 &=\frac{1}{hwc}
   \left\lVert\bm{W}\odot
    \bigl(\mathcal{D}(\tilde{\bm{z}})-\bm{X}_{\mathrm{c}}\bigr)
   \right\rVert_{\mathrm{F}}^2,\\
 \tilde{\bm{z}}_{\mathrm{g}}^{(t)}
 &=\tilde{\bm{z}}^{(t)}
   -s\nabla_{\tilde{\bm{z}}}L_{\mathrm{g}}
       (\tilde{\bm{z}}^{(t)})\text{。}
 \end{aligned}
 \label{eq:vision-diffbir-guidance}
\end{equation}
这里$s$控制引导强度，权重图$\bm{W}=\bm{1}-\bm{G}$，$\bm{G}$是把$\bm{X}_{\mathrm{c}}$分成不重叠图块后，用Sobel算子（\ref{sec:vision-transforms-sharpening}节）计算的块级梯度幅值（归一化到$[0,1]$）。平坦区域梯度小、权重大，被拉向第一阶段的平滑结果；纹理区域权重小，保留更多生成细节。更新后的清晰估计参与下一采样步。损失中的权重在平方之内，实际像素平方误差的系数是$w_{i,j,k}^2$。论文在合成数据上比较了$s=0,0.5,1$，在真实照片上取$s=0$以求更好的观感。每一步引导都要经过解码器求梯度，推理成本因此高于不引导的采样。这一步与约束最小二乘的思路相通：都在“符合先验”的候选中拉近与一个观测依据的距离，只是这里的先验来自扩散模型，观测依据换成了第一阶段的复原结果。

在固定低清人像和随机种子下，增加$s$会加强对第一阶段轮廓及平坦区域的约束，同时可能压低生成细节。被追随的仍是一个复原估计，它若已把眼睛或文字恢复错，引导也可能保留错误。皮肤纹理、身份细节与招牌字形应分别核验，不能把新的高频外观当作相机曾经记录到的事实。

三种超分辨率模型可以在同一张多次压缩的低清招牌上对照，表\ref{tab:vision-sr-compare}归纳它们改变了什么。Real-ESRGAN相对DnCNN变体改变了训练退化和目标函数，换来对真实照片的适用性和更锐利的纹理，代价是输出不再是条件均值，PSNR通常不再是它追求的指标，补出的笔画也可能与原文不符。DiffBIR相对Real-ESRGAN又改变了先验的来源和推理方式，换来更丰富的细节和可调的保真程度，代价是50步采样和一个大规模预训练扩散模型的计算，并且同一输入在不同随机种子下会给出不同细节。

\begin{table*}[tbp]
 \centering
 \small
 \begin{tabularx}{\linewidth}{@{}p{18mm}XXX@{}}
  \toprule
  比较维度 & DnCNN超分辨率变体 & Real-ESRGAN & DiffBIR\\
  \midrule
  网络规模 & 20层同尺寸卷积，在高分辨率网格上计算 & 23个RRDB，在低分辨率上计算后上采样 & 固定扩散U-Net加复制的编码器支路\\
  训练退化 & 双三次下采样，倍数2、3、4 & 两轮随机模糊、缩放、噪声、JPEG及sinc滤波 & 经典退化合成条件图；推理接各任务复原模块\\
  学习目标 & 残差的平方误差 & $L_1$、感知与对抗损失 & 条件扩散的噪声预测误差\\
  细节来源 & 训练配对的条件均值 & 生成器与判别器学到的纹理分布 & 预训练文生图扩散模型\\
  更新参数 & 全部网络 & 生成器与判别器 & 仅条件支路，主干固定\\
  推理代价 & 插值加一次前向 & 一次前向 & 复原模块加50步采样，可选逐步引导\\
  主要风险 & 输出平滑、只适用干净退化 & 补出与原文不符的纹理 & 随机种子间细节不同，可能保留第一阶段错误\\
  \bottomrule
 \end{tabularx}
 \caption{分辨率不足下三种复原模型的比较。三者面对同一个难点，即一个低清观测对应多个可能的高清图像，区别在于用什么补出平均估计中缺失的细节。}
 \label{tab:vision-sr-compare}
\end{table*}

\subsection{区域缺失下的复原}
\label{sec:vision-restoration-inpainting}
\label{sec:vision-restoration-partial-convolution}

坏点、划痕或遮挡与加性噪声不同：对应位置根本没有可信观测。以二值掩码$\bm{M}$标明有效像素，$m_{i,j}=1$表示可用，0表示缺失，并在通道维广播。\term{图像修补}（image inpainting）依据周围内容估计缺口原有内容，输入是受损图像和掩码，输出是完整图像。难点有两个：缺口形状不规则、大小不定；缺口中的占位数值没有意义，却会被普通卷积当作真实输入读取。若要求删除杯子并创造一只花瓶，则还包含新的语义目标，交由\ref{chap:vision-image-generation}章讨论，那里也介绍基于样例图块复制的经典修补方法。

普通卷积网络处理缺口时，先在缺口中填入零或图像均值，再与有效像素一起卷积；输出因占位数值而变化，常在缺口内出现颜色偏差、模糊，在缺口边缘出现人为的边缘响应，需要额外后处理。此前的学习式方法又多在图像中央的方形缺口上训练，遇到不规则缺口时表现下降。\term{部分卷积}（partial convolution）的想法是让卷积只读取有效位置，并按有效位置的数量调整响应；每算一层就更新掩码，让有效区域逐层向缺口内部推进\scite{vision-partial-convolution2018}

把一个窗口的输入、掩码和权重分别展平为$\bm{x},\bm{m},\bm{w}\in\mathbb{R}^q$，部分卷积的输出及新掩码为
\begin{equation}
 x'=
 \begin{cases}
 \dfrac{q}{\sum_i m_i}\bm{w}^{\top}(\bm{x}\odot\bm{m})+b,
   &\sum_i m_i>0,\\
 0,&\sum_i m_i=0,
 \end{cases}
 \qquad
 m'=\mathbb{I}\!\left[\sum_i m_i>0\right]\text{，}
 \label{eq:vision-partial-convolution}
\end{equation}
其中$b$是偏置。乘以$\bm{m}$屏蔽了缺口中的数值；因子$q/\sum_i m_i$把只覆盖部分有效位置的加权和放大到与完整窗口可比的尺度，它只乘加权和，不乘偏置。完全无有效输入时输出零，新掩码也为零。只要窗口内有一个有效输入，输出位置就标为有效，因此掩码中的有效区域每层向外扩展，经过足够多层后，只要原图有有效像素，掩码最终全为1。

取窗口$\bm{x}=(0.2,0,0.8)$、$\bm{m}=(1,0,1)$、$\bm{w}=(1/3,1/3,1/3)$且$b=0$，屏蔽后的加权和为$1/3$，再乘$3/2$得到0.5，$m'=1$。若把缺口占位改成100，结果仍为0.5，这正是部分卷积要保证的性质。窗口内全缺失时不能计算这种归一化，必须走零输出分支。图\ref{fig:vision-partial-mask}说明，堆叠窗口使可计算范围逐层进入缺口。

\begin{figure}[htbp]
 \centering
 % Binary masks; kernel width 3, stride 1. No image values are fabricated.
\begin{tikzpicture}[x=1mm,y=1mm,font=\figurefont\small,
 arr/.style={-{Stealth[length=2mm]},BookMuted,line width=.8pt}]
 \foreach \r/\textlabel in {0/输入掩码,1/第一层,2/第二层} {
  \node[anchor=east] at (0,{-\r*18}) {\textlabel};
  \foreach \i in {0,...,6} {
   \draw[BookMuted,line width=.6pt] ({5+\i*11},{-\r*18-4})
     rectangle ++(11,8);
  }
 }
 \foreach \r/\indices in {0/{0,1,5,6},1/{0,1,2,4,5,6},2/{0,1,2,3,4,5,6}} {
  \foreach \i in \indices {
   \fill[BookBlue!12] ({5+\i*11+.4},{-\r*18-3.6}) rectangle ++(10.2,7.2);
   \fill[BookBlue] ({10.5+\i*11},{-\r*18}) circle (.8mm);
  }
 }
 \draw[BookAmber,line width=1pt] (27,-5.5) rectangle (60,5.5);
 \draw[BookAmber,line width=1pt] (27,-23.5) rectangle (60,-12.5);
 \draw[arr] (43.5,-6)--(43.5,-12);
 \draw[arr] (43.5,-24)--(43.5,-31);
 \node[anchor=north] at (43.5,-43) {中心掩码：$0\to0\to1$};
\end{tikzpicture}

 \caption{宽度为3、步幅为1的掩码更新教学例。带点填充的格子表示已有有效输入，白格表示缺失；每层只要窗口内有一个有效位置，该层输出就标为有效。图中省略边界外输入。中间格在第二层获得可计算支持，不代表新取得了原像素观测。}
 \label{fig:vision-partial-mask}
\end{figure}

修补网络是一个U形编码器—解码器，所有卷积层都换成部分卷积层，表\ref{tab:vision-pconv-layers}按原文附录列出逐层配置。输入为$512\times512\times3$的受损图像与同尺寸掩码。编码器有8层步幅为2的部分卷积，卷积核依次为$7\times7$、两层$5\times5$和五层$3\times3$，通道数从64、128、256升到512后保持不变，空间尺寸从512逐层减半到2；除第一层外都接批归一化，激活函数为ReLU。解码器有8级，每级先用最近邻插值把特征和掩码放大2倍，再分别与同级编码器的特征和掩码沿通道拼接，然后做一层$3\times3$、步幅1的部分卷积，接批归一化和斜率0.2的LeakyReLU。最后一级拼接的是原始受损图像和原始掩码，使网络能直接复制有效像素；这一级的卷积输出3个通道，不接批归一化和激活函数。图像边界外也视为缺口，用部分卷积代替普通的零填充，避免边界处的结果受图像外数值影响。

\begin{table*}[tbp]
 \centering
 \small
 \begin{tabularx}{\linewidth}{@{}p{26mm}Xp{20mm}p{12mm}p{12mm}p{26mm}@{}}
  \toprule
  层组 & 层类型 & 卷积核与步长 & 输出通道 & 重复次数 & 输出尺寸\\
  \midrule
  编码1 & 部分卷积，ReLU & $7\times7$，2 & 64 & 1 & $256\times256\times64$\\
  编码2 & 部分卷积，批归一化，ReLU & $5\times5$，2 & 128 & 1 & $128\times128\times128$\\
  编码3 & 同上 & $5\times5$，2 & 256 & 1 & $64\times64\times256$\\
  编码4至8 & 同上 & $3\times3$，2 & 512 & 5 & $32^2$逐层减半至$2^2$，512通道\\
  解码9至12 & 最近邻上采样$\times2$，拼接编码7至4，部分卷积，批归一化，LeakyReLU & $3\times3$，1 & 512 & 4 & $4^2$逐层加倍至$32^2$，512通道\\
  解码13 & 同上，拼接编码3 & $3\times3$，1 & 256 & 1 & $64\times64\times256$\\
  解码14 & 同上，拼接编码2 & $3\times3$，1 & 128 & 1 & $128\times128\times128$\\
  解码15 & 同上，拼接编码1 & $3\times3$，1 & 64 & 1 & $256\times256\times64$\\
  解码16 & 上采样，拼接受损图像，部分卷积 & $3\times3$，1 & 3 & 1 & $512\times512\times3$\\
  \bottomrule
 \end{tabularx}
 \caption{部分卷积修补网络的逐层结构（输入$512\times512$）。每层同时输出特征和更新后的掩码，拼接时特征与掩码各自拼接。输出尺寸按步幅推算；层配置取自原文附录。}
 \label{tab:vision-pconv-layers}
\end{table*}

训练从清晰图像$\bm{X}$出发，随机施加不规则掩码，得到网络输出$\widehat{\bm{X}}_{\mathrm{net}}$和合成图$\widehat{\bm{X}}_{\mathrm{comp}}=\bm{M}\odot\bm{X}+(\bm{1}-\bm{M})\odot\widehat{\bm{X}}_{\mathrm{net}}$。总损失为
\begin{equation}
 L=L_{\mathrm{valid}}+6L_{\mathrm{hole}}+0.05L_{\mathrm{perc}}
   +120\bigl(L_{\mathrm{style}}^{\mathrm{out}}+L_{\mathrm{style}}^{\mathrm{comp}}\bigr)
   +0.1L_{\mathrm{tv}}\text{。}
 \label{eq:vision-pconv-loss}
\end{equation}
$L_{\mathrm{valid}}$和$L_{\mathrm{hole}}$分别是有效区与缺口区的$L_1$像素误差，缺口区权重更大；$L_{\mathrm{perc}}$用预训练VGG-16前三个池化层的特征比较$\widehat{\bm{X}}_{\mathrm{net}}$、$\widehat{\bm{X}}_{\mathrm{comp}}$与参考；风格项先对各层特征计算通道间的Gram矩阵，再比较其差别，约束纹理统计；$L_{\mathrm{tv}}$在缺口外扩一个像素的区域内惩罚相邻像素差，使填补内容与周围平滑衔接。这些权重由在100幅验证图上的超参数搜索确定。作者的消融显示，去掉风格项或风格权重过小，输出会出现鱼鳞状伪影；去掉感知项，网格状伪影更明显。

训练掩码来自视频相邻帧之间的遮挡估计，共55116个，训练时再随机膨胀、旋转和裁剪；图像取自ImageNet、Places2和CelebA-HQ。批归一化在缺口处会把无效像素计入均值和方差，因此训练分两步：先开启批归一化，以学习率$2\times10^{-4}$用Adam训练，批量6；再以学习率$5\times10^{-5}$微调，同时冻结编码器中批归一化的统计和参数，解码器中的批归一化保持更新，因为到解码器时缺口已基本被填满。

使用时输入受损图像和原始掩码，网络预测$\widehat{\bm{X}}_{\mathrm{net}}$，再形成
$\widehat{\bm{X}}=\bm{M}\odot\bm{Y}+(\bm{1}-\bm{M})\odot\widehat{\bm{X}}_{\mathrm{net}}$，使已知像素按要求保留。与填入占位值后直接卷积的做法相比，部分卷积消除了结果对占位值的依赖，并能处理任意形状的缺口；但它没有改变一个事实：缺口内的内容全部来自周围像素和训练先验。宽度很小的划痕常有两侧结构可参考，大缺口则允许多种内容，感知项和风格项都不能唯一确定遮挡背后的实际对象。中间掩码变为1只说明特征已有计算来源，原始观测掩码仍须保留，用来区分记录与估计。

\section{图像增强}
\label{sec:vision-restoration-enhancement}

有些图像已经保留了需要的结构，只是亮度或对比不适合观看；另一些暗图则同时缺少可靠观测。两种情况外观相似，处理方法和可以期待的结果却不同。增强的输入是一幅图像（或RAW数据），输出是适合指定用途的图像；它没有唯一正确的目标，因此每次增强都要说明希望看清什么，以及哪些颜色、测量量或局部关系必须保留。

本节按增强的对象分为两部分。曝光与亮度不足时，难点在于暗处信号本身被噪声淹没，单纯放大会同时放大噪声，所以需要回到更早的成像阶段，联合完成去噪和颜色形成。色彩与对比不当时，内容已经足够，只需重新分配数值，这时不必训练网络，可以直接选用\ref{chap:vision-image-transforms}章的点运算或直方图方法；全局曲线和局部直方图的差别在于映射是否随位置变化。

\subsection{曝光与亮度的增强}
\label{sec:vision-restoration-low-light}
\label{sec:vision-restoration-sid}

低照度时，传感器收到的光子较少，信号相对噪声较弱。按式\eqref{eq:vision-poisson-gaussian}，光子散粒噪声的标准差随信号的平方根增长，信号越弱，信噪比越低。若将像素与加性噪声一起乘上$a$，信号幅值变为$a$倍，噪声标准差也变为$a$倍，这种线性提亮没有改善二者的比例。延长真实曝光能收集更多光子，却可能引入运动模糊；对已经拍下的短曝光照片，事后乘增益不能替代重新曝光。传统相机处理管线在提亮前已依次做了白平衡、去马赛克、去噪、色彩转换和伽马校正，在极低信噪比下这些模块各自出错，结果中常见强噪声和偏色；若在管线输出上再做去噪，噪声已经被非线性处理改变了分布，一般的去噪方法难以对付，色偏也无法消除。多帧对齐融合能提高信噪比，但需要稠密对应估计，对运动场景和视频不易推广。

Learning to See in the Dark的想法是绕开传统管线，用一个全卷积网络直接把短曝光的RAW数据映射成正常亮度的RGB图像\scite{vision-sid2018}网络读取的是颜色插值和非线性处理之前的数据，可以同时完成去噪、去马赛克和颜色形成，并在训练中学习适应数据的局部降噪。为此作者构建了SID数据集：用Sony $\alpha$7S~II（Bayer阵列）和富士X-T2（X-Trans阵列）在室内外暗光场景下各拍短曝光图像和曝光时间长100到300倍的参考图像，共5094幅短曝光RAW，对应424幅参考；相机固定在三脚架上，以保证配对对齐。

以Bayer传感器为例，每个$2\times2$位置采样一红、两绿和一蓝（Bayer阵列见\ref{sec:vision-formation-acquisition}节），原始阵列为$h\times w\times1$。按位置打包后得到$\frac{h}{2}\times\frac{w}{2}\times4$数组，两个绿色位置保留为不同通道；X-Trans阵列的原始数据按$6\times6$的块排列，作者通过交换相邻元素把它打包成9个通道，而不是36个。打包只是重新排列观测，没有提前补齐每个像素的三种颜色。打包后的数据减去传感器黑电平，按有效范围归一化，再乘外部给定的放大比例。表\ref{tab:vision-sid-network}列出整个管线的形状变化。核心网络默认采用U-Net编码器—解码器（U-Net结构见\BookChapterRef{chap:convolutional-networks}{卷五“卷积神经网络”章}），作者也试过多尺度上下文聚合网络，U-Net的PSNR更高；论文没有给出U-Net的级数、各级通道数和卷积配置，只说明选择网络时要求能在显存中处理$4240\times2832$或$6000\times4000$的整幅图像，因此避开了需要切块再拼接的全连接层。网络输出$\frac{h}{2}\times\frac{w}{2}\times12$的特征，再经子像素重排，把每4个通道组成一个$2\times2$空间块，恢复成$h\times w\times3$的RGB图像。作者试过残差连接，没有观察到收益，推测原因是输入和输出处在不同的颜色空间。

\begin{table}[htbp]
 \centering
 \small
 \begin{tabularx}{\linewidth}{@{}p{16mm}Xp{28mm}@{}}
  \toprule
  阶段 & 操作 & 输出尺寸（Bayer）\\
  \midrule
  打包 & 每个$2\times2$拜耳单元拆成4个通道 & $\frac h2\times\frac w2\times4$\\
  归一化 & 减黑电平，除以有效范围，乘放大比例 & $\frac h2\times\frac w2\times4$\\
  核心网络 & U-Net编码器—解码器（层配置论文未给出） & $\frac h2\times\frac w2\times12$\\
  子像素重排 & 每4个通道排成一个$2\times2$空间块 & $h\times w\times3$\\
  \bottomrule
 \end{tabularx}
 \caption{Learning to See in the Dark的处理管线。X-Trans阵列打包为9个通道，其余步骤相同；输出直接位于sRGB空间。}
 \label{tab:vision-sid-network}
\end{table}

假设某个RAW读数为530，黑电平为512，可用上界为16383，则归一化信号为$(530-512)/(16383-512)\approx0.00113$。乘100后约为0.113。若忽略黑电平，直接用$530/16383$再乘100会超过3，数值主要来自偏置并很快被裁剪；这不是同一种输入。相机的黑电平、颜色阵列和增益约定必须与训练一致，作者也为每台相机单独训练一个网络。

训练从头开始，以输出与参考RGB之间的$L_1$距离为损失，参考图像由RAW处理库libraw从长曝光数据生成。训练时放大比例取为两次曝光时间之比，例如100、250或300倍；每次迭代随机裁取$512\times512$的图块，并做随机翻转和旋转。优化器为Adam，学习率$10^{-4}$，2000轮后降为$10^{-5}$，共训练4000轮。推理时输入完整分辨率的RAW数据，放大比例由用户给定，作用类似相机的ISO设置：比例越大，输出越亮。网络并不自动推断最佳曝光。长曝光参考仍可能有噪声，移动人物和树叶也会破坏配对；若参考错位，逐像素监督就会把位置差异也当作待修正误差。

固定同一夜景，可以比较三种做法：传统管线输出后乘增益、传统管线输出后用BM3D去噪、以及RAW到RGB的网络。前两者都在已经非线性化的RGB数据上工作，第一种保留全部噪声，第二种需要人为指定噪声水平，设小了残留噪声、设大了过度平滑，而且两种情形可能出现在同一幅图的不同区域；网络在RAW数据上工作，能按局部数据调整降噪，同时修正偏色。作者用感知实验比较了这几种做法，并说明为公平起见对基线使用了参考图像的白平衡和逐通道均值，这相当于给基线提供了额外信息。比较的结论仍有范围：网络的收益部分来自更早阶段的传感器数据，不能全部归结为更强的后处理；它只在两台相机和静态场景上训练，更换相机和处理动态场景都需要单独验证。

\subsection{色彩与对比的增强}
\label{sec:vision-restoration-tone}

当图像已有可用内容时，增强只需重新分配数值，不必训练网络。输入是一幅RGB图像，输出是同尺寸的图像，难点在于同一幅图中阴影、中间调和高光对调整的需求不同，而且任何调整都会改变一些原有的数值关系。

两条路线的区别在于映射是否随位置变化。全局曲线对整幅图使用同一个映射，相同输入值得到相同输出，结果可预测，像素之间的明暗次序不会颠倒，但难以同时照顾暗部和亮部。局部直方图方法依据每个位置周围的亮度分布改变映射，能分别拉开阴影和高光中的层次，代价是相同亮度在不同位置得到不同输出，区域之间的相对关系也随之改变。两类映射的定义与推导都在\ref{sec:vision-transforms-intensity}节，本小节只讨论在夜间招牌这类应用中怎样选用、按什么顺序使用，以及各自会带来什么副作用。

\subsubsection{基于全局曲线的增强}
\label{sec:vision-restoration-white-balance}

全局调整通常分两步：先校正颜色，再调整亮度层次。白平衡按颜色通道调整增益，使选定照明条件下的中性物体呈现为中性色，它在成像管线中的位置见\ref{sec:vision-formation-isp}节。在简化的线性RGB模型中，$\bm{c}'=\operatorname{diag}(g_r,g_g,g_b)\bm{c}$，$\bm{c}$是一个像素的颜色向量。若招牌旁的灰色墙砖被拍成$(0.20,0.16,0.10)$，取增益$(0.8,1,1.6)$会得到$(0.16,0.16,0.16)$。这依赖参考本来应是中性；夜间街景常有多种光源，钠灯下的墙面本身偏黄，误把它当灰卡会使整图颜色偏移。

亮度层次用\term{色调曲线}（tone curve）调整，即把亮度$I\in[0,1]$映射成单调递增的$T(I)$。\ref{sec:vision-transforms-intensity}节已说明，曲线斜率大于1的亮度段被拉开，小于1的被压缩；幂函数$T(I)=I^\gamma$在$\gamma<1$时拉开暗部，$\gamma>1$时压暗中间值。对夜间招牌，取$\gamma=1/2$会把0.04、0.25和0.81依次变为0.2、0.5和0.9，墙面细节变得可见，代价是亮部层次被压缩。与之对照，两倍线性提亮再裁剪会把0.6和0.9都变为1，招牌高光之间的差异完全丢失；单调曲线只是压缩，不会把不同值合并为同一值，除非量化精度不足。

使用中有两个容易出错的地方。第一，逐通道应用非线性曲线通常会改变通道比例，从而改变颜色，例如让红色招牌偏橙。若希望尽量保持线性RGB比例，可先计算亮度$I$，再在$I>0$处用$\bm{c}'=\bm{c}\,T(I)/I$统一缩放，随后处理超出色域的数值；接近零时需稳定分母并规定黑色输出。保持比例不等于所有显示条件下的感知色相严格不变，颜色空间和裁剪仍有作用。第二，曲线应作用在明确的编码上：同一个$\gamma$作用在线性光数据和已经伽马编码的sRGB数据上，效果不同，需要记录是在哪一种数值上调整的。

全局曲线只重新分配现有层次。招牌暗部可能更容易辨认，平坦暗区的噪声也可能一起显露；已经被传感器饱和或量化合并的字画，不能仅靠单调曲线恢复。它的参数只有几个，可以记录下来并在需要时近似撤销，这是它在档案和测量用途中的优点。

\subsubsection{基于局部直方图的增强}
\label{sec:vision-restoration-clahe}

一幅夜景中，招牌亮、墙面暗，同一条全局曲线拉开了墙面的层次，就会压缩招牌的层次。直方图均衡化按累计分布确定曲线，像素多的亮度段被拉开（推导见\ref{sec:vision-transforms-histogram}节）；对整幅图求一条曲线时，它仍然无法兼顾招牌和墙面。受限局部直方图均衡化（CLAHE）把图像分成若干块，每块用截断后的直方图求映射，截断限制了映射的最大斜率，相邻块的映射再经双线性插值平滑过渡（算法与数值例见\ref{sec:vision-transforms-clahe}节）。

用在招牌照片上，需要决定三件事。其一是作用在哪个通道上：若对R、G、B三个通道分别均衡，各通道的映射不同，颜色会被改变；常见做法是先转换到亮度与色度分离的空间（\ref{sec:vision-formation-color-spaces}节），只对亮度通道做CLAHE，再转回RGB。其二是块的大小和截断上限：块较小时更关注局部，但噪声及小结构更容易主导直方图；上限放宽则增加局部对比，平坦墙面上的噪声也被拉开。上限通常以相对于平均计数的倍数给出，各实现对余数的分配规则也不同，这些都须记录。其三是与去噪的先后：CLAHE会放大平坦区的微小起伏，暗光照片应先去噪再做局部均衡，否则噪声会被当作“层次”拉开。

CLAHE的副作用在招牌边缘最明显。块内有明亮招牌时，直方图的大部分计数落在暗部，暗部被拉开，招牌周围的墙面却可能因所在块的映射不同而变亮或变暗，形成环绕亮边的光晕。同一亮度在不同位置输出不同，所以不能在增强后的图像上比较两处的亮度，也不能把增强后的灰度用于测量。

两种增强在同一张夜间招牌上的差别见表\ref{tab:vision-tone-compare}。全局曲线保持了像素间的明暗次序，参数少，易于记录和复现，但无法同时照顾招牌和墙面；CLAHE能分别拉开两处的层次，代价是区域间的明暗关系可能被改变，平坦区噪声和亮边周围的光晕也会被增强。医学、档案和测量图像应先规定哪些灰度关系允许变化，再决定能否使用局部方法。

\begin{table}[htbp]
 \centering
 \small
 \begin{tabularx}{\linewidth}{@{}p{16mm}XX@{}}
  \toprule
  比较维度 & 全局曲线 & 受限局部直方图均衡化\\
  \midrule
  映射依据 & 人为设定或由全图统计确定 & 每块直方图，截断后求累计分布\\
  空间一致性 & 相同输入值得到相同输出 & 随位置变化，相邻块间插值\\
  参数 & 通道增益、$\gamma$等少数几个 & 块大小、截断上限、分配规则、作用通道\\
  使用顺序 & 白平衡在前，曲线在后 & 去噪在前，只作用于亮度通道\\
  主要风险 & 暗部与亮部难以兼顾 & 平坦区噪声放大、光晕、改变区域间关系\\
  \bottomrule
 \end{tabularx}
 \caption{色彩与对比增强两条路线的比较。两者都只重新分配已有层次，不能恢复已饱和或已被量化合并的细节。}
 \label{tab:vision-tone-compare}
\end{table}

\section{图像压缩}
\label{sec:vision-restoration-compression}

压缩的输出是可传输、可独立解码的比特流。一个网络即使能把图像映射到很小的浮点数组，也不等于已经压缩：浮点数需要多少位、哪些辅助信息要传、解码器如何复现概率和符号，都必须明确。压缩要同时控制两个量：码率$R$，即平均每幅图像（或每个像素）用多少比特，常用式\eqref{eq:vision-bpp}的每像素比特数计；失真$D$，即重建图像与原图的差别。二者此消彼长，比特越少，失真越大；在给定码率下使失真最小，或在给定失真下使码率最小，称为\term{率失真权衡}（rate-distortion trade-off）。

图像之所以能被压缩，是因为它含有冗余。第一种是编码冗余：若每个像素都用8位定长码，而某些灰度值出现得远比其他值频繁，定长码就浪费了比特。第二种是空间冗余：相邻像素的值高度相关，知道左边的像素，就能大致猜出右边的像素。第三种是与用途无关的信息：人眼对高频细节和色度的变化不如对低频亮度敏感，这部分信息去掉后观看者不易察觉。前两种冗余可以无损地去除，第三种的去除必然有损。

编码冗余的下限由熵给出。若符号$s$以概率$p(s)$独立出现，其熵为
\begin{equation}
 H=-\sum_{s}p(s)\log_2p(s)\quad\text{比特/符号}\text{，}
 \label{eq:vision-entropy}
\end{equation}
任何可唯一解码的编码，平均码长都不低于$H$；而存在前缀码使平均码长小于$H+1$（可参考\BookChapterRef{chap:information-theory-and-statistical-geometry}{卷二“概率分布”章}）。前缀码是指任何一个码字都不是另一个码字开头的编码，读到一个完整码字就能立刻确定符号。熵只针对已经给定的符号及其概率；空间冗余要先通过变换或预测转化为“符号分布更集中”，才能被熵编码利用。

一个有损编码器因此由三个环节组成：先用变换或预测去掉空间冗余，再用量化去掉与用途无关的信息，最后用熵编码去掉编码冗余。两条路线在这三个环节上的差别，在于它们是由人设计，还是由数据和率失真目标共同学习。人工设计的路线以JPEG为代表，每个环节都有明确的数学依据，便于标准化；学习的路线以尺度超先验为代表，三个环节在同一个目标下联合训练。下面先讲前者，它的每一步都为理解后者提供了对照。

\subsection{基于人工设计编码流程的压缩}
\label{sec:vision-restoration-classic-codec}

JPEG是1992年前后定型的静止图像压缩标准，至今仍是照片最常用的格式\scite{vision-wallace1992jpeg}它的基线编码流程如图\ref{fig:vision-jpeg-pipeline}：先把RGB图像转换为一个亮度和两个色度分量并对色度下采样，再把每个分量切成$8\times8$的块，每块做二维离散余弦变换，按量化表除并取整，把量化后的系数按之字形顺序排成序列，最后对序列做游程编码与Huffman编码。三个环节分别对应上文的三种冗余：离散余弦变换去掉块内的空间相关，量化表和色度下采样丢掉人眼不敏感的信息，Huffman编码去掉符号的编码冗余。只有色度下采样和量化是有损的。JPEG 2000把块变换换成小波，把Huffman编码换成上下文自适应的算术编码，解决了JPEG在低码率下的块效应。下面按“变换与量化—符号化—变长编码”的顺序展开，每一步都先讲JPEG，再说明JPEG 2000改变了什么。

\begin{figure*}[tbp]
 \centering
 % 基线JPEG编码流程；只画数据流，不表示各步的计算量。宽约150mm，跨栏使用。
\begin{tikzpicture}[font=\figurefont\small,
  stage/.style={draw=FigureBlue,fill=FigureBlueFill,rounded corners=1.8pt,
    minimum width=21mm,minimum height=9mm,align=center,inner sep=1.5mm},
  lossy/.style={stage,draw=FigureRed,fill=FigureRedFill},
  code/.style={stage,draw=FigureYellow,fill=FigureYellowFill},
  flow/.style={-{Stealth[length=1.8mm]},draw=BookMuted,line width=0.8pt},
  note/.style={text=BookMuted,font=\figurefont\footnotesize}]
 \node[stage] (rgb) at (0,0) {RGB图像};
 \node[lossy] (ycc) at (3.0,0) {YCbCr变换\\色度下采样};
 \node[stage] (blk) at (6.0,0) {$8\times8$分块\\减去128};
 \node[stage] (dct) at (9.0,0) {二维DCT};
 \node[lossy] (q) at (12.0,0) {除以量化表\\取整};
 \node[stage] (zz) at (12.0,-1.75) {之字形扫描};
 \node[code] (dc) at (9.0,-1.1) {DC：与前块\\作差分};
 \node[code] (ac) at (9.0,-2.4) {AC：零游程\\与幅值配对};
 \node[code] (hf) at (6.0,-1.75) {Huffman\\变长编码};
 \node[stage] (bs) at (3.0,-1.75) {码流\\含表与头信息};
 \draw[flow] (rgb)--(ycc);
 \draw[flow] (ycc)--(blk);
 \draw[flow] (blk)--(dct);
 \draw[flow] (dct)--(q);
 \draw[flow] (q)--(zz);
 \draw[flow] (zz)--(dc);
 \draw[flow] (zz)--(ac);
 \draw[flow] (dc)--(hf);
 \draw[flow] (ac)--(hf);
 \draw[flow] (hf)--(bs);
 \node[note] at (7.5,.85) {红框为有损步骤，其余步骤可逆或无损};
\end{tikzpicture}

 \caption{基线JPEG的编码流程。上行把图像变成量化后的系数块，下行把系数块变成比特；解码按相反顺序执行。红框标出两个有损步骤，色度下采样丢掉色度分辨率，量化丢掉系数精度；其余步骤都可以精确逆转。}
 \label{fig:vision-jpeg-pipeline}
\end{figure*}

\subsubsection{像素块的变换与量化}
\label{sec:vision-restoration-jpeg-transform}

第一步要把像素变成大多数取零的整数。JPEG先把RGB转换为YCbCr，其中亮度$Y=0.299R+0.587G+0.114B$，Cb、Cr分别是蓝色与红色相对亮度的差（\ref{sec:vision-formation-color-spaces}节）。人眼对色度的空间变化不如对亮度敏感，常用的4:2:0采样把两个色度分量在水平和垂直方向各降一半，每个像素平均的样本数从3降到$1+2\times\frac14=1.5$，在进入任何变换之前就省下一半数据。招牌的红色字画在色度上分辨率降低，边缘可能出现轻微色晕，但亮度边缘不受影响。

各分量减去128，使8位样本以零为中心，再切成$8\times8$的块。每块$f_{i,j}$（$i,j=0,\dots,7$，$i$为行，$j$为列）做二维离散余弦变换：
\begin{equation}
 F_{k,l}=\frac14C(k)C(l)\sum_{i=0}^{7}\sum_{j=0}^{7}f_{i,j}
 \cos\frac{(2i+1)k\pi}{16}\cos\frac{(2j+1)l\pi}{16}\text{，}
 \label{eq:vision-jpeg-dct}
\end{equation}
其中$C(0)=1/\sqrt2$，$k>0$时$C(k)=1$。$F_{k,l}$是垂直频率$k$、水平频率$l$的系数，$F_{0,0}$与块均值成正比，称为直流（DC）系数，其余63个称为交流（AC）系数。\ref{sec:vision-transforms-dct}节已说明，离散余弦变换是正交变换，不损失信息；对相邻像素高度相关的平滑块，能量集中在少数低频系数上，这就是它去掉空间冗余的方式。选择$8\times8$是一个折中：块太小，能量集中不充分；块太大，块内更可能同时含平坦区和边缘，计算量也更大。

\term{量化表}（quantization table）为64个系数各规定一个步长$Q_{k,l}$，量化把系数变成整数$\hat F_{k,l}=\operatorname{round}(F_{k,l}/Q_{k,l})$，解码端用$\hat F_{k,l}Q_{k,l}$近似原系数。步长越大，丢掉的精度越多；量化误差在$[-Q_{k,l}/2,Q_{k,l}/2]$内，若近似为均匀分布，方差为$Q_{k,l}^2/12$（表\ref{tab:vision-noise-models}）。标准附录给出的亮度示例表的第一行是$(16,11,10,16,24,40,51,61)$，第一列是$(16,12,14,14,18,24,49,72)$，右下角高频处的步长接近100。高频步长大，是因为人眼对高频误差不敏感，这正是在量化环节去掉与用途无关的信息。常用实现用一个质量因子整体缩放这张表，质量因子越低，步长越大，码率越低。

用招牌上一个竖直笔画边缘的块作可复算的例子。设块内每一行都是$(60,60,60,60,200,200,200,200)$，左半是暗背景、右半是亮笔画。减去128后每行为$(-68,-68,-68,-68,72,72,72,72)$。因为各行相同，式\eqref{eq:vision-jpeg-dct}对$i$的求和在$k>0$时为零（余弦在一个周期内的等距采样和为零），只剩第一行系数$F_{0,l}=\sqrt2\,C(l)\sum_jf_{0,j}\cos\frac{(2j+1)l\pi}{16}$。计算得
\[
 (F_{0,0},\dots,F_{0,7})\approx(16,\,-507.4,\,0,\,178.2,\,0,\,-119.1,\,0,\,100.9)\text{，}
\]
偶数频率为零，是因为去掉均值后这一行关于中心反对称。除以量化表第一行并取整，得到$(1,-46,0,11,0,-3,0,2)$，整个块64个系数中只剩5个非零。解码时乘回步长再做逆变换，每行重建为
\[
 (60.6,\,58.4,\,63.6,\,56.5,\,203.5,\,196.4,\,201.6,\,199.4)\text{，}
\]
最大误差约3.6。误差并不均匀：原本平坦的暗区出现了$58.4,63.6,56.5$这样的起伏，这是截断和粗量化高频系数后，在边缘附近出现的振铃。

块变换在低码率下还有一个更显眼的问题。每块独立量化，相邻两块的直流系数各自取整，平坦区跨块边界的亮度会出现台阶；交流系数被大量置零后，块内只剩平滑的低频形状，块与块之间无法平滑衔接。这种沿$8\times8$网格出现的方形边界称为块效应，就是开篇招牌放大后看到的“方块”。

JPEG 2000改变的正是变换环节\scite{vision-skodras2001}它对整幅图（或较大的图块）做多级二维离散小波变换（\ref{sec:vision-transforms-wavelet}节），每级把图像分成一个低频子带和三个方向的高频子带，再对低频子带继续分解。小波基在空间上局部、在尺度上多级，平坦区域由少数粗尺度系数表示，边缘只影响其附近的细尺度系数，既不需要人为分块，也就不会产生沿固定网格的块效应；低码率下的主要失真变为边缘附近的模糊和振铃。标准提供两种小波：可逆的整数5/3小波用于无损压缩，9/7小波用于有损压缩。量化是带死区的标量量化，每个子带一个步长。

\subsubsection{量化系数的符号化}
\label{sec:vision-restoration-symbolization}

第二步要把一块整数排成适合熵编码的符号序列。量化后的非零系数集中在左上角的低频区域，右下角大多为零。\term{之字形扫描}（zigzag scan）按频率从低到高的顺序沿反对角线来回读取64个系数，使非零值集中在序列前端、零值连成长串，扫描次序见图\ref{fig:vision-zigzag}。上例中五个非零系数在扫描序列中的位置是0、1、6、15和28。

\begin{figure}[htbp]
 \centering
 % 8x8系数块的之字形扫描次序；格内数字为扫描序号，带色格为正文算例中非零量化系数的位置（序号0,1,6,15,28）。宽约62mm。
\begin{tikzpicture}[x=7mm,y=-7mm,font=\figurefont\footnotesize]
 \foreach \j in {0,1,3,5,7} {
   \fill[FigureBlueFill] (\j,0) rectangle (\j+1,1);
 }
 \foreach \i in {0,...,7} {\foreach \j in {0,...,7} {
   \draw[BookRule,line width=.5pt] (\j,\i) rectangle (\j+1,\i+1);
 }}
 \foreach \j in {0,1,3,5,7} {
   \draw[FigureBlue,line width=.8pt] (\j,0) rectangle (\j+1,1);
 }
 \draw[BookAmber,line width=.6pt,opacity=.55] plot coordinates {(0.5,0.5)(1.5,0.5)(0.5,1.5)(0.5,2.5)(1.5,1.5)(2.5,0.5)(3.5,0.5)(2.5,1.5)(1.5,2.5)(0.5,3.5)(0.5,4.5)(1.5,3.5)(2.5,2.5)(3.5,1.5)(4.5,0.5)(5.5,0.5)(4.5,1.5)(3.5,2.5)(2.5,3.5)(1.5,4.5)(0.5,5.5)(0.5,6.5)(1.5,5.5)(2.5,4.5)(3.5,3.5)(4.5,2.5)(5.5,1.5)(6.5,0.5)(7.5,0.5)(6.5,1.5)(5.5,2.5)(4.5,3.5)(3.5,4.5)(2.5,5.5)(1.5,6.5)(0.5,7.5)(1.5,7.5)(2.5,6.5)(3.5,5.5)(4.5,4.5)(5.5,3.5)(6.5,2.5)(7.5,1.5)(7.5,2.5)(6.5,3.5)(5.5,4.5)(4.5,5.5)(3.5,6.5)(2.5,7.5)(3.5,7.5)(4.5,6.5)(5.5,5.5)(6.5,4.5)(7.5,3.5)(7.5,4.5)(6.5,5.5)(5.5,6.5)(4.5,7.5)(5.5,7.5)(6.5,6.5)(7.5,5.5)(7.5,6.5)(6.5,7.5)(7.5,7.5)};
 \foreach \i/\j/\k in {0/0/0,0/1/1,1/0/2,2/0/3,1/1/4,0/2/5,0/3/6,1/2/7,2/1/8,3/0/9,4/0/10,3/1/11,2/2/12,1/3/13,0/4/14,0/5/15,1/4/16,2/3/17,3/2/18,4/1/19,5/0/20,6/0/21,5/1/22,4/2/23,3/3/24,2/4/25,1/5/26,0/6/27,0/7/28,1/6/29,2/5/30,3/4/31,4/3/32,5/2/33,6/1/34,7/0/35,7/1/36,6/2/37,5/3/38,4/4/39,3/5/40,2/6/41,1/7/42,2/7/43,3/6/44,4/5/45,5/4/46,6/3/47,7/2/48,7/3/49,6/4/50,5/5/51,4/6/52,3/7/53,4/7/54,5/6/55,6/5/56,7/4/57,7/5/58,6/6/59,5/7/60,6/7/61,7/6/62,7/7/63} {
   \node[text=BookInk,fill=white,fill opacity=.8,text opacity=1,inner sep=.4pt] at (\j+.5,\i+.5) {\k};
 }
 \node[anchor=south] at (4,-.15) {水平频率增大 $\rightarrow$};
 \node[anchor=south,rotate=90] at (-.15,4) {$\leftarrow$ 垂直频率增大};
\end{tikzpicture}

 \caption{$8\times8$系数块的之字形扫描。格内数字是扫描序号，左上角为直流系数；带色的五格是正文例中量化后的非零系数，依次为$1,-46,11,-3,2$，它们都在第一行，扫描序号为0、1、6、15、28。}
 \label{fig:vision-zigzag}
\end{figure}

直流系数和交流系数用不同的方法符号化，分别利用两种不同的冗余。

相邻块的直流系数都反映局部平均亮度，彼此接近，JPEG只编码它与前一块直流系数之差。这是\term{预测编码}（predictive coding）的特例：用已经编码的值预测当前值，只传预测误差。一般形式称为差分脉冲编码调制（differential pulse-code modulation, DPCM），设预测为前一个值，误差为$e_i=x_i-x_{i-1}$。招牌一行像素$(100,101,101,102,103,103,105,104)$的误差序列是$(100,1,0,1,1,0,2,-1)$，除第一个值外都集中在0附近，熵比原始灰度低得多。若预测误差还要量化，预测必须基于解码端也能得到的重建值，否则误差会逐步累积。闭环形式为
\begin{equation}
 \hat e_i=Q\bigl(x_i-\tilde x_{i-1}\bigr),\qquad
 \tilde x_i=\tilde x_{i-1}+\hat e_i\text{，}
 \label{eq:vision-dpcm}
\end{equation}
其中$\tilde x_i$是重建值。两式相减得$x_i-\tilde x_i=(x_i-\tilde x_{i-1})-\hat e_i$，即每个位置的重建误差恰好等于该位置的量化误差，不随$i$累积；步长为$\Delta$的取整量化下，误差不超过$\Delta/2$。若改用原值$x_{i-1}$预测，编码端和解码端的预测值不同，量化误差会一路叠加。上例的直流系数若前一块是纯暗背景，即每个样本为$-68$，其直流系数为$8\times(-68)=-544$，量化为$-34$；当前块为1，于是只编码差值$1-(-34)=35$。

交流系数中大段的零用\term{游程编码}（run-length coding）表示：不逐个写出零，而是记录“连续几个零之后出现一个什么非零值”。JPEG把每个非零交流系数写成一个符号对（游程，位数）加上若干附加位：游程是它前面的零的个数（0至15），位数$s$是表示其幅值所需的比特数，满足$2^{s-1}\leq|v|\leq2^s-1$；附加位对正数是$v$的二进制，对负数是$v+2^s-1$的二进制。超过15个零用一个专门的符号表示16个零，块末尾剩下的零全部用块结束符号表示。上例的交流部分：$-46$前无零，幅值在32到63之间，符号为$(0,6)$，附加位为$-46+63=17$即$010001$；$11$前有4个零（序号2至5），符号为$(4,4)$，附加位$1011$；$-3$前有8个零，符号为$(8,2)$，附加位$00$；$2$前有12个零，符号为$(12,2)$，附加位$10$；之后是块结束符号。原来的63个交流系数变成5个符号和14个附加位。直流差值35的位数为6，附加位$100011$。

\subsubsection{符号的变长编码}
\label{sec:vision-restoration-huffman}

第三步要为每个符号分配码字，让常见的符号短、罕见的符号长。\term{Huffman编码}（Huffman coding）给出平均码长最短的前缀码\scite{vision-huffman1952}它从概率最小的符号开始，自底向上构造一棵二叉树，如算法\ref{alg:vision-huffman}所示。

\begin{algorithm}[Huffman编码的构造]
 \label{alg:vision-huffman}
 \begin{algorithmic}[1]
  \Require 符号集合及其概率$p(s)$，并列时的合并约定
  \Ensure 每个符号的二进制码字
  \State 为每个符号建立一个叶结点，权重为其概率，全部放入待合并集合
  \While{待合并集合中多于一个结点}
   \State 取出权重最小的两个结点$a,b$
   \State 新建父结点，权重为二者之和，两条边分别标0和1
   \State 把父结点放回待合并集合
  \EndWhile
  \State 每个符号的码字为从根到该叶结点路径上的比特串
 \end{algorithmic}
\end{algorithm}

用DPCM误差作例子。设五个误差符号$0,+1,-1,+2,-2$的概率为$0.4,0.2,0.2,0.1,0.1$（教学构造）。第一步合并$+2$与$-2$，得到权重0.2的结点；此时有$0$（0.4）、$+1$（0.2）、$-1$（0.2）和新结点（0.2）。第二步在三个0.2中合并$+1$与$-1$，得到0.4；第三步合并剩下的0.2结点与一个0.4结点，这里选$\{+1,-1\}$那一个，得到0.6；最后合并0.4与0.6。码字如图\ref{fig:vision-huffman-tree}：$0\mapsto0$，$+1\mapsto100$，$-1\mapsto101$，$+2\mapsto110$，$-2\mapsto111$。平均码长为$0.4\times1+0.6\times3=2.2$比特，而熵为$-(0.4\log_20.4+2\times0.2\log_20.2+2\times0.1\log_20.1)\approx2.12$比特，满足$H\leq L<H+1$；若用3位定长码则需3比特。并列时若改变合并选择，码长可能变为$(2,2,2,3,3)$，平均码长同样是2.2，可见Huffman码不唯一，但最优平均码长唯一。

\begin{figure}[htbp]
 \centering
 % 五个预测残差符号的Huffman树，概率为教学构造(0.4,0.2,0.2,0.1,0.1)；并列时按正文约定合并。宽约90mm。
\begin{tikzpicture}[x=1cm,y=1cm,font=\figurefont\small,
  leaf/.style={draw=FigureBlue,fill=FigureBlueFill,rounded corners=1.8pt,
    minimum width=13mm,minimum height=7mm,align=center,inner sep=1pt},
  inner/.style={draw=FigureYellow,fill=FigureYellowFill,circle,
    minimum size=8mm,inner sep=0pt},
  edge/.style={draw=BookMuted,line width=.8pt},
  bit/.style={text=BookTeal,font=\figurefont\footnotesize}]
 \node[inner] (r) at (3.2,3.6) {$1$};
 \node[leaf] (z) at (1.2,2.4) {$0$\\$0.4$};
 \node[inner] (n3) at (5.2,2.4) {$0.6$};
 \node[inner] (n2) at (3.6,1.2) {$0.4$};
 \node[inner] (n1) at (6.8,1.2) {$0.2$};
 \node[leaf] (p1) at (2.8,0) {$+1$\\$0.2$};
 \node[leaf] (m1) at (4.4,0) {$-1$\\$0.2$};
 \node[leaf] (p2) at (6.0,0) {$+2$\\$0.1$};
 \node[leaf] (m2) at (7.6,0) {$-2$\\$0.1$};
 \draw[edge] (r)--node[bit,above left] {0} (z);
 \draw[edge] (r)--node[bit,above right] {1} (n3);
 \draw[edge] (n3)--node[bit,above left] {0} (n2);
 \draw[edge] (n3)--node[bit,above right] {1} (n1);
 \draw[edge] (n2)--node[bit,left] {0} (p1);
 \draw[edge] (n2)--node[bit,right] {1} (m1);
 \draw[edge] (n1)--node[bit,left] {0} (p2);
 \draw[edge] (n1)--node[bit,right] {1} (m2);
 \node[text=BookMuted,font=\figurefont\footnotesize,anchor=north] at (1.2,1.85) {码字 0};
 \node[text=BookMuted,font=\figurefont\footnotesize,anchor=north] at (5.2,-.55) {码字 100, 101, 110, 111};
\end{tikzpicture}

 \caption{五个预测误差符号的Huffman树。方框为叶结点，标出符号及概率；圆形为合并后的内部结点，标出权重；边上的0和1连成码字。概率为教学构造，并列时按正文的约定合并。}
 \label{fig:vision-huffman-tree}
\end{figure}

Huffman码为什么最优，可以用交换论证说明思路。在最优前缀码中，概率越小的符号码字不会越短，否则交换两者的码字就能缩短平均码长；概率最小的两个符号可以取为同一父结点下码长最长的一对兄弟。把这对兄弟合并为一个概率为二者之和的符号后，原问题的平均码长等于新问题的平均码长加上这两个概率之和，后者是常数，所以原问题的最优码由新问题的最优码展开得到，这正是算法每一步所做的事。解码时从根出发，按读到的比特向下走，到达叶结点就输出一个符号并回到根；前缀性质保证这一过程不需要分隔符。

JPEG把Huffman编码用在上一步得到的符号上：直流差值的位数和交流的（游程，位数）对各用一张码表编码，附加位原样写出。码表写在文件头中，可以使用标准附录给出的示例表，也可以先统计当前图像的符号频率再生成专用表，后者需要两遍扫描，但码率更低。解码端读入码表，按前缀码解出符号，再读入相应数目的附加位，就恢复了量化系数。

Huffman编码有两个限制。码字长度必须是整数比特，每个符号至少1比特，当某个符号的概率很大时，例如块中大量出现的零游程，理想码长$-\log_2p$远小于1，整数码长造成的浪费就很明显。码表又是静态的，不能随图像局部统计的变化而调整。JPEG 2000因此在熵编码环节采用EBCOT方案\scite{vision-taubman2000}它把每个子带划分成独立编码的码块（常用$64\times64$），在每个码块内按位平面从最高位到最低位扫描，每个位平面分三遍处理，按邻近系数是否已经显著来选择上下文，再用自适应的二值算术编码器（MQ编码器）编码；算术编码不受整数码长的限制，上下文又让概率随局部内容变化（原理见\ref{sec:vision-restoration-entropy-coding}节）。按位平面编码使每个码块的比特流可以在任意位置截断，编码器在所有码块编完之后，再按率失真斜率决定各码块保留多少比特，以精确达到目标码率；同一个码流因此可以按质量或分辨率逐级解码。

两种经典编码器的差别归纳在表\ref{tab:vision-classic-codec-compare}中。JPEG 2000在低码率下避免了块效应，码率控制也更精确，代价是小波变换与位平面算术编码的计算量明显更大，而在中高码率的普通照片上，两者的观感差距并不大，这也是JPEG至今仍被广泛使用的原因之一。

\begin{table}[htbp]
 \centering
 \small
 \begin{tabularx}{\linewidth}{@{}p{14mm}XX@{}}
  \toprule
  环节 & JPEG基线 & JPEG 2000\\
  \midrule
  颜色 & YCbCr，色度4:2:0 & 不可逆或可逆的分量变换\\
  变换 & $8\times8$块离散余弦变换 & 整图或大图块的多级小波，9/7或5/3\\
  量化 & 每个系数一个量化表步长 & 每个子带一个步长，带死区\\
  符号化 & 之字形扫描，直流差分，交流游程 & 码块内按位平面三遍扫描\\
  熵编码 & 静态Huffman码表 & 上下文自适应的二值算术编码\\
  码率控制 & 质量因子间接控制 & 码块截断，按率失真斜率分配\\
  典型失真 & 块效应、振铃 & 模糊、边缘振铃\\
  \bottomrule
 \end{tabularx}
 \caption{两种经典图像编码器的比较。两者都把变换、量化、熵编码分开设计；JPEG 2000还支持无损模式和按质量或分辨率逐级解码。}
 \label{tab:vision-classic-codec-compare}
\end{table}

\subsection{基于端到端学习的压缩}
\label{sec:vision-restoration-learned-codec}

经典编码器的每个环节都由人分别设计：离散余弦变换或小波是固定的线性变换，量化表凭视觉实验确定，Huffman码表或算术编码的上下文靠人工规则。各环节单独看都有依据，合在一起却不一定最优，例如固定线性变换无法适应图像中的非线性统计结构，静态码表也不知道当前位置是平坦区还是纹理区。学习式压缩保留“变换—量化—熵编码”的框架，但把变换换成神经网络，把概率模型换成可以随内容变化的学习模型，并让它们在同一个率失真目标下联合训练。

以下以尺度超先验方法为主线\scite{vision-hyperprior2018}该方法建立在更早的端到端优化压缩框架之上\scite{vision-balle2017}一次压缩分四步：分析变换$g_{\mathrm{a}}$把图像$\bm{X}$变为潜变量$\bm{y}$，取整得到整数符号$\hat{\bm{y}}$；概率模型为每个符号给出概率$P(\hat{\bm{y}})$；熵编码器按这些概率把符号写成比特；解码端按相同概率读回符号，再由合成变换$g_{\mathrm{s}}$得到重建$\widehat{\bm{X}}=g_{\mathrm{s}}(\hat{\bm{y}})$。若熵编码接近理想，码长约为$-\log_2P(\hat{\bm{y}})$比特，这是式\eqref{eq:vision-entropy}所述下界在单个序列上的形式。于是率失真权衡可以直接写成训练目标
\begin{equation}
 L=\underbrace{\mathbb{E}_{\bm{X}}\bigl[-\log_2P(\hat{\bm{y}})\bigr]}_{R}
  +\lambda\,\underbrace{\mathbb{E}_{\bm{X}}\bigl[d(\bm{X},\widehat{\bm{X}})\bigr]}_{D}\text{，}
 \qquad \hat{\bm{y}}=\operatorname{round}\bigl(g_{\mathrm{a}}(\bm{X})\bigr)\text{，}
 \label{eq:vision-rate-distortion}
\end{equation}
其中$d$是失真度量，例如均方误差，$\lambda>0$决定权衡点，作用相当于JPEG的质量因子。码率项依赖概率模型，失真项依赖两个变换，取整同时影响两项。四个步骤因此不是各自独立设计，而是由同一个目标联合训练：变换决定符号的分布，概率模型决定码长，熵编码只负责把理想码长兑现成实际比特，解码重建则要求两端算出完全相同的概率。

这个训练目标还可以从变分自编码器的角度理解。把$g_{\mathrm{s}}$看成生成模型，把“$g_{\mathrm{a}}$加均匀噪声”看成近似后验$q(\tilde{\bm{y}}\mid\bm{X})=\prod_i\mathcal{U}(\tilde y_i\mid y_i-\tfrac12,y_i+\tfrac12)$（加噪的理由在下一小节说明），变分推断要最小化$q$与真实后验之间的KL散度的期望：
\begin{equation}
 \mathbb{E}_{\bm{X}}\,\mathbb{E}_{q}\Bigl[
 \log q(\tilde{\bm{y}}\mid\bm{X})
 -\log p(\bm{X}\mid\tilde{\bm{y}})
 -\log p(\tilde{\bm{y}})\Bigr]+\text{常数}\text{。}
 \label{eq:vision-compression-vae}
\end{equation}
第一项为零，因为单位宽均匀密度在支撑集内恒为1。若取$p(\bm{X}\mid\tilde{\bm{y}})$为均值$g_{\mathrm{s}}(\tilde{\bm{y}})$、协方差$(2\lambda)^{-1}\bm{I}$的高斯分布，第二项等于$\lambda\lVert\bm{X}-g_{\mathrm{s}}(\tilde{\bm{y}})\rVert_2^2$加常数，即加权的平方误差失真；第三项是用先验$p(\tilde{\bm{y}})$编码潜变量的交叉熵，即码率（以自然对数计，除以$\ln2$换成比特）。这说明式\eqref{eq:vision-rate-distortion}在加噪近似下就是一种变分自编码器的训练目标，区别在于$\lambda$由应用选择，而不是由数据的真实噪声决定；变分自编码器的一般原理见前文已引的卷六“概率生成模型”章。下面按编码流程依次展开这四个步骤，并在最后与经典路线逐环节对照。

\subsubsection{图像的变换与量化}
\label{sec:vision-restoration-quantization}

第一步要把像素变成便于压缩的符号。JPEG用固定的离散余弦变换去掉相关性，学习式压缩则让变换由数据和率失真目标共同决定：
\begin{equation}
 \bm{y}=g_{\mathrm{a}}(\bm{X}),\qquad
 \hat y_i=\operatorname{round}(y_i),\qquad
 \widehat{\bm{X}}=g_{\mathrm{s}}(\hat{\bm{y}})\text{。}
 \label{eq:vision-transform-coding}
\end{equation}
量化把连续数映射到有限精度的离散值，舍弃区间内的差别。此处以单位步长取整为例；步长的作用由分析与合成变换吸收，网络可以把某个分量放大，相当于对它用更细的量化，这对应JPEG量化表中逐系数不同的步长。

用两个像素作可复算的变换例。对$\bm{x}=(0.2,0.8)$，令$y_1=2(x_1+x_2)=2$、$y_2=2(x_2-x_1)=1.2$。取整得到$(2,1)$，合成按$x_1=(y_1-y_2)/4$、$x_2=(y_1+y_2)/4$得到$(0.25,0.75)$。均方误差为0.0025。若把变换系数从2改成4，$y_2=2.4$取整为2，重建为$(0.1875,0.8125)$，误差降到约$1.6\times10^{-4}$，但符号的取值范围变大，编码需要更多比特。这里分析与合成是手工线性变换，学习式网络用数据寻找更适合图像的非线性变换；量化造成不可逆信息损失的原因相同。

式\eqref{eq:vision-rate-distortion}不能直接用梯度训练，因为取整的导数几乎处处为零。原方法在训练时用加均匀噪声代替取整：$\tilde y_i=y_i+u_i$，$u_i\sim\mathcal{U}(-\tfrac12,\tfrac12)$相互独立。带波浪号的连续量用于训练，带帽号的整数用于实际编码。这个替代有一个精确的依据。设$y_i$的密度为$p_{y}$，取整后符号为$k$的概率是
\[
 P(\hat y_i=k)=\int_{k-1/2}^{k+1/2}p_{y}(v)\,\mathrm{d}v\text{；}
\]
而$\tilde y_i$的密度是$p_y$与单位宽均匀密度的卷积，$p_{\tilde y}(t)=\int_{-1/2}^{1/2}p_y(t-s)\,\mathrm{d}s=\int_{t-1/2}^{t+1/2}p_y(v)\,\mathrm{d}v$。两式在$t=k$处相同，即加噪变量的密度在整数点的值正好等于取整符号的概率。所以训练中用$-\log_2p_{\tilde y}(\tilde y_i)$估计码率，与实际编码的码长是同一个函数在相邻点上的取值。加噪后的失真$d(\bm{X},g_{\mathrm{s}}(\tilde{\bm{y}}))$也近似于取整后的失真，但二者不完全相同，交付前仍须用真正的取整、编码和解码检查。这里的均匀噪声与表\ref{tab:vision-noise-models}中的均匀噪声是同一个模型：单位步长量化误差在$[-\tfrac12,\tfrac12]$内，方差$1/12$。

尺度超先验方法中，$g_{\mathrm{a}}$由四层步幅为2的$5\times5$卷积组成，前三层后各接一个广义除法归一化层（各层配置见表\ref{tab:vision-hyperprior-layers}）。输入$h\times w\times3$的图像，前三层输出$N$个通道，最后一层输出$M$个通道，得到$\bm{y}\in\mathbb{R}^{\frac{h}{16}\times\frac{w}{16}\times M}$。$N,M$随$\lambda$选取：码率较低的5个$\lambda$用$N=128,M=192$，较高的3个用$N=192,M=320$。$g_{\mathrm{s}}$结构对称，用步幅为2的转置卷积上采样，用近似逆运算替换归一化层，最后输出3个通道。\term{广义除法归一化}（generalized divisive normalization, GDN）把每个通道的响应除以同一位置各通道响应平方的加权和：
\begin{equation}
 v_i=\frac{u_i}{\bigl(\beta_i+\sum_j\gamma_{i,j}u_j^2\bigr)^{1/2}}\text{，}
 \label{eq:vision-gdn}
\end{equation}
其中$u_i$是同一空间位置第$i$个通道的卷积输出，$\beta_i>0$与$\gamma_{i,j}\geq0$是学习参数。它与批归一化不同：批归一化训练后变成固定的仿射变换，GDN的除数随当前位置的响应而变，相当于局部增益控制。作者此前的工作表明，这种归一化能有效削弱自然图像局部统计中的相关性，使变换后的分量更接近独立\scite{vision-balle2017}与离散余弦变换相比，$g_{\mathrm{a}}$是非线性的，并且以16倍下采样覆盖比$8\times8$块大得多的范围，没有固定的块边界。

设一幅$256\times256$的图像，$M=192$，则$\bm{y}$为$16\times16\times192$，共49152个符号，比原图的196608个像素值少四分之三。但符号数减少不等于压缩完成：每个符号要多少比特，取决于下一步的概率模型。

\subsubsection{符号的概率估计}
\label{sec:vision-restoration-bitstream}
\label{sec:vision-restoration-hyperprior}

第二步要为每个整数符号给出概率，编码器与解码器必须共享这组概率。概率越准，常见符号的码长越短，平均码率越低。端到端优化压缩的第一版用因子化先验：每个通道的所有位置共享一个学习得到的一维分布，各符号独立\scite{vision-balle2017}它的缺点在训练后的潜变量中能直接看到：非零符号集中在边缘和纹理区域，平坦区域几乎全为零；也就是说，符号的幅度尺度随位置变化，且在空间上成片相关。一个对所有位置相同的分布只能折中，平坦区域的零符号概率被估低，纹理区域的大幅值符号概率也被估低，两处都多花比特。这与JPEG静态码表的局限相同。

传统编码器用辅助信息解决类似问题：例如先告诉解码器每个块用哪种划分或量化表，再按此解码。\term{辅助信息}（side information）是编码器额外发送、帮助解码器选择概率模型的信息。尺度超先验把辅助信息也交给网络学习：用一组更小的潜变量$\bm{z}$描述$\bm{y}$在各位置的尺度，先传$\bm{z}$，再以它为条件传$\bm{y}$。在给定$\bm{z}$时，$\bm{y}$的各分量被假设为相互独立，空间相关性全部由$\bm{z}$承担。

\term{尺度超先验}（scale hyperprior）由两个小变换组成，它们与主变换的逐层配置一并列在表\ref{tab:vision-hyperprior-layers}中。超分析变换$h_{\mathrm{a}}$先对$\bm{y}$取绝对值，使输入只反映幅度，再经一层步幅为1的$3\times3$卷积和两层步幅为2的$5\times5$卷积，层间接ReLU，得到$\bm{z}\in\mathbb{R}^{\frac{h}{64}\times\frac{w}{64}\times N}$，取整为$\hat{\bm{z}}$。超合成变换$h_{\mathrm{s}}$结构对称：两层步幅为2的$5\times5$转置卷积上采样，再接一层步幅为1、输出$M$个通道的$3\times3$卷积，每层后都接ReLU，因此输出的尺度$\bm{\sigma}=h_{\mathrm{s}}(\hat{\bm{z}})$非负，与$\bm{y}$同形状，每个$\sigma_i$控制一个主符号的分布。对前面的$256\times256$例子，$\hat{\bm{z}}$为$4\times4\times128$，共2048个符号，不到主符号数的二十分之一。

\begin{table*}[tbp]
 \centering
 \small
 \begin{tabularx}{\linewidth}{@{}p{22mm}Xp{24mm}p{12mm}p{12mm}p{26mm}@{}}
  \toprule
  变换 & 层类型 & 卷积核与步长 & 输出通道 & 重复次数 & 输出尺寸\\
  \midrule
  分析$g_{\mathrm{a}}$ & 卷积下采样，GDN & $5\times5$，2 & $N$ & 3 & $\frac h8\times\frac w8\times N$\\
   & 卷积下采样 & $5\times5$，2 & $M$ & 1 & $\frac h{16}\times\frac w{16}\times M$\\
  超分析$h_{\mathrm{a}}$ & 取绝对值，卷积，ReLU & $3\times3$，1 & $N$ & 1 & $\frac h{16}\times\frac w{16}\times N$\\
   & 卷积下采样，ReLU & $5\times5$，2 & $N$ & 1 & $\frac h{32}\times\frac w{32}\times N$\\
   & 卷积下采样 & $5\times5$，2 & $N$ & 1 & $\frac h{64}\times\frac w{64}\times N$\\
  超合成$h_{\mathrm{s}}$ & 转置卷积上采样，ReLU & $5\times5$，2 & $N$ & 2 & $\frac h{16}\times\frac w{16}\times N$\\
   & 卷积，ReLU & $3\times3$，1 & $M$ & 1 & $\frac h{16}\times\frac w{16}\times M$\\
  合成$g_{\mathrm{s}}$ & 转置卷积上采样，逆GDN & $5\times5$，2 & $N$ & 3 & $\frac h2\times\frac w2\times N$\\
   & 转置卷积上采样 & $5\times5$，2 & 3 & 1 & $h\times w\times3$\\
  \bottomrule
 \end{tabularx}
 \caption{尺度超先验模型的逐层结构，依据原文的网络结构图。码率较低的5个模型取$N=128,M=192$，较高的3个取$N=192,M=320$；因子化先验模型使用相同的$g_{\mathrm{a}}$与$g_{\mathrm{s}}$。}
 \label{tab:vision-hyperprior-layers}
\end{table*}

原方法以零均值高斯描述$y_i$，并将密度在单位量化区间上积分；按上文的加噪论证，这正是高斯与均匀分布卷积后在整数点的值。对整数$k$，离散概率为
\begin{equation}
 P_i(k\mid\hat{\bm{z}})
 =\int_{k-1/2}^{k+1/2}
   \frac{1}{\sqrt{2\pi}\sigma_i}
   \exp\!\left(-\frac{v^2}{2\sigma_i^2}\right)\,\mathrm{d}v
 =\Phi\!\left(\frac{k+1/2}{\sigma_i}\right)
  -\Phi\!\left(\frac{k-1/2}{\sigma_i}\right)\text{，}
 \label{eq:vision-hyperprior-probability}
\end{equation}
其中$\Phi$是标准正态累计分布函数。不能把密度在整数处的值直接当作符号概率；密度有单位，且在整数点求和不保证等于1。

算两个数可以看清尺度的作用。平坦区域$\sigma_i=0.5$时，$P_i(0)=\Phi(1)-\Phi(-1)\approx0.683$，零符号只需约0.55 bit；$P_i(1)=\Phi(3)-\Phi(1)\approx0.157$，约2.67 bit。纹理区域$\sigma_i=2$时，$P_i(0)=\Phi(0.25)-\Phi(-0.25)\approx0.197$，约2.34 bit；$P_i(1)\approx0.175$，约2.52 bit。若两处共用$\sigma=2$，平坦区每个零符号要多花约1.8 bit；共用$\sigma=0.5$，纹理区出现的大幅值又会很贵。按位置给出尺度，正是为了让每处的概率贴近该处的实际分布。平坦区零符号只需0.55 bit，低于Huffman码每个符号至少1比特的限制，这也是学习式压缩配合算术编码的原因。

辅助潜变量$\hat{\bm{z}}$本身也要编码，它采用因子化分布：同一通道的所有位置共享一个学习得到的非参数一维密度，同样与单位宽均匀密度卷积后在整数点取值，记为$P_j^{z}$。共享分布使模型能用于任意尺寸的图像。一次图像的理想估计码长为
\begin{equation}
 R_{\mathrm{est}}=
 -\sum_j\log_2 P_j^{z}(\hat z_j)
 -\sum_i\log_2 P_i(\hat y_i\mid\hat{\bm{z}})\text{。}
 \label{eq:vision-hyperprior-rate}
\end{equation}
训练时把两部分都换成加噪潜变量的连续密度，与失真共同构成目标，同时更新$g_{\mathrm{a}},g_{\mathrm{s}},h_{\mathrm{a}},h_{\mathrm{s}}$和$\bm{z}$的因子化分布参数。从变分角度看，这相当于在式\eqref{eq:vision-compression-vae}中把近似后验扩展为$\tilde{\bm{y}}$与$\tilde{\bm{z}}$的均匀分布之积，把先验写成$p(\tilde{\bm{y}}\mid\tilde{\bm{z}})p(\tilde{\bm{z}})$，KL中便多出一项编码$\tilde{\bm{z}}$的交叉熵。辅助流的项不可省略，否则网络可以把大量信息转移到$\bm{z}$而不付出代价，目标就不再衡量总压缩成本。

训练图像是从网上收集的约100万幅彩色JPEG照片，先随机缩小以减弱原有的压缩伪影，再裁取$256\times256$的图块；批量为8，用Adam训练，学习率$10^{-4}$，作者发现批归一化和学习率衰减都没有带来收益。作者用8个$\lambda$值覆盖不同码率，并分别以均方误差和MS-SSIM为失真度量，与同架构的因子化先验模型对照。在Kodak数据集上，以均方误差训练、按PSNR评价时，因子化先验模型优于JPEG但不及基于HEVC的BPG，加入超先验后接近BPG，高码率处表现更好；以MS-SSIM训练、按MS-SSIM评价时，超先验在所有码率上都带来进一步提升。每个$\lambda$训练一个模型，换一个码率就要换一套参数。

条件化是否划算，要把辅助流计入后再看。假设共享全局分布编码主符号需要100 bit，传入尺度辅助信息后主流降至72 bit。若辅助流为12 bit，总计84 bit，有16 bit收益；若辅助流需要35 bit，总计107 bit，反而增大。只有节省的主流成本超过辅助成本时，这种条件化才划算；训练目标同时含两项，网络会自己寻找这个平衡。表\ref{tab:vision-entropy-model-compare}对照了两种概率模型。

\begin{table}[htbp]
 \centering
 \small
 \begin{tabularx}{\linewidth}{@{}p{16mm}XX@{}}
  \toprule
  比较维度 & 因子化先验 & 尺度超先验\\
  \midrule
  符号分布 & 每通道一个共享分布，与位置无关 & 每个符号一个高斯尺度，由$\hat{\bm{z}}$预测\\
  空间相关 & 不建模 & 由辅助潜变量承担\\
  额外比特 & 无 & 辅助流，约占主符号数的二十分之一\\
  解码依赖 & 各符号可直接解码 & 先解$\hat{\bm{z}}$，再算尺度，再解$\hat{\bm{y}}$\\
  原文结果 & 优于JPEG，不及BPG（PSNR） & 接近BPG，高码率更好（PSNR）\\
  \bottomrule
 \end{tabularx}
 \caption{两种概率模型的比较。两者使用相同的分析与合成变换，差别只在符号概率怎样随位置变化；结果均指原论文在Kodak数据集上以均方误差训练的模型。}
 \label{tab:vision-entropy-model-compare}
\end{table}

\subsubsection{符号的熵编码}
\label{sec:vision-restoration-entropy-coding}

第三步要把整数符号和它们的概率变成实际的比特串。\term{熵编码}（entropy coding）根据共享概率把离散符号无损地转成比特；这里“无损”指整数符号能够原样恢复，前面的量化仍然有损。概率模型给出的$-\log_2P$只是理想码长，熵编码器要做的是让实际码长尽量接近它，并且能被唯一解码。上一小节的平坦区零符号只需0.55 bit，Huffman编码做不到，因为它的码字必须是整数比特。

\term{算术编码}（arithmetic coding）的一种解释是不断缩小$[0,1)$区间：每读入一个符号，就把当前区间按各符号的概率切分，保留该符号对应的一段；最后用落在最终区间内的一个二进制前缀标识整个序列。最终区间的宽度等于序列概率的乘积，所以需要的位数约为$-\log_2$（序列概率），与理想码长只差常数位。一个符号不再对应整数个比特，整个序列共享比特，因而能逼近任意小于1比特的理想码长。

设三个教学符号$A,B,C$的概率依次为$1/2,1/4,1/4$，基础区间为$[0,1/2)$、$[1/2,3/4)$、$[3/4,1)$。编码序列$ABC$时，$A$把当前区间缩为$[0,1/2)$，$B$选择其中第二段得到$[1/4,3/8)$，$C$再选择最后四分之一得到$[11/32,12/32)$。其宽度为$1/32$，对应理想码长5 bit；二进制前缀$01011$恰好表示这个区间。

若解码端已知长度为3，可用区间内代表值$11/32$恢复序列。它落在基础区间的$A$段，归一化到该段后得到$11/16$，落在$B$段；再次减去$1/2$并除以$1/4$，得到$3/4$，落在$C$段。区间统一左闭右开，因此边界也有确定归属。图\ref{fig:vision-arithmetic-coding}画出逐步缩小的区间；实际实现使用有限精度整数和归一化输出位，无需保存无限精度实数。

\begin{figure}[htbp]
 \centering
 % All intervals use the same coordinate scale; probabilities .5,.25,.25.
\begin{tikzpicture}[x=80mm,y=1mm,font=\figurefont\small]
 \foreach \row/\lab in {0/初始,1/$A$,2/$AB$,3/$ABC$} {
  \draw[BookRule,line width=.7pt] (0,{-15*\row})--(1,{-15*\row});
  \node[anchor=east] at (-.05,{-15*\row}) {\lab};
 }
 \fill[BookBlue!15] (0,-3) rectangle (1,3);
 \draw[BookBlue,line width=.8pt] (0,-3) rectangle (1,3);
 \fill[BookTeal!20] (0,-18) rectangle (.5,-12);
 \draw[BookTeal,line width=.8pt] (0,-18) rectangle (.5,-12);
 \fill[BookTeal!35] (.25,-33) rectangle (.375,-27);
 \draw[BookTeal,line width=.8pt] (.25,-33) rectangle (.375,-27);
 \fill[BookTeal] (.34375,-48) rectangle (.375,-42);
 \foreach \x/\lab in {0/0,.5/{1/2},1/1} {
  \node[anchor=south] at (\x,4) {$\lab$};
 }
 \node[anchor=west] at (.54,-15) {$[0,1/2)$};
 \node[anchor=west] at (.43,-30) {$[1/4,3/8)$};
 \node[anchor=west] at (.43,-45) {$[11/32,12/32)$};
 \draw[BookMuted,dashed,line width=.6pt] (.5,-18)--(.375,-27);
 \draw[BookMuted,dashed,line width=.6pt] (0,-18)--(.25,-27);
 \draw[BookMuted,dashed,line width=.6pt] (.25,-33)--(.34375,-42);
 \draw[BookMuted,dashed,line width=.6pt] (.375,-33)--(.375,-42);
\end{tikzpicture}

 \caption{固定概率下序列$ABC$的算术编码教学例。每行共享$[0,1)$坐标，实色区间为读入当前前缀后仍可能的范围。最终区间宽$1/32$，可由5位前缀$01011$表示；序列长度和概率表在本例中预先共享，不计入这5位。}
 \label{fig:vision-arithmetic-coding}
\end{figure}

在尺度超先验模型中，每个符号的概率由式\eqref{eq:vision-hyperprior-probability}给出，不同位置的概率表不同。原文实现采用一种简单的非自适应二值算术编码器：把每个整数符号写成二进制，从最高位到最低位逐位编码，每位的概率由该符号的离散分布直接算出。因为解码$\hat{\bm{y}}$时尺度$\hat{\bm{\sigma}}$已经由$\hat{\bm{z}}$确定，编码器不需要处理符号之间的条件依赖，也不需要像JPEG 2000的MQ编码器那样在线调整概率。更一般的实现会把式\eqref{eq:vision-hyperprior-probability}离散化成整数累计概率表，规定精度、尾部截断与逃逸机制，并使两端采用一致的符号顺序。零概率不能分配给实际出现的符号，否则无法编码。有限精度概率、流终止、填充以及尺寸和版本元数据，都会使实际字节数偏离式\eqref{eq:vision-hyperprior-rate}的估计。

真实文件应先提供足以解析的尺寸和模型约定，再存辅助流与主流，或用等价可解析的组织方式记录偏移。若一张$256\times256$图像的全部文件为4096 byte，则按式\eqref{eq:vision-bpp}，实际每像素比特数为$8\times4096/(256\times256)=0.5$，分母不再乘RGB的通道数。模型权重若由收发双方预先共享，常规每图码率不含权重；若每张图传一个专用模型，则这项成本也必须纳入所声明的传输预算。JPEG文件头中的量化表和Huffman码表也属于每图都要付出的成本。

\subsubsection{码流的解码重建}
\label{sec:vision-restoration-decoding}

第四步由解码端完成，它只拿到比特流和事先共享的模型。解码按依赖顺序执行：用共享的因子化分布从辅助流$S_z$恢复$\hat{\bm{z}}$，由$h_{\mathrm{s}}$算出尺度$\hat{\bm{\sigma}}$，据此构造每个主符号的概率表，从主流$S_y$恢复$\hat{\bm{y}}$，最后由$g_{\mathrm{s}}$得到$\widehat{\bm{X}}$。图\ref{fig:vision-hyperprior-codec}画出了两端的依赖。主流解码必须等待辅助流，因为没有尺度就没有概率表。

\begin{figure*}[tbp]
 \centering
 % Encoder/decoder dependencies. All ordinary edges are axis-aligned.
\begin{tikzpicture}[x=1mm,y=1mm,font=\figurefont\small,
 b/.style={draw=BookBlue,fill=BookBlue!8,minimum width=13mm,
   minimum height=8mm,inner sep=1mm,line width=.8pt,align=center},
 o/.style={b,draw=BookTeal,fill=BookTeal!8},
 s/.style={draw=BookMuted,minimum width=12mm,minimum height=8mm,
   inner sep=1mm,line width=.8pt},
 a/.style={-{Stealth[length=1.8mm]},BookInk,line width=.9pt},
 c/.style={a,BookAmber,rounded corners=1mm}]
 \node at (5,0) (x) {$\bm{X}$};
 \node[b] at (22,0) (ga) {$g_{\mathrm{a}}$};
 \node[b] at (42,0) (qy) {$Q_y$};
 \node[b] at (64,0) (ey) {E};
 \node[s] at (83,0) (sy) {$S_y$};
 \node[o] at (102,0) (dy) {D};
 \node[o] at (123,0) (gs) {$g_{\mathrm{s}}$};
 \node at (144,0) (xh) {$\widehat{\bm{X}}$};
 \node[b] at (22,-32) (ha) {$h_{\mathrm{a}}$};
 \node[b] at (42,-32) (qz) {$Q_z$};
 \node[b] at (64,-32) (ez) {E};
 \node[s] at (83,-32) (sz) {$S_z$};
 \node[o] at (102,-32) (dz) {D};
 \node at (123,-32) (zh) {$\hat{\bm{z}}$};
 \node[b] at (42,-16) (hse) {$h_{\mathrm{s}}$};
 \node[o] at (123,-16) (hsd) {$h_{\mathrm{s}}$};
 \foreach \u/\v in {x/ga,ga/qy,qy/ey,ey/sy,sy/dy,dy/gs,gs/xh,
   ha/qz,qz/ez,ez/sz,sz/dz,dz/zh} {\draw[a] (\u)--(\v);}
 \draw[a] (ga)--(ha);
 \draw[c] (qz)--(hse);
 \draw[c] (hse.east)--(64,-16)--(ey.south);
 \draw[c] (zh)--(hsd);
 \draw[c] (hsd.west)--(102,-16)--(dy.south);
 \node[anchor=west,text=BookAmber] at (66,-13) {$\bm{\sigma}$};
 \node[anchor=east,text=BookAmber] at (100,-13) {$\bm{\sigma}$};
 \node[anchor=south] at (33,4) {$\bm{y}$};
 \node[anchor=south] at (112.5,4) {$\hat{\bm{y}}$};
 \node[anchor=north,text=BookMuted] at (42,-40) {编码端};
 \node[anchor=north,text=BookMuted] at (112,-40) {解码端};
\end{tikzpicture}

 \caption{尺度超先验的实际编解码依赖。$Q$为取整，E与D分别为熵编码和熵解码，$S_y,S_z$为两条比特流。两端都从同一组整数$\hat{\bm{z}}$经$h_{\mathrm{s}}$产生尺度，主流解码因此必须等待辅助流。图中省略参数与尺寸元数据，未画入不可传输的训练噪声或清晰参考。}
 \label{fig:vision-hyperprior-codec}
\end{figure*}
\FloatBarrier

编码端也必须用量化后的$\hat{\bm{z}}$计算尺度，不能用未经传输的连续$\bm{z}$，否则两侧的概率表不同，算术解码会从出错的符号起全部失效。同样的道理，$h_{\mathrm{s}}$在两端的计算结果必须逐位一致：若编码端和解码端在不同硬件上做浮点运算，$\hat{\bm{\sigma}}$的末位差异可能使离散化后的概率表不同。JPEG的解码只用整数查表和逆离散余弦变换，不同实现的重建可以有微小差别，却不会因此无法解析码流；这一逐位一致的要求是学习式编解码比传统编解码多出的工程约束。

可以用一次独立解码检查上述要求：只给解码程序文件、声明的模型参数和必要配置，确认它能复现输出。原图、训练参考和编码中间浮点值都不在它的可用输入中。一个只展示分析网络和合成网络前向结果的实验，用的是加噪或连续的潜变量，尚未完成这个检查。

重建质量还要回到用途。同样的0.5 bit/像素可以优先保留平滑颜色，也可以保留文本边缘，不同失真度量会形成不同取舍：原文以MS-SSIM训练的模型按PSNR评价时表现明显较差。应把解码图像送入固定的文字读取、识别或测量流程，检查招牌字符串、车牌边界及尺寸。感知上自然的纹理不一定有助于这些用途，较小的平均像素误差也可能掩盖一个关键数字的错误。视频还会利用帧间相关，运动流、残差流和参考帧依赖在\ref{chap:vision-video-restoration}章接续。

两条路线现在可以逐环节对照，表\ref{tab:vision-codec-compare}列出学习式压缩改变了什么。变换从固定的线性块变换变成可学习的非线性变换，消除了固定的块边界；量化步长从人工量化表变成由变换吸收的缩放；概率模型从静态码表变成随位置变化、由辅助信息驱动的分布，配合算术编码可以给出低于1比特的码长；各环节由同一个率失真目标联合训练，失真度量也可以按用途更换。换来的是更低的码率，付出的是解码端必须运行神经网络、两端浮点计算须逐位一致，以及每个码率点一套参数的存储代价。

\begin{table*}[tbp]
 \centering
 \small
 \begin{tabularx}{\linewidth}{@{}p{18mm}XXX@{}}
  \toprule
  环节 & JPEG基线 & JPEG 2000 & 尺度超先验\\
  \midrule
  变换 & $8\times8$离散余弦变换，固定线性 & 多级小波，固定线性 & 四层卷积与GDN，学习的非线性变换\\
  量化 & 人工量化表，质量因子缩放 & 每个子带一个步长 & 单位取整，步长由变换吸收\\
  概率模型 & 静态Huffman码表 & 按位平面上下文自适应 & 由辅助潜变量逐位置预测高斯尺度\\
  熵编码 & Huffman，每符号至少1比特 & MQ算术编码 & 二值算术编码\\
  设计方式 & 各环节分别设计 & 各环节分别设计 & 率失真目标联合训练\\
  码率调节 & 一套参数，改质量因子 & 一套参数，截断码块 & 每个$\lambda$一套参数\\
  解码要求 & 整数查表与逆变换 & 小波逆变换与算术解码 & 神经网络前向，浮点须逐位一致\\
  \bottomrule
 \end{tabularx}
 \caption{经典编码与学习式压缩的逐环节比较。三者共享“变换—量化—熵编码”的框架，差别在于每个环节由人设计还是由数据学习。}
 \label{tab:vision-codec-compare}
\end{table*}

\section{处理效果的核验}
\label{sec:vision-restoration-evaluation}

前三节的每种处理都要回答：结果是否值得采用。这个问题有三个方面，本节按维度分开讨论：结果与参考相比有多接近、看起来是否自然；结果能否支持读字、识别和测量等后续用途；为此付出了多少计算、时间和比特。三个方面不能互相替代，同一个处理在一个维度上的改进常伴随另一个维度上的代价。

比较前要先固定条件：输入退化、参考对齐、像素范围和颜色约定。对没有可信参考的真实图像，应明确哪些结论只能来自任务结果或人工观察，不能给不存在的原图计算“恢复正确率”。

\subsection{保真与感知质量的核验}
\label{sec:vision-restoration-quality}

这一维度的输入是一对对齐的图像：处理结果和参考。难点在于“接近”有不同含义：逐像素的数值差、局部结构的相似、人眼感觉到的差别，三者可以给出不同排序，尤其在一幅低清观测对应多个可能细节时。PSNR和SSIM用固定公式衡量前两种接近，LPIPS用学习得到的特征距离逼近人的判断；后者改变的是比较空间，从像素换成深层特征，并用人类判断校准权重。

\subsubsection{PSNR与SSIM}
\label{sec:vision-restoration-psnr-ssim}

PSNR与SSIM的定义已在式\eqref{eq:vision-psnr}和式\eqref{eq:vision-ssim}给出，这里讨论它们各自衡量了什么、怎样计算才可比较。PSNR是均方误差的对数变换，二者排序相同；误差为零时PSNR为正无穷。对数尺度只是让大范围的误差便于比较，MSE减半对应PSNR提高约3 dB。前述两像素压缩例的MSE为0.0025，取$v_{\max}=1$得到约26.02 dB。对一组图像平均各图PSNR，与先合并所有像素误差再计算PSNR不是同一聚合方式，报告时要写明。

SSIM的第一个因子比较窗口均值，第二个因子同时比较方差和协方差：两窗口的变化幅度相同且同步起伏时接近1。它因此对整体亮度偏移不如PSNR敏感，对结构被破坏更敏感。窗口大小、权重、边界和颜色处理都会改变数值，应固定后再比较。

取参考$(0,0,1,1)$，平滑结果$(0,0.25,0.75,1)$的MSE为0.03125；只把第三个位置改成0的结果$(0,0,0,1)$，MSE为0.25，按固定长度比较时也可视为边缘平移一格。再取$(0,0,0.8,1)$，MSE只有0.01，但若第三位置恰是识别一笔关键字画的依据，任务错误仍可能很大。像素分数描述这些改变的平均幅度，SSIM还考察局部共同变化，二者都不知道这笔画的业务含义。

配准误差会使正确内容与错误位置比较。评价应先确认对应关系，再统一是否裁剪边缘、只评亮度还是评RGB，以及线性值或显示编码值。例如Restormer报告去雨结果时在YCbCr的亮度通道上计算PSNR和SSIM；不同论文的数字只有在同一口径下才能比较。不能为每个方法单独选择最有利的对齐和颜色口径。压缩比较还要把PSNR与实际码率配对，按式\eqref{eq:vision-bd-rate}汇总整条曲线，单个码率点上的PSNR高低不足以说明哪个编解码器更好。

\subsubsection{LPIPS}
\label{sec:vision-restoration-lpips}

PSNR和SSIM都由人设计，它们与人眼判断的一致程度有限：轻微模糊和轻微平移可能得到相近的PSNR，人却觉得前者明显更差。深度网络的中间特征对纹理和结构敏感，作者据此提出问题：能否直接用网络特征的距离衡量感知差别，并用人的判断来校准。LPIPS就是这样一种度量\scite{vision-lpips2018}

设第$\ell$层在位置$p$的特征向量经通道归一化后为$\bar{\bm{f}}_{\ell,p}(\bm{X})$，该层共有$n_\ell$个位置，学习到的通道权重为$\bm{w}_\ell$。其距离写成
\begin{equation}
 d_{\mathrm{LPIPS}}(\bm{X},\widehat{\bm{X}})
 =
 \sum_\ell\frac1{n_\ell}\sum_p
 \left\lVert
 \bm{w}_\ell\odot
 \bigl(\bar{\bm{f}}_{\ell,p}(\bm{X})
      -\bar{\bm{f}}_{\ell,p}(\widehat{\bm{X}})\bigr)
 \right\rVert_2^2\text{。}
 \label{eq:vision-lpips}
\end{equation}
两幅图像都经过同一网络与预处理。先按位置比较归一化特征，再平均空间差异并累加层，不能将整图分类向量的余弦距离直接叫作LPIPS。若某位置的两通道差为$(0.2,0.1)$，权重为$(1,2)$，该位置贡献为$0.2^2+(2\times0.1)^2=0.08$。权重全取1时，归一化特征的平方距离等价于余弦距离，学习的作用是给不同通道不同的重要性；作者把权重限制为非负，因为两幅图之间的距离不应为负。

论文比较了三种骨干网络，各自取用的层见表\ref{tab:vision-lpips-layers}。三者都只取卷积层的输出，空间尺寸随层加深逐级减小，所以式\eqref{eq:vision-lpips}的求和同时覆盖了从边缘纹理到较大结构的多个尺度。

\begin{table}[htbp]
 \centering
 \small
 \begin{tabularx}{\linewidth}{@{}p{16mm}Xp{30mm}@{}}
  \toprule
  骨干网络 & 取用的层 & 各层通道数\\
  \midrule
  AlexNet & conv1至conv5 & 64，192，384，256，256\\
  VGG & conv1\_2，conv2\_2，conv3\_3，conv4\_3，conv5\_3 & 64，128，256，512，512\\
  SqueezeNet & conv1，fire3，fire5至fire9 & 64，128，256，384，384，512，512\\
  \bottomrule
 \end{tabularx}
 \caption{LPIPS三种骨干网络的取用层。层名与通道数依据原文附录中各层学习权重的展示；每层一个与通道数等长的非负权重向量$\bm{w}_\ell$。}
 \label{tab:vision-lpips-layers}
\end{table}

权重$\bm{w}_\ell$来自人的判断。作者构建了BAPPS数据集：从图像中取$64\times64$的参考块，施加两种不同的失真，请人判断哪一个更接近参考，即两选一强制选择；失真既包括噪声、模糊、压缩等传统操作，也包括各种网络输出，数据集共有超过16.1万个图块。训练时，对每个三元组分别算出两个失真块到参考的距离$d_0,d_1$，再用一个小网络$G$把这两个距离映射为“人选择第一个”的概率$\hat h$；$G$由两层32通道的全连接层加ReLU、一层单通道全连接层和sigmoid组成。损失为交叉熵$-h\log\hat h-(1-h)\log(1-\hat h)$，$h$是人群中实际选择第一个的比例。常用的版本固定预训练骨干的参数，只学习各层的线性权重$\bm{w}_\ell$；作者也比较了微调骨干和从头训练的版本。$G$只在训练中使用，使用时只需两幅对齐的图像，一次前向即可按式\eqref{eq:vision-lpips}算出距离。

平滑砖墙可能产生较小像素误差，却改变深层纹理特征；把招牌“8”改成“3”又可能只影响很小区域，整体感知距离仍然较小。LPIPS衡量对应图像的特征差异，不能独立核验文字、身份或事实，也不能把一项分数直接解释成“恢复了多少真实内容”。

像素保真与感知质量之间存在理论上的取舍，这要用比较分布的另一个定义来说明。设$p_{\bm{X}}$为真实图像分布，复原规则$p(\widehat{\bm{X}}\mid\bm{Y})$诱导输出分布$p_{\widehat{\bm{X}}}$，以$\Delta$度量逐对失真，以$d$度量两种分布的差异。给定失真预算$D$，定义
\begin{equation}
 P(D)=
 \inf_{\substack{p(\widehat{\bm{X}}\mid\bm{Y})\\
  \mathbb{E}[\Delta(\bm{X},\widehat{\bm{X}})]\leq D}}
 d(p_{\bm{X}},p_{\widehat{\bm{X}}})\text{。}
 \label{eq:vision-perception-distortion}
\end{equation}
放宽$D$扩大可选复原规则集合，因此$P(D)$不增；若$d$对第二个分布具有适当凸性，还可得到凸的最优边界。这个框架中的感知差异比较的是分布，不能直接替换成任意一个单图美观分数\scite{vision-perception-distortion2018}

一个完全没有信息的观测可以说明这种取舍。假设原像素以相同概率为0或1，平方误差最优输出为0.5，平均失真0.25，但输出分布只有一个中间值，与真实分布完全不同。若独立地以相同概率输出0或1，边缘分布与真实像素相同，平均平方误差却为0.5。后一规则没有因此猜对每个原像素。这正对应前文的模型差别：维纳滤波和DnCNN变体靠近前一种规则，Real-ESRGAN和DiffBIR向后一种规则移动。实际方法远离最优边界时，改进模型仍可能同时降低多项误差；上述取舍不表示所有评价指标必然互相恶化。

三种指标的差别归纳在表\ref{tab:vision-quality-metric-compare}中。报告复原结果时，至少应同时给出一项像素或结构指标和一项感知指标，并说明各自的计算口径。

\begin{table}[htbp]
 \centering
 \small
 \begin{tabularx}{\linewidth}{@{}p{14mm}XXX@{}}
  \toprule
  比较维度 & PSNR & SSIM & LPIPS\\
  \midrule
  比较对象 & 逐元素差 & 局部均值、方差与协方差 & 多层归一化深度特征\\
  参数来源 & 只有峰值约定 & 窗口与常数$c_1,c_2$ & 骨干网络与人类判断校准的权重\\
  方向 & 越大越接近 & 越接近1越接近 & 越小越接近\\
  敏感于 & 平移、整体亮度偏移 & 结构破坏、对比变化 & 纹理与结构的感知差别\\
  不能判断 & 误差落在哪里、是否重要 & 内容语义 & 文字、身份等事实是否正确\\
  \bottomrule
 \end{tabularx}
 \caption{三种参考指标的比较。三者都需要对齐的参考图像，都不能回答处理结果是否保留了用途所需的具体内容。}
 \label{tab:vision-quality-metric-compare}
\end{table}

\subsection{下游用途的核验}
\label{sec:vision-restoration-downstream}

参考指标衡量与参考图的接近程度，下游核验则直接问处理结果能否完成任务。招牌任务可以固定一套文字读取器，分别输入原照片、处理结果和可信参考，再与人工核实的字符串比较。若原图识别成“18路”、处理后成“13路”、参考为“18路”，锐化后的感知改善没有帮助这次读取。还应检查处理图像本身是否改出了“3”，或字符仍为“8”但读取器产生误判，以区分图像处理与下游系统的误差。

对于测量，可固定边缘定位与标定步骤，比较长度、角点和深度；对于识别，可比较相同对象的类别及拒识结果。阈值需要在验证集上确定，不能反复查看测试标签后为某个处理器调优。若联合训练了处理器与下游模型，还应把这看作一个新系统，与使用固定下游模型的对照分别报告。压缩也一样：同样的码率下，JPEG的块效应和学习式压缩的平滑可能对文字读取造成不同的影响，应以读取正确率对码率作曲线，而不只看PSNR对码率的曲线。

当真实参考不可得时，可以保留原始RAW、多次曝光或其他视角作为额外证据，并记录生成补全区域，例如部分卷积的原始掩码或DiffBIR的引导强度与随机种子。补出的皮肤、纹理和字符是否可接受，由用途决定；需要记录保真的任务应能追溯处理前后对应关系，不能以模型确信度替代丢失观测。

\subsection{资源成本的核验}
\label{sec:vision-restoration-resources}

质量相同的两种处理，成本可能相差几个数量级。复原比较应固定输入分辨率、退化、硬件、采样步数和候选数量。经典滤波器的成本最低：中值滤波逐像素排序，维纳滤波和约束最小二乘只需每通道一次快速傅里叶变换，后者选参数时再重复几次。前馈网络与扩散流程的成本不只相差参数量：DnCNN、Restormer和Real-ESRGAN一次前向即得结果，DiffBIR需要先运行第一阶段复原模块，再进行50步采样，每步都运行一次去噪网络和条件支路，启用复原引导时每步还要经过解码器求梯度。即使同为一次前向，注意力的形式也决定了分辨率增大时成本怎样增长，这正是Restormer改用通道注意力的原因；论文报告它的计算量比此前的MPRNet少81\%。DnCNN超分辨率变体先插值再在高分辨率网格上计算，Real-ESRGAN则在低分辨率上计算，放大倍数越大，二者的差距越明显。

图像压缩则要同时统计实际文件码率、编码时间、解码时间与内存，不能只比较训练时的估计码率。JPEG的编解码只需块变换、查表和整数运算，几乎所有设备都有硬件支持；JPEG 2000的小波与位平面算术编码明显更慢。尺度超先验的解码必须先解辅助流再解主流，两段不能并行，每次解码都要运行合成网络；每个$\lambda$对应一个模型，覆盖多个码率意味着收发双方都要保存多套参数。传统编解码器常用一套参数加不同量化表覆盖全部码率，这一点上更省存储。

表\ref{tab:vision-restoration-choice}把三类目的与所需证据对应起来。一个应用可以同时占据多行，例如先从RAW复原，再调整显示并压缩发送。最终判断仍需回到整个流程的输出，前段留下的错误可能被后段放大。

\begin{table}[htbp]
 \centering
 \small
 \begin{tabularx}{\linewidth}{@{}p{18mm}XX@{}}
  \toprule
  主要目的 & 需要保留的证据 & 共同记录的成本与条件\\
  \midrule
  内容复原 & 对齐参考、退化回投影、文字或测量结果 & 输入范围、模型版本、采样步数与候选预算\\
  呈现增强 & 原始图像、亮部裁剪、颜色与局部层次变化 & 颜色空间、曲线、增益和局部参数\\
  存储传输 & 独立解码结果、实际文件、用途误差、BD-rate & 辅助比特、元数据、编码及解码资源\\
  \bottomrule
 \end{tabularx}
 \caption{按处理目的选择证据。表中各行可以组合；每项判断都同时保留输入条件、处理结果与成本，避免只用一个质量分数作选择。}
 \label{tab:vision-restoration-choice}
\end{table}

\section{本章小结}
\label{sec:vision-restoration-summary}

回到开篇的夜间招牌。让它变得可用，需要区分三件事：看清、看着自然和记录正确。本章开头把这件事拆成五个挑战，各节依次给出了回答。

对于观测信息丢失带来的解不唯一，复原一节按退化条件说明了缺失信息从哪里补回。噪声下，经典方法由噪声模型推出滤波规则：均值滤波对高斯噪声，中值与自适应中值滤波对椒盐噪声，方差稳定化变换把泊松—高斯噪声化为近似单位方差；DnCNN把同样的噪声模型放进训练数据，用残差学习一次处理一个强度范围。已知模糊核时，逆滤波、维纳滤波和约束最小二乘的差别只在弱频率上相信什么，后两者在数学上是同一形式，先验分别来自功率谱和平滑约束。分辨率不足时，同一观测对应多个高清图像，平方误差训练只能给出平滑的条件均值。区域缺失时，部分卷积保证结果不依赖占位数值，但缺口内容仍全部来自周围像素和先验。

对于退化未知且混合带来的模型失配，盲反卷积先估计核再复原，受限于先验偏向无模糊解和空间不变假设；Restormer、DnCNN-B与Real-ESRGAN改为扩大训练中的退化分布，覆盖得越广，失配越少，但合成退化仍无法穷尽真实相机和传播链路。对于低照度与目标不唯一带来的增强取舍，暗光下的提亮要回到RAW数据才能同时处理噪声和颜色；内容充足时，全局曲线保持明暗次序，CLAHE换来局部层次却改变区域间关系，选用时要规定作用通道和处理顺序。对于细节补全带来的保真与感知冲突，Real-ESRGAN和DiffBIR借助对抗训练和预训练扩散模型补出细节，补出的越多，越需要区分观测支持的内容与先验生成的内容；感知—失真边界说明这种取舍有理论依据，不是某个模型的缺陷。对于比特预算带来的码率与失真权衡，JPEG与JPEG 2000把变换、量化、熵编码分开设计，学习式压缩保留同一框架，把三个环节放进同一个率失真目标联合训练：加均匀噪声的训练近似之所以成立，是因为加噪密度在整数点正好等于取整概率；尺度超先验用少量辅助比特换取更准的概率，前提是辅助流计入总码率，且两端算出完全相同的概率表。

本章没有解决的问题仍然明显。合成退化无法穷尽真实相机和传播链路，生成先验补出的细节缺少可核验的依据，低码率下关键文字可能在平均失真很小的情况下丢失。它们分别指向退化建模、观测约束和面向用途的编码。核验部分给出的做法是同时报告参考指标、感知指标、下游任务结果和资源成本，并说明结论依赖哪类输入。平均感知质量上升若伴随关键字段错误增加，就需要改变目标或应用范围。下一章允许根据用户要求创建和修改内容，届时仍需明确哪些变化是目标、哪些内容必须保持。

\paragraph{扩展阅读}
复原方向，BM3D把图中相似的图像块堆叠成三维组，在变换域中联合收缩后再加权放回原位；它与本章的经典滤波一样需要给定噪声水平，但利用的是图像内部重复出现的结构，长期作为DnCNN等学习方法的对照基线\scite{vision-dabov2007bm3d}FFDNet把噪声水平图与下采样得到的子图像一起送入网络，一个模型即可处理不同强度乃至空间上变化的噪声，可与本章DnCNN-S和DnCNN-B之间“专用模型更准、通用模型不必知道强度”的取舍对照阅读\scite{vision-ffdnet2018}SwinIR走的是本章讨论Restormer时提到的另一条路：在局部窗口内计算空间注意力，逐层平移窗口以在相邻窗口之间传递信息，并用同一结构处理超分辨率、去噪和JPEG去块\scite{vision-swinir2021}超分辨率方面，SRCNN是较早用卷积网络学习低分辨率到高分辨率映射的工作，它的三层卷积分别对应稀疏编码方法中的块提取、非线性映射和重建，有助于理解学习式复原怎样从经典流程演变而来\scite{vision-srcnn2016}SR3为超分辨率单独训练一个条件扩散模型，以低分辨率图像为条件从噪声逐步去噪，不需要ESRGAN式的对抗训练；它与DiffBIR借用预训练文生图扩散模型的做法不同，适合比较生成先验来自何处\scite{vision-sr3-2023}

增强方向，EnlightenGAN只用互不配对的低光图像和正常光照图像训练，以全局与局部判别器约束结果的曝光外观，再用自特征保持损失约束内容不变；SID需要成对的短、长曝光RAW数据，二者对照可以看出监督来源怎样决定增强目标\scite{vision-enlightengan2021}Eilertsen等人从单张曝光的低动态范围图像预测饱和高光区内的亮度，生成高动态范围图像；它处理的是本章亮度增强之外的另一类信息丢失，即被传感器截断的高光无法靠曲线拉伸找回，只能由学到的先验补出\scite{vision-hdrcnn2017}

压缩方向，Minnen等人在尺度超先验之外加入自回归上下文模型，用已解码的相邻潜变量和辅助信息共同预测每个位置的均值与尺度，进一步降低码率，代价是解码必须逐位置串行进行\scite{vision-minnen2018}ELIC把潜变量按通道不均匀地分组，在空间上采用棋盘式的两遍上下文，在保留大部分上下文收益的同时恢复了并行解码，可用来比较概率模型精度与解码速度之间的取舍\scite{vision-elic2022}JPEG XL是JPEG委员会制定的新一代静止图像编码，在人工设计的框架中采用可变大小的块变换和模块化的无损编码，并能把已有JPEG文件无损地重新压缩，可以看作本章经典编码流程的现代延续\scite{vision-jpegxl2019}理论方面，Blau与Michaeli把感知—失真边界推广到有码率约束的情形，提出率—失真—感知函数，说明提高解码结果的感知质量要以更高码率或更大失真为代价，这为低码率下的生成式压缩提供了分析框架\scite{vision-rdp-tradeoff2019}

\chapter{图像生成与编辑}
\label{chap:vision-image-generation}

设想一位设计者手里有一只红色玩具狗的几张照片，提出四个请求：“画这只玩具狗坐在蓝色杯子左边，在雪地里”；“把这张照片转成水彩画”；“把照片里狗旁边的杯子换成花瓶”；“根据绕着它拍的一圈照片，给出从正后方看的样子”。四个请求的输出都是一幅新图像，但它们给定的东西不同，允许变化的东西也不同。第一个请求只规定内容和主体，构图、光照和背景细节都可以由生成模型补充；第二个请求有一幅输入，狗的姿态与数量要保持，笔触和色调可以整体改变；第三个请求同样有输入，却只允许杯子所在的区域改变；第四个请求的答案受真实拍摄约束，背面若没有拍到，模型补出的内容就不能算作观测结果。

本章要回答的问题是：给定这些不同的条件，怎样得到满足要求的图像，又怎样核验它确实满足了要求。生成模型的训练与采样原理已在前面各卷建立：扩散模型通过学习逐步去噪，从高斯噪声出发采样出图像，可参考\BookChapterRef{chap:pgm-generative}{卷六的“概率生成模型”章}；生成对抗网络以判别器的反馈训练生成器，条件生成对抗网络和配对、无配对的图像转换见\BookChapterRef{chap:neural-generative-models}{卷五的“基于神经网络的生成模型”章}。本卷\ref{chap:vision-multimodal-models}章又讲了LDM、unCLIP、GLIDE与流匹配等图像生成模型的结构和训练。本章不再重讲这些模型，而是研究它们在任务中怎样使用，并完整讲解只服务于某一任务的模型，例如空间控制用的ControlNet、主体适配用的DreamBooth和新视角合成用的NeRF。

四个请求依次放宽“可以变化”的范围，也依次加强“必须保持”的依据，这决定了本章的路线。我们先把四类任务的输入、输出和评价指标写清楚，并从这四个请求中归纳出贯穿全章的四个挑战：一个条件对应许多合格图像，多个对象的属性容易互相串换，局部修改与内容保持彼此牵制，有限的照片难以保证任意视角下的一致。新图像的创建只受条件约束，难点是把文字、空间图和主体参考送进生成过程，并从多个候选中选出满足要求的一幅；已有图像的转换多了一幅输入，难点变成怎样学到保持对应关系的映射；指定内容的编辑进一步要求区域外不变，这里先讲泊松图像编辑、基于样例的修补和PatchMatch等经典做法，它们只从已有图像中搬运内容，再讲按文字生成新内容的学习方法，还要在多轮请求之间维护已接受的状态；新视角合成则不再依靠生成先验，而是在多视图几何的基础上，从多幅照片拟合一个可以从任意相机渲染的场景。最后一节把这四类任务的核验放在一起，按条件遵循、输入保持、观测支持与成本三个维度比较。贯穿全章的玩具狗、杯子及各类分数和坐标都是教学构造，不代表任何模型的实测结果。

\section{图像生成与编辑的问题与挑战}
\label{sec:vision-generation-challenges}

开篇四个请求的输出都是图像，输入却各不相同。新图像的创建以文字、空间图或主体照片为条件，输出一幅或几幅满足条件的图像；已有图像的转换以一幅图像为输入，输出同一内容的另一种呈现形式；指定内容的编辑以原图、编辑请求和可能附带的范围标注为输入，输出只在请求涉及的部分发生变化的图像；新视角合成以多幅带相机参数的照片为输入，输出指定新相机下的图像。前三类任务的合格答案通常不止一个，新视角合成原则上只有一个答案，即在那个相机位置真实拍到的照片。这一差别决定了评价方式：前三类任务只能检查结果是否落在可接受的范围里，新视角合成可以逐像素与真实照片比较。

评价指标按所核验的对象分为四组。第一组比较生成图像整体的分布与真实图像的分布。\term{Fréchet Inception距离}（Fréchet Inception distance, FID）把一组真实图像和一组生成图像分别送入在ImageNet上训练的Inception-v3网络，取最后一个池化层的2048维特征，为两组特征各拟合一个高斯分布，再计算两个高斯之间的Fréchet距离\scite{vision-fid2017}：
\begin{equation}
 \mathrm{FID}
 =\lVert\bm{\mu}_{\mathrm{r}}-\bm{\mu}_{\mathrm{g}}\rVert_2^2
 +\operatorname{tr}\!\Bigl(
  \bm{\Sigma}_{\mathrm{r}}+\bm{\Sigma}_{\mathrm{g}}
  -2\bigl(\bm{\Sigma}_{\mathrm{r}}\bm{\Sigma}_{\mathrm{g}}\bigr)^{1/2}
 \Bigr)\text{，}
 \label{eq:vision-fid}
\end{equation}
其中$\bm{\mu}_{\mathrm{r}},\bm{\Sigma}_{\mathrm{r}}$和$\bm{\mu}_{\mathrm{g}},\bm{\Sigma}_{\mathrm{g}}$分别是真实特征与生成特征的均值和协方差，$(\cdot)^{1/2}$表示矩阵平方根。第一项比较两组特征的中心，第二项比较它们的分散程度和各维之间的相关；FID越小，两组特征的分布越接近。一维时公式化简为$(\mu_{\mathrm{r}}-\mu_{\mathrm{g}})^2+(\sigma_{\mathrm{r}}-\sigma_{\mathrm{g}})^2$。若真实特征的均值为0、标准差为1，生成特征的均值为0.5、标准差为2，则FID为$0.25+1=1.25$；只把生成特征的标准差调回1，FID降为0.25，说明这个指标会同时惩罚位置偏移和多样性的偏差。这里的Fréchet指标是两个高斯之间2阶Wasserstein距离的平方；一般协方差下的闭式及其一维化简可参考\BookChapterRef{chap:information-theory-and-statistical-geometry}{卷2《数学准备：概率、因果与决策》的“概率分布的分析与优化”章}。FID只比较两个分布，不检查任何一幅图是否满足它自己的条件。它还依赖Inception特征和样本数：原论文用5万幅生成图像估计均值和协方差，样本少时估计误差大，所以只有在相同样本数、相同特征提取实现下得到的FID才能互相比较。

第二组衡量结果与文字条件的一致程度。\term{CLIP分数}（CLIPScore）用CLIP ViT-B/32分别编码文字$c$和图像$\bm{X}$，取两个向量的余弦相似度\scite{vision-clipscore2021}：
\begin{equation}
 \mathrm{CLIPScore}(c,\bm{X})
 =w\max\bigl(\cos(\bm{e}_{\mathrm{T}}(c),\bm{e}_{\mathrm{I}}(\bm{X})),0\bigr)\text{，}
 \label{eq:vision-clipscore}
\end{equation}
其中$\bm{e}_{\mathrm{T}},\bm{e}_{\mathrm{I}}$是CLIP的文字与图像编码器，原论文取$w=2.5$，只用于拉伸取值范围，并在文字前加上“A photo depicts”。余弦为0.30时分数为0.75，余弦为负时记为0。原论文提出它是为了评价图像描述，生成任务借用它时比较的是同一对象：图文是否相关。它是一个整图对整句的相似度，对数量、左右关系和颜色归属并不敏感，\ref{sec:vision-generation-reranking}节会给出按它排序选错候选的例子。主体保真也用相同的余弦形式：DreamBooth把生成图像与真实主体照片的DINO特征余弦相似度在所有配对上取平均，称为DINO分数，用CLIP图像特征时称为CLIP-I，用图文特征时称为CLIP-T\scite{vision-dreambooth2023}。

第三组依靠人的判断。人工偏好评价把两个系统在同一条件下的结果并排给评审，请他们选出更好的一幅或判为相当。若$n$次比较中系统A胜$n_{\mathrm{w}}$次、平$n_{\mathrm{t}}$次，则A的胜率记为$(n_{\mathrm{w}}+n_{\mathrm{t}}/2)/n$。例如100次比较中胜55次、平10次，胜率为0.6。评审被问到的问题决定了这个数字的含义：“哪幅更真实”和“哪幅更符合文字”是两种不同的偏好，GLIDE就分别就照片真实感和与文字的相符程度作了两两比较\scite{vision-glide2022}。报告人工偏好时，应写明问题措辞、比较次数、评审来源和候选的选取方式。

第四组用于有输入或有真实答案的任务。编辑任务需要分别衡量“该改的是否改了”和“该留的是否留了”。对于明确的保留掩码$\bm{M}\in\{0,1\}^{h\times w}$，$m_{i,j}=1$表示保留区，可以计算保留区的平均变化
\begin{equation}
 E_{\mathrm{keep}}=
 \frac{\sum_{i,j,k}m_{i,j}
       |x'_{i,j,k}-x_{i,j,k}|}
 {c\sum_{i,j}m_{i,j}}\text{，}
 \label{eq:vision-edit-preservation}
\end{equation}
其中$\bm{X}$与$\bm{X}'$是原图与结果，$c$是通道数，分母要求至少有一个保留像素。原样返回原图时$E_{\mathrm{keep}}=0$，但请求根本没有执行，所以这一项必须与修改完成度一起报告。新视角合成可以在留出的测试相机下与真实照片$\bm{Y}$比较。峰值信噪比$\mathrm{PSNR}=10\log_{10}(V_{\max}^2/\mathrm{MSE})$以均方误差为基础，$V_{\max}$是像素取值上限；像素在$[0,1]$内、均方误差为0.001时，PSNR为30\,dB。结构相似性（SSIM）比较局部窗口内的均值、方差和协方差，LPIPS比较深度特征之间的距离，二者的定义与适用条件见\ref{sec:vision-restoration-psnr-ssim}节和\ref{sec:vision-restoration-lpips}节。

这些指标各自只回答一个问题。下面四个小节讨论的困难，正是让这些问题难以同时得到肯定回答的原因。

\subsection{一对多映射带来的条件遵循困难}
\label{sec:vision-generation-challenge-one-to-many}

“雪地里的红色玩具狗”对应无数幅合格的图像：狗可以坐着或站着，天空可以晴或阴。条件与图像之间不是函数关系，一个条件$c$对应一个分布$p(\bm{X}\mid c)$。若像回归问题那样训练一个确定性网络$\widehat{\bm{X}}=G(c)$并最小化平方误差，最优输出是条件均值$\mathbb{E}[\bm{X}\mid c]$。设某个像素在合格图像中一半取0（阴影里）、一半取1（雪面上），平方误差$\frac12(y-0)^2+\frac12(y-1)^2$对$y$求导得$2y-1$，在$y=0.5$处最小，网络会输出一个两种情形都不是的灰色。许多像素同时取平均，就得到模糊、缺少细节的图像。所以生成模型必须表示整个分布并从中采样；而一旦采样，每幅输出都带有随机性，有的满足条件、有的不满足，条件遵循就成了一个概率意义上的要求。

经典方法不去学习一个分布，而是直接从样例中取值。Efros与Leung的纹理合成在样例图中找出所有与当前邻域相似的位置，把它们中心像素的经验分布当作条件分布，再从中随机抽取一个值\scite{vision-efros-leung1999}。因为每次取的都是样例中真实出现过的值，输出保留了样例的锐利纹理，没有平均化。纹理本身可以用局部统计刻画，经典描述见\ref{sec:vision-features-texture}节。不过这种非参数采样只能复制样例中已有的统计规律，读不懂“红色”“雪地”这样的语义条件，它在编辑中的用法见\ref{sec:vision-generation-exemplar-inpainting}节。

学习方法用生成模型表示$p(\bm{X}\mid c)$：生成对抗网络让判别器惩罚不像真实样本的输出，扩散模型学习逐步去噪的条件分布。平均化由此得到解决，问题转成两件事：怎样让采样结果更集中地落在满足条件的区域，以及一次采样不满足时怎么办。\ref{sec:vision-generation-creation}节从条件接口、采样中的引导与采样后的重排、生成后的修订三方面回答；\ref{sec:vision-generation-pix2pix}节的pix2pix则在有配对数据的转换中，用局部判别器抵消$\ell_1$像素误差带来的平均化。

\subsection{组合语义带来的属性绑定错误}
\label{sec:vision-generation-challenge-binding}

“红色玩具狗在蓝色杯子左边”由两个对象、两种颜色和一条空间关系组成。常见的错误是画出蓝色的狗和红色的杯子：每个词都在画面中出现了，但颜色附着到了错误的对象上。\term{属性绑定}（attribute binding）指把每个属性正确地附着到它所修饰的对象上，这里的困难来自组合：对象和属性各自都不难画，难的是按句子规定的方式把它们配对，数量和空间关系也属于同一类组合要求。

以交叉注意力读取文字的生成网络中，每个空间位置都可以按学到的权重读取所有词元；结构上并没有一个机制规定“读取了‘狗’的那些位置也必须读取‘红色’而不是‘蓝色’”，这种配对只能从训练数据中学到。训练数据中少见的组合，例如蓝色的狗，学得往往更不充分。

经典图像合成用图层显式地完成绑定。设计者从一张照片里抠出红色玩具狗，从另一张照片里抠出蓝色杯子，分别放到画布上指定的位置；每个图层的颜色来自它自己的源照片，不会串换，位置也由人直接摆放。这种做法的问题在于拼接处：两张照片的光照和色调不同，直接粘贴会留下明显的边界，\ref{sec:vision-generation-poisson}节的泊松图像编辑就是为消除这种边界而设计的。它的代价是需要现成的源照片和人工摆放，生成不了任何照片里都没有的组合。

学习方法目前没有一个能保证绑定正确的结构，能做的是把绑定要求外显出来，并在生成后逐项检查。\ref{sec:vision-generation-conditions}节用空间图规定布局，用参考角色记录规定每幅参考负责哪个对象；\ref{sec:vision-generation-selection}节先按要求记录筛选候选，再排序；\ref{sec:vision-generation-unireason-refinement}节针对具体错误作一轮修订；\ref{sec:vision-generation-condition-check}节把绑定正确率作为单独的分项报告。

\subsection{局部修改带来的内容保持冲突}
\label{sec:vision-generation-challenge-preservation}

把狗旁边的杯子换成花瓶时，花瓶要立在桌面上，受光方向与周围一致，它的影子会落到杯子原来轮廓之外；杯子原来的影子也应随之消失。若严格要求杯子轮廓以外的像素一个都不变，新旧内容在边界处就会接不上，旧影子也留在原处；若允许外面变化，又可能顺带改动了不该改的狗。保持与修改之间的这种牵制，在转换中以另一种形式出现：照片转成水彩画时每个像素的颜色都要改变，需要保持的是姿态、数量和布局这些结构，而不是像素值。

经典方法的回答是只从已有图像中搬运内容。泊松图像编辑保持区域外的目标图像完全不变，区域内只保留源图像的梯度，并要求结果在区域边界上与周围的值相等，从而消除拼接缝（\ref{sec:vision-generation-poisson}节）。移除杯子之后的空洞，可以用同一幅图中别处的桌面图像块逐块填补，Criminisi等人的方法决定先填哪里（\ref{sec:vision-generation-exemplar-inpainting}节），PatchMatch则让“找到最相似的图像块”这一步快到可以交互使用（\ref{sec:vision-generation-patchmatch}节）。这些方法的局限也来自同一点：新内容只能来自现有图像，粘贴的物体在泊松编辑后整体颜色会被周围带偏，透视和光照方向也不会自动调整，桌面上本来没有的东西更是补不出来。

学习方法能按文字生成原图中不存在的内容。带掩码的修补模型学会根据可见区域生成与之协调的新内容，指令编辑模型不需要掩码，由模型自己判断改哪里，但未提及的区域也可能变化；多轮编辑还会让主体逐轮漂移。转换中，循环一致性和感知损失分别用往返重建和固定网络的特征来保持结构。\ref{sec:vision-generation-translation}节和\ref{sec:vision-generation-editing}节回答这一挑战，\ref{sec:vision-generation-change-check}节讨论怎样分别核验修改与保持。

\subsection{有限观测带来的三维一致性困难}
\label{sec:vision-generation-challenge-3d}

绕玩具狗在两个高度上各拍20张照片，要求给出从正后方偏上看的图像。新图像中的每个像素对应空间中的一条射线，它的颜色必须与所有拍到同一表面点的照片相符；相机绕狗转动时，耳朵应留在同一个三维位置，而不是每换一个视角就重新画一次。二维生成模型可以画出一幅看起来合理的背面，却没有任何机制保证换一个角度后它与前一幅、与真实照片一致。狗的底面从未被拍到，这一部分无论用什么方法都没有观测依据。

经典路线以多视图几何为基础：运动恢复结构从照片中同时估计各相机的位姿和一组稀疏的三维点，稠密重建再得到表面（\ref{sec:vision-formation-multi-view}节）；也可以不显式重建几何，而把照片看成对光线的采样，新视角直接由已采样的光线插值得到，这就是基于图像的渲染。前一条路线在毛绒、半透明和细小结构上难以得到准确的表面，后一条路线需要非常密集的拍摄。

学习方法为每个场景拟合一个可以从任意相机渲染的三维表示，并通过可微渲染让它在所有训练视图下与照片一致。一致性来自所有视图共享同一个三维表示，而不是逐幅图像各自生成。NeRF和3D Gaussian Splatting是两种代表，见\ref{sec:vision-generation-novel-view}节；没有拍到的表面仍然没有观测支持，\ref{sec:vision-generation-costs}节讨论怎样在交付时标明这一点，需要靠生成先验补全的情形留给\ref{chap:vision-multimodal-generation}章。

四个挑战并不一一对应四类任务，而是按约束由弱到强逐步叠加。创建只受条件约束，一对多和属性绑定是它的主要困难；转换和编辑在此之上多了一幅输入，保持与修改的冲突随之出现，转换在整幅图上保持结构，编辑只在局部修改；新视角合成要求结果由真实观测支持，三维一致性取代了“看起来合理”成为首要标准。所以本章按创建、转换、编辑、新视角合成的次序展开，每一类任务都继承前一类的困难；最后一节再按条件遵循、输入保持、观测支持与成本三个维度汇总核验方法，它们分别对应前两个挑战、第三个挑战和第四个挑战。

\section{新图像的创建}
\label{sec:vision-generation-creation}

创建新图像时，输入是一组条件：一段文字，可能还有姿态图、深度图或几张主体照片；输出是一幅或几幅满足条件的图像。难点有三处。第一，文字很难精确规定位置和形状，几张照片又难以通过文字表达，因此条件要按其约束的对象选用不同接口送入模型。第二，扩散模型的采样是随机的，同一组条件会给出质量和条件满足程度不同的候选，需要决定在采样中修正还是在采样后挑选。第三，生成后才能看到的错误，例如数量多了一只狗，需要一个“检查—修订”的环节。这三处难点对应创建流程中依次进行的步骤：先把请求整理成可检查的要求，再组织条件，然后生成并选择候选，最后在必要时作一次有记录的修订。这些步骤是任务流程中的职责划分，不要求模型内部恰好有四个同名模块。

\subsection{创作要求的确定}
\label{sec:vision-generation-requirements}

生成模型接受的是一段提示，验收者却需要一份可以逐项打勾的清单。考虑请求“红色玩具狗放在蓝色杯子左边”。它包含两个对象、各自的颜色、每种对象的数量以及一条左右关系。把它写成记录，就是“玩具狗，一只，红色；杯子，一只，蓝色；狗在杯子左边”。这份记录供后续生成选择和结果核验使用，并不假定生成器内部也存有同样的对象表。

把请求拆开以后，才能看清几层不同的要求。对象出现和属性绑定是两层：画面里同时有红色和蓝色，不说明红色属于狗、蓝色属于杯子；画面里有狗也有杯子，不说明数量和左右关系正确。若还要求杯子上写“OPEN”，字符应作为单独一项检查，因为整体图文相似度对几个字母的拼写并不敏感。

输出用途还决定哪些部分可以变化。插图允许材质、光照和背景有多种合理选择；产品展示要求主体形状固定；排版素材还要规定留白和裁剪范围。条件没有指定的部分可以自由变化，已经指定的部分不能因为“画面自然”而放宽。一个请求通常有许多可接受结果，因此要求记录应写出必须满足的项和允许变化的项，而不是把某一幅参考图的全部像素当作唯一答案。

\subsection{创作条件的组织}
\label{sec:vision-generation-conditions}

有了要求记录，下一步是把它变成模型能读取的条件。玩具狗例子里的要求分属三类信息：“雪地、杯子、红色”这类语义适合用文字表达；“狗抬起左前爪、坐在画面左侧”这类几何关系，用文字描述既冗长又不精确；“就是照片里这只玩具狗”这类身份信息，文字几乎无法表达，只能从参考照片中取得。条件组织要解决的问题是：每类信息用什么形式编码，在生成过程的哪一步进入，以及几类条件冲突时怎样判断结果。

按条件的形式，本节分四种情形。文字条件下的生成直接使用\ref{chap:vision-multimodal-models}章的LDM和unCLIP，条件经文字编码器进入去噪网络，不需要额外训练；空间图条件下的生成需要把一幅与输出对齐的图接入预训练模型，ControlNet用一条可训练的复制支路完成这件事，原模型参数不动；单一主体参考下的生成需要让模型记住一个具体主体，DreamBooth把主体信息写进模型参数本身；多参考图像下的生成则要再回答每幅参考负责什么。从文字到多参考，条件携带的信息越来越具体，代价也从“只改提示”增加到“训练支路”“微调整个模型”和“需要显式规定参考角色”。

\subsubsection{文字条件下的生成}
\label{sec:vision-generation-latent-entry}

只有文字时，最直接的工具是\ref{sec:vision-models-ldm}节的LDM。它先用自编码器把图像压缩成空间潜网格，在潜空间中训练以文字为条件的去噪网络，文字经编码器形成词元序列，再通过交叉注意力被每个空间位置读取；采样时从高斯噪声出发反复去噪，最后由解码器还原像素\scite{pgm-generative-rombach2022}。\ref{sec:vision-models-unclip}节的unCLIP则先由先验模型把文字映射为CLIP图像向量，再由解码器根据这个向量生成图像；给定一幅参考图时，也可以直接编码它的CLIP向量，跳过先验，生成语义相近的变体\scite{vision-unclip2022}。两者的模型结构和训练在该章已讲，这里只讨论使用时的输入组织、超参数和失败情形。

使用LDM时，一次生成由以下量共同决定：提示文本、目标分辨率、随机种子、采样器与步数、引导强度。分辨率先换算成潜网格大小。若自编码器在高宽方向各压缩8倍，$512\times512$的目标对应$64\times64$的潜网格，通道数由所选自编码器规定，例如4。种子决定初始高斯噪声$\bm{z}^{(K)}\in\mathbb{R}^{64\times64\times4}$，采样器按噪声日程调用去噪网络$K$次，每次读取文字特征，把当前状态更新为下一状态；最后解码得到$\bm{X}\in\mathbb{R}^{512\times512\times3}$，再按训练时的像素归一化反变换到显示范围。潜数组不是可交付的RGB图片；改变长宽比也必须相应改变潜网格，只在最后拉伸图像会让采样条件与报告不符。

unCLIP的参考图路径与LDM的空间潜网格不同。CLIP图像向量是一个整图汇总向量，参考图经它进入解码器时，被保留的是“这是一只红色毛绒玩具狗”一类语义，耳朵的精确形状、身上的标志和原背景并不保证保留。所以unCLIP适合产生与参考风格或语义相近的多个变体，不适合要求“同一只玩具”的身份保持。使用时需要同时保存先验、解码器、上采样器及其预处理版本，只保留CLIP编码器并不能重现生成路径。

文字条件的常见失败有三类。第一，属性绑定错误，例如狗被画成蓝色、杯子被画成红色，原因是交叉注意力并没有为“红色”指定一个对象框，词语与区域的对应完全靠训练数据学到。第二，数量与空间关系不准，例如“两个杯子”画成三个。第三，小字和细小结构在潜空间压缩中丢失。这些失败不能靠更长的提示保证消除，后文的空间条件、候选选择和修订分别针对其中一部分。实际比较时应固定模型、提示、分辨率和候选数，例如对同一描述用四个种子生成，观察姿态、背景和细节的变化；相同种子只有在模型、采样器和实现都一致时才可复现。

\subsubsection{空间图条件下的生成}
\label{sec:vision-generation-controlnet}

文字“狗抬起左前爪，坐在画面左三分之一处”很难让LDM把爪子放到准确位置。若手头有一幅姿态关键点图、边缘图或深度图，它与输出图像逐像素对齐，直接规定了布局。问题是，已有的文字生成模型只有文字条件接口，没有读取空间图的入口。从头训练一个同时读取文字和空间图的模型代价很高；直接在已有模型上微调全部参数，又可能在只有几万对条件图的数据上破坏原模型学到的生成能力。ControlNet解决的就是“在不损坏预训练模型的前提下，为它增加空间条件接口”\scite{vision-controlnet2023}。

ControlNet的主要想法是：冻结原模型，把原模型编码部分复制一份作为可训练支路，由支路读取空间条件，再通过初始化为零的卷积把支路输出加回原模型。复制保证支路从一开始就具备原模型的特征提取能力，而不是从随机权重学起；零初始化保证训练开始时支路的输出为零，原模型的行为与微调前完全相同，训练只会逐渐加入控制信号。\term{零卷积}（zero convolution）指权重和偏置都初始化为零的$1\times1$卷积，它是支路与主网络之间的唯一连接。

先看一个模块上的计算。设预训练模块为$F_{\bm{\theta}}$，复制模块为$F_{\bm{\theta}_{\mathrm{c}}}$，输入特征为$\bm{h}$，条件特征为$\bm{c}$，两个零卷积为$Z_1,Z_2$，则模块输出为
\begin{equation}
 \bm{h}'=F_{\bm{\theta}}(\bm{h})+
 Z_2\!\left(F_{\bm{\theta}_{\mathrm{c}}}
       \bigl(\bm{h}+Z_1(\bm{c})\bigr)\right)\text{。}
 \label{eq:vision-controlnet}
\end{equation}
$Z_1$把条件特征注入复制模块的输入，$Z_2$把复制模块的输出加回主路径。训练开始时两者输出都为零，$\bm{h}'=F_{\bm{\theta}}(\bm{h})$。零输出并不意味着支路没有梯度：$Z_2$的权重梯度等于上游误差与它的输入$F_{\bm{\theta}_{\mathrm{c}}}(\cdot)$的乘积，后者不为零，所以$Z_2$第一次更新后就变成非零，梯度随后才能继续传到复制模块和$Z_1$。图\ref{fig:vision-controlnet}标出了两条路径的关系。

\begin{figure}[htbp]
 \centering
 % Single-module sketch of ControlNet Eq.(2); no complete U-Net is implied.
\begin{tikzpicture}[x=1mm,y=1mm,font=\figurefont\small,
 b/.style={draw=BookBlue,fill=BookBlue!10,line width=.8pt,
  minimum width=16mm,minimum height=9mm,inner sep=1mm,align=center},
 op/.style={circle,draw=BookInk,fill=white,line width=.8pt,
  minimum size=5mm,inner sep=0},
 a/.style={-{Stealth[length=2mm]},BookInk,line width=.9pt,
  rounded corners=1mm}]
 \node at (5,0) (h) {$\bm{h}$};
 \node[b,draw=BookMuted,fill=BookPaper,minimum width=25mm] at (42,0)
  (base) {$F_{\bm{\theta}}$};
 \node[op] at (82,0) (sum) {$+$};
 \node at (101,0) (out) {$\bm{h}'$};
 \node[op] at (30,-25) (cond) {$+$};
 \node[b,minimum width=23mm] at (55,-25) (copy) {$F_{\bm{\theta}_{\mathrm{c}}}$};
 \node[b,minimum width=12mm] at (82,-25) (z2) {$Z_2$};
 \node at (5,-45) (c) {$\bm{c}$};
 \node[b,minimum width=12mm] at (30,-45) (z1) {$Z_1$};
 \draw[a] (h)--(base);
 \draw[a] (base)--(sum);
 \draw[a] (sum)--(out);
 \draw[BookInk,line width=.9pt] (14,0)--(14,-25);
 \draw[a] (14,-25)--(cond);
 \fill[BookInk] (14,0) circle (.6mm);
 \draw[a] (cond)--(copy);
 \draw[a] (copy)--(z2);
 \draw[a] (z2)--(sum);
 \draw[a] (c)--(z1);
 \draw[a] (z1)--(cond);
 \node[anchor=south,text=BookMuted] at (42,7) {冻结};
 \node[anchor=north,text=BookBlue] at (55,-32) {可训练副本};
 \node[anchor=north,text=BookBlue] at (82,-32) {从零连接};
\end{tikzpicture}

 \caption{ControlNet的模块级教学示意。主模块参数$\bm{\theta}$冻结，复制模块及零初始化连接可训练；相同输入进入两侧，条件只通过支路注入。零初始化让控制输出从零开始，训练后经加法调节主路径。完整网络在多个尺度接入控制特征，图中仅展开一处。}
 \label{fig:vision-controlnet}
\end{figure}

零卷积启动的顺序可以手算。把外层连接简化为标量$a$乘以支路输出$q$，主路径输出为$f$，目标为$y$，损失为$(f+aq-y)^2/2$。取$a=0,q=2,f=0.2,y=0.6$，则$\partial L/\partial a=(f+aq-y)q=-0.8$，而$\partial L/\partial q=(f+aq-y)a=0$。以学习率0.1更新一次得到$a=0.08$；下一步$\partial L/\partial q$就不再为零。这个教学例只说明梯度怎样逐层启动，不表示完整网络的收敛步数。

把这个模块放回整个网络。原论文以Stable Diffusion为基础，它的去噪U-Net在$64\times64\times4$的潜网格上工作，由编码部分、中间块和带跳连的解码部分组成；编码与解码部分各有12个块，分布在$64,32,16,8$四种分辨率上，每种分辨率3个块。ControlNet复制这12个编码块和1个中间块，复制支路的13个输出各经一个零卷积，分别加到U-Net解码部分的12条跳连和中间块的输出上；解码部分本身不复制，保持冻结。论文没有列出各块的通道数，表\ref{tab:vision-controlnet-structure}的通道与层内组成取自Stable Diffusion v1公开的配置文件：基础通道320，四个分辨率的通道倍率为$1,2,4,4$，前三个分辨率的块含空间Transformer，文字由冻结的CLIP文字编码器编码为77个768维词元。

\begin{table*}[tbp]
 \centering
 \small
 \begin{tabularx}{\linewidth}{@{}p{30mm}Xp{24mm}p{22mm}p{12mm}p{20mm}@{}}
  \toprule
  阶段或层组 & 层类型 & 卷积核与步长 & 输出通道 & 重复次数 & 输出尺寸\\
  \midrule
  条件编码器（论文） & 卷积，ReLU & $4\times4$，步长2 & 16、32、64、128 & 4 & 论文称$64\times64$\\
  条件编码器（官方代码） & 卷积，SiLU；末层为零卷积 & $3\times3$，其中3层步长2 & 16、16、32、32、96、96、256，末层320 & 8 & $64\times64$\\
  输入卷积 & 卷积 & $3\times3$，步长1 & 320 & 1 & $64\times64$\\
  编码分辨率1 & 残差块，空间Transformer & $3\times3$，步长1 & 320 & 2 & $64\times64$\\
  下采样 & 卷积 & $3\times3$，步长2 & 320 & 1 & $32\times32$\\
  编码分辨率2 & 残差块，空间Transformer & $3\times3$，步长1 & 640 & 2 & $32\times32$\\
  下采样 & 卷积 & $3\times3$，步长2 & 640 & 1 & $16\times16$\\
  编码分辨率3 & 残差块，空间Transformer & $3\times3$，步长1 & 1280 & 2 & $16\times16$\\
  下采样 & 卷积 & $3\times3$，步长2 & 1280 & 1 & $8\times8$\\
  编码分辨率4 & 残差块 & $3\times3$，步长1 & 1280 & 2 & $8\times8$\\
  中间块 & 残差块，空间Transformer，残差块 & $3\times3$，步长1 & 1280 & 1 & $8\times8$\\
  零卷积（仅ControlNet） & 卷积 & $1\times1$，步长1 & 与所接的块相同 & 13 & 与所接的块相同\\
  解码部分（冻结，不复制） & 拼接跳连后的残差块（除$8\times8$外含空间Transformer），分辨率之间上采样 & $3\times3$，步长1 & 1280、1280、640、320 & 每个分辨率3个块 & $8\to64$\\
  输出层 & 组归一化，SiLU，卷积 & $3\times3$，步长1 & 4 & 1 & $64\times64$\\
  \bottomrule
 \end{tabularx}
 \caption{ControlNet所依附的Stable Diffusion v1去噪U-Net与新增部分。编码部分从“输入卷积”到“编码分辨率4”共12个块，ControlNet把它们和中间块各复制一份；“零卷积”一行是新增的13个连接。条件编码器给出论文文字与官方代码两种描述，二者不一致，见正文。通道与层内组成取自公开配置，论文正文只给出块数与分辨率。}
 \label{tab:vision-controlnet-structure}
\end{table*}

表中的两种模块按以下顺序计算。残差块对输入依次作组归一化、SiLU、$3\times3$卷积，加上由噪声步编码而来的向量，再作组归一化、SiLU、$3\times3$卷积，最后与输入相加（通道数变化时输入先经$1\times1$卷积）。空间Transformer块把$h\times w$个位置展平为词元序列，依次经过自注意力、读取文字词元的交叉注意力和前馈层，每一步都带残差连接，再恢复为特征图。文字条件只经交叉注意力进入，冻结主网络和复制支路都读取同一份文字词元。

空间条件图的分辨率是$512\times512$，而支路在$64\times64$的潜网格上工作，因此先由一个小卷积网络$E$把条件图编码成$64\times64$的特征，再经$Z_1$加到复制支路输入卷积的输出上。论文文字说$E$由4层$4\times4$、步长2的卷积组成，通道依次为16、32、64、128，激活为ReLU；但4次步长为2的下采样会把512降到32，与它所说的64不符。作者公开的实现用8层$3\times3$卷积，其中3层步长为2，激活为SiLU，最后一层是输出320通道的零卷积，恰好把512降到64并与输入卷积的通道数对齐。复现时应以所用代码为准，并在报告中写明。数据流可以概括为：带噪潜变量与文字进入冻结U-Net；同一带噪潜变量经复制的输入卷积后加上编码后的条件图，进入复制编码器；复制编码器各块的输出经零卷积加到冻结U-Net解码端的相应位置；最后由冻结U-Net输出噪声预测。论文报告，在单张40GB A100上，带ControlNet的训练每次迭代比单独微调Stable Diffusion多用约23\%的显存和34\%的时间，因为冻结部分不需要计算参数梯度。

训练数据是“目标图像、文字、条件图”三元组，条件图通常由检测器从目标图像自动提取，例如用边缘检测得到边缘图、用姿态估计得到关键点图。训练目标仍是扩散模型的噪声预测误差。设目标图像的潜表示为$\bm{z}_i$，条件图为$\bm{C}_i$，提示为$c_i$；在小批量中为每个样本抽噪声步$k_i$和噪声$\bm{\epsilon}_i$，构造带噪状态$\bm{z}_i^{(k_i)}$，则
\begin{equation}
 L_{\mathrm{ctrl}}=
 \frac{1}{m}\sum_{i=1}^{m}
 \left\lVert\bm{\epsilon}_i-
 \epsilon_{\bm{\theta},\bm{\psi}}
  (\bm{z}_i^{(k_i)},k_i,c_i,\bm{C}_i)
 \right\rVert_2^2\text{，}
 \label{eq:vision-controlnet-training}
\end{equation}
其中$m$是小批量大小，$\bm{\psi}$汇集复制支路、条件编码器和零卷积的参数，只有$\bm{\psi}$被更新，$\bm{\theta}$保持冻结。原论文训练时以50\%的概率把文字替换为空串，迫使模型从条件图本身读出语义，而不是只依赖文字。论文还报告了一种“突然收敛”现象：模型在一段训练中几乎不遵循条件，随后在不到一万步内突然开始遵循。每种条件类型（边缘、深度、姿态等）分别训练一个支路。

推理时同时给出文字和条件图，在每个去噪步注入支路特征。支路输出可以乘以一个控制强度系数：强度为零时退回原模型，强度增大时输出更贴近条件图的结构。多种条件可以同时使用，原论文把多个ControlNet的输出直接相加。但若深度图规定前爪在身体后方、姿态图又规定前爪伸向相机，两种条件没有一致的几何解，相加并不能消解矛盾。常见失败还包括条件提取器本身出错，例如姿态估计把左右爪弄反，ControlNet会忠实地画出错误姿态。因此需要先核验条件图，再记录支路版本、条件分辨率和控制强度。空间条件只规定布局，不规定“这是哪一只玩具狗”，后者需要下面的主体适配。

\subsubsection{单一主体参考下的生成}
\label{sec:vision-generation-dreambooth}

文字“玩具狗”指一个类别，ControlNet只规定布局，二者都无法让模型画出照片里那只特定的玩具狗，例如它左耳的缝线和背上的补丁。unCLIP的图像向量又压缩掉了这些细节。要让模型在新场景、新姿态中重现一个具体主体，就需要把这个主体的外观写进模型。DreamBooth的做法是用该主体的3到5张照片微调整个文字生成模型，并把主体绑定到提示里的一个特殊标识上\scite{vision-dreambooth2023}。

主要想法分两部分。其一，选择一个在词表中很少出现、本身几乎没有语义的词元作为标识，记作“[V]”，训练提示写成“a [V] toy dog”。因为标识原本没有含义，微调可以把主体外观绑定到它上面，同时类别词“toy dog”提供了类别先验，让模型不必从零学习“玩具狗长什么样”。其二，只在几张照片上微调会带来两种退化：模型会把“toy dog”这个类别词也缩窄成这一只玩具，即\term{语言漂移}（language drift）；输出还会趋同于几张训练照片的姿态和视角。\term{类别先验保持}（class-specific prior preservation）用来抑制这两种退化：在微调前，先用原模型以提示“a toy dog”生成约1000幅一般玩具狗图像，微调时让模型继续能生成这些图像，从而保持类别词的原有含义。

DreamBooth不改变模型结构，输入输出都与原文字生成模型相同：输入是提示与带噪图像（或潜表示），输出是去噪预测。所以它的“结构”就是所微调的基础模型的结构，需要交代的是微调了其中哪些部分。原论文在Imagen和Stable Diffusion上实验，二者的组成不同。Imagen是级联模型：一个以T5-XXL文字编码为条件、输出$64\times64$图像的基础扩散模型，接两个文字条件的超分辨率扩散模型，分别完成$64\to256$和$256\to1024$。原论文发现，主体保真的最好结果来自微调全部层，包括以文字编码为条件的层；只微调基础模型时，未经微调的超分辨率模型会给主体添上错误的高频纹理，所以$64\to256$的超分辨率模型对多数主体都必须一起微调，$256\to1024$的模型只对细节特别多的主体有帮助。超分辨率模型训练时会给低分辨率输入加噪声并以噪声强度为条件，称为噪声条件增强；原论文在微调$64\to256$模型时把这一噪声强度从训练时的$10^{-3}$降到$10^{-5}$，否则主体的细密花纹会被抹平。Stable Diffusion只有一个在潜空间工作的去噪U-Net（结构见表\ref{tab:vision-controlnet-structure}），微调的是这个U-Net，文字编码器可以一并微调，自编码器的解码器保持不变。标识的选择需要经过分词器检查：原论文对Imagen从T5-XXL词表中编号5000到10000、长度不超过3个字符的词元中随机选取，随意拼出的一串字母可能被切成多个常用词元，反而带有已有含义。

训练目标由两项组成：
\begin{equation}
 L=L_{\mathrm{subj}}+\lambda_{\mathrm{pr}}L_{\mathrm{class}}\text{，}
 \label{eq:vision-dreambooth}
\end{equation}
原论文取$\lambda_{\mathrm{pr}}=1$。原论文以预测干净图像的形式写去噪目标。设主体照片为$\bm{x}$，其提示为$c_{\mathrm{subj}}$（“a [V] toy dog”），加噪状态为$\alpha_k\bm{x}+\sigma_k\bm{\epsilon}$，主体项为
\begin{equation}
 L_{\mathrm{subj}}=
 \mathbb{E}_{\bm{x},k,\bm{\epsilon}}
 \left[
 w_k\left\lVert
 f_{\bm{\theta}}(\alpha_k\bm{x}+\sigma_k\bm{\epsilon},
                  k,c_{\mathrm{subj}})-\bm{x}
 \right\rVert_2^2\right]\text{，}
 \label{eq:vision-dreambooth-subject}
\end{equation}
其中$\alpha_k,\sigma_k$是噪声日程系数，$w_k$是噪声步权重，$f_{\bm{\theta}}$输出对干净图像的估计。类别项把$(\bm{x},c_{\mathrm{subj}})$替换为预先生成的类别图像$\bm{x}_{\mathrm{pr}}$与普通提示$c_{\mathrm{pr}}$（“a toy dog”），另抽噪声和噪声步。主体项教模型“[V]是这一只”，类别项教模型“toy dog仍然是一般玩具狗”，二者监督不同的提示。若基础模型输出噪声预测，可由$\hat{\bm{x}}=(\bm{z}^{(k)}-\sigma_k\hat{\bm{\epsilon}})/\alpha_k$换算，得$\lVert\hat{\bm{x}}-\bm{x}\rVert_2^2=(\sigma_k^2/\alpha_k^2)\lVert\hat{\bm{\epsilon}}-\bm{\epsilon}\rVert_2^2$，所以改写成噪声误差时要同步变换噪声步权重。

训练时每个小批量同时抽主体照片和类别图像，分别计算两项后加权求和，反传更新全部被微调的参数。两种基础模型上的训练配置见表\ref{tab:vision-dreambooth-config}。论文没有给出批大小和优化器设置。训练结束后保存微调权重、标识和类别提示。推理时输入“a [V] toy dog sitting in the snow”，即可在雪地中生成这只玩具狗，不再需要参考照片。

\begin{table}[htbp]
 \centering
 \small
 \begin{tabularx}{\linewidth}{@{}p{22mm}XX@{}}
  \toprule
  项目 & Imagen & Stable Diffusion\\
  \midrule
  基础模型组成 & $64\times64$基础模型，$64\to256$与$256\to1024$两级超分辨率 & 潜空间去噪U-Net，自编码器，文字编码器\\
  微调的部分 & 全部层，含读取文字编码的层；$64\to256$超分辨率模型一并微调，$256\to1024$可选 & 去噪U-Net，文字编码器可选；解码器不变\\
  主体照片 & 3到5张 & 3到5张\\
  类别图像 & 约1000幅，也可更少 & 约1000幅，也可更少\\
  $\lambda_{\mathrm{pr}}$ & 1 & 1\\
  学习率 & $10^{-5}$ & $5\times10^{-6}$\\
  迭代次数 & 约1000 & 约1000\\
  训练时间 & 单个TPUv4约5分钟 & 单张A100约5分钟\\
  其他 & 标识取自T5-XXL词表编号5000到10000、不超过3个字符的词元；超分辨率噪声条件增强由$10^{-3}$降至$10^{-5}$ & 论文未给出标识的选取范围\\
  \bottomrule
 \end{tabularx}
 \caption{DreamBooth在两种基础模型上的微调配置，数值取自原论文。DreamBooth不改动网络结构，表中列出的是被微调的部分与训练设置；批大小与优化器论文未给出。}
 \label{tab:vision-dreambooth-config}
\end{table}

DreamBooth的边界也来自它的训练数据。若只有正面照片，生成背面时的纹理来自类别先验，而不是这只玩具；增大$\lambda_{\mathrm{pr}}$保持了类别多样性，却不能补出从未拍到的外观。训练过久会让模型复制参考照片的背景；多个主体各自微调后再组合，又会出现标识绑定混淆。原论文用DINO特征相似度和CLIP图像相似度衡量主体保真，用CLIP图文相似度衡量提示遵循，二者需要同时报告，因为单纯复制参考照片会让前者很高、后者很低。

\subsubsection{多参考图像下的生成}
\label{sec:vision-generation-reference-roles}

现在请求变成：“让第一张照片里的玩具狗，坐在第二张照片里那种蓝釉杯子旁边，画成第三张图的水彩风格”。三幅参考分别负责主体身份、局部外观和整体风格。上面三种情形都只有一种条件，此时新出现的难点是参考之间的职责分配：若只写“参考这三张图”，模型可能把杯子的蓝色涂到狗身上，把第一张的桌面背景一起复制，或者把水彩风格只用在杯子上。

解决这个问题的第一步不在模型里，而在要求记录中：为每幅参考写明它对应的目标对象和应取用的属性，例如“参考1：玩具狗的身份（耳朵、斑纹）；参考2：杯子的釉色与形状；参考3：全图笔触与色调；布局由文字规定”。我们把这种对应称为参考角色绑定。有了它，生成后才能分项检查是否发生串换，例如狗是否变成蓝色、杯子是否保留了形状。

第二步是选用支持多图输入的模型接口，并确认它确实能区分各参考的角色。可选的做法分两类：像DreamBooth那样为主体单独微调，再用其他条件提供其余参考；或使用在推理时直接读取参考图像的接口。\ref{sec:vision-generation-kontext-history}节要讲的FLUX.1 Kontext把参考图像编码为词元序列，与待生成图像的词元一起送入同一个Transformer，并为不同来源的图像分配不同的位置编号；但其原论文主要以单张上下文图像训练，多张参考同时输入不属于它已经验证的能力。所以角色绑定是一份要求规格，不能为只接收单图的模型凭空增添多参考能力。

冲突应在采样前暴露。若参考2是蓝杯子，提示却要求杯子是红色，就要在记录里写明“取杯型、改颜色”；若两处条件要求杯子分别在狗的左右两侧，必须先选定一个。保留属性和可变属性分别记录以后，才能判断输出中的颜色变化是按要求修改还是串换。参考未展示的属性仍由生成先验补充，精确文字仍需读取核验，整体参考相似度不能代替这些分项判断。

四种条件的比较见表\ref{tab:vision-generation-condition-comparison}。固定任务为“在雪地里画这只玩具狗”，它们的差别在于条件提供的信息、模型要付出的适配成本以及典型的失败方式。文字条件不需要任何训练，但只能约束语义；ControlNet用每类条件一次支路训练换来像素级布局控制；DreamBooth用每个主体一次全模型微调换来身份保持；多参考接口省去逐主体训练，却要求模型确实支持多图角色区分，并把角色规格交给使用者写明。

\begin{table}[htbp]
 \centering
 \small
 \begin{tabularx}{\linewidth}{@{}l X X X X@{}}
  \toprule
  条件情形 & 约束的信息 & 需要的适配 & 每次使用的输入 & 典型失败\\
  \midrule
  文字（LDM、unCLIP） & 语义、类别、风格 & 无 & 提示 & 属性串换、数量与位置不准\\
  空间图（ControlNet） & 逐像素对齐的布局 & 每类条件训练一条支路，原模型冻结 & 提示与条件图 & 条件图错误被忠实画出，多条件矛盾\\
  单一主体（DreamBooth） & 具体主体外观 & 每个主体微调整个模型 & 含标识的提示 & 复制背景、未见外观来自类别先验\\
  多参考图像 & 各参考的指定属性 & 依所用接口而定 & 提示、多幅参考与角色记录 & 属性在参考间串换\\
  \bottomrule
 \end{tabularx}
 \caption{四种创作条件在同一创建任务下的比较。“需要的适配”指为得到该能力而额外进行的训练；推理成本另随采样步数和候选数变化。}
 \label{tab:vision-generation-condition-comparison}
\end{table}

\subsection{候选图像的生成与选择}
\label{sec:vision-generation-selection}

条件组织好以后，采样会给出一幅具体图像，但不保证它满足要求记录。固定条件和种子，问题变为怎样提高最终交付图像满足条件的概率。有两种时机可以干预。一是在采样过程中修正每一步的去噪方向，使轨迹更偏向条件，这改变了候选本身；二是先生成多个候选，再用打分器挑选，这不改变任何候选，只改变交付哪一个。两者都可能用到图文相似度，但作用的对象、时机和成本不同，下面依次讨论。

\subsubsection{采样中的引导修正}
\label{sec:vision-generation-glide-guidance}

\ref{sec:vision-models-glide}节已经推导了GLIDE中的两种采样引导\scite{vision-glide2022}。无分类器引导在每一步分别计算同一网络的有条件和无条件噪声预测，按式\eqref{eq:vision-glide-cfg-interface}以强度$g$外推；CLIP引导则用一个在带噪图像上训练过的CLIP模型，按式\eqref{eq:vision-glide-clip-interface}把图文匹配分数对当前图像的梯度乘以强度$g_{\mathrm{clip}}$加到反向均值上。这里不再重复推导，只讨论使用时最常遇到的三个问题：强度怎样选择，代价多大，以及怎样核验它的效果。

先看强度对结果的影响。$g=1$等于不加引导，$g>1$时每一步都沿“有条件减无条件”的方向多走一段。GLIDE报告，提高无分类器引导强度会让样本更贴合文字、画面更像照片，同时样本多样性下降，二者此消彼长；在人工比较中，无分类器引导的结果在真实感和图文相符两方面都比CLIP引导更受偏好。原因之一是CLIP引导沿着一个固定判别模型的梯度修改像素，生成器可能找到让该CLIP模型打高分、人看来却不合理的图像，这类似于针对该模型的对抗样本。强度过大时还会出现颜色过饱和、对比过强和构图单调。所以强度没有“越大越好”的方向，通常在一个中等范围内按任务确定。

强度的选择应作为一次有记录的实验。固定模型、提示、种子集合和采样步数，例如对“红色玩具狗在蓝色杯子左边”用8个种子，分别取$g=1,3,5,7.5,10$生成，得到40幅图。每个强度下按要求记录逐项统计满足率（颜色归属、数量、左右关系），同时统计同一强度下8幅图之间的差异，例如背景和姿态是否趋同。若满足率在某个强度之后不再上升而多样性继续下降，更大的强度就只带来代价。选定的强度只对这组模型、采样器和任务有效，换一种采样器或换成编辑任务都要重新确定。

再看成本。无分类器引导在每一步需要两次网络前向，一次有条件、一次无条件，可以拼成一个批次计算，计算量约为不加引导的两倍。CLIP引导除了去噪网络，每一步还要前向并反传一个噪声感知CLIP模型，而且这个模型需要单独训练，普通CLIP只在清晰图像上训练过，直接接到带噪状态上不能沿用GLIDE的结论。报告引导效果时，应把每步网络调用次数和额外模型一并写明。

核验时要记住，引导提高的是模型自身的条件概率，或者某个CLIP模型的匹配分数，不是要求记录中的每一项。更强的引导可能让“红色”更显著，却同时让狗与杯子的颜色更容易串换。因此引导强度的效果只能由逐项核验判断，不能以平均图文相似度上升代替。

\subsubsection{生成后的候选重排}
\label{sec:vision-generation-reranking}

重排发生在候选已经形成之后。对$k$幅候选$\bm{X}_1,\ldots,\bm{X}_k$计算打分并排序，它不修改任何候选的像素，也不能从全部错误的候选中造出正确答案。最常用的打分是CLIP图文相似度：把每个候选和提示分别编码成单位向量，取内积。GLIDE对照的DALL-E就采用了这种生成后的CLIP重排\scite{vision-glide2022}；不同系统使用的候选总数不同，比较时必须连同选择协议一起阅读。

全局图文相似度只衡量粗粒度的相关性，按它排序会选错。表\ref{tab:vision-generation-candidates}给出四个构造候选：若只取相似度最高者，会选中出现两只狗的A；若先按要求记录逐项筛选，再在通过者中按质量排序，只有D满足全部要求。这里的分项检查由人工或独立的检测、读取工具完成，不是CLIP分数提供的保证。

\begin{table}[htbp]
 \centering
 \small
 \begin{tabularx}{\linewidth}{@{}c c X c@{}}
  \toprule
  候选 & 构造相似度 & 按请求核验的主要结果 & 保留\\
  \midrule
  A & 0.34 & 出现两只玩具狗 & 否\\
  B & 0.32 & 狗为蓝色，杯子为红色 & 否\\
  C & 0.30 & 属性正确，但狗在右侧 & 否\\
  D & 0.28 & 数量、颜色归属与左右关系均满足 & 是\\
  \bottomrule
 \end{tabularx}
 \caption{全局相似度与逐项条件检查的教学对照。分数人为设置，用于展示两种选择规则的区别；不代表任何模型的实测排序。}
 \label{tab:vision-generation-candidates}
\end{table}

由此得到一个两阶段的选择流程：先用硬条件筛选，例如用检测器数狗和杯子的个数、读取它们的位置和颜色；再在通过者中按图文相似度或审美分数排序。重排的成本与候选数成正比：生成$k$幅图就要付出$k$次完整采样，再加上$k$次打分。若一个请求生成了100幅图、人工选出1幅，报告时应写明100幅的生成成本和筛选规则；按“最佳候选”计分与按“一次随机生成”计分，衡量的是两个不同的系统。没有候选通过硬条件时，应返回“条件未满足”，或在明确预算内继续生成或修订，不能把最高分候选自动称为成功。

两种干预的关系可以这样概括。引导在采样中改变每个候选，代价是每步多一次前向（或一次CLIP反传），影响所有候选，并以多样性为代价；重排不改变候选，代价随候选数线性增加，而且只能在已有候选中选择。二者可以组合，例如以中等引导强度生成8个候选，再按硬条件筛选。比较不同系统时应固定引导强度和候选总数，否则无法判断差异来自模型还是选择过程。

\subsection{生成结果的反思修订}
\label{sec:vision-generation-unireason-refinement}

引导和重排都在“一次生成”内起作用。若8个候选都把杯子画在了狗的左侧，重排无能为力，提高引导强度也未必能修正一条具体的空间关系。此时更直接的办法是看着已有的草图，指出哪里不符合要求，再按这条意见修改。生成后的反思修订就是这样一个“生成—检查—编辑”的环节。\ref{sec:vision-models-unireason}节介绍的UniReason 1.0把它训练进一个统一模型：模型先生成草图，再读取草图与原请求，用文字写出偏差说明，最后以草图和偏差说明为条件生成修订图像\scite{vision-unireason-2026}。它的结构、训练数据和联合目标都在那一节讲过，这里只讨论在创作中怎样使用：编辑请求如何形成，修订预算如何规定，修订结果如何验收。

编辑请求的质量决定修订能否成功。偏差说明必须依据实际草图指出“哪个对象、哪个属性或哪条关系需要怎样改变”，并说明哪些部分应保持。例如草图出现两只狗、杯子在狗左侧，合格的请求是“只保留一只红色玩具狗；把蓝杯子移到狗的右侧；保持狗的斑纹和背景”；而“请更符合原提示”只是重复原要求，模型无从知道改哪里。若草图数量和位置都正确，只有阴影方向与光源矛盾，请求就应只涉及阴影，不应让模型重画主体。使用时可以把模型自己写出的偏差说明与要求记录逐项对照：它漏掉的偏差不会被修正，它误报的偏差可能引入新的错误。

修订预算需要事先规定。UniReason原论文的实现只作一轮修订，不能外推为“修订越多次越好”。每一轮都要额外付出一次图像编码、一段文字生成、一次完整图像采样和解码，计算量大约是一次直接生成的两倍以上。更重要的是，修订可能在纠正数量的同时引起主体漂移，例如狗的斑纹改变。所以合理的规则是设定修订轮数上限（例如一轮），并保存原草图；若修订结果在已满足的项上退步，就拒绝修订、退回草图，或者转为下一节的局部编辑。

验收应固定原条件，把直接生成、加入文字推理和再加一轮修订的三个结果并排比较，逐项检查数量、颜色归属、左右关系和阴影，然后才看自然度。若偏差说明全部正确而修订图仍画错关系，应记为图像实施失败，而不能因文字推理正确就算任务完成。UniReason训练时用外部大模型核验草图并构造修订数据，验收时若仍用同类模型裁判打分，结果可能只是更迎合这个裁判；因此还要加入区域检测、字符读取或人工抽查等独立检查。

\section{已有图像的转换}
\label{sec:vision-generation-translation}

转换以一幅已有图像为输入，输出另一种风格、呈现或成像形式的图像，例如把玩具狗照片转成水彩画、把边缘草图转成照片、把白天街景转成夜景。与创建相比，输入图像带来了一个新的约束：输出必须与输入保持某种对应，例如狗的位置、姿态和数量不变；与下一节的编辑相比，转换作用于整幅图像，变化的方式由事先确定的“域”决定，而不是由每次的请求临时规定。难点在于怎样学到一个既改变外观、又保持对应的映射，而能够学到什么样的对应，取决于训练数据提供了什么。

\subsection{转换要求的规定}
\label{sec:vision-generation-translation-requirements}

转换要求要写清两件事：输出属于什么目标形式，以及哪些内容必须与输入对应。把玩具狗照片转成水彩画时，纹理、笔触和色调可以改变，狗的姿态、数量和画面布局要保持；把边缘图转成照片时，每条边缘的位置都必须保留，颜色和材质可以由模型补充；若把一种成像方式转换成另一种，例如把卫星图转成地图，还要求每个位置的结构与实际地物对应。三种用途即使用相同网络实现，允许的误差也不同。

训练数据必须与这种要求相符。对齐的“边缘图—照片”配对，逐像素地告诉模型输入的每个位置应该变成什么；只有一组照片和一组水彩画、二者并不对应时，数据只说明两个域各自的外观；只有一幅指定风格的画时，数据提供的是纹理统计，而不是任何一幅图的目标。若把没有逐位置对应的图像当作像素监督，损失就会迫使网络拟合本不存在的对应关系。因此，转换要求、可用数据和下面的模型选择是联动的。

\subsection{转换映射的学习}
\label{sec:vision-generation-translation-mapping}

转换映射是一个前馈网络$G$，读入一幅图像，一次前向输出转换结果。三种模型的差别在于用什么监督来约束“输出与输入的对应”。pix2pix有逐像素配对的目标图像，直接用目标图约束输出，并用判别器约束局部外观；CycleGAN没有配对，只能要求转换后再转回来能恢复输入，用这种往返关系代替配对目标；Johnson等人的感知损失网络面对的是固定风格转换，既没有配对目标也不需要另一个域的数据集，它用一个固定的预训练网络在特征空间中分别度量“内容是否保留”和“风格是否接近”。这三条路线依次放宽了对训练数据的要求，也依次减弱了对输出的约束。

\subsubsection{pix2pix}
\label{sec:vision-generation-pix2pix}

卷五的“基于神经网络的生成模型”章已经完整讲了pix2pix的结构与目标：它是条件生成对抗网络，生成器是带跳连的U-Net，判别器是读入“输入—输出”对的局部块判别器，目标是条件对抗项加上$\lambda$倍的$\ell_1$像素误差\scite{vision-pix2pix2017}。这里回顾它要解决的问题，然后侧重在转换任务中怎样组织数据、选择超参数和使用。

pix2pix的动机是：对于有配对数据的转换，只用$\ell_1$像素误差训练，网络会对“同一输入可能对应的多种纹理”取平均，输出模糊；只用对抗项，输出逼真但可能不保留输入结构。它的主要想法是两者兼用：$\ell_1$项把输出拉向配对目标，保证低频结构对应；判别器只看局部图像块，专门惩罚模糊和不真实的纹理。判别器同时读入输入，使它判断的是“这个输出是否像这个输入的真实对应”，而不仅是“像不像照片”。

以$256\times256$的边缘图转玩具狗照片为例，逐层结构见表\ref{tab:vision-pix2pix-structure}。论文附录用两个记号描述各层：C$k$表示“卷积—批归一化—ReLU”，有$k$个卷积核；CD$k$在批归一化之后多一个丢弃率50\%的dropout。所有卷积都是$4\times4$，编码器和判别器中步长为2、每层把分辨率减半，解码器中每层把分辨率加倍。编码器的ReLU都是斜率0.2的leaky ReLU，第一层C64不加批归一化；解码器使用普通ReLU。

\begin{table*}[tbp]
 \centering
 \small
 \begin{tabularx}{\linewidth}{@{}p{30mm}Xp{24mm}p{22mm}p{12mm}p{20mm}@{}}
  \toprule
  阶段或层组 & 层类型 & 卷积核与步长 & 输出通道 & 重复次数 & 输出尺寸\\
  \midrule
  \multicolumn{6}{@{}l}{\textit{生成器（U-Net），输入$256\times256\times3$}}\\
  编码层1 & 卷积，leaky ReLU（无批归一化） & $4\times4$，步长2 & 64 & 1 & $128\times128$\\
  编码层2 & C128 & $4\times4$，步长2 & 128 & 1 & $64\times64$\\
  编码层3 & C256 & $4\times4$，步长2 & 256 & 1 & $32\times32$\\
  编码层4至8 & C512 & $4\times4$，步长2 & 512 & 5 & $16\times16\to1\times1$\\
  解码层1至3 & 上采样卷积，批归一化，dropout，ReLU，再拼接同分辨率编码特征 & $4\times4$，步长$1/2$ & 512，拼接后1024 & 3 & $2\times2\to8\times8$\\
  解码层4 & 同上，无dropout & $4\times4$，步长$1/2$ & 512，拼接后1024 & 1 & $16\times16$\\
  解码层5至7 & 同上，无dropout & $4\times4$，步长$1/2$ & 256、128、64，拼接后512、256、128 & 3 & $32\times32\to128\times128$\\
  输出层 & 上采样卷积，tanh & $4\times4$，步长$1/2$ & 3 & 1 & $256\times256$\\
  \midrule
  \multicolumn{6}{@{}l}{\textit{$70\times70$判别器，输入为输入图与输出图沿通道拼接的$256\times256\times6$}}\\
  C64 & 卷积，leaky ReLU（无批归一化） & $4\times4$，步长2 & 64 & 1 & $128\times128$\\
  C128 & 卷积，批归一化，leaky ReLU & $4\times4$，步长2 & 128 & 1 & $64\times64$\\
  C256 & 卷积，批归一化，leaky ReLU & $4\times4$，步长2 & 256 & 1 & $32\times32$\\
  C512 & 卷积，批归一化，leaky ReLU & $4\times4$，步长1（代码） & 512 & 1 & $31\times31$\\
  输出层 & 卷积，sigmoid & $4\times4$，步长1（代码） & 1 & 1 & $30\times30$\\
  \bottomrule
 \end{tabularx}
 \caption{pix2pix的生成器与$70\times70$判别器。层序列、通道与激活取自论文附录：编码器C64-C128-C256-C512-C512-C512-C512-C512，U-Net解码器按论文记号为CD512-CD1024-CD1024-C1024-C1024-C512-C256-C128，记号中的数字是各解码层的输入通道数，即上一层输出与跳连拼接之后的通道数，最后的C128对应输出层的输入。论文写所有卷积步长为2，但那样判别器的感受野不是$70\times70$；表中判别器最后两层的步长1与$30\times30$输出取自作者公开的代码。}
 \label{tab:vision-pix2pix-structure}
\end{table*}

数据流如下。输入$\bm{X}\in\mathbb{R}^{256\times256\times3}$经编码器8层步长为2的卷积，逐步缩小到$1\times1\times512$；解码器对称地8次上采样，第$i$个解码层的输出与第$8-i$个编码层的输出沿通道拼接后送入下一层，所以论文把解码层的通道写成1024、512等拼接后的数目。最终输出$\widehat{\bm{Y}}\in\mathbb{R}^{256\times256\times3}$，经tanh落在$[-1,1]$。跳连使边缘位置可以直接从编码器传到解码器，而不必经过$1\times1$的瓶颈；没有跳连的编码—解码版本，解码器为CD512-CD512-CD512-C512-C256-C128-C64，所有信息都必须穿过瓶颈。判别器把$\bm{X}$与$\widehat{\bm{Y}}$（或真实$\bm{Y}$）沿通道拼成6通道输入，输出一个$30\times30$的分数图。每个分数的感受野可以逐层累加：对$4\times4$卷积，第$\ell$层使感受野增加$3\times(\text{前面各层步长之积})$，五层的步长依次为$2,2,2,1,1$，感受野为$1+3\times(1+2+4+8+8)=70$。所以每个分数只判断输入上一个$70\times70$邻域内的配对是否真实。论文还比较了感受野为$1\times1$、$16\times16$和$286\times286$的判别器，它们只是增减了C层的数目，例如$16\times16$判别器为C64-C128。

训练时交替更新：判别器以真实对为正、生成对为负，在分数图的每个位置上计算对数损失再平均；生成器使用非饱和对抗项，并加上$\lambda\lVert\bm{Y}-\widehat{\bm{Y}}\rVert_1$，原论文取$\lambda=100$。优化判别器时把目标除以2，减慢判别器相对生成器的学习；两者都用Adam，学习率$2\times10^{-4}$，$\beta_1=0.5,\beta_2=0.999$，权重从均值0、标准差0.02的高斯分布初始化，全部从头训练。数据增强先把$256\times256$的图像放大到$286\times286$，再随机裁回$256\times256$。可以手算一次：设两个局部块的真实分数为$(0.8,0.6)$、生成分数为$(0.2,0.4)$，判别损失为$-\{\log0.8+\log0.6+\log0.8+\log0.6\}/2\approx0.734$（自然对数）；生成器的非饱和对抗项为$-(\log0.2+\log0.4)/2\approx1.263$。像素项是求和还是求平均必须与$\lambda$配套说明，否则权重不可比较。数据增强必须对输入与目标做相同的裁剪和翻转，保证配对位置不变。

使用时只运行生成器，判别器和目标图都不再需要。原实现在测试时保留dropout，并用当前测试批次计算批归一化统计（批大小为1时相当于实例归一化），复现时要保存这些设置。常见失败有三种：训练配对存在系统性错位时，网络会学会这种错位；同一输入有多种合理纹理时，输出缺乏多样性，原论文发现加入随机噪声会被网络忽略，dropout带来的变化也很有限；局部块逼真不保证整幅图的结构和身份正确。

\subsubsection{CycleGAN}
\label{sec:vision-generation-cyclegan}

若要把玩具狗照片转成水彩画，几乎不可能找到同一只狗、同一姿态的“照片—水彩”配对，pix2pix的$\ell_1$项就无从计算。CycleGAN处理的正是只有两组互不对应的图像的情形\scite{vision-cyclegan2017}；其目标和推导同样已在卷五该章讲过。它的主要想法是：同时训练照片到水彩的$G:A\to B$和水彩到照片的$F:B\to A$，用目标域判别器要求转换结果像目标域图像，再用\term{循环一致性}（cycle consistency）要求$F(G(\bm{x}))\approx\bm{x}$、$G(F(\bm{y}))\approx\bm{y}$。往返能恢复输入，说明转换没有丢掉输入的信息，这就替代了缺失的配对目标。

循环损失为
\begin{equation}
 L_{\mathrm{cyc}}=
 \mathbb{E}_{\bm{x}\sim p_A}
 [\lVert F(G(\bm{x}))-\bm{x}\rVert_1]
 +
 \mathbb{E}_{\bm{y}\sim p_B}
 [\lVert G(F(\bm{y}))-\bm{y}\rVert_1]\text{。}
 \label{eq:vision-cycle-consistency}
\end{equation}
总目标是两个方向的对抗项加上$\lambda L_{\mathrm{cyc}}$，原论文取$\lambda=10$。架构上，两个生成器沿用Johnson等人的转换网络（见下一小节），逐层结构见表\ref{tab:vision-cyclegan-structure}。论文附录沿用Johnson代码库的命名：c7s1-$k$是$7\times7$、步长1、$k$个卷积核的“卷积—实例归一化—ReLU”；d$k$是$3\times3$、步长2的同类层，用于下采样；R$k$是含两层$3\times3$卷积、两层卷积核数相同的残差块；u$k$是$3\times3$、步长$1/2$的分数步长卷积层，用于上采样。卷积使用反射填充以减少边界伪影。对$256\times256$及更高分辨率的训练图像，生成器为c7s1-64、d128、d256、9个R256、u128、u64、c7s1-3；对$128\times128$的图像，残差块减为6个。实例归一化对每幅图、每个通道单独计算均值和方差，与批大小1的训练相配。两个判别器都是$70\times70$的局部块判别器，层序列为C64-C128-C256-C512，C$k$为$4\times4$的“卷积—实例归一化—leaky ReLU（斜率0.2）”，第一层不加实例归一化，最后再用一层卷积输出单通道分数图。与pix2pix不同，它不读入输入图像，只判断一幅图是否属于本域。

\begin{table*}[tbp]
 \centering
 \small
 \begin{tabularx}{\linewidth}{@{}p{30mm}Xp{24mm}p{22mm}p{12mm}p{20mm}@{}}
  \toprule
  阶段或层组 & 层类型 & 卷积核与步长 & 输出通道 & 重复次数 & 输出尺寸\\
  \midrule
  \multicolumn{6}{@{}l}{\textit{生成器（9个残差块），输入$256\times256\times3$}}\\
  c7s1-64 & 卷积，实例归一化，ReLU & $7\times7$，步长1 & 64 & 1 & $256\times256$\\
  d128 & 卷积，实例归一化，ReLU & $3\times3$，步长2 & 128 & 1 & $128\times128$\\
  d256 & 卷积，实例归一化，ReLU & $3\times3$，步长2 & 256 & 1 & $64\times64$\\
  R256 & 残差块，两层卷积 & $3\times3$，步长1 & 256 & 9 & $64\times64$\\
  u128 & 分数步长卷积，实例归一化，ReLU & $3\times3$，步长$1/2$ & 128 & 1 & $128\times128$\\
  u64 & 分数步长卷积，实例归一化，ReLU & $3\times3$，步长$1/2$ & 64 & 1 & $256\times256$\\
  c7s1-3 & 卷积 & $7\times7$，步长1 & 3 & 1 & $256\times256$\\
  \midrule
  \multicolumn{6}{@{}l}{\textit{$70\times70$判别器，输入$256\times256\times3$}}\\
  C64至C512 & 卷积，实例归一化（C64除外），leaky ReLU & $4\times4$，见表\ref{tab:vision-pix2pix-structure} & 64、128、256、512 & 4 & $128\times128\to31\times31$\\
  输出层 & 卷积 & $4\times4$，步长1（代码） & 1 & 1 & $30\times30$\\
  \bottomrule
 \end{tabularx}
 \caption{CycleGAN的生成器与判别器，层序列取自论文附录。判别器与pix2pix的$70\times70$判别器结构相同，只是把批归一化换成实例归一化、输入只有3通道；各层步长与输出尺寸的说明同表\ref{tab:vision-pix2pix-structure}。附录对最后一层c7s1-3的输出激活没有单独说明。}
 \label{tab:vision-cyclegan-structure}
\end{table*}

论文只说明残差块含两层$3\times3$卷积；按作者公开的实现，单个残差块内的计算顺序是：$3\times3$卷积、实例归一化、ReLU、$3\times3$卷积、实例归一化，再与块的输入相加。9个残差块都在$64\times64\times256$上工作，转换主要在这个分辨率上完成；下采样把感受野扩大，上采样再恢复到原分辨率。由于整个网络都是卷积，训练后可以处理与训练尺寸不同的输入。

训练时从两个域独立抽样，批大小为1。对抗项采用最小二乘形式：以域$B$为例，判别器最小化$\frac12\mathbb{E}[(D_B(\bm{y})-1)^2]+\frac12\mathbb{E}[D_B(G(\bm{x}))^2]$，生成器最小化$\mathbb{E}[(D_B(G(\bm{x}))-1)^2]$，平方误差在判别器很自信时仍能给生成器提供梯度。判别器用最近生成的50幅图像组成的缓冲池更新，减弱它只追随最新生成器的振荡。原实现用Adam、学习率$2\times10^{-4}$训练100轮，再在接下来100轮线性衰减到零。对照片转绘画这类任务，原论文还加入“输入已属于目标域时尽量不改动”的恒等映射项，以保持颜色构成。使用时只运行需要的方向$G$，反向生成器和两个判别器只在训练中使用。

循环一致性有一个根本限制：它约束信息不丢失，却不唯一确定语义对应。设两个域各只有一幅狗图和一幅猫图，概率各为$1/2$。若$G$把狗映成猫、猫映成狗，$F$执行逆映射，那么两个输出分布都正确、循环误差为零，但语义完全交换。图\ref{fig:vision-cycle-ambiguity}把这个反例与保持类别的映射并排显示。实际训练中，卷积网络的局部性倾向于保持布局，这类交换较少出现，但这是网络结构带来的倾向，不是目标函数的保证。

\begin{figure}[htbp]
 \centering
 % Both maps are bijections of two equiprobable categories.
% Target order is changed in the right panel to avoid line crossings.
\begin{tikzpicture}[x=1mm,y=1mm,font=\figurefont\small,
 b/.style={draw=BookBlue,fill=BookBlue!8,line width=.8pt,
  minimum width=12mm,minimum height=8mm,inner sep=1mm},
 a/.style={{Stealth[length=1.7mm]}-{Stealth[length=1.7mm]},
  BookInk,line width=.9pt}]
 \node[font=\figurefont\small\bfseries] at (20,14) {保持类别};
 \node[font=\figurefont\small\bfseries] at (79,14) {交换类别};
 \node[anchor=south,text=BookMuted] at (5,5) {域$A$};
 \node[anchor=south,text=BookMuted] at (35,5) {域$B$};
 \node[anchor=south,text=BookMuted] at (64,5) {域$A$};
 \node[anchor=south,text=BookMuted] at (94,5) {域$B$};
 \node[b] at (5,0) (ad) {狗};
 \node[b] at (35,0) (bd) {狗};
 \node[b] at (5,-17) (ac) {猫};
 \node[b] at (35,-17) (bc) {猫};
 \node[b] at (64,0) (ad2) {狗};
 \node[b] at (94,0) (bc2) {猫};
 \node[b] at (64,-17) (ac2) {猫};
 \node[b] at (94,-17) (bd2) {狗};
 \foreach \u/\v in {ad/bd,ac/bc,ad2/bc2,ac2/bd2} {
  \draw[a] (\u)--node[above=1mm] {$G/F$}(\v);
 }
 \node[anchor=north] at (20,-25) {循环误差为0};
 \node[anchor=north] at (79,-25) {循环误差也为0};
\end{tikzpicture}

 \caption{循环一致性不能唯一确定语义对应。两域各有等概率的狗、猫两项。左侧保持类别，右侧交换类别；每条往返路径都回到原项，且目标域边缘分布相同。反例是离散教学构造，用来说明循环与域判别约束仍有多解。}
 \label{fig:vision-cycle-ambiguity}
\end{figure}

CycleGAN的常见失败也由此而来。转换需要大幅改变几何形状时，例如把猫转成狗，局部纹理改变容易而形状改变困难，结果往往只改了纹理；生成器还可能把输入信息以人眼不可见的细微信号藏在输出中，使往返误差很小而中间图像并不忠实。因此在测量或医学等需要真实对应的转换中，无配对训练不能提供保真保证。

\subsubsection{Johnson感知损失网络}
\label{sec:vision-generation-perceptual-style}

第三种情形是：只有一幅指定的水彩画，希望任何输入照片都能转成它的风格。Gatys等人的风格迁移已经表明，在预训练卷积网络的特征空间中，用特征本身度量内容、用特征的通道相关统计度量风格，就能把一幅图的内容与另一幅图的风格组合起来\scite{vision-gatys2016}；但那种方法对每幅新照片都要从噪声出发、用数百次梯度迭代直接优化像素，每幅图都要付出一次完整的优化。Johnson等人的动机是保留这种特征空间的度量，却把逐图优化改成一次前向计算\scite{vision-perceptual2016}。

主要想法是区分两个网络。\term{感知损失}（perceptual loss）指在固定预训练网络的特征空间中计算的差异，而不是像素差异。可训练的图像变换网络$G_{\bm{\theta}}$负责生成图像；固定的损失网络（ImageNet上预训练的VGG-16）只用于计算感知损失，不更新参数。用许多内容图训练$G_{\bm{\theta}}$，使它的输出在损失网络看来内容接近输入、风格接近指定画作；训练完成后，转换一幅新图只需$G_{\bm{\theta}}$的一次前向，原论文报告与逐图优化的质量相当，而速度快约三个数量级。

风格转换用的变换网络逐层结构见表\ref{tab:vision-johnson-structure}，取自论文的补充材料。输入$256\times256\times3$的照片，先经一层$9\times9$卷积得到32通道，再经两次步长为2的$3\times3$卷积下采样到$64\times64\times128$；接着5个残差块在这个分辨率上处理；然后两次步长为$1/2$的$3\times3$卷积恢复到$256\times256$，最后一层$9\times9$卷积输出3通道。除输出层外，每个非残差卷积层后面都接空间批归一化和ReLU；输出层改用缩放的tanh，把像素限制在$[0,255]$。网络没有池化层，下采样和上采样都由带步长的卷积完成。全部为卷积层，所以训练后可以处理任意尺寸的输入。

\begin{table}[htbp]
 \centering
 \small
 \begin{tabularx}{\linewidth}{@{}p{16mm}Xp{15mm}p{9mm}p{6mm}p{15mm}@{}}
  \toprule
  阶段或层组 & 层类型 & 卷积核与步长 & 输出通道 & 重复次数 & 输出尺寸\\
  \midrule
  输入层 & 卷积，批归一化，ReLU & $9\times9$，步长1 & 32 & 1 & $256\times256$\\
  下采样 & 卷积，批归一化，ReLU & $3\times3$，步长2 & 64、128 & 2 & $128\times128\to64\times64$\\
  残差块 & 两层卷积，相加后不接ReLU & $3\times3$，步长1 & 128 & 5 & $64\times64$\\
  上采样 & 分数步长卷积，批归一化，ReLU & $3\times3$，步长$1/2$ & 64、32 & 2 & $128\times128\to256\times256$\\
  输出层 & 卷积，缩放tanh & $9\times9$，步长1 & 3 & 1 & $256\times256$\\
  \bottomrule
 \end{tabularx}
 \caption{Johnson等人用于风格转换的变换网络，取自论文补充材料。按层数计共16层（5个残差块各含2层）。同一论文用于超分辨率的网络结构不同：没有下采样，残差块为64通道，再用2或3次步长$1/2$的卷积放大4倍或8倍。}
 \label{tab:vision-johnson-structure}
\end{table}

残差块采用Gross与Wilber的设计：$3\times3$卷积、批归一化、ReLU、$3\times3$卷积、批归一化，与输入相加，相加之后不再接ReLU。补充材料对比了把每个残差块换成同样两层卷积的非残差网络，二者最终的训练损失相近，残差网络收敛更快；作者推测残差连接对更深的网络更关键，这个16层的网络较浅。

损失由三项组成。设损失网络第$\ell$层的特征为$\bm{F}_\ell(\bm{X})\in\mathbb{R}^{c_\ell\times n_\ell}$，$c_\ell$为通道数，$n_\ell$为空间位置数。内容项比较输出$\widehat{\bm{X}}$与输入$\bm{X}$在relu2\_2层的特征：$\frac{1}{c_\ell n_\ell}\lVert\bm{F}_\ell(\widehat{\bm{X}})-\bm{F}_\ell(\bm{X})\rVert_{\mathrm{F}}^2$。风格项比较输出与风格图$\bm{S}$的Gram矩阵
\begin{equation}
 \bm{\Gamma}_\ell(\bm{X})
 =\frac{1}{c_\ell n_\ell}
   \bm{F}_\ell(\bm{X})\bm{F}_\ell(\bm{X})^\top\text{，}
 \label{eq:vision-style-gram}
\end{equation}
在relu1\_2、relu2\_2、relu3\_3、relu4\_3四层上累加$\lVert\bm{\Gamma}_\ell(\widehat{\bm{X}})-\bm{\Gamma}_\ell(\bm{S})\rVert_{\mathrm{F}}^2$。第三项是总变差正则，惩罚输出相邻像素的过大跳变，使图像平滑。三项加权求和，权重决定内容与风格的折中。

为什么内容项用特征本身、风格项用Gram矩阵？Gram矩阵对空间位置求和，只记录哪些通道同时激活，不记录它们出现在哪里。用一个两通道、两位置的例子验证：内容特征$\bm{F}=\left(\begin{smallmatrix}1&0\\0&2\end{smallmatrix}\right)$给出$\bm{\Gamma}=\operatorname{diag}(1/4,1)$；把两个位置对调得到$\bm{F}'=\left(\begin{smallmatrix}0&1\\2&0\end{smallmatrix}\right)$，Gram矩阵不变，但内容特征的平方差和为10，按$c_\ell n_\ell=4$归一化为2.5。所以风格项只要求纹理统计相近，允许笔触出现在任何位置；内容项才把结构锁定在输入的相应位置。

训练数据是约8万幅MS-COCO图像，缩放为$256\times256$，与一幅固定风格图配合。每个小批量（大小为4）只抽内容图，经$G_{\bm{\theta}}$得到输出，再让损失网络分别处理输出与内容图，风格图的Gram矩阵可预先计算。梯度经过损失网络反传到输出图像，再传到$G_{\bm{\theta}}$，但损失网络本身不更新。原论文用Adam、学习率$10^{-3}$训练4万次迭代，约两轮；总变差正则的权重在$10^{-6}$到$10^{-4}$之间按每种风格交叉验证选取，不用权重衰减和dropout，单张GTX Titan X上训练约4小时。这种网络一个模型只对应一种风格，换一幅风格画要重新训练；它也不保证徽标、细小文字等精确结构的保留。

三种转换模型的比较见表\ref{tab:vision-generation-translation-comparison}。固定任务为“把玩具狗照片转成某种绘画”，并以“姿态与布局是否保持、是否像目标形式”为评价口径。CycleGAN相对pix2pix改变的是监督来源：用往返重建替代配对目标，换来对无配对数据的适用性，代价是语义对应不再唯一，还要多训练一个反向生成器和一个判别器。Johnson网络又把“另一个域的数据集”换成“一幅风格图加一个固定特征网络”，彻底去掉了判别器和对抗训练，训练稳定，但只能表达纹理统计层面的风格，不能完成需要改变语义或几何的转换。

\begin{table}[htbp]
 \centering
 \small
 \begin{tabularx}{\linewidth}{@{}l X X X@{}}
  \toprule
   & pix2pix & CycleGAN & Johnson感知损失网络\\
  \midrule
  训练数据 & 逐像素对齐的输入—目标配对 & 两个域互不对应的图像集 & 内容图集与一幅风格图\\
  对应关系的来源 & 配对目标的$\ell_1$项 & 往返重建的循环项 & 固定网络的内容特征项\\
  目标外观的来源 & 条件局部块判别器 & 目标域局部块判别器 & Gram矩阵风格项\\
  训练时的网络 & 生成器、判别器 & 两个生成器、两个判别器 & 变换网络、固定损失网络\\
  使用时的网络 & 生成器 & 一个方向的生成器 & 变换网络\\
  主要局限 & 依赖对齐配对，输出缺乏多样性 & 语义对应不唯一，难改几何形状 & 一个模型一种风格，只改纹理统计\\
  \bottomrule
 \end{tabularx}
 \caption{三种转换映射在同一转换任务下的比较。三者使用时都是一次前向计算，差别主要在训练数据与监督。}
 \label{tab:vision-generation-translation-comparison}
\end{table}

\subsection{转换结果的核验}
\label{sec:vision-generation-translation-check}

转换结果要同时满足两个方面：具有目标形式，并保持规定的对应。输出仍像照片，说明风格没有充分改变；输出是自然的水彩画但狗变成了猫，说明语义保持失败。两方面各需独立检查，一方的好结果不能抵消另一方的失败。CycleGAN的循环误差很小，只回答了往返约束，不能代替对中间结果的检查；Johnson网络的感知损失很低，也只说明在VGG特征下接近，不说明徽标和文字保留。

具体做法是固定一组未参与训练的输入，逐项记录对象、轮廓、文字与颜色的变化，按用途区分允许和不允许的改变。轮廓保持可以用边缘图或分割掩码在输入与输出之间比较；对象身份可以用\ref{chap:vision-image-recognition}章的检测器检查类别与数量；目标形式可以用目标域的分类器或人工判断。若发现系统性错位，应先检查训练配对；若没有配对依据却要求严格的成像保真，就需要增加相应监督，或者收窄用途。转换由预先确定的域关系决定变化方式；若每次请求要求改变的内容不同，就进入下一节的编辑。

\section{指定内容的编辑}
\label{sec:vision-generation-editing}

编辑的输入是一幅原图和一个请求，例如“把狗旁边的杯子换成花瓶”，输出是只在请求涉及的部分发生变化的图像。它与转换的区别在于变化由当前请求规定，而且通常只覆盖局部；与创建的区别在于区域外的内容是必须保持的约束。难点有三处：确定改哪里，在改动区域内生成与周围协调的新内容同时保持其余部分，以及在多轮请求中不丢失已经接受的修改。本节按编辑流程的步骤展开，最后讨论核验。

\subsection{编辑范围的确定}
\label{sec:vision-generation-editing-scope}

“把玩具狗旁边的杯子换成花瓶”至少包含三个要素：待替换的对象、替换后的内容和应保留的周围物体。若画面里有两个杯子，仅凭类别词无法确定改哪一个。用户可以直接给出框、掩码或参考编号，也可以从\ref{chap:vision-image-recognition}章的检测与分割结果中选择；对“左边那个蓝色杯子”这类指代的自动解析由\ref{chap:vision-multimodal-understanding}章展开，不是编辑算法的前提。

范围通常以掩码表示。沿用\ref{chap:vision-image-restoration}章的约定，$\bm{M}\in\{0,1\}^{h\times w}$中$m_{i,j}=1$表示保留区，$m_{i,j}=0$表示允许生成区。若花瓶比原杯子高，只遮住杯子轮廓就没有容纳新物体的空间，还会留下杯子原来的阴影；合理的范围应覆盖新物体可能占据的位置、接触区和阴影，并说明玩具狗是否允许被遮挡。掩码是任务要求的一部分，生成器不会自动推断出正确的编辑边界。

文字也可能不足以说明形状或位置。视觉指令把箭头、圈选、骨架或草图作为请求的一部分：圈住杯子并写“换成花瓶”时，圈选用于选择目标；画出一个更高的花瓶轮廓时，草图规定形状；画一个侧向箭头表示光线时，要求的是相应的受光与阴影，而箭头本身不应出现在结果中。VIBE基准按指示定位、形态修改和因果关系三层组织这类编辑任务，用多模态模型裁判按任务专用问题判分\scite{vision-vibe-2026}。它是评价材料，不是编辑模型。用它检验时，可以让文字相同而箭头分别指向左右两个杯子，检查被替换的对象是否随之改变；若只看整图是否出现花瓶，两次请求可能都得高分，视觉指令是否起作用却没有被检验。

位置标注必须跟随前处理变换。原图为$1000\times800$像素，目标点为$(600,400)$；若等比缩小一半，点变为$(300,200)$；若再从左侧裁去50像素，模型输入中的点是$(250,200)$。输出恢复时要逆向补回裁剪偏移并还原比例，掩码与草图也作相同变换。只缩放图像而保留原坐标，会把编辑送到错误位置。

\subsection{单轮编辑的实施}
\label{sec:vision-generation-editing-methods}

有了原图、请求和范围，单轮编辑要回答两个问题：新内容从哪里来，以及怎样让它与周围衔接。经典路线只从已有图像中搬运内容。泊松图像编辑把另一幅照片中现成的花瓶贴进范围，在梯度域中求解，使边界处没有接缝；Criminisi样例修补用同一幅图中别处的图像块，填补移除杯子后留下的空洞，并按优先级决定先填哪里；PatchMatch解决这两类块方法共同的计算瓶颈，即为每个图像块快速找到最相似的块。学习路线按文字生成原图中不存在的内容。GLIDE修补模型要求用户给出掩码，只在掩码内生成与文字和周围协调的内容；InstructPix2Pix只要一句指令，由模型自己判断哪里该改、改成什么。沿着这个顺序，新内容的来源从“现成图像”变成“生成模型”，范围的给出者从用户变成模型：经典路线对保留区的约束最严格，却只能重组已有内容；学习路线能创造新内容，但需要大量训练数据，指令路线还可能改动未指定的区域。

\subsubsection{泊松图像编辑}
\label{sec:vision-generation-poisson}

假设手头有一张现成的花瓶照片，想把花瓶放到玩具狗照片里杯子的位置。直接把花瓶区域的像素复制过去，两张照片的曝光和色温不同，花瓶边缘会出现一圈明显的接缝；在边界附近用渐变的权重混合两幅图，接缝变得模糊，却会留下一圈半透明的光晕，花瓶自身的细节也被冲淡。Pérez等人的出发点是，人眼对局部的明暗变化比对整体亮度更敏感：只要花瓶内部的明暗起伏保留下来，整体亮度略有偏移并不显眼\scite{vision-poisson2003}。于是问题改写为：区域内的结果应尽量保持源图像的梯度，区域边界上的值必须与目标图像相等。

记目标图像（玩具狗照片）为$f^*$，源图像（花瓶照片，已对齐到目标图像的坐标）为$g$，要粘贴的区域为$\Omega$，其边界为$\partial\Omega$；未知的是$\Omega$内的结果$f$。各颜色通道分别处理，下面只写一个通道。给定区域内的一个向量场$\bm{v}$作为引导，称为引导场（guidance field），泊松图像编辑求
\begin{equation}
 \min_{f}\iint_{\Omega}\lVert\nabla f-\bm{v}\rVert_2^2\,\mathrm{d}x\,\mathrm{d}y,
 \qquad f|_{\partial\Omega}=f^*|_{\partial\Omega}\text{。}
 \label{eq:vision-poisson-variational}
\end{equation}
把花瓶无缝地贴进来，取$\bm{v}=\nabla g$，这种用法称为无缝克隆。$\bm{v}=\bm{0}$时，问题变成在固定边界值下求最平缓的填充，解满足拉普拉斯方程$\Delta f=0$，其中$\Delta f=\partial^2f/\partial x^2+\partial^2f/\partial y^2$。这正是Dirichlet能量最小化：在边界值固定的函数中，使$\iint|\nabla f|^2$最小的函数满足拉普拉斯方程，可参考\BookChapterRef{chap:functional-analysis-and-approximation}{卷一“函数空间”章}。加入引导场以后，推导步骤相同，只多出一个源项。

设$f$是最小解，$\eta$是$\Omega$上任意一个光滑且在$\partial\Omega$上为0的函数。对任意实数$\epsilon$，$f+\epsilon\eta$仍满足边界条件，因此$J(\epsilon)=\iint_{\Omega}\lVert\nabla f+\epsilon\nabla\eta-\bm{v}\rVert_2^2$在$\epsilon=0$处取最小，$J'(0)=0$。展开平方并对$\epsilon$求导得
\begin{equation*}
 J'(0)=2\iint_{\Omega}(\nabla f-\bm{v})\cdot\nabla\eta\,\mathrm{d}x\,\mathrm{d}y=0\text{。}
\end{equation*}
对向量场$\bm{u}=\nabla f-\bm{v}$，乘积法则给出$\operatorname{div}(\eta\bm{u})=\nabla\eta\cdot\bm{u}+\eta\operatorname{div}\bm{u}$；对左边用散度定理，区域积分化为边界上的通量$\oint_{\partial\Omega}\eta\,\bm{u}\cdot\bm{n}\,\mathrm{d}s$，其中$\bm{n}$是边界的外法向，而$\eta$在边界上为0，这一项为0。于是
\begin{equation*}
 \iint_{\Omega}\nabla\eta\cdot(\nabla f-\bm{v})
 =-\iint_{\Omega}\eta\,\bigl(\Delta f-\operatorname{div}\bm{v}\bigr)=0
\end{equation*}
对所有这样的$\eta$成立。若连续函数$\Delta f-\operatorname{div}\bm{v}$在某点不为0，取一个只在该点附近非零、与它同号的$\eta$，积分就不为0，矛盾；所以最小解满足带源项的泊松方程
\begin{equation}
 \Delta f=\operatorname{div}\bm{v}\quad\text{在}\ \Omega\ \text{内},
 \qquad f|_{\partial\Omega}=f^*|_{\partial\Omega}\text{。}
 \label{eq:vision-poisson-equation}
\end{equation}
无缝克隆时$\operatorname{div}\bm{v}=\Delta g$，记修正量$\tilde f=f-g$，方程变为$\Delta\tilde f=0$，边界值为$f^*-g$。这给出了一个直观的解释：结果等于源图像加上一层平缓的修正，修正量在边界上恰好抹平两幅图的差，再向内部平滑过渡。花瓶内部的细节由$g$保留，整体颜色则被边界上的差带偏，这既是接缝消失的原因，也是被贴物体颜色会改变的原因。这里用到的多元变分推导与带源项的泊松方程，前两卷只讲了无源项的特例，本章给出的推导已足够当前使用。

实际计算在像素网格上进行。设$\Omega$是一组像素，$N_p$是像素$p$的4个邻居，$\partial\Omega$是不属于$\Omega$但与$\Omega$相邻的像素。对每对相邻像素$(p,q)$规定引导量$v_{p,q}$，满足$v_{q,p}=-v_{p,q}$，无缝克隆取$v_{p,q}=g_p-g_q$。离散问题是
\begin{equation*}
 \min_{\{f_p\}_{p\in\Omega}}
 \sum_{\langle p,q\rangle\cap\Omega\neq\varnothing}
 (f_p-f_q-v_{p,q})^2,
 \qquad f_q=f^*_q\ \ (q\in\partial\Omega)\text{，}
\end{equation*}
求和遍历至少有一端在$\Omega$内的相邻像素对，每对只算一次。对$p\in\Omega$求偏导：含$f_p$的项正好是$p$与它的各个邻居组成的对，每一项$(f_p-f_q-v_{p,q})^2$对$f_p$的导数是$2(f_p-f_q-v_{p,q})$（若这一对在求和中写成$(q,p)$，由$v_{q,p}=-v_{p,q}$得到同一结果）。令偏导为0，把已知的边界值移到右边，得到
\begin{equation}
 |N_p|\,f_p-\sum_{q\in N_p\cap\Omega}f_q
 =\sum_{q\in N_p\cap\partial\Omega}f^*_q
 +\sum_{q\in N_p}v_{p,q},
 \qquad p\in\Omega\text{。}
 \label{eq:vision-poisson-discrete}
\end{equation}
左边是离散拉普拉斯算子的相反数：内部像素$|N_p|=4$，即“4倍中心值减去四邻之和”，图像锐化中用到的就是这个模板（\ref{sec:vision-transforms-sharpening}节）。式\eqref{eq:vision-poisson-discrete}对$\Omega$中每个像素写一个方程，构成一个稀疏线性方程组。它的系数矩阵对称；若$\Omega$的每个连通部分都与边界相邻，矩阵还是正定的，因为对应的二次型等于$\Omega$内相邻像素差的平方和加上与边界相邻像素值的平方和，只有全零向量能使它为0。对称正定的稀疏方程组可以用共轭梯度法求解，可参考\BookChapterRef{chap:linear-algebra-and-matrix-analysis}{卷一“线性空间”章}；原论文用的是带逐次超松弛的Gauss–Seidel迭代或V循环多重网格。Gauss–Seidel迭代按扫描顺序逐个更新未知量，每次都使用当前最新的邻居值，见算法\ref{alg:vision-poisson-gauss-seidel}。

\begin{algorithm}[泊松图像编辑的Gauss–Seidel求解]
 \label{alg:vision-poisson-gauss-seidel}
 \begin{algorithmic}[1]
  \Require 目标图像$f^*$，源图像$g$，区域$\Omega$，松弛系数$\omega\in[1,2)$，扫描次数$T$
  \Ensure 区域内的结果$\{f_p\}_{p\in\Omega}$
  \State 初始化：对$p\in\Omega$令$f_p\gets g_p$；对相邻像素令$v_{p,q}\gets g_p-g_q$
  \For{$t=1,\ldots,T$}
   \For{按扫描顺序取$p\in\Omega$}
    \State $u\gets\Bigl(\sum_{q\in N_p\cap\Omega}f_q+\sum_{q\in N_p\cap\partial\Omega}f^*_q+\sum_{q\in N_p}v_{p,q}\Bigr)\big/|N_p|$
    \State $f_p\gets(1-\omega)f_p+\omega u$
   \EndFor
  \EndFor
  \State \Return $\{f_p\}_{p\in\Omega}$，区域外取$f^*$
 \end{algorithmic}
\end{algorithm}

用一维的例子可以手算全过程，见图\ref{fig:vision-poisson-1d}。像素编号0到5，$\Omega=\{1,2,3,4\}$，目标图像在两个边界像素上的值为$f^*_0=10$、$f^*_5=20$；源图像为$g=(0,0,4,8,4,0)$，像一个小鼓包。一维时每个像素只有左右两个邻居，式\eqref{eq:vision-poisson-discrete}给出
\begin{equation*}
 \begin{aligned}
 2f_1-f_2&=10+(0-0)+(0-4)=6,\\
 2f_2-f_1-f_3&=(4-0)+(4-8)=0,\\
 2f_3-f_2-f_4&=(8-4)+(8-4)=8,\\
 2f_4-f_3&=20+(4-8)+(4-0)=20\text{。}
 \end{aligned}
\end{equation*}
解为$f=(12,18,24,22)$，代回四个方程可以逐一验证。直接粘贴得到$(10,0,4,8,4,20)$，两端分别跳变$-10$和$+16$；泊松结果的相邻差为$(2,6,6,-2,-2)$，源图像的相邻差为$(0,4,4,-4,-4)$，每一处都多了2。边界上两幅图的总落差是$(20-10)-(0-0)=10$，被平均分摊到5个相邻差上，鼓包的形状完整保留，只是整体被抬高并倾斜。修正量$\tilde f=f-g=(12,14,16,18)$在10与20之间线性变化，这正是一维的$\Delta\tilde f=0$。若用$\omega=1$的Gauss–Seidel从$f=g$出发，第一遍扫描依次得到$f_1=(4+10-4)/2=5$、$f_2=(5+8+0)/2=6.5$、$f_3=(6.5+4+8)/2=9.25$、$f_4=(9.25+20+0)/2=14.625$，离解还很远：边界信息每遍只向内传播有限的距离，区域越大需要的扫描越多，这就是原论文同时采用多重网格的原因。

\begin{figure}[htbp]
 \centering
 % 一维泊松融合教学例：像素0-5，Omega={1,2,3,4}，边界值f*_0=10、f*_5=20，
% 源图像g=(0,0,4,8,4,0)；直接粘贴(10,0,4,8,4,20)，泊松解(10,12,18,24,22,20)。宽约93mm。
\begin{tikzpicture}[x=13mm,y=1.6mm,font=\figurefont\small]
 \fill[BookPaper] (.5,-1) rectangle (4.5,27);
 \node[anchor=north west,text=BookMuted] at (.55,27) {区域$\Omega$};
 \draw[BookMuted,line width=.7pt,->] (-.3,0)--(5.5,0) node[right] {像素};
 \draw[BookMuted,line width=.7pt,->] (-.3,0)--(-.3,28) node[above] {值};
 \foreach \y in {10,20} {
  \draw[BookRule,line width=.5pt] (-.3,\y)--(5.3,\y);
  \node[anchor=east] at (-.32,\y) {$\y$};
 }
 \node[anchor=east] at (-.32,0) {$0$};
 \foreach \x in {0,...,5} {
  \draw[BookMuted,line width=.5pt] (\x,0)--(\x,-.8);
  \node[below] at (\x,-.8) {$\x$};
 }
 \draw[BookBlue,dashed,line width=1pt] plot[mark=*,mark size=1.3pt,
  mark options={solid}] coordinates {(0,10)(1,0)(2,4)(3,8)(4,4)(5,20)};
 \draw[BookTeal,line width=1pt] plot[mark=*,mark size=1.3pt] coordinates {
  (0,10)(1,12)(2,18)(3,24)(4,22)(5,20)};
 \fill[BookMuted] (-.08,9.2) rectangle (.08,10.8);
 \fill[BookMuted] (4.92,19.2) rectangle (5.08,20.8);
 \draw[BookBlue,dashed,line width=1pt] (.2,-8.5)--(.6,-8.5);
 \node[anchor=west] at (.65,-8.5) {直接粘贴};
 \draw[BookTeal,line width=1pt] (2.2,-8.5)--(2.6,-8.5);
 \node[anchor=west] at (2.65,-8.5) {泊松融合};
 \fill[BookMuted] (4.05,-9.5) rectangle (4.25,-7.5);
 \node[anchor=west] at (4.3,-8.5) {边界值};
\end{tikzpicture}

 \caption{一维泊松融合的教学例。灰色方块是目标图像在边界像素0与5上的值，蓝色折线是直接粘贴源图像的结果，在两端各有一次跳变；青色折线是泊松融合的结果，两端与目标图像相接，内部保留了源图像的鼓包形状，每个相邻差都比源图像多2。数值为教学构造。}
 \label{fig:vision-poisson-1d}
\end{figure}

引导场不一定取源图像的梯度。把带孔洞或半透明的物体贴到有纹理的墙上时，可以在每对相邻像素上比较两幅图的变化，取较强的一个：$|f^*_p-f^*_q|>|g_p-g_q|$时取$v_{p,q}=f^*_p-f^*_q$，否则取$g_p-g_q$。原论文称之为混合梯度，它让墙面纹理透过字母之间的空隙保留下来。

泊松图像编辑成立的前提是源区域与目标区域在边界附近内容相容。若$\Omega$的边界横穿一条强边缘，例如桌沿，边界上的落差会沿着修正量扩散进花瓶，形成一片色斑；被贴物体的颜色总会被周围带偏，无法同时保持“边界无缝”和“颜色忠实于源图像”；它也不调整透视和光照方向，更不会为花瓶补上投在桌面的影子。“新内容在边界处与周围一致”这一要求在学习方法中仍然存在，只是不再显式解方程，而由修补模型读取可见区域、在训练中学会满足。

\subsubsection{Criminisi样例修补}
\label{sec:vision-generation-exemplar-inpainting}

泊松编辑需要一张现成的源图像。若请求只是“把杯子去掉”，杯子移除后留下的区域没有任何源图像可贴，只能用同一幅图中的桌面和背景补上。较早的修补方法沿着边界上的等照度线（亮度相等的曲线）把已知值扩散进来，适合细长的划痕，但填补大区域时会把纹理抹平；纹理合成能补出逼真的纹理，却不理会桌沿这样的直线结构\scite{vision-criminisi2004}。Criminisi等人的方法把纹理合成的采样方式与一个决定填充顺序的优先级结合起来，下面先看它所继承的纹理合成。

Efros与Leung把纹理看作马尔可夫随机场：给定一个像素周围$w\times w$窗口内的邻域，这个像素与图像其余部分无关\scite{vision-efros-leung1999}。窗口宽度$w$是唯一由用户设定的参数，应与纹理中最大的规则结构相当。合成一个像素$p$时，它的窗口$\omega(p)$中只有部分像素已知，算法在样例图中寻找最相似的窗口：距离$d$是两个窗口在已知像素上的平方差之和，用二维高斯核加权，使靠近中心的像素更重要，再除以已知像素的个数。设最小距离为$d_{\mathrm{best}}$，所有满足$d<(1+\epsilon)d_{\mathrm{best}}$的窗口组成候选集，原论文取$\epsilon=0.1$；从候选集中随机取一个窗口，用它的中心像素作为$p$的值。整幅纹理从一个随机取自样例的$3\times3$种子向外一层层生长，修补空洞时则从空洞边缘向内生长，见算法\ref{alg:vision-efros-leung}。

\begin{algorithm}[Efros–Leung非参数纹理合成]
 \label{alg:vision-efros-leung}
 \begin{algorithmic}[1]
  \Require 样例图$I_{\mathrm{smp}}$，待填像素集合，窗口宽度$w$，阈值系数$\epsilon=0.1$
  \Ensure 填好的图像
  \While{仍有待填像素}
   \State 取紧邻已知区域的一层待填像素
   \For{这一层中的每个像素$p$}
    \State 对样例图中每个窗口，按已知像素计算高斯加权、按已知像素数归一化的平方差$d$
    \State $d_{\mathrm{best}}\gets$最小的$d$；候选集$\gets\{$窗口$:d<(1+\epsilon)d_{\mathrm{best}}\}$
    \State 从候选集中随机取一个窗口，把它的中心像素值赋给$p$
   \EndFor
  \EndWhile
 \end{algorithmic}
\end{algorithm}

这一步为什么不会像回归那样取平均，可以用一维例子看清。设样例是一行像素$(0,1,1,0,1,1,0,1,2)$，窗口只看左边两个像素，要合成上下文为$(0,1)$之后的像素。样例中$(0,1)$出现在三处，后面分别跟着1、1、2，三处都完全匹配，$d_{\mathrm{best}}=0$，候选集就是这三处（$d_{\mathrm{best}}=0$时只保留完全匹配）。于是以$2/3$的概率取1、$1/3$的概率取2，每次取到的都是样例中真实出现过的值；若取三者的平均$4/3$，得到的是样例中从未出现的灰度。

逐像素合成很慢，且每个像素独立选择，大结构容易断裂。Efros与Freeman的图像缝合改为以$B\times B$的块为单位：相邻块之间重叠$B/6$的宽度，新块从样例中随机选取，要求它在重叠区的误差与最佳块相差不超过最佳误差的0.1倍；两块之间的接缝不取直线，而是沿重叠区误差最小的路径切开\scite{vision-quilting2001}。设重叠区第$i$行、第$j$列的平方差为$e_{i,j}$，累积误差按动态规划递推
\begin{equation*}
 E_{i,j}=e_{i,j}+\min\bigl(E_{i-1,j-1},E_{i-1,j},E_{i-1,j+1}\bigr),
\end{equation*}
最后一行$E$的最小值给出最小误差路径的终点，再逐行回溯。动态规划的一般原理可参考\BookChapterRef{chap:graph-optimization}{卷一“图”部分的“图上的约束优化”章}。以3行2列的重叠区为例，取$e=\left(\begin{smallmatrix}1&4\\3&1\\2&5\end{smallmatrix}\right)$，则第一行$E=(1,4)$，第二行$E=(3+1,1+1)=(4,2)$，第三行$E=(2+2,5+2)=(4,7)$；终点在第一列，回溯得到逐行的列号$(1,2,1)$，路径误差为$1+1+2=4$。

纹理合成把空洞当作一片纹理，按同心层的顺序填补，桌沿这样的直线结构在填到中间时往往错位。Criminisi等人指出，问题出在填充顺序：应当优先填补那些既有较多可靠邻居、又正好位于强边缘延长线上的位置。设待填区域为$\Omega$，已知的源区域为$\Phi$，二者之间的分界线称为填充前沿$\delta\Omega$。对前沿上的每个点$p$，以它为中心取一个块$\Psi_p$，原论文默认$9\times9$，实际应略大于纹理中最大的可辨元素。$p$的优先级定义为
\begin{equation}
 P(p)=C(p)D(p),\quad
 C(p)=\frac{\sum_{q\in\Psi_p\cap(I-\Omega)}C(q)}{|\Psi_p|},\quad
 D(p)=\frac{|\nabla I_p^{\perp}\cdot\bm{n}_p|}{\alpha}\text{。}
 \label{eq:vision-criminisi-priority}
\end{equation}
置信项$C(p)$是块内已知像素的置信度之和除以块的面积$|\Psi_p|$，初始化时$\Omega$内的像素置信度为0，其余为1；它偏向那些被已知像素包围较多的位置。数据项$D(p)$中，$\nabla I_p^{\perp}$是与梯度垂直、沿等照度线方向的向量，长度等于梯度幅值，$\bm{n}_p$是前沿在$p$处的单位法向，$\alpha$是归一化常数，灰度图取255；等照度线越强、越垂直地撞上前沿，$D(p)$越大，于是边缘被优先沿原方向延伸进空洞。选出优先级最高的$\hat p$以后，在源区域中找与$\Psi_{\hat p}$最相似的块$\Psi_{\hat q}$，距离是两块已知像素在CIE Lab颜色空间中的平方差之和；把$\Psi_{\hat q}$中对应位置的像素复制到$\Psi_{\hat p}$的未知像素上，并把这些新填像素的置信度设为$C(\hat p)$，见算法\ref{alg:vision-criminisi}。

\begin{algorithm}[Criminisi样例修补]
 \label{alg:vision-criminisi}
 \begin{algorithmic}[1]
  \Require 图像$I$，待填区域$\Omega$，块大小（默认$9\times9$）
  \Ensure 填好的图像
  \State 对$p\in\Omega$令$C(p)\gets0$，对其余像素令$C(p)\gets1$
  \While{$\Omega\neq\varnothing$}
   \State 找出填充前沿$\delta\Omega$，按式\eqref{eq:vision-criminisi-priority}计算前沿上每点的$P(p)$
   \State $\hat p\gets\arg\max_{p\in\delta\Omega}P(p)$
   \State 在源区域中找$\Psi_{\hat q}=\arg\min_{\Psi_q\subset\Phi}d(\Psi_{\hat p},\Psi_q)$，$d$只在$\Psi_{\hat p}$的已知像素上计算
   \State 把$\Psi_{\hat q}$的对应像素复制到$\Psi_{\hat p}\cap\Omega$
   \State 对新填像素$q$令$C(q)\gets C(\hat p)$，并把它们移出$\Omega$
  \EndWhile
 \end{algorithmic}
\end{algorithm}

用两个前沿点比较优先级。点$a$位于桌沿的延长线上，块内81个像素中已知54个，$C(a)=54/81\approx0.667$；桌沿的梯度幅值为120，等照度线垂直撞上前沿，$D(a)=120/255\approx0.471$，$P(a)\approx0.314$。点$b$位于平坦的桌面上，块内已知63个，$C(b)\approx0.778$；梯度幅值只有10，$D(b)\approx0.039$，$P(b)\approx0.030$。尽管$a$周围已知像素更少，它的优先级仍高出约10倍，桌沿会先被延伸进空洞。这些数值是教学构造。

这种贪心的逐块复制有明确的边界。它只能使用图像中已有的块，空洞里若应出现一个图中别处没有的物体，例如被杯子挡住的狗爪，就无法补出；某一步复制错误后，新填像素会成为后续匹配的依据，错误会沿前沿继续传播；块太小抓不住纹理元素，块太大又难以找到匹配。每次填块都要在整个源区域中搜索最相似的块，这一步的代价由下面的PatchMatch处理。

\subsubsection{PatchMatch}
\label{sec:vision-generation-patchmatch}

纹理合成、图像缝合和样例修补都反复执行同一个操作：对一个区域中的每个块，在另一个区域中找最相似的块。Barnes等人把这种对应关系称为\term{最近邻域场}（nearest-neighbor field, NNF）：对图像$A$中每个块的坐标$\bm{a}$，记它在图像$B$中最相似块的坐标为$\bm{b}$，则$f(\bm{a})=\bm{b}-\bm{a}$是一个偏移量，所有偏移量排成一张与$A$同尺寸的二维数组\scite{vision-patchmatch2009}。若区域含$M$个块位置、每块$m$个像素，穷举搜索的代价是$O(mM^2)$。百万像素的图像用$7\times7$的块，约为$49\times10^{12}$次像素运算，远不能交互使用；kd树等近似最近邻结构在高维的块空间中仍然慢，还要占用较多内存。

PatchMatch依据三点观察。其一，块空间的维数很高，但偏移量只是二维的，在偏移空间中搜索更省时省内存。其二，自然图像是连贯的：若某个块在$B$中的最佳匹配偏移了$\bm{v}$，它右边的块很可能也偏移$\bm{v}$，所以好的偏移可以传给邻居。其三，在大量随机猜测中，总有一部分恰好猜对。设$A$与$B$都含$M$个块位置，随机初始化时某个指定位置猜对的概率只有$1/M$，但至少有一个位置猜对的概率是$1-(1-1/M)^M$，$M$很大时趋于$1-1/e\approx0.63$；若把落在正确偏移附近$C$个像素内的猜测也算作好的猜测，这个概率是$1-(1-C/M)^M\approx1-e^{-C}$，$C=5$时约为0.993。只要连贯区域中有一个位置猜对，传播就能把它扩散到整个区域。

算法先给每个块随机分配一个偏移，然后反复迭代，见算法\ref{alg:vision-patchmatch}。每次迭代按扫描顺序（从左到右、从上到下）访问每个块，对它依次做传播和随机搜索。传播时，用$D(\bm{v})$表示$A$中$(x,y)$处的块与$B$中$(x,y)+\bm{v}$处的块之间的距离，在$\{f(x,y),f(x-1,y),f(x,y-1)\}$中取使$D$最小的偏移；偶数次迭代按逆扫描顺序进行，改用$f(x+1,y)$与$f(x,y+1)$，使好的偏移也能向左、向上传播。随机搜索在当前偏移$\bm{v}_0=f(x,y)$周围按指数缩小的范围试探
\begin{equation}
 \bm{u}_i=\bm{v}_0+w\alpha^i\bm{R}_i,\qquad i=0,1,2,\ldots
 \label{eq:vision-patchmatch-search}
\end{equation}
其中$\bm{R}_i$在$[-1,1]\times[-1,1]$中均匀随机，$w$是最大搜索半径，原论文取图像的最大边长，$\alpha=1/2$；半径$w\alpha^i$小于1个像素时停止，候选落到$B$外时截到边界上，凡是更好的候选都替换当前偏移。原论文的结果都只迭代4到5次，此时偏移场几乎不再变化。

\begin{algorithm}[PatchMatch近似最近邻域场]
 \label{alg:vision-patchmatch}
 \begin{algorithmic}[1]
  \Require 图像$A,B$，块距离$D$，迭代次数$T$（4到5），搜索半径$w$，缩小比例$\alpha=1/2$
  \Ensure 偏移场$f$
  \State 对$A$中每个块位置，从$B$的范围内均匀随机抽取$f(x,y)$
  \For{$t=1,\ldots,T$}
   \If{$t$为奇数}
    \State 扫描顺序为从左到右、从上到下，邻居集合$S(x,y)=\{(x-1,y),(x,y-1)\}$
   \Else
    \State 扫描顺序反向，$S(x,y)=\{(x+1,y),(x,y+1)\}$
   \EndIf
   \For{按扫描顺序取每个块位置$(x,y)$}
    \State 传播：$f(x,y)\gets\arg\min_{\bm{v}\in\{f(x,y)\}\cup\{f(\bm{s}):\bm{s}\in S(x,y)\}}D(\bm{v})$
    \State 随机搜索：对$i=0,1,\ldots$直到$w\alpha^i<1$，按式\eqref{eq:vision-patchmatch-search}生成$\bm{u}_i$，若$D(\bm{u}_i)<D(f(x,y))$则$f(x,y)\gets\bm{u}_i$
   \EndFor
  \EndFor
  \State \Return $f$
 \end{algorithmic}
\end{algorithm}

用一维例子手算传播。设$A$中位置0到3的四个块，在$B$中的真实匹配都在10到13，即正确偏移都是$+10$；随机初始化得到$(+3,+10,-1,+5)$，只有位置1碰巧猜对。第一次迭代从左到右：位置0没有左邻，假设随机搜索没有命中；位置1已正确；位置2比较自身的$-1$与左邻的$+10$，后者距离为0，取$+10$；位置3同理取$+10$。第二次迭代从右到左：位置3、2、1不变，位置0比较自身的$+3$与右邻的$+10$，取$+10$。两次迭代后全部正确，而初始时只有一个位置猜对。随机搜索的候选数也可以算出：若$w=1000$，半径依次为$1000,500,\ldots,1000/2^9\approx1.95$，下一个$1000/2^{10}<1$，所以每个块每次迭代试探10个候选。图\ref{fig:vision-patchmatch-steps}画出了初始化和两次传播后的偏移场。

\begin{figure}[htbp]
 \centering
 % PatchMatch一维教学例：A中位置0-3的正确偏移均为+10（匹配B中10-13）。
% 初始化(+3,+10,-1,+5)；第一次正向传播后(+3,+10,+10,+10)；第二次逆向传播后全部为+10。
% 三个面板上下排列，每格5mm，宽约107mm、高约108mm。
\begin{tikzpicture}[x=5mm,y=1mm,font=\figurefont\small,
 cellA/.style={draw=BookBlue,fill=BookBlue!10,line width=.6pt},
 cellB/.style={draw=BookMuted,fill=white,line width=.5pt},
 good/.style={-{Stealth[length=1.6mm]},BookTeal,line width=.9pt},
 bad/.style={-{Stealth[length=1.6mm]},BookBlue,dashed,line width=.8pt},
 prop/.style={-{Stealth[length=1.4mm]},BookAmber,line width=.8pt}]
 \foreach \panel/\ytop/\title/\offs in {
   0/0/{初始化}/{3,10,-1,5},
   1/-34/{第一次迭代：正向传播}/{3,10,10,10},
   2/-68/{第二次迭代：逆向传播}/{10,10,10,10}} {
  \node[anchor=west] at (-4.6,\ytop+5) {\title};
  \node[anchor=east] at (-.2,\ytop-3) {$A$};
  \node[anchor=east] at (-.2,\ytop-19) {$B$};
  \foreach \i in {0,...,3} {
   \draw[cellA] (\i,\ytop-6) rectangle (\i+1,\ytop);
   \node at (\i+.5,\ytop-3) {\footnotesize\i};
  }
  \foreach \j in {0,...,15} {
   \draw[cellB] (\j,\ytop-22) rectangle (\j+1,\ytop-16);
   \node[text=BookMuted] at (\j+.5,\ytop-19) {\scriptsize\j};
  }
  \foreach \o [count=\k from 0] in \offs {
   \pgfmathtruncatemacro{\tgt}{\k+\o}
   \ifnum\o=10
    \draw[good] (\k+.5,\ytop-6.3) -- (\tgt+.5,\ytop-15.7);
   \else
    \draw[bad] (\k+.5,\ytop-6.3) -- (\tgt+.5,\ytop-15.7);
   \fi
  }
 }
 \draw[prop] (1.6,-34+0.8) to[bend left=50] (2.4,-34+0.8);
 \draw[prop] (2.6,-34+0.8) to[bend left=50] (3.4,-34+0.8);
 \draw[prop] (1.4,-68+0.8) to[bend right=50] (0.6,-68+0.8);
 \draw[good] (3,-96)--(4,-96); \node[anchor=west] at (4.1,-96) {正确偏移};
 \draw[bad] (7.5,-96)--(8.5,-96); \node[anchor=west] at (8.6,-96) {错误偏移};
 \draw[prop] (12,-96) to[bend left=40] (13,-96); \node[anchor=west] at (13.1,-96) {传播方向};
\end{tikzpicture}

 \caption{PatchMatch传播过程的一维教学例。每个面板上排是$A$中位置0至3的块，下排是$B$中的位置，箭头从$A$的块指向当前偏移所指的$B$中位置，实线为正确偏移$+10$，虚线为错误偏移；弯箭头表示本次迭代中好偏移的传播方向。初始化时只有位置1碰巧正确；第一次迭代按扫描顺序把$+10$传给位置2和3，第二次迭代逆向扫描再传给位置0。图中省略了随机搜索。数值为教学构造。}
 \label{fig:vision-patchmatch-steps}
\end{figure}

在原论文的修补中，空洞记为$T$，图像其余部分记为$S$，修补按类似期望最大化的方式交替进行：先用PatchMatch为$T$中每个块在$S$中找最近邻，每个匹配块把自己的像素“投给”它覆盖的位置；再让每个像素取所得投票的平均值，得到新的$T$。为避免陷入差的局部解，计算从低分辨率开始，逐级放大并在每级继续迭代。论文在最大约200万像素的测试图像上验证了收敛性，并报告CPU实现比带PCA降维的kd树快20到100倍，典型参数下的时间复杂度为$O(mM\log M)$，额外内存为$O(M)$，可以交互使用。

PatchMatch只加速了搜索，没有改变样例方法的根本局限：新内容仍只能来自已有的块。它的偏移场也只按块距离优化，不包含“这是一只杯子”之类的语义，空洞里应出现的内容若在图中找不到相似的块，结果只能是貌似合理的纹理拼贴。学习方法中的注意力也计算一个位置与所有位置的相似度，再按相似度从其他位置取值，可以看成一种可学习的软性最近邻；而要生成原图中不存在的花瓶，则需要下面以文字为条件的生成模型。

\subsubsection{GLIDE修补模型}
\label{sec:vision-generation-glide-inpainting}

最简单的修补做法是：用普通文字生成模型采样，在每一步把掩码外的像素替换为原图相应噪声水平的像素。这样能强制保留已知区域，但模型在训练中从未见过“部分像素已知”的输入，它不知道已知区域里有什么，生成的花瓶常与周围光照和边界不协调。GLIDE修补模型的动机是让模型在每一步都能直接看到清晰的保留区域，从而学会利用上下文\scite{vision-glide2022}。

主要想法是把可见图像和掩码作为额外输入通道，对\ref{sec:vision-models-glide}节的文字条件GLIDE作修补专用的微调。模型在每个去噪步同时读到当前带噪图像、清晰的保留区域和掩码，因此可以根据桌面纹理、光照方向来生成缺失区域。

架构上，GLIDE基础模型在$64\times64$的像素空间工作。修补模型的输入是三部分沿通道拼接：当前带噪图像（3通道）、可见图像$\bm{V}=\bm{M}\odot\bm{X}$（3通道）和掩码$\bm{M}$（1通道），共$64\times64\times7$。只有第一层卷积的输入通道从3增加到7，新增4个通道对应的权重初始化为零，其余权重从文字生成检查点继承，所以微调开始时模型的行为与原模型相同。文字仍经Transformer编码后接入去噪网络。上采样模型把分辨率提到$256\times256$，它读取完整的低分辨率图像，以及高分辨率下的未遮挡区域。掩码通道不可省略：可见图像中被遮住的位置填0，而真实图像中也可能有黑色像素，只有掩码能区分二者。

各组件的规模见表\ref{tab:vision-glide-inpaint-config}。基础模型沿用Dhariwal与Nichol的ImageNet $64\times64$扩散模型（ADM）的结构\scite{vision-adm2021}，只把宽度扩大到512通道，视觉部分约23亿参数；文字先编码为$K$个词元，送入一个24个残差块、宽2048的Transformer，约12亿参数。这个Transformer的输出有两种用法：最后一个词元的向量代替ADM中的类别嵌入，最后一层的$K$个词元向量则分别投影到ADM每个注意力层的维数，拼接到该层的注意力上下文中，让图像特征可以注意到文字。GLIDE论文没有逐层列出ADM的结构，各分辨率的残差块数与注意力位置以ADM论文为准，本章不再展开。

\begin{table}[htbp]
 \centering
 \small
 \begin{tabularx}{\linewidth}{@{}p{24mm}Xp{20mm}@{}}
  \toprule
  组件 & 结构与输入输出 & 规模\\
  \midrule
  文字编码器 & Transformer，24个残差块，宽2048；输出$K$个词元向量 & 约12亿参数\\
  基础扩散模型 & ADM结构，宽512通道；$64\times64$像素空间 & 视觉部分约23亿参数\\
  修补模型的输入层 & 第一层卷积由3通道扩为7通道，新增4通道的权重初始化为0 & —\\
  上采样扩散模型 & $64\to256$，ADM上采样结构，基础通道384；文字编码器宽1024 & 约15亿参数\\
  训练 & 基础模型250万次迭代、批大小2048；上采样模型160万次迭代、批大小512；再以20\%空文字微调以支持无分类器引导 & —\\
  修补微调 & 随机擦除区域，以可见图像与掩码为附加输入；完整模型的修补微调步数论文未给出 & —\\
  采样 & 修补样例用100个扩散步，其余展示样例用150步 & —\\
  \bottomrule
 \end{tabularx}
 \caption{GLIDE及其修补模型的配置，数值取自GLIDE论文。基础模型与文字编码器合计约35亿参数；论文公开的小规模过滤版本在110万次预训练之后，另用50万次迭代同时微调无分类器引导与修补。}
 \label{tab:vision-glide-inpaint-config}
\end{table}

训练数据就是原有的图文对，无需额外标注。每个样本随机擦去图像的一些区域得到$\bm{M}$，构造$\bm{V}$，再对完整图像加噪，以完整图像的噪声预测误差微调模型。由于每次抽到的掩码不同，模型学会在各种可见范围下补全。推理时固定用户给出的掩码、原图和文字“花瓶”，从噪声出发反复去噪；可见条件在整个采样中保持不变，只有带噪图像随步更新。

这里要分清训练目标和使用目标。训练时，模型补全的是被擦去的真实内容；使用时，文字要求的花瓶原图中并不存在，这属于创作，模型借助的是训练中学到的“依上下文补全”的能力。若只是要恢复被遮挡的杯子，那是\ref{sec:vision-restoration-inpainting}节的复原目标。掩码条件不保证掩码外的每个像素都不变，因为生成和上采样都在全图上进行。需要严格保持时，可以按掩码把原图回填，这是一个额外的后处理步骤，回填后要重新检查边界。

\subsubsection{InstructPix2Pix}
\label{sec:vision-generation-instruct-pix2pix}

用户常常只给一句“把杯子换成花瓶”，不画掩码，也不写完整的新场景描述。这时需要一个读入原图和编辑指令、直接输出编辑结果的模型。训练这样的模型需要大量“原图、指令、编辑后图像”三元组，而人工制作这样的配对几乎不可行。InstructPix2Pix的动机就是自动合成这些训练数据，再用它们训练一个指令编辑模型\scite{vision-instruct-pix2pix2023}。

主要想法分两步。第一步合成数据：用语言模型从图像描述生成编辑指令和编辑后的描述，再用文字生成模型把编辑前后两段描述分别生成图像。第二步训练一个以原图和指令为条件的扩散模型。合成数据的关键困难是，独立生成的两幅图会在姿态、背景上全面不同，不是一对编辑前后的图像。InstructPix2Pix用Prompt-to-Prompt方法解决：两幅图使用相同的初始噪声和相同的采样噪声，并在前一部分去噪步中让第二幅图沿用第一幅图的注意力权重，使两幅图的布局保持一致，只在描述不同的地方变化\scite{vision-prompt2prompt2022}。Prompt-to-Prompt原本按编辑类型替换交叉注意力；InstructPix2Pix的附录说明，它改为对所有编辑统一替换第二幅图的自注意力权重。

数据流程的细节如下。先人工撰写700条三元组（输入描述、编辑指令、输出描述），用默认训练参数把GPT-3微调一轮；再让微调后的GPT-3为大量图像描述生成指令和新描述。对每对描述，用Stable Diffusion配合Prompt-to-Prompt生成100对图像，共享注意力的步数比例在0.1到0.9之间随机抽取；然后用CLIP相似度筛选，要求两图的CLIP图像相似度超过0.75、图文相似度超过0.2，以及“图像变化方向与描述变化方向”的方向相似度超过0.2，通过的配对按方向相似度排序，每对描述最多保留4对。生成图像用Stable Diffusion v1.5的指数滑动平均权重、Euler祖先采样器和100个去噪步。最终得到超过45万个训练样本。共享注意力减少了无关变化，筛选去掉了明显失败的配对，但并不保证每条指令都被正确执行。

模型以Stable Diffusion v1.5初始化。原图经冻结的自编码器编码为潜表示$\bm{z}_I$；训练时目标图像编码后加噪得到$\bm{z}^{(t)}$；两者都是$h'\times w'\times4$，沿通道拼接为8通道输入去噪U-Net的第一层，新增4个通道的权重初始化为零。指令$c_T$经文字编码器和交叉注意力进入网络，不拼进像素通道。训练目标为
\begin{equation}
 L=\mathbb{E}\!\left[
 \left\lVert\bm{\epsilon}-
 \epsilon_{\bm{\theta}}
   \bigl(\bm{z}^{(t)},t,\bm{z}_I,c_T\bigr)
 \right\rVert_2^2\right]\text{，}
 \label{eq:vision-instruct-training}
\end{equation}
其中$\bm{\epsilon}$是加到目标潜表示上的噪声。原图条件在整个训练和使用中都保持干净，只有目标在加噪。为支持下面的双条件引导，训练时以5\%的概率只丢弃图像条件、5\%只丢弃文字条件、5\%同时丢弃二者。训练配置见表\ref{tab:vision-instruct-config}。论文说明全部可用权重都从预训练检查点初始化，其余训练设置沿用公开的Stable Diffusion代码库，但没有单独说明自编码器和文字编码器是否参与更新。

\begin{table}[htbp]
 \centering
 \small
 \begin{tabularx}{\linewidth}{@{}p{24mm}Xp{20mm}@{}}
  \toprule
  组件 & 结构与输入输出 & 规模\\
  \midrule
  去噪U-Net & Stable Diffusion v1.5的U-Net（逐层结构见表\ref{tab:vision-controlnet-structure}），从其指数滑动平均权重初始化 & 论文未给出\\
  输入层 & 第一层卷积读取带噪目标潜表示与原图潜表示，共8通道；新增4通道的权重初始化为0 & —\\
  条件接入 & 指令经原有的文字条件机制（交叉注意力）进入；原图只经输入通道拼接进入 & —\\
  训练数据 & 超过45万个由GPT-3与Stable Diffusion合成的“原图、指令、编辑后图像”样本 & —\\
  训练 & $256\times256$，1万步，总批大小1024，学习率$10^{-4}$且无预热；随机水平翻转，随机缩放到256至288再裁回256；8张40GB A100约25.5小时 & —\\
  推理 & $512\times512$，Euler祖先采样器100步，每步三次前向；单张A100约9秒一幅 & —\\
  \bottomrule
 \end{tabularx}
 \caption{InstructPix2Pix的配置，数值取自原论文正文与附录。模型在$256\times256$上训练，论文报告它在推理时能推广到$512\times512$。}
 \label{tab:vision-instruct-config}
\end{table}

推理时输入原图和指令，从随机噪声开始采样，不需要目标图像。每步计算三种噪声预测：$\bm{\epsilon}_0=\epsilon_{\bm{\theta}}(\bm{z}^{(t)},t,\varnothing,\varnothing)$、$\bm{\epsilon}_I=\epsilon_{\bm{\theta}}(\bm{z}^{(t)},t,\bm{z}_I,\varnothing)$和$\bm{\epsilon}_{I,T}=\epsilon_{\bm{\theta}}(\bm{z}^{(t)},t,\bm{z}_I,c_T)$，再组合为
\begin{equation}
 \hat{\bm{\epsilon}}
 =\bm{\epsilon}_0
  +s_I(\bm{\epsilon}_I-\bm{\epsilon}_0)
  +s_T(\bm{\epsilon}_{I,T}-\bm{\epsilon}_I)\text{。}
 \label{eq:vision-instruct-guidance}
\end{equation}
这是无分类器引导在两个条件上的推广：第二项强调原图，第三项在已有原图条件下强调指令带来的变化。以两个分量示意，令$\bm{\epsilon}_0=(0.1,0.2)$、$\bm{\epsilon}_I=(0.2,0.1)$、$\bm{\epsilon}_{I,T}=(0.4,-0.1)$，取$s_I=1.5,s_T=5$，得$(1.25,-0.95)$；两项强度都取1时恢复$\bm{\epsilon}_{I,T}$。原论文指出$s_T$在5到10、$s_I$在1到1.5之间通常效果较好，并需按例调整：$s_I$越大越接近原图，$s_T$越大修改越强。每步三次前向使推理成本约为普通采样的三倍。

由于没有显式掩码，未提及的区域也可能改变，例如狗的颜色随杯子一起变化。数量、空间关系类的指令以及需要大幅改变几何的指令容易失败，这部分错误继承自合成数据：生成数据的文字生成模型和筛选器本身就不擅长这些情形。

五种单轮编辑方法的比较见表\ref{tab:vision-generation-editing-comparison}。固定任务为“把狗旁边的杯子换成花瓶”，评价口径为替换是否完成、杯子以外是否保持，以及新内容与周围是否协调。三种经典方法都不训练模型，新内容来自现成图像：泊松图像编辑需要一张花瓶照片，在保留区严格不变的前提下用梯度域求解消除接缝，代价是花瓶颜色被带偏；Criminisi样例修补与PatchMatch只能完成“去掉杯子”这一半，用图中已有的桌面补洞，前者规定填充顺序，后者把块搜索加速到可交互，二者都补不出图中没有的花瓶。GLIDE修补模型把新内容的来源换成以文字为条件的生成模型，掩码仍由用户给出，可见区域作为输入让模型学会与周围协调，代价是需要大规模图文数据训练，且掩码外不保证逐像素不变。InstructPix2Pix相对GLIDE修补模型又改变了两处设计：范围由用户掩码改为模型从指令中推断，条件接入由“可见图像加掩码”改为“完整原图潜表示”。它换来了只需一句话的使用方式，代价是失去对保留区域的显式约束，并依赖合成数据的质量。

\begin{table*}[tbp]
 \centering
 \small
 \begin{tabularx}{\linewidth}{@{}p{26mm}p{30mm}XXX@{}}
  \toprule
  方法 & 用户输入 & 新内容的来源 & 范围外的保持 & 典型失败\\
  \midrule
  泊松图像编辑 & 原图、区域、源图像 & 源图像的梯度，经求解泊松方程融合 & 区域外严格不变 & 边界横穿强边缘时出现色斑，被贴物体颜色偏移，不处理透视与阴影\\
  Criminisi样例修补 & 原图、待填区域 & 同一图中别处的图像块，按优先级逐块复制 & 区域外严格不变 & 图中没有的内容补不出，复制错误沿前沿传播\\
  PatchMatch & 原图、待填区域 & 同一图中的图像块，由随机化最近邻搜索与投票合成 & 区域外严格不变 & 与样例修补相同；只加速搜索，不含语义\\
  GLIDE修补模型 & 原图、掩码、对新内容的描述 & 文字条件扩散模型，读取可见区域与掩码 & 由掩码约束，可回填强制保持 & 掩码过小留下旧阴影、边界不协调\\
  InstructPix2Pix & 原图、编辑指令 & 以原图潜表示和指令为条件的扩散模型 & 无显式约束，靠$s_I$调节 & 未提及区域改变，数量与几何指令失败\\
  \bottomrule
 \end{tabularx}
 \caption{五种单轮编辑方法在同一替换任务下的比较。经典方法不需要训练，每次编辑的计算量随区域大小增长；两种学习方法需要预先在大规模数据上训练，单次编辑的代价是一次完整的扩散采样，InstructPix2Pix每步还需三次前向。}
 \label{tab:vision-generation-editing-comparison}
\end{table*}

\subsection{多轮编辑的状态维护}
\label{sec:vision-generation-kontext-history}

一次替换之后，用户还可能继续要求“把花瓶改成蓝色”“再把它移到狗的右边”。每一轮都有新的变化，同时已经接受的修改和始终不应变化的主体都要保持。若只把当前图像送入下一轮，模型能看到上一轮的结果，却无法恢复已被前几轮采样削弱的狗的斑纹；若每轮都从最初的原图开始，又会丢掉已接受的花瓶。因此多轮编辑需要维护一份状态：原始依据、每轮的已接受版本以及修改记录，并写明每一项是保持还是被本轮覆盖。这份状态由模型的输入能力和外部的版本管理共同完成，二者应分开描述。

模型一侧，需要一个能把已有图像作为上下文、按指令生成新图的接口，而且最好能在多轮中保持主体。FLUX.1 Kontext是这样的模型\scite{vision-kontext-2025}。它的动机是：InstructPix2Pix沿通道拼接原图和目标，要求二者尺寸相同，并需要为编辑单独训练；而Kontext希望用同一个模型完成文字生成和图像编辑，让条件图与目标图的尺寸可以不同。主要想法是沿序列方向拼接：上下文图像经冻结自编码器编码并分块后成为$n_y$个词元，待生成图像成为$n_x$个词元，二者拼成长度$n_x+n_y$的序列送入以流匹配训练的Transformer，文字词元经另一条流与之交换注意力。为区分两幅图中相同空间坐标的词元，三维旋转位置编码给目标词元赋$(0,h,w)$，给上下文词元赋$(1,h,w)$，首坐标只是图像来源编号。

训练从FLUX.1文字生成检查点开始，样本为“上下文图像$\bm{Y}$、指令$c$、目标图像$\bm{X}$”，纯文字生成样本令$\bm{Y}$为空。设$\bm{x}=E(\bm{X})$为目标潜表示，抽标准高斯噪声$\bm{\epsilon}$与$t\in(0,1)$，目标为
\begin{equation}
 \begin{aligned}
 \bm{z}(t)&=(1-t)\bm{x}+t\bm{\epsilon},\\
 L_{\mathrm{ctx}}&=
 \mathbb{E}\!\left[
 \left\lVert v_{\bm{\theta}}
  (\bm{z}(t),t,E(\bm{Y}),c)
   -(\bm{\epsilon}-\bm{x})\right\rVert_2^2
 \right]\text{。}
 \end{aligned}
 \label{eq:vision-kontext-flow}
\end{equation}
误差只计在目标词元上，上下文词元保持干净、不作为输出。这里$t=0$是干净目标、$t=1$是噪声，与\ref{sec:vision-models-flow-matching}节的方向相反，所以采样沿$t$递减积分。自编码器冻结，更新的是生成主干。论文没有公开全部关系数据构造与训练预算；公开的dev版本还经过引导蒸馏，推理步数与完整模型不同。

Kontext的上下文保持干净，并不意味着输出在对应区域逐像素不变：编码、全图生成和解码都会改变细节，模型也不会因为序列拼接而自动保存编辑历史。论文的形式可以扩展到多幅上下文图像，但所发布的版本主要以单幅上下文训练。所以外部版本管理不可省略：保存原图$\bm{X}^{(0)}$，接受杯子替换得到$\bm{X}^{(1)}$，花瓶变蓝得到$\bm{X}^{(2)}$，移动花瓶得到$\bm{X}^{(3)}$。每轮以当前已接受的图像作为单幅上下文，同时记录模型、提示、种子、尺寸和候选数。第三轮要同时检查花瓶的位置、是否仍为蓝色、狗的耳朵和斑纹、未要求改变的背景。若第三轮颜色退回，应回到$\bm{X}^{(2)}$重做，而不是把新结果当成下一轮的正确起点。

历史要求也有不同类型。ImgEdit把多轮编辑的评价分成三类：早期设定的全局要求在后续轮次中仍然有效，例如“所有结果保持木质纹理”；后一轮请求依赖前文指代，例如“把它改成蓝色”中的“它”；以及回到某个历史版本的操作\scite{vision-imgedit-2025}。这三类要求都依赖保存的状态，而不是单轮的“原图—结果”配对。随着轮数增加，编码、采样和解码误差会累积，即使某轮没有要求改动主体，主体也可能漂移；Kontext论文同样展示了多轮后出现伪影和指令遗漏的失败例。外部回退可以阻止错误继续传播，却补不回原始参考从未提供的细节。

\subsection{编辑结果的核验}
\label{sec:vision-generation-editing-check}

核验时应把原图、编辑结果、要求记录和允许变化的范围放在一起。杯子仍在，是目标未修改；花瓶出现在另一侧，是改错位置；狗的斑纹变化，是保留失败；花瓶与桌面之间出现断边，是边界问题。把这些错误合成一个“编辑质量”分数，就无法判断应该修改流程中的哪一步。

对于明确的保留掩码$\bm{M}$，用式\eqref{eq:vision-edit-preservation}计算保留区的平均变化。若直接返回原图，$E_{\mathrm{keep}}=0$，但杯子没有变成花瓶，所以这一项必须与修改完成度一起报告。若请求允许全图色调变化，逐像素的保留误差就要换成与要求相符的判断。泊松图像编辑是一个需要特别说明的例子：它在保留区内严格不变，$E_{\mathrm{keep}}=0$，但粘贴进来的花瓶整体颜色会被调整，这一变化要按“新内容是否保持源物体外观”另行检查。

为强制保持范围外的内容，可以合成$\bm{X}_{\mathrm{out}}=\bm{M}\odot\bm{X}+(1-\bm{M})\odot\bm{X}'$。它保证保留像素严格相等，却可能截断新花瓶的阴影；若使用羽化掩码，就要说明过渡带也会改变。多轮编辑时，每轮都应与最初的原图和保留要求对照，而不只与上一轮比较，否则身份漂移会逐轮积累而不被发现。缩小范围、调整条件强度和重新生成都要记录轮数与候选预算，最终成功的结果应连同得到它的成本一起报告。

\section{观测场景的新视角合成}
\label{sec:vision-generation-novel-view}

回到开篇的第四个请求：根据绕玩具狗拍摄的一圈照片，给出从正后方看的样子。\term{新视角合成}（novel view synthesis）根据已拍摄的视图，渲染同一场景在未拍摄相机位置下的图像。它与前面三类任务的根本区别在于，结果应由真实观测支持，而不是由生成先验补充。输入是多幅照片及其相机参数，输出是任意给定相机下的图像，难点在于从有限的二维观测中拟合一个能够解释所有视图的三维表示。本节按“组织相机与观测—拟合场景表示—渲染新视角并核验”的步骤展开。只有一幅图、需要靠先验补出背面的情形，属于\ref{chap:vision-multimodal-generation}章的多视角生成。

\subsection{相机与观测的组织}
\label{sec:vision-generation-cameras}

拟合场景之前，每幅照片都要有已知的相机参数。这里用到\ref{chap:vision-image-formation}章建立的针孔相机模型：内参矩阵$\bm{K}$（焦距、主点）决定像素与相机坐标系中方向的关系，外参即旋转$\bm{R}$与平移$\bm{t}$决定相机在世界中的位置和朝向，世界点$\bm{x}$投影到像素的齐次坐标为$\bm{K}(\bm{R}\bm{x}+\bm{t})$，再除以第三个分量（\ref{sec:vision-formation-camera-intrinsics}节）。把这个过程反过来，像素$(u,v)$对应一条从相机中心$\bm{o}=-\bm{R}^{\top}\bm{t}$出发、方向为$\bm{d}\propto\bm{R}^{\top}\bm{K}^{-1}(u,v,1)^{\top}$的射线，这是后面所有渲染计算的出发点。例如$\bm{K}$的焦距为500像素、主点为$(320,240)$，相机位于世界原点且不旋转，则像素$(570,240)$对应方向$(0.5,0,1)$，即向右偏离光轴约26.6度。

相机参数通常由运动恢复结构从照片本身估计（\ref{sec:vision-formation-multi-view}节）。它先在照片之间匹配局部特征，由两视图的对极几何得到相对位姿（\ref{sec:vision-formation-two-view}节），再逐步加入更多照片，用三角化得到三维点，最后用光束法平差同时调整所有相机和点，使三维点在各照片上的重投影误差最小。输出是每幅照片的$\bm{K},\bm{R},\bm{t}$和一组稀疏的三维点。NeRF对真实场景使用COLMAP运动恢复结构软件估计相机\scite{vision-nerf2020}，3D Gaussian Splatting还用它给出的稀疏点初始化场景。相机也可以用\ref{chap:vision-image-recognition}章“相机与场景几何的估计”中的前馈几何模型估计。多视图几何本身只给出相机和稀疏点；要从任意新相机渲染出逼真图像，还需要一个能表示所有可见表面外观的场景表示，这是下一小节的问题。

拍摄的覆盖范围决定哪些表面有多视角证据。绕玩具狗一圈、在两个高度上各拍若干张，可以覆盖大部分表面；若只拍了正面，背面在任何拟合中都没有观测支持。拍摄前还应留出一些相机作为测试视图，它们不参与拟合，用于后面的核验。原始NeRF和3D Gaussian Splatting都要求场景静止、曝光一致、相机准确：玩具被碰动、光照变化或标定误差，都会让同一空间点需要解释互相矛盾的像素，拟合过程可能用错误的密度或外观吸收这些矛盾，训练视图看起来变好，几何却更不准确。

\subsection{场景表示的拟合}
\label{sec:vision-generation-scene-fitting}

有了带相机的照片，需要决定用什么表示场景，才能从新相机渲染出图像。基于图像的渲染给出了最早的一种回答：不重建几何，直接把照片看成对光线的采样。在没有遮挡物的自由空间中，沿一条射线传播的光强不变，所以可以用射线与两个平行平面的交点$(u,v)$和$(s,t)$来索引它，全部射线构成一个四维函数，称为光场；新视角的每个像素对应一条射线，只需在已拍到的射线中查找并插值\scite{vision-lightfield1996}。这种做法不需要估计任何深度，但只有在相机排布得足够密集、相邻射线足够接近时，插值才可靠；只拍了一圈几十张照片时，大多数新射线离已拍射线很远，插值结果会出现重影。另一条经典路线先用多视图立体重建出网格，再把照片作为纹理贴上去，对表面清晰的物体很有效，但毛绒、半透明和细小结构难以得到准确的网格。

学习路线把场景表示成一组可优化的参数，要求渲染过程可微，从而用“渲染图与照片的差异”直接优化表示，不必先得到准确的表面。两种代表性的路线对“场景用什么表示、怎样渲染”给出了不同回答。NeRF用一个神经网络隐式表示连续的密度和颜色场，渲染时沿每条像素射线采样并积分；3D Gaussian Splatting用数百万个显式的三维高斯椭球表示场景，渲染时把它们投影到图像平面并按深度合成。两者都对每个场景单独优化，换一个场景要重新拟合；它们的差别主要体现在渲染速度、存储和细节表达上。

\subsubsection{NeRF}
\label{sec:vision-generation-nerf}

在NeRF之前，新视角合成多依赖显式的网格、点云或离散体素。网格需要先重建出准确的表面，对毛绒、半透明和细小结构很困难；离散体素的存储随分辨率的立方增长，细节受限。\term{神经辐射场}（neural radiance field, NeRF）的动机是用一个连续函数表示场景，让存储与分辨率无关，并用可微的体渲染从照片直接优化它\scite{vision-nerf2020}。

主要想法是：用一个多层感知机把三维位置$\bm{x}\in\mathbb{R}^3$和观察方向$\bm{d}$映射为体密度$\sigma(\bm{x})\geq0$和颜色$\bm{c}(\bm{x},\bm{d})\in[0,1]^3$。密度描述该位置有多少“物质”会截获光线，只依赖位置；颜色还依赖方向，以表达高光随视角的变化。渲染一个像素时，沿它的射线累加各处的颜色贡献，这种按密度与遮挡加权累加的计算称为\term{体渲染}（volume rendering）。整个渲染过程对网络参数可微，因此可以用照片像素直接监督。

先写出连续形式。相机像素对应射线$\bm{r}(s)=\bm{o}+s\bm{d}$，$\bm{o}$是相机中心，$s$是沿射线的距离。在近远界$s_{\mathrm{n}},s_{\mathrm{f}}$之间，像素颜色为
\begin{equation}
 \begin{aligned}
 \bm{C}(\bm{r})
 &=\int_{s_{\mathrm{n}}}^{s_{\mathrm{f}}}
    T(s)\sigma(\bm{r}(s))
    \bm{c}(\bm{r}(s),\bm{d})\,\mathrm{d}s,\\
 T(s)&=\exp\!\left(
     -\int_{s_{\mathrm{n}}}^{s}\sigma(\bm{r}(u))\,\mathrm{d}u
    \right)\text{。}
 \end{aligned}
 \label{eq:vision-nerf-continuous}
\end{equation}
$T(s)$称为\term{透射率}（transmittance），表示光从相机出发到达距离$s$之前没有被截获的比例；近处密度越大，远处颜色的可见贡献越小。$\sigma$的单位是长度的倒数，使指数无量纲。

数值计算时沿射线取$n$个有序样本，假定每段内密度和颜色不变。设第$i$段长为$\delta_i$，得到
\begin{equation}
 \begin{aligned}
 \alpha_i&=1-\exp(-\sigma_i\delta_i),&
 T_i&=\prod_{j<i}(1-\alpha_j),\\
 w_i&=T_i\alpha_i,&
 \widehat{\bm{C}}&=\sum_{i=1}^{n}w_i\bm{c}_i
       +T_{n+1}\bm{c}_{\mathrm{bg}}\text{。}
 \end{aligned}
 \label{eq:vision-nerf-discrete}
\end{equation}
$\alpha_i$是光在第$i$段被截获的概率，空乘积取1，$\bm{c}_{\mathrm{bg}}$是背景色。权重$w_i$由“前面各段都没有截获”和“本段截获”两部分相乘得到，所以远处的颜色不能只按自身的$\alpha_i$相加。各$w_i$与背景权重$T_{n+1}$之和为1。

用一个两段的例子手算。取$\alpha_1=0.5,\alpha_2=0.75$，颜色分别为红$(1,0,0)$和蓝$(0,0,1)$，背景为白$(1,1,1)$。第一段权重为0.5；第二段只能使用剩下的0.5，权重为$0.5\times0.75=0.375$；背景权重为$0.5\times0.25=0.125$。合成颜色为$(0.625,0.125,0.5)$。若忘记透射率，直接用0.5和0.75加权，就会重复计算被前景遮住的贡献，权重和也超过1。图\ref{fig:vision-ray-compositing}画出了这个顺序和权重。

\begin{figure}[htbp]
 \centering
 % Two intervals, alpha .5 and .75; background transmittance .125.
\begin{tikzpicture}[x=1mm,y=1mm,font=\figurefont\small]
 \draw[BookMuted,line width=.9pt,->] (5,0)--(100,0) node[right] {$s$};
 \draw[BookInk,line width=.9pt] (5,-3)--(0,-6)--(0,6)--(5,3)--cycle;
 \node[anchor=south] at (4,9) {相机};
 \fill[BookBlue!12] (24,-7) rectangle (41,7);
 \draw[BookBlue,line width=.9pt] (24,-7) rectangle (41,7);
 \fill[BookTeal!15] (58,-7) rectangle (75,7);
 \draw[BookTeal,line width=.9pt] (58,-7) rectangle (75,7);
 \draw[BookMuted,line width=1pt] (91,-7)--(91,7);
 \node[anchor=south] at (32.5,8) {$\alpha_1=0.5$};
 \node[anchor=south] at (66.5,8) {$\alpha_2=0.75$};
 \node[anchor=north] at (32.5,-9) {$T_1=1$};
 \node[anchor=north] at (66.5,-9) {$T_2=0.5$};
 \node[anchor=south] at (91,8) {背景};
 \node[anchor=east] at (7,-31) {$w$};
 \fill[BookBlue!20] (10,-35) rectangle (50,-27);
 \fill[BookTeal!20] (50,-35) rectangle (80,-27);
 \fill[BookPaper] (80,-35) rectangle (90,-27);
 \draw[BookInk,line width=.7pt] (10,-35) rectangle (90,-27);
 \draw[BookInk,line width=.7pt] (50,-35)--(50,-27);
 \draw[BookInk,line width=.7pt] (80,-35)--(80,-27);
 \node at (30,-31) {$0.5$};
 \node at (65,-31) {$0.375$};
 \node[anchor=north] at (30,-37) {第一段};
 \node[anchor=north] at (65,-37) {第二段};
 \node[anchor=north] at (85,-37) {$0.125$};
\end{tikzpicture}

 \caption{两段射线的体渲染教学例。近处段先取得0.5的贡献权重，后段只能使用剩余0.5中的0.75，背景使用最终透射率0.125。下方长度条表示实际权重，三段和为1。颜色值人为设定，用于复算，图中形状不表示真实物体测量。}
 \label{fig:vision-ray-compositing}
\end{figure}

网络本身的结构如下。多层感知机直接读取三维坐标时倾向于学到低频函数，难以表示细纹理，因此先对输入作多频正余弦编码：每个坐标分量$p$映射为$(\sin2^0\pi p,\cos2^0\pi p,\ldots,\sin2^{L-1}\pi p,\cos2^{L-1}\pi p)$，位置取$L=10$得到60维，方向取$L=4$得到24维。这种位置编码（positional encoding）使网络能够拟合高频细节。编码前，位置的三个坐标先归一化到$[-1,1]$，方向取单位向量，各分量本身也在$[-1,1]$内。网络的逐层结构见表\ref{tab:vision-nerf-structure}：位置编码经8层宽256的全连接层（ReLU激活），按DeepSDF的做法，第5层的输出再与输入编码拼接后送入第6层；之后一层同时输出经ReLU保证非负的密度$\sigma$和一个不加激活的256维特征；该特征与方向编码拼接，经一层宽128的全连接层（ReLU）和一层sigmoid输出层得到RGB颜色。密度不读取方向，保证几何在各视角下一致。

\begin{table}[htbp]
 \centering
 \small
 \begin{tabularx}{\linewidth}{@{}p{16mm}Xp{15mm}p{9mm}p{6mm}p{15mm}@{}}
  \toprule
  阶段或层组 & 层类型 & 卷积核与步长 & 输出通道 & 重复次数 & 输出尺寸\\
  \midrule
  位置编码 & $\gamma(\bm{x})$，$L=10$ & — & 60 & 1 & 每个采样点\\
  主干1至5 & 全连接，ReLU & — & 256 & 5 & 每个采样点\\
  跳连拼接 & 第5层输出与$\gamma(\bm{x})$拼接 & — & 316 & 1 & 每个采样点\\
  主干6至8 & 全连接，ReLU & — & 256 & 3 & 每个采样点\\
  密度与特征 & 全连接；$\sigma$经ReLU，特征无激活 & — & 1与256 & 1 & 每个采样点\\
  方向分支 & 特征与$\gamma(\bm{d})$（$L=4$，24维）拼接，全连接，ReLU & — & 128 & 1 & 每个采样点\\
  颜色输出 & 全连接，sigmoid & — & 3 & 1 & 每个采样点\\
  \bottomrule
 \end{tabularx}
 \caption{NeRF的多层感知机，按原论文正文与附录图7整理。网络逐点计算，没有卷积，故卷积核一列为空、输出尺寸为单个采样点。原论文同时优化“粗”“细”两个网络，参数各自独立；论文只给出一种网络结构，没有区分二者。论文报告单个场景的网络权重约5\,MB。}
 \label{tab:vision-nerf-structure}
\end{table}

逐点查询网络很昂贵，而射线上大部分点位于空处或被遮挡。NeRF因此用粗细两级采样：粗网络在射线上分层随机取64个点，计算权重$w_i$，把它们归一化成沿射线的分段分布；再按这个分布多采128个点，细网络在全部192个点上重新计算颜色。粗细两个网络的颜色都与真实像素比较。训练时每次从所有训练照片中随机抽4096条像素射线，损失是粗、细渲染颜色与真实颜色的平方误差之和，用Adam优化，学习率从$5\times10^{-4}$指数衰减到$5\times10^{-5}$，其余参数取默认值$\beta_1=0.9,\beta_2=0.999,\epsilon=10^{-7}$；原论文一个场景训练10万到30万次迭代，在一张V100上需要1到2天。

渲染新视角时，为新相机的每个像素构造射线，重复上述采样与积分。一幅图有$p$条射线、每条查询$n$个点，计算量随$pn$增长。渲染时每条射线经粗网络查询64个点、经细网络查询192个点，共256次；原论文的合成场景每幅图有64万条射线，真实场景有76.2万条，每幅图需要1.5亿到2亿次网络查询，在一张V100上约30秒一帧，远不能实时渲染。图像损失本身不唯一确定几何：未拍到的表面、镜面反射和错误的相机都可能让新视角失败，而训练视图仍然拟合得很好。

\subsubsection{3D Gaussian Splatting}
\label{sec:vision-generation-gaussians}

NeRF的瓶颈在于每个像素都要沿射线查询上百次网络，渲染速度远不能实时，训练也要一两天。大部分查询落在空处，却都要付出完整的网络计算。3D Gaussian Splatting的动机是用显式的几何单元只在有物体的地方表示场景，并用光栅化代替射线采样，从而实现实时渲染和更快的训练\scite{vision-gaussian2023}。

主要想法是：用一组三维高斯椭球表示场景，每个高斯有位置、形状、不透明度和随方向变化的颜色；渲染时把每个高斯投影成图像上的二维椭圆，按深度从前到后做透明度合成。这个合成公式与NeRF的离散体渲染具有相同的结构，但只在高斯覆盖的像素上计算，不需要沿射线查询网络。

每个高斯的参数包括：中心$\bm{\mu}\in\mathbb{R}^3$；协方差$\bm{\Sigma}=\bm{R}_{g}\bm{S}\bm{S}^{\top}\bm{R}_{g}^{\top}$，其中旋转$\bm{R}_{g}$由单位四元数（4个数）表示、$\bm{S}$是由3个正尺度组成的对角阵，这种分解保证协方差始终半正定；不透明度$o\in(0,1)$，经sigmoid得到；颜色用3阶\term{球谐函数}（spherical harmonics）系数表示，即在方向球面上用16个基函数的组合表示每个颜色通道随方向的变化，共48个系数。一个高斯共有$3+4+3+1+48=59$个参数，一个场景通常有一百万到五百万个高斯。论文没有把59这个总数写出来，它是按上述参数化逐项相加得到的；场景规模“一百万到五百万个高斯”是论文对所测全部场景的报告。各参数的维数、约束方式与初始化见表\ref{tab:vision-gaussian-parameters}。

\begin{table}[htbp]
 \centering
 \small
 \begin{tabularx}{\linewidth}{@{}p{20mm}p{9mm}XX@{}}
  \toprule
  参数 & 维数 & 约束方式 & 初始化\\
  \midrule
  中心$\bm{\mu}$ & 3 & 无约束；学习率按指数衰减 & 运动恢复结构的稀疏点\\
  旋转$\bm{q}$ & 4 & 单位四元数，使用前归一化 & 论文未给出\\
  尺度$\bm{s}$ & 3 & 经指数函数得到正尺度 & 各向同性，取到最近三个点的平均距离\\
  不透明度$o$ & 1 & 经sigmoid落在$[0,1)$ & 论文未给出\\
  球谐系数 & 48 & 3阶，每个颜色通道16个；先只优化0阶，每1000次迭代开放一阶 & 论文未给出\\
  \bottomrule
 \end{tabularx}
 \caption{3D Gaussian Splatting中单个高斯的参数，取自原论文。59个参数在每个高斯上独立优化，场景中高斯的数目由自适应密度控制在训练中改变。}
 \label{tab:vision-gaussian-parameters}
\end{table}

渲染时，先把三维高斯变换到相机坐标，再在中心附近把透视投影线性化，得到二维协方差
\begin{equation}
 \bm{\Sigma}_{2}
 =\bm{J}\bm{R}_{c}\bm{\Sigma}
    \bm{R}_{c}^{\top}\bm{J}^{\top}\text{，}
 \label{eq:vision-gaussian-projection}
\end{equation}
其中$\bm{R}_{c}$把世界方向转到相机坐标，$\bm{J}\in\mathbb{R}^{2\times3}$是该点处投影的Jacobian矩阵。透视投影不是线性的，所以这只是局部近似。设投影中心为$\bm{\mu}_{2,i}$，第$i$个高斯在像素$\bm{u}$处的有效不透明度为
\begin{equation}
 a_i(\bm{u})=o_i\exp\!\left[
 -\frac12(\bm{u}-\bm{\mu}_{2,i})^\top
 \bm{\Sigma}_{2,i}^{-1}(\bm{u}-\bm{\mu}_{2,i})
 \right]\text{。}
 \label{eq:vision-gaussian-alpha}
\end{equation}
把覆盖该像素的高斯按深度从近到远排列，颜色为$\sum_i a_i\bm{c}_i\prod_{j<i}(1-a_j)$，必要时加上剩余背景。与式\eqref{eq:vision-nerf-discrete}对照，$a_i$对应$\alpha_i$，$\prod_{j<i}(1-a_j)$对应透射率$T_i$；区别在于$a_i$来自投影椭圆和学到的不透明度，而不是对密度的分段积分。

为了在GPU上高效完成，原方法把屏幕分成$16\times16$像素的图块，每个高斯按它覆盖的图块复制若干份，用“图块编号加深度”作为键做一次全局基数排序，然后每个图块并行地按深度合成。排序按高斯中心的深度进行，同一像素内的次序只是近似，不做精确的表面求交。

初始化来自运动恢复结构给出的稀疏点：每个点成为一个高斯，初始为各向同性，尺度取到最近三个点的平均距离。损失为$L=(1-\lambda)L_1+\lambda L_{\text{D-SSIM}}$，$\lambda=0.2$，其中D-SSIM是由结构相似性构造的差异项，公开实现取1减SSIM。为了稳定，优化先在缩小为四分之一的图像分辨率上进行，第250次和第500次迭代后各放大一倍。优化同时更新所有高斯的位置、尺度、旋转、不透明度和球谐系数；球谐阶数从0开始，每1000次迭代增加一阶，直到3阶。

稀疏点不足以表达所有细节，因此需要\term{自适应密度控制}（adaptive density control），在训练中改变高斯的数目。优化热身之后每100次迭代检查一次：视图空间位置梯度的平均幅值超过阈值0.0002的高斯，说明该处重建不足。若高斯较小，就复制一份并沿梯度方向移动，用于覆盖缺失的结构；若高斯较大，就分裂成两个，尺度各除以1.6，新位置从原高斯中采样，用于细化过大的区域。不透明度低于阈值$\epsilon_\alpha$的高斯被删除（论文未给出$\epsilon_\alpha$的数值），世界空间中过大或在视图中覆盖范围过大的高斯也定期删除；每3000次迭代把所有不透明度重置为接近零，让多余的高斯在后续优化中自然淘汰。例如玩具狗的耳朵边缘若被一个大高斯覆盖，耳朵与背景的颜色会被混在一起；分裂后两个较小的高斯可以分别拟合耳朵和背景，但这仍需后续优化，分裂本身不代表测得了新的表面点。原论文训练3万次迭代，在Mip-NeRF360数据集的场景上约需40分钟。

渲染新视角只需重新投影、排序和合成，原论文在这些场景上以超过100帧每秒的速度渲染，而对照的Mip-NeRF360每帧需要数秒以上。代价是存储：一个场景模型约数百兆字节，而NeRF类模型只需几兆字节。高斯集合也不是带拓扑的三角网格，没有物理材质，渲染逼真不说明它可以直接用于碰撞检测或精确编辑。

两种场景表示的比较见表\ref{tab:vision-generation-scene-comparison}。固定任务为“从绕玩具狗拍摄的照片渲染新视角”，评价口径为保留视图上的画质、训练与渲染时间以及存储。3D Gaussian Splatting相对NeRF改变了表示（隐式网络改为显式高斯）和渲染方式（射线采样改为投影排序），换来的是实时渲染和以分钟计的训练，付出的是数百倍的存储和对初始点云的依赖；两者都无法为没有拍到的表面提供观测支持。

\begin{table}[htbp]
 \centering
 \small
 \begin{tabularx}{\linewidth}{@{}l X X@{}}
  \toprule
   & NeRF & 3D Gaussian Splatting\\
  \midrule
  场景表示 & 多层感知机隐式表示的密度与颜色场 & 百万量级的显式三维高斯\\
  渲染方式 & 每条射线查询网络两百余次 & 投影、按图块排序、透明度合成\\
  初始化 & 随机初始化网络 & 运动恢复结构的稀疏点\\
  训练时间（原论文设置） & 单场景1到2天 & 单场景约40分钟\\
  渲染速度 & 逐射线查询，远低于实时 & 每秒百帧以上\\
  存储 & 几兆字节量级的网络权重 & 数百兆字节的高斯参数\\
  \bottomrule
 \end{tabularx}
 \caption{两种场景表示在同一新视角合成任务下的比较。时间与存储取自各自原论文报告的设置，硬件与数据集不同，只表示数量级差别。两者都需逐场景优化，要求静态场景与准确相机，且不能为未观测表面提供证据。}
 \label{tab:vision-generation-scene-comparison}
\end{table}

\subsection{新视角的渲染与核验}
\label{sec:vision-generation-rendering-check}

渲染一个新视角时，给出目标相机的内参和外参，按上面的方法得到图像。核验的关键是使用在拟合前就留出的测试视图：在这些相机下渲染，与真实照片比较PSNR、SSIM或LPIPS（见\ref{chap:vision-image-restoration}章），同时检查轮廓、细节和反光。测试视图的像素不能用于调节场景参数；若相机标定使用了测试图像，也应在协议中说明。

单个测试视图还可能漏掉不连续现象。沿一条连续的相机路径渲染视频，可以观察遮挡变化时是否出现漂浮的斑点、外观是否突然跳变。路径上没有真实照片的位置只能做一致性观察，无法计算误差。几何精度还需要点云、深度等独立测量，逼真连贯的渲染动画不能直接证明表面准确。对玩具狗的背面，若拍摄没有覆盖，两种表示都可能把错误拟合延伸过去。若用途只是展示某些可见视角，应注明适用的视角范围；若需要完整的三维资产，就需要补拍背面并核验几何。动态场景的时间变化由\ref{chap:vision-multimodal-generation}章扩展。

\section{生成结果的评价}
\label{sec:vision-generation-evaluation}

前四节分别给出了各自任务的核验方法，本节把它们放在同一组维度下比较。四类任务的参照不同：创建主要对照条件，转换和编辑还要对照输入，新视角要对照真实观测。相应地，评价分为三个维度：条件是否被遵循，输入中应保持的内容是否被保持，以及结果有多少受观测支持、付出了多少成本。\ref{sec:vision-generation-challenges}节定义的FID、CLIP分数、人工偏好、保留区变化和PSNR等指标分别落在这三个维度之下，本节讨论怎样组合使用它们。比较模型时还要固定提示、参考、分辨率、候选数量和修订预算，否则无法判断差异来自模型本身还是选择过程。

\subsection{条件遵循的核验}
\label{sec:vision-generation-condition-check}

回到玩具狗与杯子的要求记录，逐项检查存在、数量、颜色归属和左右关系；对主体适配，还要与参考照片比较身份。可以用检测器统计对象和位置，用区域特征比较主体，用文字读取器检查杯子上的字符。DreamBooth的评价是一个例子：它用DINO特征相似度衡量主体保真、用CLIP图文相似度衡量提示遵循，两个指标同时报告，防止“复制参考照片”这种退化在单一指标上得高分。这些检查工具本身也会漏检或误读，自动分数应注明工具的输入和输出，并抽样人工复核，不能把工具的错误算到生成器头上。

整体图文相似度适合衡量粗粒度相关性，不能单独证明复杂条件全部成立。汇总时除平均分外，应保留关键分项和失败率，例如对象数量正确率、属性绑定正确率。多个属性中一个绑定错误，即使只影响一小块区域，也会使整个请求失败。对使用视觉指令的编辑，还要检查结果是否随指令改变，而不只是检查目标对象是否出现。

\subsection{输入保持的核验}
\label{sec:vision-generation-change-check}

转换和编辑都需要同时检查“该改的改了”和“该留的留了”。原样返回原图，可以使式\eqref{eq:vision-edit-preservation}的保留误差为零，却没有执行请求；把杯子换成花瓶的同时改变了狗的斑纹，完成了目标却破坏了保留条件。两类错误应分别记录。多轮编辑还应分别报告当前轮修改的完成率、历次已接受修改的保持情况，以及相对最初主体的漂移。

保持的含义随任务而变。照片转成绘画后，像素距离大是预期的现象，应该比较轮廓、对象数量和身份，而不是逐像素误差；LPIPS这类感知距离也要结合允许的风格变化来解释。新视角中，同一物体因投影而在画面中移动，也不是编辑漂移；评价必须与目标相机下的真实照片比较，而不是与原相机的图像逐像素相减。

\subsection{观测支持与制作成本的核验}
\label{sec:vision-generation-costs}

观测支持回答“结果中的哪些内容有证据”。新视角合成的输出应标明哪些区域被训练视图覆盖；DreamBooth生成的背面若没有参考照片，就只是类别先验的推测；修补区域中的花瓶是按请求创作的内容，而不是被遮挡物体的恢复。三者都可以产生逼真图像，但真实性的依据不同，交付时应分别说明。

成本则回答“得到这一幅结果花了多少”。一项可交付成果的成本包括条件制作（画姿态图、写掩码）、模型适配（ControlNet支路训练、DreamBooth微调）、视图采集与标定、场景拟合、候选生成与筛选以及修订轮次。一次渲染只需几毫秒，不表示完成整个场景也只需一次前向；从一百幅候选中选出一幅，也不表示单次生成就稳定满足要求。报告时应把这些成本与结果一起列出，使不同系统在相同预算下比较。

\section{本章小结}
\label{sec:vision-generation-summary}

开篇的四个请求，现在可以分别回答。它们共同面对四个挑战：一个条件对应许多合格图像，多个对象的属性容易串换，局部修改与内容保持互相牵制，有限的照片难以保证三维一致；评价时则需要分别用FID、CLIP分数、人工偏好、保留区变化和PSNR等指标回答不同的问题。在雪地里画这只玩具狗，属于新图像的创建：文字经LDM或unCLIP约束语义，ControlNet通过冻结主干、复制支路和零卷积接入逐像素对齐的空间图，DreamBooth用少量照片和类别先验保持把具体主体写进模型参数，多参考生成则需要显式规定每幅参考的角色。采样中的引导以多样性和计算为代价提高条件贴合度，生成后的重排只能在已有候选中选择，一轮反思修订可以纠正某些空间关系，但需要预算上限和独立验收。转成水彩画属于已有图像的转换：pix2pix依赖配对目标，CycleGAN用循环一致性替代配对却不能唯一确定语义，Johnson感知损失网络用固定特征网络把风格转换变为一次前向计算。把杯子换成花瓶属于指定内容的编辑：泊松图像编辑在梯度域中把现成的花瓶无缝贴入，Criminisi样例修补按置信度与等照度线强度决定填充顺序、用图中已有的块补洞，PatchMatch用随机初始化、传播和随机搜索把最近邻块的搜索加速到可交互，这三种经典方法让保留区严格不变，却只能重组已有内容；GLIDE修补模型靠掩码与可见区域约束范围，InstructPix2Pix用合成的编辑三元组和双条件引导实现只需一句指令的编辑，多轮编辑还要在模型之外维护版本状态。给出正后方视图属于新视角合成：多视图几何先给出相机与稀疏点，光场渲染需要密集拍摄才能直接插值光线；NeRF用隐式场与体渲染、3D Gaussian Splatting用显式高斯与投影合成拟合同一组观测，后者以存储换来实时渲染。

评价贯穿这四类任务：条件遵循、输入保持、观测支持与成本应分别核验，不能用一个自然度分数或全局相似度代替。仍有一些问题本章没有解决。多主体和多参考的角色绑定还缺少经过验证的通用模型；精确数量、文字和复杂空间关系仍是文字条件生成的常见失败；多轮编辑的主体漂移只能靠外部回退限制，不能从模型内部消除；缺少视角的观测场景无法由本章的拟合方法补全，需要\ref{chap:vision-multimodal-generation}章引入生成先验，那时又必须重新区分哪些内容有观测依据。视频生成还会增加时间连续性和音画同步的约束，本章关于条件、保持与观测证据的区分在那里仍然适用。

\paragraph{扩展阅读}
创作条件方向，SDXL在LDM的基础上扩大去噪网络，改用两个文字编码器，并把训练图像的原始尺寸和裁剪位置也作为条件输入，再用一个细化模型在同一潜空间中补充细节，可用来观察本章的LDM结构在规模扩大后还需要补充哪些条件\scite{vision-sdxl2024}T2I-Adapter同样向冻结的文生图模型接入草图、深度或姿态等空间图，但只训练一个轻量的适配网络，不像ControlNet那样复制编码器支路，多个适配器的输出还可以相加组合\scite{vision-t2i-adapter2024}IP-Adapter为参考图像单独设置一组与文字交叉注意力并行的交叉注意力，使参考图像不经微调即可作为条件；Textual Inversion则冻结整个模型，只从少量照片中学习一个新词的嵌入向量。二者与DreamBooth一样要把具体主体或风格带入生成，但分别把信息放在推理时的条件和词嵌入中，而不是写进模型参数，保真程度和成本也随之不同\scite{vision-ip-adapter2023,vision-textual-inversion2023}

编辑方向，SDEdit先把用户的笔画草图或待改图像加噪到某个中间时刻，再用扩散模型去噪回来；加噪越多，结果越自然，与输入的对应却越弱。它不需要配对数据或额外训练，可与本章pix2pix和CycleGAN的学习式转换对照\scite{vision-sdedit2022}本章的InstructPix2Pix只在两幅生成图像之间使用Prompt-to-Prompt；Null-text Inversion先用确定性采样把真实照片反演为一条去噪轨迹，再逐步优化无分类器引导中空文本条件的嵌入，使加了引导的采样仍能准确重建原图，从而把同样的注意力控制用于真实照片的编辑\scite{vision-null-text-inversion2023}新视角合成方向，Mip-NeRF把每个像素对应的射线换成锥体，用积分位置编码描述锥体截面覆盖的区域，缓解NeRF在拍摄距离和分辨率变化时的混叠；Instant-NGP用多分辨率哈希表存储可学习的特征，把NeRF类模型的训练从一两天缩短到秒级至分钟级，从隐式表示一侧回应了3D Gaussian Splatting所针对的速度问题\scite{vision-mip-nerf2021,vision-instant-ngp2022}

评价方向，Kynkäänniemi等人发现，只要让生成图像在ImageNet类别上的分布与真实图像对齐，即使画质没有改善，FID也会明显下降，说明FID在相当程度上反映的是Inception特征中的类别信息\scite{vision-fid-imagenet-classes2023}Jayasumana等人认为Inception特征和正态分布假设都不适合评价现代文生图模型，改用CLIP特征和不假设分布形式的最大均值差异，提出CMMD\scite{vision-cmmd2024}GenEval用目标检测器核对生成图像中的物体是否出现，以及数量、颜色和相对位置是否符合提示，把本章所说的属性绑定和计数错误变成可以自动统计的检查项\scite{vision-geneval2023}

